\documentclass[a4paper,10pt]{article}
\usepackage[top=25mm,bottom=24mm,left=27mm,right=27mm,headheight=15pt,headsep=8mm,footskip=12mm]{geometry}
\usepackage{fontspec,amsmath,amsthm,unicode-math}
\usepackage[english,french]{babel}
\usepackage{microtype,xcolor,fancyhdr,titlesec,booktabs,longtable,array,tabularx}
\usepackage{graphicx,tikz,enumitem,float,needspace,fvextra}
\usetikzlibrary{arrows.meta,calc,positioning,decorations.pathreplacing,matrix}
\usepackage[unicode,hidelinks]{hyperref}
\usepackage[nameinlink,noabbrev]{cleveref}
\setmainfont{STIXTwoText-Regular.otf}[BoldFont=STIXTwoText-Bold.otf,ItalicFont=STIXTwoText-Italic.otf,BoldItalicFont=STIXTwoText-BoldItalic.otf]
\setmathfont{STIXTwoMath-Regular.otf}
\hypersetup{pdftitle={Structure constitutive du vide et propagation électromagnétique},pdfauthor={Oumar Haidara Fall; Rubén Ramos Balsa},pdfsubject={Dérivation de la constante de Catalan, de la perméabilité et de la permittivité}}
\newcommand{\Autores}{Oumar Haidara Fall \textperiodcentered\ Rubén Ramos Balsa}
\newcommand{\RR}{\mathbb R}\newcommand{\ZZ}{\mathbb Z}\newcommand{\QQ}{\mathbb Q}\newcommand{\CC}{\mathbb C}\newcommand{\FF}{\mathbb F}
\newcommand{\Id}{\operatorname{id}}\newcommand{\tr}{\operatorname{tr}}
\newcommand{\Span}{\operatorname{span}}\newcommand{\Cyl}{\operatorname{Cyl}}
\newcommand{\square}{\mdlgwhtsquare}
\newcommand{\TPK}{\mathrm{TPK}}\newcommand{\APP}{\mathrm{APP}}\newcommand{\TRIT}{\{-1,0,+1\}}
\newcommand{\HMT}{\mathrm{HMT}}\newcommand{\Tr}{\operatorname{Tr}}\newcommand{\dd}{\mathrm d}
\newtheorem{theorem}{Théorème}[section]
\newtheorem{lemma}[theorem]{Lemme}
\newtheorem{proposition}[theorem]{Proposition}
\newtheorem{corollary}[theorem]{Corollaire}
\theoremstyle{definition}\newtheorem{definition}[theorem]{Définition}\newtheorem{axiom}[theorem]{Axiome}
\theoremstyle{remark}\newtheorem{remark}[theorem]{Remarque}
\numberwithin{equation}{section}
\setlength{\parindent}{1em}\setlength{\parskip}{3pt plus 1pt minus .5pt}
\setlength{\emergencystretch}{2em}
\clubpenalty=10000
\widowpenalty=10000
\displaywidowpenalty=10000
\renewcommand{\normalsize}{\fontsize{10.5}{13.4}\selectfont}\normalsize
\setlist{nosep,topsep=5pt}
\titleformat{\section}{\fontsize{13}{16}\selectfont\bfseries}{\thesection}{.65em}{}
\titleformat{\subsection}{\fontsize{11.2}{14}\selectfont\bfseries}{\thesubsection}{.65em}{}
\titleformat{\subsubsection}{\normalsize\bfseries}{\thesubsubsection}{.65em}{}
\AddToHook{cmd/section/before}{\clearpage}
\AddToHook{cmd/subsection/before}{\Needspace{7\baselineskip}}
\AddToHook{cmd/subsubsection/before}{\Needspace{6\baselineskip}}
\setcounter{tocdepth}{2}
\titlespacing*{\section}{0pt}{17pt}{9pt}\titlespacing*{\subsection}{0pt}{13pt}{6pt}
\pagestyle{fancy}\fancyhf{}
\fancyhead[L]{\fontsize{8}{10}\selectfont Structure constitutive du vide et propagation électromagnétique}
\fancyhead[R]{}
\fancyfoot[C]{\thepage}\renewcommand{\headrulewidth}{.2pt}
\raggedbottom
% BEGIN NUCLEO_ADITIVO_20260921_PREAMBULO
\makeatletter
\@ifundefined{example}{\theoremstyle{definition}\newtheorem{example}[theorem]{Exemple}}{}
\@ifundefined{notatabla}{\newcommand{\notatabla}[1]{\par\smallskip{\footnotesize #1\par}}}{}
\makeatother
\providecommand{\FF}{\mathbb F}

% END NUCLEO_ADITIVO_20260921_PREAMBULO
\newif\ifHMTDeltaEnglish\HMTDeltaEnglishfalse
\newcommand{\HMTLang}[2]{\ifHMTDeltaEnglish#2\else#1\fi}
\makeatletter
\renewcommand*{\l@subsection}{\@dottedtocline{2}{1.5em}{3.3em}}
\makeatother
% Apertura independiente del índice: corrección quirúrgica 20260928.
\let\HMTIndiceOriginal\tableofcontents
\renewcommand{\tableofcontents}{\clearpage\begingroup\setlength{\parskip}{0pt}\fontsize{9.2}{10.5}\selectfont\HMTIndiceOriginal\endgroup\clearpage}
\begin{document}
\thispagestyle{plain}
\begingroup\centering
\setlength{\parskip}{0pt}
{\fontsize{10}{12}\selectfont Manuscrit de compilation pour le transfert académique\par}
\vspace{6mm}
{\fontsize{17}{20.5}\selectfont\bfseries Structure constitutive du vide et propagation électromagnétique\par}
\vspace{6pt}
{\fontsize{11}{13.5}\selectfont Dérivation de la constante de Catalan, de la perméabilité et de la permittivité\par}
\vspace{8pt}
{\fontsize{11.5}{14}\selectfont\Autores\par}
\vspace{3pt}
{\fontsize{8}{10}\selectfont Coautorat sans préséance de contribution\par}
\vspace{3pt}
{\fontsize{8}{10}\selectfont Intégration et compilation documentaires générées par ChatGPT - Astra.\par}
\par\endgroup

\begin{center}\bfseries Résumé\end{center}
\phantomsection
\addcontentsline{toc}{section}{Résumé}
\begingroup
\fontsize{9.5}{10.7}\selectfont
\setlength{\parskip}{1.5pt plus .5pt minus .5pt}
% Composición local: conserva resumen y palabras clave en la apertura.
\fontsize{10.2}{11.6}\selectfont
\setlength{\parskip}{1pt plus .5pt minus .5pt}
\enlargethispage{2\baselineskip}
\noindent
On construit une réponse constitutive positive et inversible du vide
électromagnétique : deux canaux orientés déterminent conjointement les
coordonnées adimensionnelles de permittivité, de perméabilité, d’impédance et de
vitesse de propagation. On démontre leurs relations exactes et la
récupération unique de la paire angulaire qui les produit. La même réponse
permet de reconstruire la forme de l’ellipse d’action et d’exprimer la
fonctionnelle de Barbero--Immirzi au moyen de ses deux valeurs propres étiquetées.
Le résultat relie ainsi la détermination des valeurs à
l’information structurelle conservée par leur évaluation.

La construction part d’un noyau formel mathématique appelé
holographie modulaire triadique, ci-après HMT. \emph{All Possible Paths}
(APP) est la structure arithmétique des chemins sur le support toroïdal
$(\mathbb Z/9\mathbb Z)^2$ ; ses déplacements cardinaux forment le
graphe de Cayley additif, et ses feuillets conservent les évaluations somme et
produit avec résidu et quotient. TRIT est la signature orientée
$\{-1,0,+1\}$ qui distingue les régimes locaux. Le
\emph{Topological Phase Kernel} (TPK) compose la sélection, le transport,
la mise à jour de mémoire et le retour sur cet état enrichi. Le
prolongement nonadique et la clôture dodécaphasique engendrent les coordonnées
corrélées et les registres dont proviennent l’action, la paire
angulaire et les incidences utilisées par la réponse.

Les secteurs $90$ et $120$ du transport $T$ fixent
$\mathcal V=(I-T^{90})(I-T^{120})^{-1}$. Sa fonction scalaire
$f(s)=(1+s+s^2)/(1+s+s^2+s^3)$ est une bijection décroissante de $(0,1)$
sur $(3/4,1)$ : cette propriété établit l’inversion cubique et son
domaine. L’échange de feuillets permute la permittivité et la perméabilité,
inverse l’impédance et conserve la vitesse. Dans la même coordinatisation dimensionnelle,
on démontre en outre l’identité entre la vitesse constitutive et la
vitesse cinématique de l’horloge.

La réalisation sur le cycle à douze phases apporte le défaut de trace
$-1/12$, une représentation de $f$ comme quotient de déterminants et un quart de tour
de carré $-I$. La double intégration logarithmique de sa résolvante
produit le moment de Catalan. Dans la section de Pell $2-\sqrt3$, on
démontre une identité exacte avec la constante de Catalan et
$\log(2+\sqrt3)$ ; les deux arguments constitutifs maintiennent leur
détermination angulaire propre. L’inversion de $f$ compose ensuite la
réponse du vide avec la fonctionnelle orientée de Barbero--Immirzi.
La restitution intégrale démontre également que la fonction constitutive
détermine le moment orienté complet sous la condition
$\mathcal C_2(s)\to0$ à l’origine ; sa limite en $s=1$ retrouve
la constante de Catalan. On conserve ainsi une correspondance fonctionnelle
avec son orientation et ses arguments, outre ses évaluations.

La polarisation discrète s’obtient à partir du bilan borné des vacances,
et le préfacteur électronique $\sqrt3/4$ à partir de la norme du commutateur des
projections associées aux fibres de somme et de produit du support
arithmétique tridimensionnel. La section
positive de charge, construite avec $\alpha$, l’action et l’impédance,
détermine les relations de von Klitzing, de conductance, de flux magnétique
et de Josephson, y compris leurs identités croisées et leur covariance sous
le retour de l’action. Les réalisations dimensionnelles conservent les
sections d’action, de charge et de vitesse et leurs changements d’unités. L’évaluation
numérique correspond aux coordonnées internes ;
sa comparaison physique conserve explicitement la représentation et la
coordinatisation dimensionnelle utilisées.

La séparation positive du mode commun et du mode orienté spécifie
l’information conservée par la normalisation unimodulaire. La résolution
de Fourier de l’incidence dodécaphasique identifie quatre modes et
leurs poids exacts, en distinguant le défaut de rang, les multiplicités
signées et le coefficient singulier du retour.

L’inversion induit également une loi associative de composition des
réponses, avec une coordonnée angulaire additive. Dans la réalisation
électrofaible à l’arbre, les lecteurs de jauge et de Higgs retrouvent
conjointement les canaux et l’impédance, avec un domaine explicite.

\smallskip
\noindent\textbf{Mots-clés :} détermination structurelle des constantes ;
vide électromagnétique ; réponse constitutive ; inversion spectrale ;
holographie modulaire triadique ; orientation des feuillets ; constante de Catalan ;
unité de Pell ; Barbero--Immirzi ; polarisation ; quantification électrique.
\par\endgroup

\tableofcontents
\clearpage
\begingroup\setlength{\parskip}{2pt}\fontsize{10}{11.8}\selectfont
\clearpage
\section{Introduction}
\label{iii:sec:introduccion}

Une relation constitutive organise la réponse d’un système, mais sa
forme fonctionnelle et ses coefficients posent des questions différentes. Une
fois les coefficients connus, l’équation permet de calculer des réponses ;
déterminer pourquoi ces coefficients ont conjointement leurs valeurs
nécessite de décrire la structure qui les fixe. Ce travail aborde cette
seconde question au moyen d’une construction mathématique à deux canaux
orientés. Sa thèse est que la structure précédente détermine
l’équation et que l’évaluation conserve une partie récupérable de cette
structure. Le vide électromagnétique constitue ici une réalisation
précise de cette thèse : la permittivité, la perméabilité, l’impédance et la vitesse
s’obtiennent comme quatre lectures reliées d’une même réponse.

Le résultat principal comprend tant la construction directe que
l’inverse. Dans le domaine contractant, l’opérateur constitutif est positif,
ses valeurs propres appartiennent à un intervalle déterminé et ses quatre
coordonnées satisfont des identités exactes. La stricte monotonie de la
fonction de réponse fournit un inverse unique ; avec les feuillets
étiquetés, cet inverse retrouve les canaux et la paire angulaire. Cette
récupération permet de composer trois descriptions : la réponse du
vide, la forme de l’ellipse d’action et la fonctionnelle de
Barbero--Immirzi. La comparaison ne se limite donc pas à des valeurs
numériques : elle identifie les opérations qui transmettent l’orientation et
l’information spectrale entre ces descriptions.

\subsection{Noyau générateur et provenance de la réponse}

Le modèle se construit à partir d’un noyau formel mathématique appelé
holographie modulaire triadique, ci-après HMT. Son premier objet,
\emph{All Possible Paths} (APP), réunit des chemins et des évaluations
arithmétiques sur les $81$ positions du support
$X_9=(\mathbb Z/9\mathbb Z)^2$. Les translations de générateurs
$\{(1,0),(-1,0),(0,1),(0,-1)\}$ définissent son graphe de Cayley additif,
un réseau cyclique bidimensionnel de réalisation toroïdale. Les
évaluations somme et produit agissent sur ce support ; chacune conserve
son résidu et son quotient. La structure des chemins et les évaluations
ont ainsi des fonctions distinctes au sein d’APP : les déplacements
spécifient les chemins, tandis que les feuillets arithmétiques déterminent les
données qu’ils transportent.

Le deuxième objet, TRIT, est la signature ternaire orientée
$\{-1,0,+1\}$, obtenue au moyen de la représentation équilibrée des
classes modulo $3$. Sa réalisation quadratique distingue les types
hyperbolique, parabolique et elliptique, respectivement. Le troisième objet
est le \emph{Topological Phase Kernel}, ci-après TPK : il compose la
sélection de phase, le transport cardinal et la mise à jour de mémoire.
La dynamique agit sur les positions, les feuillets, les quotients, l’orientation,
les préfixes et les données de frontière. Le retour d’une phase conserve
l’histoire accumulée, de sorte qu’un tour peut fermer une coordonnée
sans restituer l’état complet.

Le prolongement nonadique transmet cette structure au moyen de préfixes et de
frontières compatibles. Dans son développement conjoint sont engendrées les
coordonnées $\pi_{\rm HMT}$, $\varphi_{\rm HMT}$,
$e_{\rm HMT}$ et $\alpha_{\rm HMT}$. Le registre dodécaphasique $K$
et son évaluation périodique interviennent dans la détermination de
$\alpha$ au moyen de la clôture dodécaphasique ; la représentation analytique complète exprime cette même
coordonnée dans la classe de queues construite. On obtient ensuite
les sections d’action et la paire angulaire. L’utilisation de ces résultats
comme antécédents conserve leurs démonstrations au sein de l’article.
Cette séquence fixe la généalogie des canaux avant d’évaluer la
réponse constitutive.

L’orientation des feuillets et leur incidence aréale--volumétrique remplissent
des rôles complémentaires. La première distingue le retour d’un chemin
de la restauration de son orientation ; la seconde provient du calendrier
des modes et de sa matrice d’incidence. La fréquence chronologique, la
mesure du système symbolique et la marque régionale de frontière maintiennent
leurs définitions propres. Ces distinctions, développées dans la
section~\ref{iii:sec:hojas}, précisent l’information qui parvient au
transport angulaire et évitent de la remplacer par un rapport scalaire isolé.

\clearpage
\subsection{Détermination constitutive et récupération de la forme}

Avec $T=q_+P_++q_-P_-$ et $0<q_\pm<1$, les incidences temporelles
$90$ et $120$ déterminent
\[
\mathcal V=(I-T^{90})(I-T^{120})^{-1}.
\]
En écrivant $S=T^{30}$, l’équation
$\mathcal V(I+S+S^2+S^3)=I+S+S^2$ montre la règle qui fixe la
réponse pour ce transport. La section~\ref{iii:sec:constitutiva}
démontre son existence et son unicité et établit la fonction scalaire
$f:(0,1)\to(3/4,1)$. Ses valeurs $r_\pm=f(q_\pm^{30})$
produisent
\[
\widehat\varepsilon=r_+^2,\qquad
\widehat\mu=r_-^2,\qquad
\widehat Z=r_-/r_+,\qquad
\widehat c=(r_+r_-)^{-1}.
\]
La positivité et les lectures de produit et de quotient déterminent les deux
carrés une fois les relations constitutives fixées. L’unicité
affirmée appartient à ce système explicite de règles et à son domaine.

L’inversion cubique retrouve $q_\pm$ puis les coordonnées
$x=-\frac12\log(q_+q_-)$,
$y=\frac12\log(q_-/q_+)$. Avec les étiquettes de feuillet conservées,
l’impédance et la vitesse suffisent à cette reconstruction.
La proposition~\ref{iii:prop:forma-LC} obtient ensuite
l’anisotropie de l’ellipse d’action et l’impédance inductive--capacitive
relative. Les références d’action, de fréquence et d’impédance complètent
sa réalisation dimensionnelle. La proposition~\ref{iii:prop:identidad-velocidad}
prouve, dans la même coordinatisation, $c_{\rm int}=c_{\rm vac}$ et relie la
longueur de trajectoire à la durée élémentaire. L’inversion
reconstruit ainsi des antécédents spectraux déterminés, tandis que le chemin
et la mémoire restent dans l’état enrichi dont ils proviennent.

Cet inverse induit également une loi de composition entre réponses :
la multiplication des transports s’exprime au moyen de l’opération
transportée et admet une coordonnée additive. La
section~\ref{amp-higgs-seccion} compose ensuite les mêmes canaux avec
les signatures angulaires de jauge et de Higgs. L’inverse conjoint, son
domaine exact et la reconstruction de l’impédance conservent les conditions
d’échelle, de raffinement et de normalisation de cette réalisation.

\subsection{Cycle, constantes et fonctionnelles orientées}

La fonction $f$ admet une réalisation sur le cycle à douze phases.
Ses sous-espaces de dimensions $3$ et $4$ produisent le quotient de
déterminants ; leur différence de trace normalisée donne $-1/12$.
Dans un plan orienté du même cycle, un quart de tour $J$ satisfait
$J^2=-I$ et sa résolvante détermine $s/(1+s^2)$. Deux intégrations
logarithmiques, avec leurs conditions à l’origine, engendrent le moment
$\mathcal C_2(s)$ et sa valeur à l’extrémité, la constante de Catalan.
Le théorème~\ref{iii:thm:catalan-pell} relie exactement
l’évaluation en $2-\sqrt3$ à cette valeur à l’extrémité et à
$\log(2+\sqrt3)$. La section algébrique de Pell et les arguments
angulaires $q_\pm^{30}$ sont des évaluations distinctes de la même famille.

Les traces des puissances du transport déterminent en outre un tableau
périodique qui reconstruit le logarithme du quotient constitutif.
Sa série de moments, la représentation de Mellin et les bornes du reste
sont démontrées dans la section~\ref{amp-afin-sec}. La composition diagonale
et la tensorielle maintiennent explicite la profondeur conservée par chaque
opération ; le retour spectral demeure distinct du retour de l’état.

Le théorème~\ref{iii:thm:restitucion-funcional-catalan} complète la
relation au moyen d’une restitution intégrale : la fonction constitutive
fixe le moment orienté entier et la condition de limite nulle à
l’origine détermine sa normalisation. La preuve distingue la récupération
de cette famille de l’inversion de ses deux évaluations de canal.
Les deux opérations conservent la lecture de vitesse
$\widehat c=(\det\mathcal V)^{-1}$ et sa trace normalisée, qui relient
le transport constitutif à la durée et à la longueur du même état.

La composition avec Barbero--Immirzi conserve une information supplémentaire :
le signe relatif des canaux. L’incidence hexade--octade et le
calendrier fixent les termes de la fonctionnelle angulaire ; l’inverse de
$f$ permet de l’exprimer sur $\mathcal V$. Le
théorème~\ref{iii:thm:barbero-factorization} démontre cette
factorisation. Un échange de canaux par rapport à une marque
d’orientation fixe inverse le signe de la fonctionnelle ; une conjugaison
simultanée de l’opérateur et de la marque le conserve. De cette manière, la réponse
constitutive transporte l’information nécessaire à sa lecture angulaire.

% BEGIN SINTESIS_ADITIVA_20260922
\par
La section~\ref{delta22:iii:restitucion} réunit le potentiel spectral de la réponse constitutive, la restitution globale des canaux au moyen de la vitesse normalisée et de Barbero--Immirzi, ainsi que ses bornes de stabilité. Sur l’image engendrée, avec les étiquettes et les bases conservées, la composition retrouve la paire angulaire, la section d’action et la lecture périodique du registre dodécaphasique.
% END SINTESIS_ADITIVA_20260922

La section~\ref{iii:sec:incidencias-espectrales} rend explicites les
normalisations de cette composition. La factorisation positive sépare
le facteur commun de la réponse et sa composante unimodulaire orientée ;
les deux données permettent de restituer les quatre coordonnées constitutives.
La représentation de Fourier distingue ensuite les modes des
incidences $90/120$, avec les poids signés de la chaîne et une
normalisation de transformée déclarée. Cette organisation permet de
comparer les réalisations antécédentes à l’opérateur utilisé
dans l’article, en conservant pour chacune sa fonction et son domaine.

L’appendice~\ref{iii:hist:antecedentes} développe en outre la réalisation
finie sur les positions, les modes et la phase : il prouve son équivariance et la
décomposition en vingt-sept composantes, conserve les phases de ses
valeurs propres et évalue sa sensibilité paramétrique. Les lecteurs angulaires
antécédents sont reproduits avec leurs opérations et données déclarées.
Leur comparaison aux lecteurs constitutifs du corps du texte distingue les
poids, les normalisations et les domaines de chaque réalisation.

\subsection{Polarisation, charge et réalisation dimensionnelle}

Le développement électrique s’appuie également sur deux résultats des
feuillets arithmétiques. Le décalage entre capacités nonadique et décimale
produit une discordance cumulée uniformément bornée ; sa différence
temporelle définit l’accroissement de polarisation. D’autre part,
l’appariement des fibres additive et multiplicative dans $D_9^3$
détermine le mode tritique et la norme $\sqrt3/4$ du commutateur de leurs
projections. La section~\ref{iii:sec:vacancias-prefactor} conserve
les preuves et les domaines des deux constructions. Leur relation avec
la limite $3/4$ de la réponse temporelle précise la donnée qu’elles partagent
sans identifier leurs espaces opératoriels.

La section~\ref{iii:sec:carga} compose l’impédance, $\alpha$ et la
constante de Planck au moyen de l’unique section positive satisfaisant
$Z_{\rm vac}e_a^2=2\alpha h_a$. On en obtient la résistance de
von Klitzing, le quantum de conductance, le flux magnétique et la constante de
Josephson, avec des identités croisées exactes. Le retour de l’action
laisse invariants les deux premiers et transforme les deux derniers par
des puissances opposées du quotient d’action. Ce contenu électrique
est une conséquence des sections précédentes, non une liste de
constantes indépendantes.

L’évaluation finale maintient les droites d’action, de charge et de vitesse,
ainsi que les transformations de leurs coordonnées lors du changement d’unités.
Le tableau de la section~\ref{iii:sec:evaluacion} présente les coordonnées
internes obtenues par évaluation. Une comparaison physique se formule ensuite, avec la coordinatisation
dimensionnelle et le régime de couplage déclarés. L’argument
s’organise autour d’une question commune : quelle structure détermine la
réponse, quelles équations elle impose de satisfaire et quelle information peut
être retrouvée à partir des valeurs ainsi obtenues.

\par\endgroup
\clearpage
% Núcleo común de I conservado literalmente; registro ampliado por inclusión.
\begingroup
\makeatletter
\def\input@path{{base_articulo_I/}}
\makeatother
\clearpage\section{Noyau formel de l'holographie modulaire triadique}
\label{sec:nucleo}

L'holographie modulaire triadique, abrégée HMT, commence par une structure finie de positions et d'opérations. L'ordre de construction importe : les positions reçoivent des évaluations arithmétiques, les transitions acquièrent une signature ternaire et la dynamique conserve l'information nécessaire pour les prolonger. Ce n'est qu'ensuite que les réalisations numériques sont formées. Une constante reconnue à la fin de ce processus ne remplace ni la règle qui a produit sa région ni l'histoire qui détermine son évaluation.

La dénomination \emph{All Possible Paths}, APP, exprime le rôle des chemins sur le support. Elle est adoptée comme dénomination de la construction, avec une référence conceptuelle à la formulation par chemins de la mécanique quantique ; elle ne suppose pas que sa loi arithmétique soit l'intégrale de chemins de Feynman. Le \emph{Topological Phase Kernel}, TPK,\footnote{Les auteurs dédient également les initiales TPK à Tan Pengkiang. Cette dédicace est indépendante de la définition mathématique.} désigne la composition de la sélection, du transport, de la mise à jour de la mémoire et du prolongement. Les sigles ne sont maintenus qu'après la définition des objets qu'ils représentent.

\subsection{APP : support positionnel et évaluations arithmétiques}

Considérons les neuf marques intérieures d'une droite graduée entre zéro et dix,
\[
 D_9=\{1,2,3,4,5,6,7,8,9\}.
\]
Leur disposition horizontale et verticale produit les quatre-vingt-une positions de \(D_9^2\). Les marques intérieures ne doivent pas être confondues avec les dix intervalles unitaires de la droite de référence. Cette disposition fixe un alphabet et deux directions ; elle n'utilise pas encore un développement réel pour décider d'une trajectoire.

L'identification cyclique des positions conduit au support
\[
 X_9=(\ZZ/9\ZZ)^2.
\]
La coordinatisation par représentants positifs conserve le symbole neuf pour la classe nulle. Pour tout entier positif, on définit
\begin{equation}
 \rho_9(n)=1+((n-1)\bmod9),\qquad
 q_9(n)=\frac{n-\rho_9(n)}9,
 \qquad n=\rho_9(n)+9q_9(n).
 \label{eq:residuo-cociente}
\end{equation}
La racine numérique positive est donc un choix de représentant de la classe résiduelle. Elle ne conserve pas à elle seule l'entier antérieur à la réduction ; la coordonnée \(q_9\) fournit précisément l'information manquante.

À chaque position \((i,j)\), l'accumulation de \(i\) répétée \(j\) fois produit \(ij\). Les deux évaluations entières et leurs projections sont
\[
 S(i,j)=i+j,\qquad P(i,j)=ij,\qquad
 \Sigma(i,j)=\rho_9(i+j),\quad \Pi(i,j)=\rho_9(ij).
\]
La symétrie \((i,j)\mapsto(j,i)\) échange les directions d'accumulation et conserve les deux évaluations. C'est une propriété d'un unique support, et non une identification ultérieure entre deux matrices indépendantes.

\begin{definition}[Évaluation arithmétique relevée]
L'évaluation d'APP qui conserve les deux feuillets est
\begin{equation}
 \mathcal A(i,j)=
 \bigl((R^+_{ij},q^+_{ij}),(R^\times_{ij},q^\times_{ij})\bigr)
 =\bigl((\rho_9(i+j),q_9(i+j)),(\rho_9(ij),q_9(ij))\bigr).
 \label{eq:app-levantada}
\end{equation}
Le terme \emph{feuillet} indique ici une évaluation sur le même support : additive ou multiplicative. Les deux coordonnées de quotient appartiennent à l'état arithmétique et ne sont pas éliminées lors de la réduction des résidus.
\end{definition}

\input{figures/app.tex}

\begin{proposition}[Reconstruction des évaluations]
Les deux couples de \eqref{eq:app-levantada} déterminent exactement \(i+j\) et \(ij\). Les résidus isolés ne déterminent ni ces évaluations entières ni leurs quotients.
\end{proposition}
\begin{proof}
La première affirmation est \eqref{eq:residuo-cociente} appliquée à chaque feuillet. Pour la seconde, \(10=1+9\) et \(19=1+2\cdot9\) ont le même résidu et des quotients distincts. Sur le support bidimensionnel, les positions \((5,5)\) et \((8,2)\) donnent la même somme et le même résidu du produit, mais les produits \(25=7+18\) et \(16=7+9\). La projection résiduelle identifie des données que l'évaluation relevée distingue.
\end{proof}

\subsubsection{Graphe de Cayley, topologie toroïdale et degré d'enroulement}

Les translations cardinales sont données, en coordonnées de ligne et de colonne, par
\[
 \nu_N=(-1,0),\quad\nu_E=(0,1),\quad
 \nu_S=(1,0),\quad\nu_O=(0,-1),\qquad T_d(x)=x+\nu_d.
\]
La première coordonnée augmente vers le sud et la seconde vers l'est. Le graphe
\[
 \mathscr G_9=\operatorname{Cay}(X_9,\{\nu_N,\nu_E,\nu_S,\nu_O\})
\]
est le réseau cyclique bidimensionnel : il possède quatre-vingt-un sommets et quatre arêtes orientées sortantes par sommet. C'est un graphe de Cayley pour la structure \emph{additive} du support. La multiplication des résidus possède des diviseurs de zéro et ne transforme pas toutes les lignes d'APP en un groupe multiplicatif.

Un mot cardinal \(w=d_1\cdots d_r\) et une origine \(x_0\) déterminent
\[
 x_k=x_0+\sum_{j=1}^k\nu_{d_j}\pmod9.
\]
Il existe \(81\cdot4^r\) chemins orientés de longueur \(r\), retours et retours en arrière compris. Lorsque le chemin se ferme, son déplacement entier avant réduction définit
\begin{equation}
 \operatorname{wind}(w)=\frac19
 \bigl(\#S(w)-\#N(w),\#E(w)-\#O(w)\bigr)\in\ZZ^2.
\end{equation}
L'inversion du chemin change le signe de ce vecteur. Le retour à la même position résiduelle n'impose donc pas que l'enroulement soit nul. C'est le premier exemple d'une différence qui réapparaîtra dans le retour de phase : une coordonnée peut se refermer tandis qu'une autre enregistre une évolution non triviale.

\subsubsection{Quotient de transport et identité de cocycle}

Soit \(s_9:\ZZ/9\ZZ\to D_9\) la section des représentants positifs. Le quotient de transport associé à la composition additive, également appelé \emph{retenue} en arithmétique positionnelle, est
\[
 c_9(a,b)=\frac{s_9(a)+s_9(b)-s_9(a+b)}9.
\]
Le caractère compositionnel de cette mémoire s'exprime par une identité exacte.

\begin{lemma}[Identité de cocycle]
Pour \(a,b,c\in\ZZ/9\ZZ\),
\[
 c_9(a,b)+c_9(a+b,c)=c_9(b,c)+c_9(a,b+c).
\]
\end{lemma}
\begin{proof}
Les deux sommes sont égales à
\[
 \frac{s_9(a)+s_9(b)+s_9(c)-s_9(a+b+c)}9.
\]
L'associativité de la somme des représentants relevés requiert précisément ce terme supplémentaire lors de la composition dans le quotient.
\end{proof}

Avec des représentants positifs, ce cocycle n'est pas normalisé en l'identité : \(c_9(0,a)=1\). Si \(s_0\) prend ses représentants dans \(\{0,\ldots,8\}\), sa retenue normalisée \(c_0\) satisfait
\[
 c_9(a,b)=c_0(a,b)+\mathbf1_{a=0}+\mathbf1_{b=0}-\mathbf1_{a+b=0}.
\]
La différence est le changement de section ; elle n'altère pas l'identité de cocycle. La décomposition ultérieure du calendrier en phase et fenêtre utilisera les représentants de zéro à huit.

Les totaux des deux évaluations distinguent la contribution visible et celle conservée dans le quotient :
\[
 \sum S=810,\quad\sum P=2025,\quad
 \sum\Sigma=405,\quad\sum\Pi=459,
\]
\[
 \sum q^+=45,\qquad\sum q^\times=174.
\]
On en obtient
\begin{equation}
 2025-810=(459-405)+9(174-45)=54+9\cdot129.
 \label{eq:defecto-integrado}
\end{equation}
La différence somme--produit ne se réduit pas au défaut visible \(54\). L'identité conserve également le défaut de quotient \(129\). Le supprimer modifierait l'arithmétique transportée, même si le tableau résiduel demeurait identique.

\subsubsection{Réduction ternaire et radical multiplicatif}

L'anneau \(R=\ZZ/9\ZZ\) a pour groupe des unités \(\{1,2,4,5,7,8\}\) et pour idéal maximal
\[
 \mathfrak m=(3)=\{0,3,6\},\qquad\mathfrak m^2=0.
\]
Toute ligne additive est une permutation des neuf représentants. Une ligne multiplicative l'est exactement lorsque son indice est une unité. Les lignes trois et six contiennent trois valeurs répétées ; la ligne neuf contient uniquement le représentant de la classe nulle. Le repliement multiplicatif distingue donc déjà, avant toute géométrie continue, les unités du secteur nilpotent.

L'application \(x\mapsto3x\) a un noyau et une image égaux à \(\mathfrak m\). Elle induit un isomorphisme d'espaces vectoriels sur \(\FF_3\),
\[
 R/\mathfrak m\longrightarrow\mathfrak m,\qquad [x]\longmapsto3x,
\]
mais pas un isomorphisme d'anneaux. En particulier, la suite visible \(3,6,9\) admet deux lectures différentes : sa réduction directe modulo trois est \(0,0,0\), tandis que l'extraction du facteur trois suivie de la réduction de sa coordonnée interne donne \(1,2,0\). Cette seconde opération conserve une coordonnée de l'idéal ; ce n'est pas la même projection que la réduction directe du nombre original.

\subsubsection{Réalisation latine de l'évaluation additive}

La différence entre les deux évaluations peut également s'exprimer au moyen d'
opérateurs. Soit \((e_r)_{r\in\ZZ/9\ZZ}\) la base orthonormale de \(\CC^9\).
L'affectation \(v_{ab}=e_{a+b}\) produit une base orthonormale dans chaque ligne
et chaque colonne. Les projecteurs \(P_{ab}=v_{ab}v_{ab}^{*}\) satisfont
\[
 \sum_bP_{ab}=I_9=\sum_aP_{ab},\qquad P_{ab}^2=P_{ab}.
\]
En effet, chaque translation \(b\mapsto a+b\) permute les neuf résidus. On obtient une réalisation par des projecteurs de rang un d'un carré latin, dans la terminologie des carrés latins quantiques \cite{mustovicary,delascuevas2023}. Cette construction part du feuillet additif déjà défini et explicite l'application vers sa représentation opératorielle.

L'évaluation multiplicative conserve un autre contenu : pour \(a=3\), les
vecteurs \(e_{3b}\) répètent les résidus \(0,3,6\), et
\[
 \sum_{b\in\ZZ/9\ZZ}|e_{3b}\rangle\langle e_{3b}|
 =3\bigl(|e_0\rangle\langle e_0|+|e_3\rangle\langle e_3|
                    +|e_6\rangle\langle e_6|\bigr).
\]
La multiplicité est enregistrée dans l'opérateur et caractérise le radical
multiplicatif. Le support toroïdal, la propriété latine et la
résolution opératorielle de l'identité sont donc des propriétés qui exigent
des définitions différentes. Les carrés magiques quantiques étudiés par
De las Cuevas, Drescher et Netzer exigent des entrées positives et des sommes de lignes
et de colonnes égales à l'identité ; leur théorie distingue en outre les
dilatations à entrées commutatives\cite{delascuevas2020}. Cette référence
fournit un langage précis pour comparer les réalisations d'APP.

\subsection{TRIT : orientation et régimes quadratiques}

\begin{definition}[Signature ternaire et relèvement local]
La signature TRIT utilise l'identification équilibrée
\[
 \FF_3\longrightarrow\TRIT,\qquad 0\mapsto0,\quad1\mapsto+1,\quad2\mapsto-1.
\]
Pour \(n\in\ZZ\), son relèvement conserve \(n=r+3c\), où \(r\in\TRIT\) et \(c\in\ZZ\). Dans une transition d'amplitude locale bornée, les trois états \((0,-1),(0,0),(0,+1)\) peuvent présenter le même résidu nul et des retenues distinctes.
\end{definition}

La signature ne remplace pas la trajectoire. Un zéro résiduel ne signifie pas l'absence d'information d'orientation, de phase ou de quotient. Il n'introduit pas non plus à lui seul une grandeur physique. Sa fonction est de distinguer algébriquement les classes de transport qui seront réalisées ci-après.

\subsubsection{Trois régimes quadratiques et une exponentielle commune}

Une fois la signature ternaire obtenue, sa réalisation quadratique est définie par
\begin{equation}
 \mathbb A_\tau=\RR[X]/(X^2+\tau),\qquad \tau\in\TRIT.
 \label{eq:algebras-trit}
\end{equation}
Si \(J_\tau\) est la classe de \(X\), alors \(J_\tau^2=-\tau I\). La structure distingue le type elliptique \(\tau=+1\), le type parabolique \(\tau=0\) et le type hyperbolique \(\tau=-1\). La réalisation réelle ne génère pas la signature : elle exprime les types que la réduction orientée a déjà produits.

\begin{proposition}[Exponentiation des trois types]
Dans chaque algèbre \eqref{eq:algebras-trit},
\begin{equation}
 \exp(tJ_\tau)=C_\tau(t)I+S_\tau(t)J_\tau,
\end{equation}
où
\[
 C_\tau(t)=\sum_{n\ge0}\frac{(-\tau)^nt^{2n}}{(2n)!},\qquad
 S_\tau(t)=\sum_{n\ge0}\frac{(-\tau)^nt^{2n+1}}{(2n+1)!}.
\]
Les trois réalisations sont
\[
 \begin{array}{c|c|c}
 \tau&J_\tau^2&\exp(tJ_\tau)\\\hline
 +1&-I&\cos t\,I+\sin t\,J_\tau\\
 0&0&I+tJ_\tau\\
 -1&I&\cosh t\,I+\sinh t\,J_\tau.
 \end{array}
\]
\end{proposition}
\begin{proof}
Séparer les termes pairs et impairs de la série exponentielle et substituer \(J_\tau^{2n}=(-\tau)^nI\) et \(J_\tau^{2n+1}=(-\tau)^nJ_\tau\) donne les formules. Pour \(\tau=0\), tous les termes de degré au moins deux s'annulent. Dans les deux autres cas, on retrouve les séries trigonométriques ou hyperboliques.
\end{proof}

L'exponentielle est commune aux trois régimes : elle n'est pas attribuée exclusivement à la branche parabolique. Les invariants
\[
 C_\tau(t)^2+\tau S_\tau(t)^2=1
\]
se déduisent de \(\exp(tJ_\tau)\exp(-tJ_\tau)=I\). Le même principe algébrique conserve l'unité à travers des réalisations qui diffèrent par leur type quadratique.

Dans le régime hyperbolique, \(P_\pm=(I\pm J_{-1})/2\) sont des projecteurs complémentaires et
\[
 \exp(tJ_{-1})=e^tP_++e^{-t}P_-.
\]
La croissance et la contraction apparaissent conjointement dans deux directions propres. Dans le régime elliptique, le paramètre décrit une rotation ; dans le régime parabolique, une translation nilpotente. L'interprétation temporelle adoptée par HMT associe \(+1\) au retour à mémoire, \(0\) au seuil et \(-1\) à l'ouverture bidirectionnelle. Cette interprétation exige un ordre de paramétrage ; elle n'ajoute pas un quatrième type algébrique.

\input{sections/trit_desarrollo.tex}

\input{figures/regimenes_trit.tex}

\newpage
\subsection{TPK : noyau topologique de phase}

La dynamique utilise le sélecteur de phase
\[
 M_{\rm ph}(r)=
 \begin{cases}
 +1,&r\in\{1,4,7\},\\
 -1,&r\in\{2,5,8\},\\
 0,&r\in\{3,6,9\}.
 \end{cases}
\]
Sa valeur décide quelle évaluation opère. Le transport cardinal déplace le curseur correspondant ; le registre conserve la valeur précédente, le nouveau quotient, l'orientation et la transition. La sélection et le transport sont donc des fonctions distinctes : un signe ternaire n'est pas un déplacement du tore.

Nous désignerons par \(\widehat X\) l'espace des états arithmétiques avec mémoire de trajectoire. Un élément contient les positions et les directions des deux curseurs, le mode actif, la signature TRIT, la phase, les résidus et les quotients des deux évaluations, les retenues, le préfixe construit et les données de frontière. Lorsqu'un cylindre de raffinement est introduit, il fait partie de l'état. La notation évite de remplacer l'objet par une version réduite qui ne conserve que le symbole visible.

\begin{definition}[Mise à jour du TPK]
Soient \(\mathsf{Sel}_t\) la sélection du mode au moyen de \(M_{\rm ph}\), \(\mathsf{Tra}_t\) le transport cardinal relevé et \(\mathsf{Mem}_t\) la mise à jour des registres. La transition est
\begin{equation}
 \mathcal U_t=\mathsf{Mem}_t\circ\mathsf{Tra}_t\circ\mathsf{Sel}_t.
 \label{eq:actualizacion}
\end{equation}
Chaque application conserve les champs qu'elle ne modifie pas et met à jour les autres au moyen de la division euclidienne et de la concaténation de la trajectoire.
\end{definition}

La notation \eqref{eq:actualizacion} exprime l'ordre de composition et ne remplace pas les règles effectives. La section suivante spécifie complètement sa réalisation finie à deux curseurs. Ensuite, le prolongement ne s'obtiendra pas en effaçant sa mémoire et en répétant la première émission, mais en transportant les préfixes et leurs frontières compatibles. Ainsi, APP fournit les évaluations, TRIT distingue leurs régimes et TPK compose leur évolution.

\input{sections/tpk_desarrollo_integrado.tex}

\input{sections/memoria_resolvente.tex}

\clearpage% BEGIN NUCLEO_ADITIVO_20260921_CENSO
\clearpage
\section{Noyau topologique de phase et génération du catalogue de germes}
\label{act21:tpk:capitulo}

Les deux évaluations de l'APP acquièrent une efficacité générative lorsque l'on précise
quelle évaluation agit, en quelle position elle s'effectue, comment cette position se déplace
et quelle information chaque transition conserve. Le \emph{Topological Phase Kernel},
abrégé en TPK et dénommé en espagnol \emph{núcleo topológico de fase}, réunit
ces opérations dans une dynamique sur des états enrichis. Sa définition
comprend la sélection, le transport, la mise à jour arithmétique et la conservation du
registre. La construction finie développée dans ce chapitre fournit
les germes sur lesquels agiront les relèvements régionaux et le prolongement
coinductif.

Le résultat principal de cette étape est un catalogue généré et classé :
les conditions initiales des deux curseurs produisent des émissions entières ;
leur lecture ternaire détermine des mots orientés ; l'action diédrale organise
ces mots en orbites. Chaque étape conserve ses préimages et les relations
qui permettent de revenir à la description précédente. La chaîne de cardinaux
\[
 104\,976\longrightarrow468\longrightarrow243\subset729,
 \qquad 243\longrightarrow43,
\]
désigne sous une forme abrégée des objets différents. Chaque objet est ensuite construit et chaque dénombrement
est justifié, avant que l'un d'eux soit utilisé comme donnée de
la sélection régionale.

\subsection{Sélection de phase, transport et état enrichi}
\label{act21:tpk:operaciones}

\subsubsection{Sélecteur tritique de phase opérationnelle}

La phase nonadique prend ses valeurs en \(D_9=\{1,\ldots,9\}\). Le sélecteur  
tritique est la fonction
\begin{equation}
 M_{\mathrm{ph}}:D_9\longrightarrow\{-1,0,+1\},\qquad
 M_{\mathrm{ph}}(g)=
 \begin{cases}
 +1,&g\in\{1,4,7\},\\
 -1,&g\in\{2,5,8\},\\
 0,&g\in\{3,6,9\}.
 \end{cases}
 \label{act21:tpk:selector}
\end{equation}
Dans la réalisation émettrice, la valeur \(+1\) active l'évaluation additive,
\(-1\) active l'évaluation multiplicative et \(0\) enregistre l'événement
neutre. L'orientation cardinale du curseur est une autre coordonnée de l'état.
Cette séparation permet de fixer simultanément quelle opération est lue et dans quelle
direction est transporté son point d'évaluation.

Pour les transitions numérotées par \(t\geq1\), la phase précédant la lecture est
\begin{equation}
 g_t=1+[(t-1)]_9,\qquad \epsilon_t=M_{\mathrm{ph}}(g_t),
 \label{act21:tpk:fase}
\end{equation}
où \([n]_b\) désigne le représentant de \(n\) modulo \(b\) dans
\(\{0,\ldots,b-1\}\). Une fenêtre de neuf transitions contient la suite
\((+1,-1,0,+1,-1,0,+1,-1,0)\).

\subsubsection{Curseurs orientés et transport cardinal}

Des indices \(I_9=\{0,\ldots,8\}\) sont utilisés pour les positions et des étiquettes positives \((a,b)=(i+1,j+1)\) pour les évaluations APP. Un curseur fini est
\[
 c=(i,j;d)\in\mathcal S:=I_9^2\times\mathcal D,
 \qquad \mathcal D=\{N,E,S,O\}.
\]
Les quatre directions agissent par
\[
 v_N=(-1,0),\quad v_E=(0,1),\quad
 v_S=(1,0),\quad v_O=(0,-1).
\]
L'application cardinale correspondant à \(d\) est
\[
 T_d(i,j)=\bigl([i+(v_d)_1]_9,[j+(v_d)_2]_9\bigr).
\]
Ainsi, \(M_{\mathrm{ph}}\) et \(T_d\) ont des domaines, des codomaines et des actions différents. L'activation tritique sélectionne l'un des deux curseurs ; le transport cardinal modifie sa position dans la direction que ce curseur conserve déjà.

Le relèvement entier du curseur est
\(\widetilde c=(x,y;d)\in\mathbb Z^2\times\mathcal D\), avec position finie
\(([x]_9,[y]_9)\). Ses nombres de tours sont récupérés en
\begin{equation}
 x=9\lfloor x/9\rfloor+[x]_9,\qquad
 y=9\lfloor y/9\rfloor+[y]_9.
 \label{act21:tpk:vueltas}
\end{equation}
Ces deux divisions concernent le transport spatial. La division d'une
évaluation additive ou multiplicative constitue un registre arithmétique
distinct, qui sera conservé conjointement avec elles.

\subsubsection{Composition de la transition}

L'état de la réalisation finie comprend les deux curseurs relevés,
leurs directions, la phase, les accumulateurs de la fenêtre active, la liste
ordonnée des blocs déjà émis et le registre des transitions. Chaque registre
local contient les positions antérieures, les évaluations entières des deux
feuillets, leurs restes et quotients, le régime actif et les orientations.
Les champs de préfixe et de frontière qui apparaîtront dans un prolongement seront
incorporés au même état avec leurs règles de compatibilité.

La transition s'ordonne comme
\begin{equation}
 \mathcal U_t=\mathsf{Mem}_t\circ\mathsf{Tra}_t\circ\mathsf{Sel}_t.
 \label{act21:tpk:composicion}
\end{equation}
La sélection détermine le régime et conserve l'échantillon antérieur à tout déplacement. Le transport applique le mouvement actif et l'inversion cardinale prescrite par le calendrier. L'actualisation incorpore cet échantillon aux accumulateurs et au registre ; à la fin d'une fenêtre, elle émet son bloc. La composition utilise donc un échantillon antérieur au transport bien que son archivage définitif s'effectue à la fin de la transition.

\subsection{Conditions initiales et énumération réversible}
\label{act21:tpk:semillas}

Les feuillets additif et multiplicatif disposent de curseurs indépendants. Le
domaine des conditions initiales est le produit ordonné
\[
 \Omega_{\mathrm{seed}}=\mathcal S_+\times\mathcal S_\times,
 \qquad |\mathcal S|=9^2\cdot4=324.
\]
Par conséquent,
\begin{equation}
 |\Omega_{\mathrm{seed}}|=324^2=104\,976.
 \label{act21:tpk:censo-inicial}
\end{equation}
Le nombre \(81\) compte des positions ; \(324\) compte des positions orientées d'un curseur ; \(104\,976\) compte des paires ordonnées de curseurs. Le choix d'une position en APP conserve uniquement une partie d'une condition initiale de l'émetteur.

On fixe la numérotation
\[
 \operatorname{ind}(N)=0,\quad\operatorname{ind}(E)=1,\quad
 \operatorname{ind}(S)=2,\quad\operatorname{ind}(O)=3,
\]
et l'on définit
\begin{equation}
 \nu(i,j;d)=4(9i+j)+\operatorname{ind}(d).
 \label{act21:tpk:indice-semilla}
\end{equation}

\begin{proposition}[Énumération exacte des conditions initiales]
\label{act21:tpk:enumeracion}
L'application \(\nu\) est une bijection de \(\mathcal S\) sur
\(\{0,\ldots,323\}\). Le couple d'indices se code de manière réversible
par \(324\nu_++\nu_\times\in\{0,\ldots,104\,975\}\).
\end{proposition}
\begin{proof}
La division par quatre restaure
\[
 \operatorname{ind}(d)=[\nu]_4,\qquad
 j=[\lfloor\nu/4\rfloor]_9,\qquad i=\lfloor\nu/36\rfloor.
\]
Les bornes du domaine garantissent que ces trois quantités appartiennent à
leurs ensembles d'indices et reconstruisent \(\nu\) au moyen de
\eqref{act21:tpk:indice-semilla}. Pour le couple, la division euclidienne par
\(324\) restitue le quotient \(\nu_+\) et le reste \(\nu_\times\).
\end{proof}

L'énumération parcourt le domaine complet avant tout regroupement
régional. Les indices servent à localiser les préimages d'une
émission, en conservant explicitement la condition qui a produit chaque trajectoire.

\subsection{Évaluation des feuillets et registre arithmétique}
\label{act21:tpk:lecturas}

En une position positive \((a,b)\), les évaluations entières sont
\(A_+(a,b)=a+b\) et \(A_\times(a,b)=ab\). Pour \(n>0\), on conserve
\begin{equation}
 \rho_9(n)=1+[n-1]_9,\qquad
 q_9(n)=\lfloor(n-1)/9\rfloor,\qquad
 n=\rho_9(n)+9q_9(n).
 \label{act21:tpk:division-evaluacion}
\end{equation}
L'activation additive incorpore \(\rho_9(A_+)\) à l'accumulateur correspondant. L'activation multiplicative utilise un observable du résidu \(\rho_9(A_\times)\), défini ci-dessous.

\subsubsection{Exposant discret dans les unités modulo neuf}

Les puissances de \(2\) modulo \(9\) parcourent
\[
 2^0,2^1,\ldots,2^5\equiv1,2,4,8,7,5\pmod9,
 \qquad 2^6\equiv1\pmod9.
\]
Le logarithme discret dans ce groupe de six unités se représente par
\begin{equation}
 \ell_9(1)=0,\quad\ell_9(2)=1,\quad\ell_9(4)=2,
 \quad\ell_9(8)=3,\quad\ell_9(7)=4,\quad\ell_9(5)=5.
 \label{act21:tpk:logaritmo}
\end{equation}
La loi émettrice étend l'observable par le biais de \(\ell_9(3)=\ell_9(6)=\ell_9(9)=0\). Pour conserver son origine, l'échantillon enregistré inclut le résidu original et l'indicateur d'appartenance au groupe d'unités. Ainsi, une valeur observée nulle peut provenir de l'unité \(1\) ou de l'un des trois résidus non inversibles, et son origine demeure récupérable.

\subsubsection{Ordre de lecture et mouvement}

Les évaluations s'effectuent aux positions antérieures à la transition.
Si \(\epsilon_t=+1\), on accumule la lecture additive et on déplace son curseur.
Si \(\epsilon_t=-1\), on accumule l'exposant discret multiplicatif et on
déplace le second curseur. Si \(\epsilon_t=0\), on incrémente le compteur
neutre et on conserve les deux positions.

À la fin des étapes \(27\) et \(54\), les deux directions sont remplacées
par leurs opposées. Les fenêtres s'enchaînent en conservant leurs positions et
orientations. Au début de chaque fenêtre, seuls ses trois accumulateurs locaux sont remis à zéro ; le registre précédent reste disponible.

\begin{proposition}[Inversion du transport relevé]
\label{act21:tpk:inversion}
Soient \(a,b\in\{0,1\}\) les indicateurs d'activation et d'inversion d'un
curseur. Si \(\operatorname{opp}\) échange \(N\leftrightarrow S\) et
\(E\leftrightarrow O\), la transformation
\[
 F_{a,b}(z,d)=\bigl(z+a v_d,\operatorname{opp}^{\,b}(d)\bigr)
\]
admet pour inverse
\[
 F_{a,b}^{-1}(z',d')=
 \bigl(z'-a v_{\operatorname{opp}^{\,b}(d')},
       \operatorname{opp}^{\,b}(d')\bigr).
\]
La réduction spatiale modulo neuf entrelace cette transformation avec le
transport fini.
\end{proposition}
\begin{proof}
Appliquer deux fois l'opposition restaure la direction initiale. Une fois
cette direction récupérée, la soustraction de \(a v_d\) annule le déplacement.
Pour la projection spatiale, il suffit de l'égalité
\([z+a v_d]_9=[[z]_9+a v_d]_9\), composante par composante. L'inversion de
direction conserve cette égalité.
\end{proof}

La phase détermine les indicateurs utilisés par l'inverse. Le registre
spatial et l'arithmétique peuvent ainsi être conservés conjointement, même en
traversant le bord du représentant \(I_9^2\).

\subsection{Six fenêtres et émission décimale par blocs}
\label{act21:tpk:emision}

\subsubsection{Accumulation dans une fenêtre nonadique}

La fenêtre \(m\), avec \(1\leq m\leq6\), comprend les étapes
\(9m-8,\ldots,9m\). Soient \(S_m\) la somme de ses résidus positifs
additifs actifs, \(E_m\) la somme de ses exposants discrets actifs et
\(Z_m\) son nombre d'événements neutres.

\begin{lemma}[Multiplicités d'activation]
\label{act21:tpk:tres-eventos}
Chaque fenêtre contient trois activations additives, trois multiplicatives et
trois neutres. En particulier, \(Z_m=3\).
\end{lemma}
\begin{proof}
Les phases d'une fenêtre parcourent exactement \(1,\ldots,9\). Chacun
des trois ensembles de \eqref{act21:tpk:selector} a pour cardinal trois.
\end{proof}

Les six fenêtres forment \(54\) transitions et produisent six blocs.  
Cette longueur appartient à l'émetteur initial défini ici. La continuation  
en profondeur agit ensuite sur les mots et leurs registres compatibles.

\subsubsection{Règle de codification décimale}

La loi émettrice spécifie les chiffres
\begin{align}
 d_{m,1}&=[S_m]_{10},\notag\\
 d_{m,2}&=[S_m+E_m]_{10},\notag\\
 d_{m,3}&=[3S_m+5E_m+7Z_m]_{10},\notag\\
 U_m&=100d_{m,1}+10d_{m,2}+d_{m,3}.
 \label{act21:tpk:bloque-decimal}
\end{align}
Les poids \(100,10,1\) sont les positions des centaines, des dizaines et des unités.
L'indice \(10\) indique le reste décimal, et chaque bloc satisfait
\(0\leq U_m\leq999\). Le mot émis est la suite ordonnée
\(\mathbf U=(U_1,\ldots,U_6)\), de largeur fixe de trois chiffres par bloc.

Les coefficients \(3,5,7\) constituent les poids explicites de cette loi
émettrice. Une fois la règle fixée, chaque chiffre se calcule à partir des
accumulateurs, et les résultats de ce chapitre se déduisent de cette règle.
Le nombre \(1\) qui apparaît sous sa forme réduite a une origine
additionnelle concrète : \(7Z_m=21\equiv1\pmod{10}\).

Si \(s_m=[S_m]_{10}\) et \(e_m=[E_m]_{10}\), on obtient
\begin{equation}
 G(s_m,e_m)=100s_m+10[s_m+e_m]_{10}
                   +[3s_m+5e_m+1]_{10}=U_m.
 \label{act21:tpk:acoplamiento-firmas}
\end{equation}
La lettre \(e_m\) désigne ici une composante de la signature d'exposants discrets. Sa définition appartient à l'émetteur fini.

\begin{proposition}[Récupération des deux signatures]
\label{act21:tpk:recuperacion-firmas}
L'application \(G:\{0,\ldots,9\}^2\to\{0,\ldots,999\}\) est injective.
Pour un bloc \(U=100a+10b+c\) de son image, l'inverse est
\[
 s=a,\qquad e=[b-a]_{10},\qquad
 c=[3s+5e+1]_{10}.
\]
En conséquence, le mot \(\mathbf U\) détermine les deux signatures de six composantes \(\mathbf s\) et \(\mathbf e\).
\end{proposition}
\begin{proof}
Les deux derniers termes de \eqref{act21:tpk:acoplamiento-firmas} appartiennent à
\(\{0,\ldots,99\}\), de sorte que la centaine récupère \(s\). La dizaine
est \([s+e]_{10}\) ; en soustrayant \(s\) et en réduisant modulo dix, on récupère
\(e\). L'unité vérifie la troisième congruence. La reconstruction s'applique
indépendamment aux six positions du mot.
\end{proof}

Les totaux entiers \(S_m,E_m\) se retrouvent à partir des lectures enregistrées ;
les centaines et les dizaines retrouvent leurs restes. Par exemple, un total
\(S_m=15\) produit \(s_m=5\). Cette différence identifie exactement
quelle information chaque niveau de codage conserve.

\subsubsection{Codification positionnelle du mot de six blocs}

Les six blocs admettent une codification entière à largeur fixe,
\begin{equation}
 \mathscr N(\mathbf U)=\sum_{m=1}^{6}U_m1000^{6-m},
 \qquad0\leq\mathscr N(\mathbf U)<1000^6.
 \label{act21:tpk:entero-base1000}
\end{equation}
Le premier bloc occupe la position de plus haut poids et le sixième, celle des unités de base mille. Un bloc \(058\) occupe une position complète avec une valeur entière \(58\) ; sa largeur de trois chiffres fait partie de la présentation décimale de cette position.

\begin{proposition}[Récupération positionnelle des émissions]
\label{act21:tpk:inversa-base1000}
La codification \eqref{act21:tpk:entero-base1000} détermine le mot complet
par le biais de
\[
 U_m=\left[\left\lfloor
 \frac{\mathscr N(\mathbf U)}{1000^{6-m}}\right\rfloor\right]_{1000},
 \qquad1\leq m\leq6.
\]
\end{proposition}
\begin{proof}
En divisant par \(1000^{6-m}\), les blocs précédents produisent des multiples entiers de \(1000\), le bloc \(m\) produit \(U_m\), et la contribution des suivants appartient à \([0,1)\). La partie entière élimine cette dernière contribution et le résidu modulo \(1000\) élimine les premières.
\end{proof}

Cette codification conserve les six émissions. Son domaine possède des positions ordonnées de base mille, tandis que le domaine APP possède des positions spatiales modulo neuf. La composition maintient les deux indices : \(m\) identifie une fenêtre d'émission et \((i,j)\) identifie la cellule visitée lors d'une transition de cette fenêtre.

\subsection{Calcul complet d'une émission qui commence par 501}
\label{act21:tpk:ejemplo501}

On fixe les conditions initiales
\[
 c_{+,0}=(0,4;N),\qquad c_{\times,0}=(2,3;N),
 \qquad \nu_+=16,\quad\nu_\times=84.
\]
Les positions positives initiales sont \((1,5)\) et \((3,4)\). L'indice
chronologique précédent est \(t=0\), la première phase lue est \(g_1=1\) et le
registre est vide. Ces coordonnées sont les entrées
du calcul ; les blocs s'obtiennent en exécutant la récurrence.

\begin{table}[htbp]
\caption{Première fenêtre : lectures précédentes et accumulation active.}
\label{act21:tpk:tabla-501}
\small
\begin{tabular}{@{}rcccrrr@{}}
\toprule
\(t\)&\(\epsilon_t\)&position \(+\)&position \(\times\)&
\(\Delta S\)&\(\Delta E\)&\((S,E,Z)\)\\
\midrule
1&\(+1\)&\((1,5)\)&\((3,4)\)&6&0&\((6,0,0)\)\\
2&\(-1\)&\((9,5)\)&\((3,4)\)&0&0&\((6,0,0)\)\\
3&0&\((9,5)\)&\((2,4)\)&0&0&\((6,0,1)\)\\
4&\(+1\)&\((9,5)\)&\((2,4)\)&5&0&\((11,0,1)\)\\
5&\(-1\)&\((8,5)\)&\((2,4)\)&0&3&\((11,3,1)\)\\
6&0&\((8,5)\)&\((1,4)\)&0&0&\((11,3,2)\)\\
7&\(+1\)&\((8,5)\)&\((1,4)\)&4&0&\((15,3,2)\)\\
8&\(-1\)&\((7,5)\)&\((1,4)\)&0&2&\((15,5,2)\)\\
9&0&\((7,5)\)&\((9,4)\)&0&0&\((15,5,3)\)\\
\bottomrule
\end{tabular}
\notatabla{Les positions sont des étiquettes positives des cellules antérieures au mouvement. Les événements neutres augmentent uniquement \(Z\).}
\end{table}

Les trois évaluations additives actives sont
\[
 1+5=6,\qquad9+5=14=5+9,\qquad8+5=13=4+9,
\]
de telle sorte que \(S_1=6+5+4=15\). Les trois évaluations multiplicatives actives
sont \(3\cdot4=12\), \(2\cdot4=8\) et \(1\cdot4=4\). Leurs résidus
positifs sont \(3,8,4\), avec des observables \(0,3,2\) ; par conséquent,
\(E_1=5\). Enfin \(Z_1=3\). La règle décimale donne
\[
 d_{1,1}=[15]_{10}=5,\quad
 d_{1,2}=[15+5]_{10}=0,\quad
 d_{1,3}=[45+25+21]_{10}=1,
\]
et on obtient \(U_1=501\).

La seconde lecture additive précédente a pour quotient arithmétique \(1\),
tandis que son curseur relevé se trouve en \((-1,4;N)\), avec un nombre de tours
verticaux \(-1\). Ce sont deux données différentes du même registre. Toutes deux
interviennent dans la récupération intégrale de l'évaluation et de sa position.

\begin{table}[htbp]
\caption{Les six fenêtres de la condition initiale \((16,84)\).}
\label{act21:tpk:tabla-seis501}
\begin{tabular}{@{}rrrrrrrr@{}}
\toprule
\(m\)&pas&\(S_m\)&\(E_m\)&\(Z_m\)&\(s_m\)&\(U_m\)&\([-U_m]_3\)\\
\midrule
1&1--9&15&5&3&5&501&0\\
2&10--18&6&5&3&6&614&1\\
3&19--27&24&5&3&4&498&0\\
4&28--36&21&5&3&1&169&2\\
5&37--45&12&5&3&2&272&1\\
6&46--54&12&5&3&2&272&1\\
\bottomrule
\end{tabular}
\end{table}

Le mot résultant et ses signatures sont
\begin{align*}
 \mathbf U&=(501,614,498,169,272,272),\\
 \mathbf s&=(5,6,4,1,2,2),&
 \mathbf e&=(5,5,5,5,5,5).
\end{align*}
L'inversion cardinale après l'étape \(27\) modifie les parcours des
trois fenêtres finales. Le bloc répété \(272\) correspond à deux
fenêtres successives possédant leurs propres positions et leur propre registre chronologique.

Comme deuxième cas, la condition initiale minimale \((\nu_+,\nu_\times)=(0,0)\)
produit
\[
 \mathbf U=(252,109,252,903,850,850),\quad
 \mathbf s=(2,1,2,9,8,8),\quad
 \mathbf e=(3,9,3,1,7,7).
\]
Dans ce cas, les deux curseurs retournent à leur position et orientation initiales
après \(54\) transitions. La phase nonadique effectue six tours, le
compteur chronologique avance de \(54\) unités et le registre contient
\(54\) échantillons. La lecture positionnelle du retour et
l'histoire enregistrée décrivent des propriétés distinctes de l'exécution.

\subsection{Factorisation du recensement fini}
\label{act21:tpk:censo}

Les signatures \(f_+(\nu)=\mathbf s\) et \(f_\times(\nu)=\mathbf e\) se
calculent en exécutant le curseur correspondant pendant les six fenêtres.
L'indépendance des positions et des directions entre les deux curseurs donne
\begin{equation}
 T_{\mathrm{em}}(\nu_+,\nu_\times)
   =G^{(6)}\bigl(f_+(\nu_+),f_\times(\nu_\times)\bigr),
 \label{act21:tpk:factorizacion}
\end{equation}
où \(G^{(6)}\) applique \eqref{act21:tpk:acoplamiento-firmas} à chaque position.
La factorisation permet d'énumérer d'abord \(324+324\) trajectoires et
de combiner ensuite les signatures distinctes, en conservant la multiplicité de
leurs préimages.

\par\noindent\begin{minipage}{\linewidth}
\subsubsection{Images des deux signatures}

Les tableaux suivants comprennent toutes les signatures. Leur multiplicité est le
nombre de conditions initiales qui les produisent ; l'indice minimal localise
l'une de ces conditions. La préimage complète est définie par
\[
 f_\sigma^{-1}(a)=\{\nu\in\{0,\ldots,323\}:f_\sigma(\nu)=a\}.
\]
La récurrence des sections précédentes détermine chaque élément de cet
ensemble par un calcul entier fini.

\begin{table}[H]
\centering\fontsize{9.5}{11.4}\selectfont
\setlength{\abovecaptionskip}{8pt}\setlength{\belowcaptionskip}{6pt}
\renewcommand{\arraystretch}{1.12}
\caption{Les dix-huit signatures additives.}
\label{act21:tpk:firmas-aditivas}
\begin{tabular}{@{}lrr@{\qquad}lrr@{}}
\toprule
signature&multiplicité&\(\nu_{\min}\)&signature&multiplicité&\(\nu_{\min}\)\\
\midrule
\texttt{122564}&18&17&\texttt{564122}&18&16\\
\texttt{122898}&18&24&\texttt{645212}&18&4\\
\texttt{212645}&18&5&\texttt{654889}&18&33\\
\texttt{212988}&18&0&\texttt{889221}&18&13\\
\texttt{221456}&18&29&\texttt{889654}&18&32\\
\texttt{221889}&18&12&\texttt{898122}&18&25\\
\texttt{456221}&18&28&\texttt{898465}&18&20\\
\texttt{465898}&18&21&\texttt{988212}&18&1\\
\texttt{546988}&18&9&\texttt{988546}&18&8\\
\bottomrule
\end{tabular}

\par\vspace{12pt}

\caption{Les vingt-six signatures multiplicatives.}
\label{act21:tpk:firmas-multiplicativas}
\begin{tabular}{@{}lrr@{\qquad}lrr@{}}
\toprule
signature&multiplicité&\(\nu_{\min}\)&signature&multiplicité&\(\nu_{\min}\)\\
\midrule
\texttt{000000}&108&8&\texttt{555555}&24&84\\
\texttt{177393}&8&1&\texttt{555717}&8&14\\
\texttt{177555}&8&54&\texttt{555771}&8&18\\
\texttt{177771}&8&11&\texttt{555933}&8&24\\
\texttt{339555}&8&6&\texttt{717555}&8&12\\
\texttt{339771}&8&15&\texttt{717717}&8&21\\
\texttt{339933}&8&59&\texttt{717933}&8&19\\
\texttt{393177}&8&0&\texttt{771177}&8&9\\
\texttt{393393}&8&45&\texttt{771339}&8&13\\
\texttt{393555}&8&40&\texttt{771555}&8&16\\
\texttt{555177}&8&52&\texttt{933339}&8&57\\
\texttt{555339}&8&4&\texttt{933555}&8&26\\
\texttt{555393}&8&41&\texttt{933717}&8&17\\
\bottomrule
\end{tabular}
\end{table}
\end{minipage}\par


\begin{proposition}[Image et multiplicités de l'émetteur]
\label{act21:tpk:468}
L'image \(\mathcal U_6=\operatorname{im}T_{\mathrm{em}}\) contient
\(468\) mots distincts. Ses préimages se répartissent en \(432\)
fibres de cardinal \(144\), \(18\) de cardinal \(432\) et \(18\) de
cardinal \(1944\).
\end{proposition}
\begin{proof}
L'évaluation exhaustive de l'algorithme de la section ~\ref{act21:tpk:algoritmo} sur les indices \(0,\ldots,323\) donne les deux tableaux ci-dessus. La première partition a \(18\) classes de cardinal \(18\), et la seconde \(24\) classes de cardinal \(8\), une de cardinal \(24\) et une de cardinal \(108\). Les sommes \(18\cdot18=324\) et \(24\cdot8+24+108=324\) vérifient la masse totale des deux partitions. La proposition ~\ref{act21:tpk:recuperacion-firmas} rend injectif le couplage de ses images ; par conséquent, il y a \(18\cdot26=468\) émissions. Par \eqref{act21:tpk:factorizacion}, la préimage d'un couple de signatures est le produit de ses deux préimages. On obtient \(18\cdot24=432\) produits de cardinal \(18\cdot8=144\), \(18\) de cardinal \(18\cdot24=432\) et \(18\) de cardinal \(18\cdot108=1944\). Enfin,
\[
 432\cdot144+18\cdot432+18\cdot1944=104\,976.
\]
La partie énumérative est une vérification finie exacte de la récurrence ;
la factorisation explique pourquoi ce décompte couvre toutes les paires.
\end{proof}

L'émission de l'exemple \ref{act21:tpk:ejemplo501} a une fibre de cardinal
\(18\cdot24=432\). L'exemple d'indices \((0,0)\) a une fibre de
cardinal \(18\cdot8=144\). Dans les deux cas, le mot émis récupère
les signatures, tandis que l'identification d'une condition initiale particulière
utilise son indice ou son registre de trajectoire.

\subsection{Lecture ternaire orientée et conservation des fibres}
\label{act21:tpk:reduccion-ternaria}

La réduction orientée agit séparément sur chaque bloc :
\begin{equation}
 \Psi:\mathcal U_6\longrightarrow\mathbb F_3^6,\qquad
 \Psi(U_1,\ldots,U_6)=([-U_1]_3,\ldots,[-U_6]_3).
 \label{act21:tpk:psi}
\end{equation}
Son signe fixe une orientation du catalogue. L'identification équilibrée
\(0\mapsto0\), \(1\mapsto+1\), \(2\mapsto-1\) s'applique ensuite,
composante par composante.

Le module neuf organise les évaluations et positions d'APP ; les restes modulo dix forment chaque chiffre ; la base mille regroupe trois positions décimales ; le module trois détermine la signature ternaire de chaque bloc. Ces quatre usages ont des objets et des fonctions définis. En particulier, \(\Psi\) lit six blocs et produit six trits ; une conversion positionnelle de l'entier formé par dix-huit chiffres serait une autre opération.

Pour les deux conditions calculées, on obtient
\begin{align*}
 \Psi(501,614,498,169,272,272)&=(0,1,0,2,1,1),\\
 \Psi(252,109,252,903,850,850)&=(0,2,0,0,2,2).
\end{align*}
L'émetteur décimal est défini avant cette réduction. La fonction de
\(\Psi\) est de produire l'objet ternaire sur lequel agissent la symétrie
orbitale et les sélecteurs régionaux. Cette fonction explique sa position dans la
généalogie : elle relie une émission déjà calculée à une classification structurale.

L'image \(\mathcal W_6:=\Psi(\mathcal U_6)\) contient exactement
\(243\) mots au sein de l'ensemble ambiant de \(3^6=729\) mots.
L'énumération de l'algorithme final donne l'histogramme
\begin{equation}
\begin{array}{c|rrrrrr}
 |\Psi^{-1}(w)|&1&2&3&4&5&6\\\hline
 \text{nombre de mots }w&132&66&18&3&6&18.
\end{array}
\label{act21:tpk:fibras-ternarias}
\end{equation}
Les sommes des deux lectures sont
\[
132+66+18+3+6+18=243,
\qquad132+2\cdot66+3\cdot18+4\cdot3+5\cdot6+6\cdot18=468.
\]
La première compte les mots ternaires ; la seconde récupère le nombre d'émissions réparties entre ses fibres.

\begin{example}[Deux émissions avec la même signature ternaire]
\label{act21:tpk:colision}
Les émissions
\[
 (498,501,614,252,252,109),\qquad
 (870,810,923,252,252,109)
\]
sont distinctes et toutes deux ont pour image \((0,0,1,0,0,2)\) par \(\Psi\).
Chacune a \(144\) conditions initiales. Le mot ternaire conserve
une classe de lecture ; les émissions et leurs préimages distinguent les données
qui convergent vers cette classe.
\end{example}

La conservation des fibres s'exprime concrètement par
\[
 \mathcal F_w=\{\mathbf U\in\mathcal U_6:\Psi(\mathbf U)=w\},
 \qquad
 \Omega_{\mathbf U}=T_{\mathrm{em}}^{-1}(\mathbf U).
\]
Ainsi, la préimage complète de \(w\) dans le domaine initial est l'union disjointe des ensembles \(\Omega_{\mathbf U}\), avec \(\mathbf U\in\mathcal F_w\). Le registre enrichi peut conserver l'émission, ses signatures et la condition initiale ainsi que sa lecture ternaire. L'inverse exacte se réfère alors à la donnée enregistrée avec son type ; le mot de six trits pris isolément conserve la portée précise de \eqref{act21:tpk:psi}.

\subsection{Action diédrale et les quarante-trois orbites}
\label{act21:tpk:orbitas}

Sur six coordonnées se définissent les permutations
\begin{align}
 r(u_1,u_2,u_3,u_4,u_5,u_6)
   &=(u_3,u_1,u_2,u_5,u_6,u_4),\\
 s(u_1,u_2,u_3,u_4,u_5,u_6)
   &=(u_4,u_5,u_6,u_1,u_2,u_3).
 \label{act21:tpk:generadores-diedrales}
\end{align}
La première fait tourner les deux groupes de trois coordonnées dans des sens conjugués ;
la seconde échange les deux groupes. Leur composition satisfait
\(r^3=s^2=1\) et \(srs=r^{-1}\). Elles engendrent ainsi une action du groupe diédral
\(D_3\), à six éléments.

\begin{proposition}[Équivariance et stabilité du catalogue]
\label{act21:tpk:equivarianza}
La réduction \(\Psi\) commute avec l'action de \(D_3\). Les ensembles
\(\mathcal U_6\) et \(\mathcal W_6\) sont stables sous leurs générateurs.
\end{proposition}
\begin{proof}
L'égalité \(\Psi(g\mathbf U)=g\Psi(\mathbf U)\) se vérifie en chaque
coordonnée parce que \(g\) permute les positions et \(\Psi\) applique la même
réduction à chaque position. La stabilité du catalogue constitue une autre
vérification : en générant ses \(468\) éléments par
\eqref{act21:tpk:factorizacion}, on vérifie que \(r\mathbf U\) et
\(s\mathbf U\) appartiennent à nouveau à cet ensemble. L'algorithme de la
section~\ref{act21:tpk:algoritmo} effectue ces vérifications d'appartenance.
L'équivariance transmet ensuite la stabilité à son image ternaire.
\end{proof}

\begin{proposition}[Décomposition orbitale complète]
\label{act21:tpk:43}
L'ensemble \(\mathcal W_6\) se décompose en \(43\) orbites : \(38\)
de cardinal six et \(5\) de cardinal trois.
\end{proposition}
\begin{proof}
Pour chaque \(w\in\mathcal W_6\) on forme
\[
 \mathcal O(w)=\{w,rw,r^2w,sw,srw,sr^2w\}.
\]
On choisit comme représentant le mot lexicographiquement minimal de cet ensemble. Deux mots ont le même représentant exactement lorsqu'ils appartiennent à la même orbite. L'algorithme génère la liste des représentants du tableau~\ref{act21:tpk:representantes} ; la somme de leurs cardinaux vaut \(38\cdot6+5\cdot3=243\). Comme chaque mot du catalogue participe à cette opération et que le catalogue est stable, la liste couvre toute l'image.
\end{proof}

\begin{table}[htbp]
\caption{Représentants minimaux et cardinaux des quarante-trois orbites.}
\label{act21:tpk:representantes}
\small
\begin{tabular}{@{}lr@{\qquad}lr@{\qquad}lr@{}}
\toprule
représentant&cardinal&représentant&cardinal&représentant&cardinal\\
\midrule
\texttt{001002}&6&\texttt{002112}&6&\texttt{012222}&6\\
\texttt{001011}&6&\texttt{002211}&6&\texttt{021022}&6\\
\texttt{001012}&6&\texttt{002220}&6&\texttt{021121}&6\\
\texttt{001021}&6&\texttt{011021}&6&\texttt{021202}&6\\
\texttt{001022}&6&\texttt{011112}&6&\texttt{022022}&3\\
\texttt{001100}&3&\texttt{011120}&6&\texttt{022112}&6\\
\texttt{001101}&6&\texttt{011201}&6&\texttt{022121}&6\\
\texttt{001110}&6&\texttt{011211}&6&\texttt{022202}&3\\
\texttt{001112}&6&\texttt{011220}&6&\texttt{022211}&6\\
\texttt{001202}&6&\texttt{011222}&6&\texttt{022222}&6\\
\texttt{001210}&6&\texttt{012112}&6&\texttt{112112}&3\\
\texttt{001211}&6&\texttt{012121}&6&\texttt{112211}&3\\
\texttt{001220}&6&\texttt{012202}&6&\texttt{112222}&6\\
\texttt{001222}&6&\texttt{012210}&6&&\\
\texttt{002012}&6&\texttt{012220}&6&&\\
\bottomrule
\end{tabular}
\end{table}

Les orbites héritent des invariants entiers de leurs fibres. Pour chaque mot,
\(n(w)=|\mathcal F_w|\) compte les émissions, tandis que
\[
 M(w)=\sum_{\mathbf U\in\mathcal F_w}|\Omega_{\mathbf U}|
\]
compte les conditions initiales. Le recensement des symboles
\(c(w)=(n_0,n_1,n_2)\) enregistre combien de fois chaque classe ternaire apparaît.
Ces données, conjointement avec l'orientation et le support des parcours,
fournissent les prédicats de la sélection régionale ultérieure. La
classification conserve ainsi plus d'information que la forme d'un mot
représentant.

\subsection{Algorithme autonome d'énumération et de contrôle}
\label{act21:tpk:algoritmo}

L'énumération utilise les entiers, les listes finies et les opérations déjà définies.  
Le procédé suivant précise son exécution ; \(\sigma\) prend les valeurs  
\(+\) et \(\times\).

\begin{enumerate}
\item Pour chaque \(\nu=0,\ldots,323\), récupérer \((i,j;d)\) par
l'inverse de \eqref{act21:tpk:indice-semilla}. Initialiser une liste vide de signatures.
\item Pour chaque fenêtre \(m=1,\ldots,6\), initialiser l'accumulateur \(A=0\).
Parcourir \(t=9m-8,\ldots,9m\), calculant \(\epsilon_t\) par
\eqref{act21:tpk:fase}.
\item Si \(\sigma=+\) et \(\epsilon_t=+1\), ajouter
\(\rho_9(i+j+2)\) à \(A\) et mettre à jour
\((i,j)\leftarrow T_d(i,j)\). Si \(\sigma=\times\) et
\(\epsilon_t=-1\), ajouter
\(\ell_9(\rho_9((i+1)(j+1)))\) et effectuer le même transport.
\item À la fin de \(t=27\) ou \(t=54\), remplacer \(d\) par son opposée.  
À la conclusion de la fenêtre, ajouter \([A]_{10}\) à la liste des signatures.  
Conserver la position et la direction pour la fenêtre suivante.
\item Grouper les \(324\) indices par égalité de leurs listes de six
composants. Ces deux groupements donnent \(f_+\), \(f_\times\) et toutes
leurs préimages.
\item Pour chaque paire de signatures distinctes, appliquer \(G^{(6)}\).
Attribuer à l'émission la multiplicité produit des deux préimages.
Vérifier la récupération des signatures dans chaque bloc et la somme des
multiplicités \(104\,976\).
\item Appliquer \(\Psi\) à chaque émission et regrouper par égalité de mots.
Compter les \(243\) valeurs et les multiplicités de
\eqref{act21:tpk:fibras-ternarias}.
\item Appliquer \(r,s\) à chaque émission et vérifier son appartenance au catalogue.
Pour chaque mot ternaire, former ses six images diédriques, éliminer les
répétitions et enregistrer son minimum lexicographique. Compter les \(43\)
orbites et les deux tailles du tableau~\ref{act21:tpk:representantes}.
\end{enumerate}

L'exécution se termine : le premier tronçon évalue \(648\) trajectoires de \(54\) transitions ; le deuxième combine \(468\) couples de signatures ; le troisième regroupe un ensemble de \(243\) mots. Une vérification directe supplémentaire peut exécuter les \(104\,976\) couples et vérifier la factorisation \eqref{act21:tpk:factorizacion} ; sa portée coïncide avec le même domaine fini.

\subsection{Conservation du registre et fonction des germes}
\label{act21:tpk:conclusion}

Soit \(\mathsf{Sample}(q_t)\) l'échantillon antérieur à la transition et \(\mathcal M_t\) le registre accumulé. La mise à jour de la mémoire est
\begin{equation}
 \mathcal M_{t+1}=\mathcal M_t\mathbin{\Vert}\mathsf{Sample}(q_t),
 \label{act21:tpk:archivo}
\end{equation}
où \(\Vert\) désigne la concaténation ordonnée. La récurrence donne
\[
 \mathcal M_n=\mathcal M_0\mathbin{\Vert}
 (\mathsf{Sample}(q_0),\ldots,\mathsf{Sample}(q_{n-1})).
\]
Les accumulateurs se reconstruisent en additionnant les contributions actives de
chaque segment de neuf enregistrements. Leurs quotients, les tours des
curseurs et les orientations restent associés aux mêmes transitions.
L'exécution détermine ce registre au moyen de la récurrence ; le fichier
préserve la provenance concrète des lectures.

À la fin de cette étape sont construits le domaine complet des conditions
initiales, l'émission décimale, ses deux signatures récupérables, l'image ternaire,
ses fibres et sa classification orbitale. Les recensements \(104\,976\), \(468\),
\(243\), \(729\) et \(43\) ont désormais un domaine et un mécanisme d'obtention.
Les sélections régionales recevront ces structures avec leurs invariants,
et les élévations agiront sur des mots orientés dont la provenance est
conservée. Cette continuité fait du catalogue fini un antécédent
constructif du développement coinductif.

% END NUCLEO_ADITIVO_20260921_CENSO
\section{Développement du noyau formel : conditions initiales, géométrie et raffinement}
\label{sec:extension}

Le support de quatre-vingt-une positions ne coïncide pas avec le domaine des conditions initiales de la dynamique. Une condition initiale doit également indiquer une orientation ; les feuillets additif et multiplicatif disposent en outre de curseurs indépendants. Cette distinction permet de reconstruire exhaustivement l’émission sans sélectionner au préalable une région numérique.

\subsection{Les états initiaux et la récurrence d’émission}

Pour l’énumération, nous utilisons les indices \(I_9=\{0,\ldots,8\}\), reliés aux étiquettes positives par \(i\mapsto i+1\). Soit
\[
 \mathcal S=I_9^2\times\{N,E,S,O\},\qquad
 \Omega_{\rm seed}=\mathcal S_+\times\mathcal S_\times.
\]
Chaque condition initiale \(\omega=(s_+,s_\times)\) détermine deux positions et deux directions. Le domaine comprend tous les couples ordonnés, de sorte que
\begin{equation}
 |\mathcal S|=324,\qquad |\Omega_{\rm seed}|=324^2=104\,976.
 \label{eq:semillas}
\end{equation}
L’annexe indépendante énumère toutes ces conditions ; le corps du texte présente la loi qui les transforme. Les lignes d’une table d’émissions sont des fibres de ce domaine de conditions initiales.

La réalisation finie utilise six fenêtres de neuf transitions. À chaque instant \(t\), la signature de phase est \(M_{\rm ph}(1+((t-1)\bmod9))\). Si elle vaut \(+1\), on lit le résidu positif \(\Sigma(i+1,j+1)=\rho_9(i+j+2)\) du premier curseur, puis on déplace ce curseur ; si elle vaut \(-1\), on lit le résidu positif \(\Pi(i+1,j+1)=\rho_9((i+1)(j+1))\) du second curseur, puis on déplace le second. La valeur zéro incrémente le compteur neutre. Au terme des transitions vingt-sept et cinquante-quatre, les orientations cardinales sont inversées. Les évaluations entières et leurs quotients sont conservés dans le registre, mais l’émetteur utilise les lectures résiduelles spécifiées ici.

Dans le secteur des unités, le lecteur multiplicatif utilise le logarithme discret de base deux dans \((\ZZ/9\ZZ)^\times\) :
\[
 \ell_9(1)=0,\quad\ell_9(2)=1,\quad\ell_9(4)=2,\quad
 \ell_9(8)=3,\quad\ell_9(7)=4,\quad\ell_9(5)=5.
\]
Pour l’émission finie, il est étendu avec la valeur zéro au secteur non unitaire. Cette extension est un observable et non une identification des éléments \(3,6,9\) à l’unité multiplicative. Le feuillet et son appartenance au secteur radical demeurent dans le registre de trajectoire.

Soient \(S_m\) la somme des trois résidus positifs additifs \(\Sigma\) de la fenêtre \(m\), \(E_m\) la somme des trois exposants discrets des résidus multiplicatifs \(\Pi\) et \(Z_m=3\) le nombre d’événements neutres. Ainsi, \(S_m\) ne somme pas les évaluations entières \(S=i+j\), dont les quotients demeurent dans le registre. Avec \([a]_{10}\in\{0,\ldots,9\}\), la règle d’émission décimale est
\begin{align}
 d_{m,1}&=[S_m]_{10},\notag\\
 d_{m,2}&=[S_m+E_m]_{10},\notag\\
 d_{m,3}&=[3S_m+5E_m+7Z_m]_{10},\notag\\
 U_m&=100d_{m,1}+10d_{m,2}+d_{m,3}.
 \label{eq:emisor-seis}
\end{align}
Les coefficients entiers et le calendrier de cette récurrence font explicitement partie de la loi d’émission déclarée. Aucune valeur de \(\pi,\varphi,e\) ni de \(\alpha\) n’est reçue. La distinction entre une règle d’émission et une constante reconnue à partir de sa sortie évite que la généalogie soit dissimulée dans une table de décimales.

Si \(s_m=[S_m]_{10}\) et \(e_m=[E_m]_{10}\), la même règle s’écrit
\begin{equation}
 U_m=100s_m+10[s_m+e_m]_{10}+[3s_m+5e_m+1]_{10}.
 \label{eq:acoplamiento}
\end{equation}
Le terme final égal à un est le résidu de \(7Z_m=21\) modulo dix. La réduction ternaire orientée du catalogue adopte la convention
\begin{equation}
 \Psi(U_1,\ldots,U_6)=([-U_1]_3,\ldots,[-U_6]_3).
 \label{eq:psi-catalogo}
\end{equation}
Le signe de cette coordinatisation est explicite : le changer modifie les mots et exige de transporter les orientations ultérieures. L'identification équilibrée de chaque composante de \(\FF_3\) au TRIT s'effectue après \eqref{eq:psi-catalogo}.

\begin{proposition}[Factorisation injective de l’émission]
Le mot \(U=(U_1,\ldots,U_6)\) détermine les deux signatures \(s=(s_1,\ldots,s_6)\) et \(e=(e_1,\ldots,e_6)\). L’image de l’émission possède \(18\cdot26=468\) éléments. La projection ternaire possède \(243\) valeurs distinctes dans l’espace ambiant de \(729\) mots.
\end{proposition}
\begin{proof}
Les centaines de \(U_m\) permettent de retrouver \(s_m\). Si \(b_m\) est son chiffre des dizaines, alors \(e_m=[b_m-s_m]_{10}\). Le chiffre des unités vérifie la troisième congruence de \eqref{eq:acoplamiento}. Le couplage est donc injectif sur le produit des signatures.

L’énumération des \(324\) conditions d’un curseur par la récurrence précédente produit dix-huit signatures additives, chacune ayant dix-huit antécédents ; et vingt-six signatures multiplicatives, dont vingt-quatre ont huit antécédents, une en a vingt-quatre et une en a cent huit. L’annexe conserve ces signatures et leurs conditions initiales, ce qui permet de répéter l’énumération finie sans catalogue numérique de référence. L’injectivité du couplage donne \(18\cdot26\) émissions et les multiplicités
\[
 \underbrace{432}_{18\cdot24}\text{ fibres de }144,
 \qquad18\text{ fibres de }432,
 \qquad18\text{ fibres de }1944.
\]
La somme est \(432\cdot144+18\cdot432+18\cdot1944=104\,976\). Appliquer \eqref{eq:psi-catalogo} aux \(468\) émissions produit exactement \(243\) mots. Ce dernier nombre est le cardinal de l’image, non celui de l’espace \(\FF_3^6\).
\end{proof}

Cette factorisation explique la différence entre exhaustivité finie et profondeur illimitée. Le domaine \eqref{eq:semillas} est complet pour le support et l’émetteur déclarés ; il ne contient pas par anticipation chaque préfixe de chaque prolongement. La profondeur appartient au système d’extensions compatibles construit sur ces conditions.

\subsection{Clôture du calendrier et progression de la mémoire}

Le calendrier de cent huit transitions se compose de douze fenêtres nonadiques. Sa structure de groupe est plus précise qu’un produit informel de deux périodes :
\begin{equation}
 0\longrightarrow C_{12}\xrightarrow{\iota}C_{108}
 \xrightarrow{p}C_9\longrightarrow0,
 \quad\iota([b])=[9b],\quad p([t])=[t]_9.
 \label{eq:extension108}
\end{equation}
Si \(t=9b+r\) avec \(0\le r<9\), la somme de deux phases produit la retenue
\[
 (b,r)+(b',r')=
 \left(b+b'+\left\lfloor\frac{r+r'}9\right\rfloor,\ [r+r']_9\right),
\]
où la première coordonnée est réduite modulo douze. La composition conserve le terme qui relie les deux coordonnées.

\begin{proposition}[Non-scindement du calendrier]
La suite \eqref{eq:extension108} est exacte et n’admet pas de section qui soit un homomorphisme de groupes.
\end{proposition}
\begin{proof}
Le noyau de la réduction modulo neuf est le sous-groupe des multiples de neuf, image de \(\iota\). Si l’extension était scindée, la commutativité donnerait \(C_{108}\cong C_{12}\times C_9\). Le premier groupe est d’exposant \(108\), tandis que le second est d’exposant \(\operatorname{mcm}(12,9)=36\), d’où une contradiction.
\end{proof}

Neuf transitions rétablissent la phase nonadique et font avancer d’une fenêtre ; cent huit rétablissent le calendrier fini. Aucun de ces faits n’impose d’effacer la trajectoire ou de rétablir le préfixe d’un prolongement. L’holonomie nonadique \(\Gamma_9\) désigne le transport d’un tour sur des états avec mémoire, de sorte que
\begin{equation}
 g(\Gamma_9x)=g(x),\qquad M(\Gamma_9x)=M(x)+1.
 \label{eq:holonomia-memoria}
\end{equation}
L’égalité de phase et l’incrément de mémoire sont compatibles et constituent deux observations du même transport.

\input{sections/temporalidad_trit.tex}

La représentation monodromique de ce retour est développée dans la
section~\ref{mono-ampl:conexion}. L’holonomie de
l’état y est distinguée de sa représentation sur la phase et le compteur. La suspension
symbolique de la section~\ref{mono-ampl:cristal} définit le modèle de
cristal temporel apériodique, et la section~\ref{mono-ampl:euler}
compose cette temporalité avec les trois réalisations d’Euler.

\subsection{Compression observable et conservation orthogonale}

Une réalisation isométrique permet d’exprimer quantitativement ce que perd une observation réduite. Soit \(U:\mathcal H\to\mathcal H\) une réalisation isométrique du transport d’un tour, \(U^*U=I\), et \(J:\mathcal H_{\rm vis}\to\mathcal H\) une inclusion isométrique. Nous définissons
\[
 T=J^*UJ,\qquad D=(I-JJ^*)UJ.
\]
Alors
\begin{equation}
 UJ=JT+D,\qquad J^*D=0,\qquad I-T^*T=D^*D.
 \label{eq:conservacion-ortogonal}
\end{equation}
La dernière égalité se déduit en prenant le produit adjoint de la première et en utilisant l’orthogonalité. La contraction de la composante visible ne constitue pas à elle seule une destruction de l’information de l’état total.

\begin{proposition}[Conservation à horizon fini]
Si \(I-T^*T=D^*D\), alors, pour tout \(n\ge1\),
\begin{equation}
 I=\sum_{r=0}^{n-1}T^{*r}D^*DT^r+T^{*n}T^n.
 \label{eq:telescopia}
\end{equation}
\end{proposition}
\begin{proof}
Multiplier \(D^*D=I-T^*T\) par \(T^{*r}\) et \(T^r\), puis sommer, produit une somme télescopique. Les termes intérieurs s’annulent et il reste \(I-T^{*n}T^n\).
\end{proof}

Le dernier terme conserve l’état terminal. Le supprimer avant de justifier une limite modifierait l’égalité. Cette distinction est le contenu opératoire de la conservation de l’unité dans la compression : la norme se répartit entre la composante visible, la mémoire et le terminal, sans être remplacée par un résidu unique.

\subsection{Développement local : commutation, signature lorentzienne et courbure discrète}

Le bloc diagonal adjacent dont les diagonales sont égales modulo neuf est unique. En effet,
\[
 k^2\equiv(k+1)^2\pmod9\quad\Longleftrightarrow\quad2k+1\equiv0\pmod9
\]
donne \(k=4\). Sa matrice visible est
\[
 B_c=\begin{pmatrix}7&2\\2&7\end{pmatrix}=7I+2R,
 \quad R=\begin{pmatrix}0&1\\1&0\end{pmatrix},\quad R^2=I.
\]
Avec \(P_\pm=(I\pm R)/2\), on obtient \(B_c=9P_++5P_-\). L’unité de la normalisation stochastique et la conservation d’une forme indéfinie sont deux réalisations de ce bloc :
\[
 \frac{B_c}{9}=P_++\frac59P_-,\qquad
 \frac{B_c}{\sqrt{45}}=\exp(\eta R),\quad\eta=\frac12\log\frac95.
\]
Pour \(Q=\operatorname{diag}(1,-1)\), la seconde matrice satisfait \(B_c^{\mathsf T}QB_c=45Q\). Il en résulte une réalisation locale dans \(SO^+(1,1)\) ; il n’est pas nécessaire d’introduire au préalable un espace-temps à quatre dimensions pour la démontrer.

La moyenne uniforme sur les classes \(C_1=\{1,4,7\}\), \(C_2=\{2,5,8\}\), \(C_0=\{3,6,9\}\) donne
\[
 M_\Pi=\begin{pmatrix}4&5&6\\5&4&6\\6&6&9\end{pmatrix},\quad
 M_\Sigma=\begin{pmatrix}5&6&4\\6&4&5\\4&5&6\end{pmatrix}.
\]
Le polynôme caractéristique de \(M_\Pi\) est
\[
 (\lambda+1)((\lambda-9)^2-72),
\]
d’où la signature \((2,1)\) de sa forme quadratique. L’écart spectral local \(9-5=4\), le coefficient \(2\), la différence agrégée \(51-45=6\) et le défaut intégré \(9\cdot6=54\) décrivent des niveaux distincts de la même comparaison arithmétique. Ce ne sont pas des indices interchangeables d’un même opérateur.

La courbure discrète peut être décrite directement avant cette réalisation spectrale. Dans les résidus, soient \(S(x,y)=x+y\), \(P(x,y)=xy\) et \(\Delta_x,\Delta_y\) les différences progressives. Les variables de liaison sont des différences de potentiel, \(A_x^T=\Delta_xT\), \(A_y^T=\Delta_yT\), et l’on fixe
\[
 F(T_x,T_y)=\Delta_x\Delta_yT_y-\Delta_y\Delta_xT_x.
\]
Puisque \(\Delta_xS=\Delta_yS=1\), \(\Delta_xP=y\) et \(\Delta_yP=x\),
\begin{equation}
 F(S,S)=F(P,P)=0,\qquad F(S,P)=+1,\quad F(P,S)=-1\pmod9.
\end{equation}
La somme et le produit, pris séparément, induisent des connexions plates ; leur mélange ordonné change le signe de la courbure de plaquette. En sommant sur un rectangle de côtés \(a,b\), les arêtes intérieures s’annulent et la somme orientée de bord vaut \(\pm ab\) modulo neuf. Une relation exacte apparaît ainsi entre l’ordre de composition et l’orientation de la frontière.

\subsection{Complétion cubique et deux notions de frontière}

L’extension du support cubique de côté neuf à celui de côté dix admet une stratification positionnelle exacte :
\begin{equation}
 10^3=9^3+3\cdot9^2+3\cdot9+1=729+243+27+1.
 \label{eq:cubo}
\end{equation}
Le premier terme correspond aux trois coordonnées intérieures ; le deuxième, à une nouvelle coordonnée ; le troisième, à deux ; et le dernier, au sommet comportant trois nouvelles coordonnées. La couronne d’extension contient \(271\) positions. En revanche, la frontière intérieure du cube discret de côté neuf en contient \(9^3-7^3=386\). La différence entre ces dénombrements est une différence de domaine, non un conflit de valeurs.

Diviser \eqref{eq:cubo} par \(1000\) donne une partition normalisée de l’unité. L’incidence du cube distingue six faces, douze arêtes et huit sommets. Le lien ultérieur avec les hexades et les octades de Steiner requiert les applications combinatoires correspondantes ; les cardinaux du cube ne s’y substituent pas. Cette complétion apporte ici la relation entre les bases \(729\) et \(1000\), tandis que le prolongement exceptionnel sera développé après la génération numérique et l’électron.

\input{sections/revision_emrh.tex}

\input{sections/geometria_modular.tex}
\input{sections/aritmetica_historias.tex}

\clearpage\section{Génération coinductive de \texorpdfstring{\(\pi,\varphi,e\)}{pi, phi, e} à profondeur arbitraire}
\label{sec:generacion}

L'émission de six blocs est une étape finie de la construction, et non un développement décimal complet. Sa fonction consiste à produire des régions munies de multiplicités, d'orientations et de règles de transport. Le prolongement doit conserver ces données et déterminer une collection compatible de préfixes de longueur arbitraire. Le mot \emph{génération} désigne cette opération ; \emph{généalogie} désigne la reconstruction de ses dépendances. Les deux notions interviennent, mais remplissent des fonctions distinctes.

\subsection{Orbites, régions et caractères de clôture, de propagation et d'auto-échelle}

Sur un mot \(w=(w_1,\ldots,w_6)\), nous définissons
\[
 r(w)=(w_3,w_1,w_2,w_5,w_6,w_4),\qquad
 s(w)=(w_4,w_5,w_6,w_1,w_2,w_3).
\]
On a \(r^3=s^2=I\) et \(srs=r^{-1}\) ; ces permutations réalisent donc le groupe diédral \(D_3\). L'action conserve le catalogue et les multiplicités et commute avec \(\Psi\). Le quotient des \(243\) mots visibles possède \(43\) orbites : trente-huit de cardinal six et cinq de cardinal trois.

La classe de clôture se reconnaît par un prédicat de fibre : ses six orientations ont cinq émissions par mot, une multiplicité totale \(1008\) et une distribution \(432+4\cdot144\). Il existe une seule orbite possédant ces propriétés dans le catalogue. Pour préciser les deux autres sélections, nous écrivons \(f(w)=\#\Psi^{-1}(w)\), \(m(w)=\sum_{U\in\Psi^{-1}(w)}\operatorname{mult}(U)\) et \(\nu(w)=(n_0,n_1,n_2)\). La classe diagonale additive d'une émission est \(d(U)=[i+j]_9\) dans les indices de zéro à huit ; sa signature additive conserve cette classe et l'orientation \(ES\) ou \(NO\). Soit \(\mathcal D(\mathcal O)\) l'ensemble de ces classes pour toutes les émissions sur une orbite.

Parmi les orbites de taille six, le prédicat conjoint
\[
 f(w)=1,\qquad m(w)=144\quad(w\in\mathcal O),\qquad
 \mathcal D(\mathcal O)=\{0,3,6\}
\]
laisse exactement six orbites. Dans ce secteur, le recensement \(\nu=(2,3,1)\), égal à celui de clôture, sélectionne uniquement la propagation ; le recensement \(\nu=(2,2,2)\) sélectionne uniquement l'auto-échelle. Le support diagonal ne peut pas être omis : sans lui, il y a quatre orbites à un seul point équilibrées de multiplicité \(144\), et non une seule. Enfin, l'existence d'une émission avec \(d=6\) et une orientation \(ES\) sélectionne un représentant de chaque orbite et conduit à
\begin{equation}
 w^{\rm cl}=010211,\qquad w^{\rm pr}=201101,\qquad w^{\rm au}=121200.
 \label{eq:palabras-regionales}
\end{equation}
Ces symboles désignent d'abord des régions du catalogue. L'identification scalaire sera démontrée au moyen de leurs lecteurs ; ils ne sont pas choisis parce qu'ils coïncident avec des chiffres d'une constante connue.

Le mot de clôture \(010211\) possède les cinq préimages suivantes :
\begin{center}\small
\begin{tabular}{@{}lll r@{}}\toprule
Émission de six blocs&Signature multiplicative&Rôle&Multiplicité\\\midrule
\(501,614,498,169,272,272\)&\(555555\)&transversal&432\\
\(810,923,870,169,272,272\)&\(339555\)&orientation positive&144\\
\(810,983,810,169,272,272\)&\(393555\)&orientation positive&144\\
\(870,923,810,169,272,272\)&\(933555\)&orientation positive&144\\
\(870,923,810,418,674,521\)&\(933717\)&retour orienté&144\\\bottomrule
\end{tabular}
\end{center}
La multiplicité \(432\) et la signature constante \(555555\) caractérisent la région transversale. Elles n'identifient pas la position centrale \((5,5)\) d'APP. Confondre une signature de six fenêtres avec une position du support éliminerait la distinction entre région numérique et structure électronique.

La microfibre possède en outre une morphologie bipartite précise. Chaque émission est un produit d'une classe de curseur additif et d'une classe multiplicative. Les trois régions positives partagent la même classe additive de dix-huit conditions ; ensemble, elles contiennent trois classes multiplicatives de huit conditions chacune. La région transversale possède un produit \(18\times24\) ; les trois positives réunies, un autre \(18\times24\) ; et le retour, un produit \(18\times8\). La décomposition par émissions comporte cinq parties, tandis que la décomposition en composantes connexes du graphe biparti en comporte trois. Toutes deux conservent les mêmes \(1008\) germes sans réduire la forme à un cardinal.

\subsection{Transport fini et frontière de prolongement}

Les endomorphismes \(L_0,L_1\) de l'espace des mots ternaires agissent avant la réalisation archimédienne. Leur détermination commence par les transitions de l'état arithmétique, et non par la prescription de deux matrices. Si \(U_t\) est l'étape composée de sélection, de transport et de mise à jour, le bloc observé satisfait
\[
 b_j^\chi=\Pi_6(x_j^\chi),\qquad
 x_{j+1}^\chi=U_{9j+9}\cdots U_{9j+1}(x_j^\chi),
 \qquad \chi\in\{\mathrm{cl},\mathrm{pr},\mathrm{au}\}.
\]
On forme six transitions pour chaque régime : deux consécutives dans chacune des trois régions. Les couples de matrices d'origines et de destinations obtenus dans le registre sont
\[
\begin{array}{c|c|c|c}
X_0&Y_0&X_1&Y_1\\\hline
010211&012222&010211&002111\\
012222&010211&002111&110221\\
201101&121221&102011&012222\\
121221&102011&012222&102011\\
121200&112202&121020&010210\\
112202&121020&010210&010200
\end{array}
\]
Chaque chaîne de six chiffres représente une ligne dans \(\FF_3^6\). Puisque \(\det X_0=1\) et \(\det X_1=-1\) dans \(\FF_3\), les équations \(X_aL_a=Y_a\) déterminent de manière unique \(L_a=X_a^{-1}Y_a\). Dans cette coordinatisation par vecteurs-lignes, on obtient
\[
 L_0=\begin{pmatrix}
 2&2&2&1&2&1\\2&1&2&2&1&1\\1&0&0&1&0&2\\
 0&0&1&1&1&0\\2&2&2&0&2&1\\2&1&2&1&0&0
 \end{pmatrix},\quad
 L_1=\begin{pmatrix}
 2&1&2&0&2&2\\1&2&1&1&2&0\\1&2&1&1&1&2\\
 0&1&0&0&2&1\\2&0&2&1&0&1\\0&2&2&2&1&1
 \end{pmatrix}\quad(\bmod3).
\]
Le calendrier \((L_0,L_0,L_1,L_1)\) produit quatre nouveaux blocs de six composantes. Leurs concaténations forment
\[
 w_6\longrightarrow w_{12}\longrightarrow w_{18}
 \longrightarrow w_{24}\longrightarrow w_{30}.
\]
Les sous-indices de \(w\) indiquent la longueur accumulée. La frontière \(R_{36}\) conserve la structure terminale et ouvre la relation de prolongement ; elle n'est pas renommée comme un dernier bloc périodique qui pourrait se répéter indéfiniment.

L'identité \(L=X^{-1}Y\) démontre l'unicité une fois les transitions fixées ; elle ne remplace pas la production de ces transitions par \(U_t\). La provenance de leur registre et l'examen des seize calendriers binaires sont documentés dans le supplément technique. Le calendrier indiqué est le seul de ces seize calendriers qui reproduise conjointement les trois suites de cinq blocs.

\subsection{Cinq régions de clôture et compensation angulaire}

La région de clôture contient une composante transversale et quatre accompagnatrices. Le lecteur symétrique agit dans l'espace de leurs cinq positions au moyen de
\[
 P_5=\frac15\mathbf1\mathbf1^{\mathsf T},\qquad
 \mathbf1=(1,1,1,1,1)^{\mathsf T}.
\]
L'invariance \(P_5\lambda=\lambda\) et la conservation \(\mathbf1^{\mathsf T}\lambda=1\) imposent \(\lambda=\mathbf1/5\). Cette normalisation répartit l'unité entre les positions du lecteur ; elle ne remplace pas les multiplicités des germes \(432,144,144,144,144\) par des fréquences uniformes. La coordinatisation elliptique du lecteur réalise chaque coordonnée par \(1+\lambda_jJ\) ; sa pente est \(\lambda_j\), d'où la pente primitive \(1/5\). La règle qui fait passer d'une coordonnée normalisée à une pente est ainsi explicitée, au lieu d'identifier les deux notions sans application. Le lecteur compose les quatre positions associées selon la loi
\begin{equation}
 a\oplus b=\frac{a+b}{1-ab}.
 \label{eq:suma-pendientes}
\end{equation}
Cette loi s'obtient en multipliant \((1+aJ)(1+bJ)\) dans le régime \(J^2=-I\) et en divisant par sa composante scalaire. La composition des pentes est donc une réalisation de la structure quadratique de TRIT, et non une somme ordinaire de quatre fractions.

\begin{proposition}[Compensateur de la clôture]
La composition de quatre pentes \(1/5\) est \(120/119\). La condition de retour à une pente un détermine un unique compensateur positif de la forme \(1/q\), avec \(q>1\), et donne \(q=239\).
\end{proposition}
\begin{proof}
La première duplication donne \((1/5)\oplus(1/5)=5/12\) ; la seconde, \((5/12)\oplus(5/12)=120/119\). La condition
\[
 \frac{120/119-1/q}{1+(120/119)/q}=1
\]
équivaut à \((120-119)q=120+119\), d'où \(q=239\).
\end{proof}

Le caractère de clôture résultant a pour réalisation
\begin{equation}
 \Pi_{\rm cl}=4\left(4\arctan\frac15-\arctan\frac1{239}\right).
 \label{eq:lector-pi}
\end{equation}
L'identité angulaire de Machin reconnaît \(\Pi_{\rm cl}=\pi\). Le sens de la construction est précis : la pente et le nombre de compositions sont déclarés par le lecteur de la pentafibre ; l'équation de compensation impose \(239\). Une fois ce caractère produit, son évaluation par des séries alternées ne sélectionne pas à nouveau l'orbite des germes.

Pour \(0<x<1\), les sommes alternées
\[
 A_n(x)=\sum_{j=0}^n\frac{(-1)^jx^{2j+1}}{2j+1}
\]
encadrent \(\arctan x\) entre deux troncatures consécutives, dont la différence est \(x^{2n+3}/(2n+3)\). Appliquées aux deux pentes de \eqref{eq:lector-pi}, elles fournissent des extrémités rationnelles d'erreur explicite. La valeur numérique est ainsi calculée comme conséquence du caractère exact, sans introduire de préfixe cible.

\subsection{Lecteur de propagation et normalisation de l'unité}

La région de propagation se réalise au moyen d'une collection formelle \(P_t\in\mathcal A[[t]]\), où \(\mathcal A\) est une algèbre commutative unitaire sur \(\QQ\). Sa composition satisfait \(P_{t+s}=P_tP_s\), d'unité \(P_0=1\), et le transport élémentaire orienté fixe le générateur scalaire \(P'_0=+1\). La condition fixe son signe, et non sa seule grandeur. En écrivant
\[
 P_t=\sum_{n\ge0}a_nt^n,
\]
la normalisation détermine \(a_0=a_1=1\). La comparaison du coefficient de \(s^nt\) dans \(P_{s+t}=P_sP_t\) donne \((n+1)a_{n+1}=a_na_1\) ; de manière équivalente, \(P'_t=P_t\). Ainsi,
\begin{equation}
 (n+1)a_{n+1}=a_n,\qquad a_n=\frac1{n!}.
 \label{eq:propagacion}
\end{equation}
Le caractère en une unité de paramètre est \(\Pi_{\rm pr}=P_1\), reconnu ultérieurement comme \(e\).

La normalisation du générateur importe. La loi de semi-groupe sans cette donnée admettrait \(P_t=e^{ct}\) pour différents \(c\) ; c'est pourquoi on n'omet pas la condition qui fixe l'unité du paramètre dans le lecteur. Cette condition est antérieure à la comparaison décimale et appartient à l'exposition de la règle opératoire, et non au choix d'une valeur expérimentale.

\begin{proposition}[Bornes rationnelles de propagation]
Pour \(n\ge1\), soit \(E_n=\sum_{j=0}^n1/j!\). Alors
\[
 E_n<\Pi_{\rm pr}<E_n+\frac{n+2}{(n+1)(n+1)!}.
\]
\end{proposition}
\begin{proof}
Le premier terme omis est \(1/(n+1)!\). Chaque quotient ultérieur est au plus égal à \(1/(n+2)\). La somme géométrique majorant la queue est
\[
 \frac1{(n+1)!}\frac1{1-1/(n+2)}
 =\frac{n+2}{(n+1)(n+1)!}.
\]
L'inégalité stricte découle de la décroissance des quotients successifs.
\end{proof}

\subsection{Incidence modale et génération de l'auto-échelle}

Le sélecteur tritique produit le cycle \((+1,-1,0)\). Dans les modes \(\Sigma,\Pi\), la règle est \(+1:m\mapsto\Sigma\), \(-1:m\mapsto\Pi\), \(0:m\mapsto m\). Avec fermeture cyclique, le mode observé est \((\Sigma,\Pi,\Pi)\). Ses arêtes sont \(\Sigma\to\Pi\), \(\Pi\to\Pi\) et \(\Pi\to\Sigma\). Son incidence, dans cet ordre des modes, est donc
\begin{equation}
 F_{\rm av}=\begin{pmatrix}0&1\\1&1\end{pmatrix}.
 \label{eq:fav}
\end{equation}
Les quatre entrées s'obtiennent à partir du calendrier, et non d'une équation qui aurait préalablement reçu la valeur de \(\varphi\). La répétition du calendrier donne les dénombrements \(3F_{\rm av},9F_{\rm av},36F_{\rm av}\) aux longueurs neuf, vingt-sept et cent huit. Ces dénombrements ne sont pas les puissances de la matrice.

\begin{theorem}[Caractère positif d'auto-échelle]
L'incidence \eqref{eq:fav} possède un unique rayon propre strictement positif. En le normalisant sous la forme \((1,x)^{\mathsf T}\), sa coordonnée \(x\) satisfait \(x^2-x-1=0\), \(x>0\), et vaut \(\varphi=(1+\sqrt5)/2\).
\end{theorem}
\begin{proof}
L'équation propre donne \(x=\lambda\) et \(1+x=\lambda x\). L'élimination de \(\lambda\) produit le polynôme. Une seule de ses deux racines est positive. Comme \(F_{\rm av}^2\) a toutes ses entrées positives, le rayon positif est unique et simple. L'expression radicale reconnaît la coordonnée obtenue ; elle ne participe pas à la production de \(F_{\rm av}\).
\end{proof}

Avec \(F_0=0,F_1=1,F_{n+1}=F_n+F_{n-1}\), l'incidence induit, pour \(n\geq1\),
\[
 F_{\rm av}^n=\begin{pmatrix}F_{n-1}&F_n\\F_n&F_{n+1}\end{pmatrix}.
\]
Les quotients consécutifs \(F_{n+1}/F_n\) alternent autour de \(\varphi\). L'identité de déterminant
\[
 F_{n+1}^2-F_nF_{n+2}=(-1)^n
\]
donne une largeur exacte \(1/(F_nF_{n+1})\) entre deux quotients consécutifs. Sa croissance rend cet encadrement rationnel arbitrairement petit.

La mesure sur toutes les histoires compatibles ne coïncide pas avec la fréquence des modes du calendrier à trois événements. La fréquence chronologique est \((1/3,2/3)\). La mesure d'entropie maximale de l'enveloppe symbolique définie par \(F_{\rm av}\) se construit à partir de ses vecteurs propres et a pour rapport de poids \(\varphi^{-2}\). Cette dernière intervient dans la lecture surfacique--volumétrique ; elle ne s'obtient pas en remplaçant simplement une fréquence de comptage par une autre.

\input{sections/retorno_areal_volumetrico.tex}

Une fois les caractères générés, les deux orientations du lecteur
\[
 R(x,y)=\frac{x}{y^2}
\]
fournissent
\begin{equation}
 \rho_A=\frac\pi{\varphi^2},\qquad
 \rho_E=\frac\varphi{\pi^2},\qquad
 \varphi^3\rho_A^2\rho_E=1.
 \label{eq:areal-electron}
\end{equation}
La seconde expression apparaîtra dans l'exposant électronique. La première s'obtient également comme limite de la lecture marquée \(\pi F_{n-1}/F_{n+1}\). L'échange des arguments les relie, mais ne transforme pas un reparamétrage de l'équation électronique en une deuxième dérivation indépendante de tous ses coefficients.

\input{sections/euler_trit_articulacion.tex}

\subsection{Cylindres compatibles et profondeur arbitraire}

La réalisation d'un bloc ternaire \(b=(b_1,\ldots,b_6)\) est l'entier
\[
 \nu(b)=\sum_{j=1}^6b_j3^{6-j}\in\{0,\ldots,728\}.
\]
Une suite de blocs définit des préfixes entiers par \(N_{n+1}=729N_n+\nu(b_{n+1})\), en commençant en \(N_0=m\), où \(m\) est la partie entière déterminée par le lecteur. Les fourchettes précédentes situent les caractères de clôture, de propagation et d'auto-échelle dans \((3,4)\), \((2,3)\) et \((1,2)\), respectivement ; leurs valeurs initiales sont donc \(m=3,2,1\). Les blocs suivants raffinent la partie fractionnaire. Le cylindre correspondant est
\begin{equation}
 I_n=\left[\frac{N_n}{729^n},\frac{N_n+1}{729^n}\right],
 \qquad |I_n|=729^{-n}.
 \label{eq:cilindro}
\end{equation}
La convention semi-ouverte est utilisée pour sélectionner un représentant lorsqu'une valeur tombe sur une frontière. Pour le passage à la limite, on utilise les fermetures de ces intervalles et on identifie les deux développements de frontière équivalents. Cette précision empêche de supposer que toute intersection d'intervalles semi-ouverts emboîtés est automatiquement non vide.

\begin{theorem}[Détermination de la réalisation scalaire]
Une histoire de blocs compatible détermine une unique limite des extrémités gauches de \eqref{eq:cilindro}. Si un caractère exact dispose d'encadrements rationnels de largeur arbitrairement petite et se sépare de la frontière de chaque cylindre décimal examiné, chaque préfixe décimal de longueur finie est déterminé après un calcul fini. Les préfixes ainsi obtenus sont compatibles entre eux.
\end{theorem}
\begin{proof}
L'extension d'un bloc conserve le préfixe précédent par division euclidienne. Les cylindres fermés sont emboîtés et leur diamètre tend vers zéro ; leurs extrémités gauches forment une suite de Cauchy et possèdent une unique limite. Pour une longueur décimale fixée, une valeur séparée des frontières de la partition possède une distance positive à celles-ci. Une fourchette suffisamment étroite est contenue dans une unique cellule et détermine son préfixe. L'unicité de la cellule et l'emboîtement des partitions prouvent la compatibilité par troncature. Aux points à développement terminal, il faut ajouter la décision exacte de frontière et sa convention de représentation ; cette décision ne se déduit pas d'une largeur numérique seule.
\end{proof}

\begin{corollary}[Obtention finie et compatibilité des préfixes régionaux]
\label{gen:prefijos-regionales}
Pour chacune des réalisations \(\pi,\varphi,e\) et pour toute longueur
décimale finie, les encadrements construits déterminent le préfixe correspondant
par un calcul fini. Les préfixes obtenus à différentes profondeurs
sont compatibles par troncature.
\end{corollary}
\begin{proof}
Les caractères de clôture, de propagation et d'auto-échelle possèdent les encadrements précédents. Après leur reconnaissance, l'irrationalité de \(\pi,e,\varphi\) garantit leur séparation des frontières rationnelles de chaque partition décimale finie. On applique le théorème précédent à chaque caractère.
\end{proof}
La profondeur arbitraire exprime le quantificateur du corollaire : pour tout
horizon fini, il existe un calcul qui détermine son préfixe. L'objet
mathématique est la collection compatible complète ; chaque exécution en obtient une
projection finie.

\subsection{Involution spéculaire et synchronisation des régions}

La connexion nonadique transporte conjointement les trois caractères. Dans la coordinatisation de la neuvième phase, l'axe d'auto-échelle et la somme ternaire de clôture et de propagation restent fixés sous l'échange spéculaire. L'inversion hélicoïdale change le sens du transport ; la réflexion des deux canaux change leur orientation relative. Ce sont des involutions différentes, bien qu'elles agissent sur le même système de trajectoires.

Les coïncidences de recensement illustrent cette coordination :
\begin{center}
\begin{tabular}{@{}c rrr l@{}}\toprule
Phase&\(N_\pi\)&\(N_e\)&\(N_\varphi\)&Égalité des dénombrements\\\midrule
4&207&207&122&\(N_\pi=N_e\)\\
5&39&151&151&\(N_e=N_\varphi\)\\
6&110&110&69&\(N_\pi=N_e\)\\
7&81&29&81&\(N_\varphi=N_\pi\)\\
8&59&59&59&synchronisation triple\\
9&43&43&43&synchronisation triple\\\bottomrule
\end{tabular}
\end{center}
Ces entiers comptent des coordonnées compatibles des régions. Ils n'affirment pas une égalité entre \(\pi\), \(e\) et \(\varphi\) en tant que nombres réels. Si
\[
 \partial(a,b,c)=(a-b,b-c,c-a),
\]
son image appartient au réseau \(A_2=\{(u,v,w)\in\ZZ^3:u+v+w=0\}\). Aux phases huit et neuf, la composante relative s'annule, mais la composante diagonale passe de \((59,59,59)\) à \((43,43,43)\). La synchronisation relative coexiste avec un déplacement de \(-16\) par canal. La fermeture qui conduira à \(\alpha\) doit conserver les deux informations : coïncidence relative et évolution de la mémoire.

\subsection{Suspension hélicoïdale sturmienne}

La relation entre les bases neuf et dix de la complétion définit le paramètre adimensionnel
\begin{equation}
 \delta=\frac{\log(10/9)}{\log10}.
 \label{eq:delta-sturm}
\end{equation}
Elle ne s'obtient pas à partir d'une constante physique cible. C'est l'échelle logarithmique relative de deux bases déjà présentes dans la construction. Son emploi comme lecture symbolique du raffinement a une conséquence exacte.

\begin{lemma}[Irrationalité de la pente]
Le nombre \(\delta\) est irrationnel.
\end{lemma}
\begin{proof}
Si \(\delta=p/q\), \(q>0\), alors \((10/9)^q=10^p\) et \(9^q=10^{q-p}\). La factorisation du premier membre contient le nombre premier trois et celle du second uniquement les nombres premiers deux et cinq. L'égalité est impossible.
\end{proof}

Le mot mécanique inférieur, d'ordonnée à l'origine \(\theta\), est
\[
 z_n(\theta)=\lfloor(n+1)\delta+\theta\rfloor-\lfloor n\delta+\theta\rfloor\in\{0,1\}.
\]
On l'obtient en observant la rotation \(\theta\mapsto\theta+\delta\pmod1\) au moyen de l'intervalle \([1-\delta,1)\). La suite d'événements a pour densité \(\delta\), et la somme sur toute fenêtre de longueur \(L\) vaut \(\lfloor L\delta\rfloor\) ou \(\lceil L\delta\rceil\). Puisque \(21<1/\delta<22\), les séparations entre événements sont vingt et un ou vingt-deux. Le mot qui indique quels sont les intervalles courts est un autre mot mécanique de pente \(\sigma=22-1/\delta\) ; ses événements sont séparés par six ou sept intervalles primaires. La seconde échelle est induite par la première et n'est pas un second paramètre libre.

\input{sections/vacancias_capacidad.tex}

\begin{theorem}[Complexité des lectures de phase]
Pour \(m\ge1\), le mot mécanique a une complexité factorielle \(p(m)=m+1\). En conservant la phase \(n\bmod9\), la complexité est \(p_9(m)=9(m+1)\) ; en conservant seulement sa classe ternaire, elle est \(p_3(m)=3(m+1)\). Les trois entropies topologiques sont nulles.
\end{theorem}
\begin{proof}
Les prédécesseurs des deux extrémités de l'intervalle de codage forment, pour les facteurs de longueur \(m\), \(m+1\) points distincts du cercle : la coïncidence entre extrémités translatées élimine les doublons et l'irrationalité empêche toute coïncidence supplémentaire. Ces points déterminent \(m+1\) intervalles à itinéraires constants et distincts. Toute orbite irrationnelle visite chaque intervalle, d'où \(p(m)=m+1\).

Une phase modulo neuf étant fixée, les positions de cette phase parcourent une rotation de pas \(9\delta\), également irrationnel. Par densité, tous les facteurs apparaissent dans chaque phase, et la première coordonnée distingue leurs neuf copies. Le même argument avec le pas \(3\delta\) donne trois copies après réduction ternaire. Enfin, \(m^{-1}\log(c(m+1))\to0\) pour \(c=1,3,9\).
\end{proof}

Le produit de phases \(C_9\times\mathbb T\), translaté par \((1,\delta)\), possède des caractères de fréquence
\[
 \frac a9+k\delta\pmod1,\qquad a\in\ZZ/9\ZZ,\ k\in\ZZ.
\]
Le module spectral n'est pas \(\log10\) : c'est le module additif engendré par \(1/9\) et \(\delta\), modulo les entiers. Le facteur \(\log10\) appartient à la normalisation de \eqref{eq:delta-sturm}.

Une présentation opératorielle rend la mémoire explicite. Dans \(\ell^2(\ZZ)\), soient \(T|n\rangle=|n+1\rangle\), \(V_9|n\rangle=\zeta_9^n|n\rangle\), \(W_\delta|n\rangle=e^{2\pi i n\delta}|n\rangle\) et \(D_M|n\rangle=\lfloor n/9\rfloor|n\rangle\). Sur les vecteurs à support fini,
\begin{equation}
 V_9T=\zeta_9TV_9,\qquad
 W_\delta T=e^{2\pi i\delta}TW_\delta,\qquad
 [D_M,T^9]=T^9.
 \label{eq:weyl-relojes}
\end{equation}
Les deux premières relations sont des covariances de Weyl ; la troisième exprime qu'un tour nonadique augmente le degré de mémoire. La rotation de base est abélienne, tandis que l'algèbre des observables et des translations ne l'est pas. L'inversion \(|n\rangle\mapsto|-n\rangle\) conjugue l'avancée avec le recul. Cette réalisation fournit une suspension hélicoïdale apériodique à retour fini de phase, et non un mot périodique de longueur neuf.

\input{figures/holonomia.tex}

Pour ce dénombrement, nous fixons l'intercept \(\theta=0\) et les extrémités cumulées \((N_0,\ldots,N_4)=(0,729,972,999,1000)\), correspondant à l'ordre cubique \((729,243,27,1)\). Les dénombrements inférieurs sont \(\lfloor N_j\delta\rfloor-\lfloor N_{j-1}\delta\rfloor\) et donnent \((33,11,1,0)\). Les supérieurs sont \(\lceil N_j\delta\rceil-\lceil N_{j-1}\delta\rceil\), avec la même frontière initiale nulle, et donnent \((34,11,1,0)\). Ces valeurs dépendent de l'intercept et de l'ordre des strates ; elles ne sont pas affirmées pour une phase initiale arbitraire. Leurs totaux sont quarante-cinq et quarante-six. L'incrément correspond à une unique unité de frontière, située dans la première strate de cette partition ordonnée. Les nombres trente-quatre et onze apparaissent ainsi dans un même dénombrement antérieur à toute coordinatisation d'unités ; les identifier à des exposants physiques exige en outre la dérivation dimensionnelle correspondante.

La dénomination \emph{cristal apériodique} est employée ici pour cet ordre symbolique, dont la complexité, le spectre et le retour sont spécifiés. L'entropie topologique nulle n'équivaut pas à elle seule à une dissipation physique nulle : une réalisation thermodynamique exige une dynamique énergétique et une opération d'effacement définies. De même, cette lecture sturmienne ne remplace pas tous les opérateurs du TPK et ne démontre pas à elle seule la périodicité ou l'apériodicité de la sortie complète de \(\alpha\). Sa fonction est d'expliciter une structure d'ordre que le simple retour du calendrier ne décrit pas.

\input{sections/revision_monodromia.tex}

\input{sections/primos_retornos.tex}

\clearpage% BEGIN NUCLEO_ADITIVO_20260921_SELECCION
\clearpage
% Ampliación demostrativa DF003--DF007. No modifica la fuente integral.
% Propietario común de los archivos Lean citados en estos comentarios:
% output/PAQUETE_ARTICULO_I_INTEGRACION_ESTADO_CAMPOS_20260919/.
% Residencia causal de la primera sección: después de los lectores regionales
% coinductivos y antes de emplear sus registros como entrada de la selección K.
% Residencia causal de la segunda sección: después de la construcción marcada
% Paley--Witt--Golay--Leech y antes de la reconstrucción de la estructura VOA.
% Los dos cortes expositivos conservan el mismo origen APP--TRIT--TPK.

\section[Génération régionale et raffinement]{Compatibilité intégrale de la génération régionale, du raffinement cylindrique et de l'incidence}
\label{act21:sec:df-lectores-cilindros}

La prolongation régionale produit des suites ordonnées de blocs, ainsi que
leurs profondeurs de lecture. La projection archimédienne et l'incidence
ternaire se calculent sur ces mêmes suites. Le lien entre ces deux
réalisations peut s'exprimer entièrement par des opérations entières :
sélection par intervalles rationnels, division euclidienne, élévation de six
trits, lecture de retenue et transformation de Witt. Cette formulation permet
de suivre chaque coefficient depuis son bloc générateur jusqu'à son emploi ultérieur.

La construction utilise les trois lecteurs régionaux déjà obtenus de l'état APP--TRIT--TPK. Nous les désignerons par les indices \(c\in\{\mathrm P,\mathrm E,\mathrm F\}\), correspondant respectivement aux régions de clôture, de propagation et d'auto-échelle. Leur réalisation réelle s'écrira \(v_c\). Les identifications ultérieures \(v_{\mathrm P}=\pi_{\mathrm{HMT}}\), \(v_{\mathrm E}=e_{\mathrm{HMT}}\) et \(v_{\mathrm F}=\varphi_{\mathrm{HMT}}\) conservent cet ordre causal. L'opération qui émet un bloc inspecte les intervalles rationnels du lecteur régional ; la valeur réelle intervient dans la démonstration de correction de cette opération.

\subsection{Émission positionnelle par séparation rationnelle de cellules}
\label{act21:subsec:df-emision-racional}

% Fuente: lean/biblioteca/RegionalPublicationComposition.lean,
% produced_brackets, eventual_publication, joint_eventual_publication,
% search_terminates, stopped_cell_correct, publications_compatible.

Pour chaque région, on dispose de deux suites rationnelles
\(\ell_c(j),u_c(j)\), construites par son lecteur, qui satisfont
\[
 \ell_c(j)\leq v_c\leq u_c(j).
\]
La propriété de séparation des cellules démontrée pour ces lecteurs affirme
que, pour chaque entier \(s>0\), il existe \(J\) tel que, pour tout \(j\geq J\),
\begin{equation}
 \lfloor s\ell_c(j)\rfloor
 =\lceil su_c(j)\rceil-1
 =\lfloor sv_c\rfloor.
 \label{act21:eq:df-separacion-celdas}
\end{equation}
L'évitement des frontières rationnelles par les valeurs régionales est l'hypothèse des théorèmes précédents qui garantit cette séparation. Ici, on emploie sa conclusion opérationnelle, avec les lecteurs et les intervalles qui la produisent.

\begin{definition}[Indice de séparation et préfixe émis]
Pour \(s>0\), soit
\[
 j_c(s)=\min\{j\in\mathbb N:
 \lfloor s\ell_c(j)\rfloor=\lceil su_c(j)\rceil-1\}.
\]
Pour une base entière \(b\geq2\) et une profondeur \(n\geq0\), on définit
\begin{equation}
 P_c(b,n)=\lfloor b^n\ell_c(j_c(b^n))\rfloor.
 \label{act21:eq:df-publicacion-prefijo}
\end{equation}
Le préfixe contient la partie entière et les premières \(n\) positions en base \(b\) ; par conséquent, \(P_c(b,0)\) est la partie entière régionale.
\end{definition}

\begin{theorem}[Terminaison et correction de l'émission]
\label{act21:thm:df-emision-correcta}
L'indice \(j_c(s)\) existe pour tout \(s>0\). Les préfixes émis satisfont
\begin{equation}
 P_c(b,n)=\lfloor b^n v_c\rfloor,
 \qquad
 P_c(b,n)\leq b^n v_c<P_c(b,n)+1.
 \label{act21:eq:df-prefijo-correcto}
\end{equation}
Pour chaque échelle \(s\) il est possible en outre de choisir un unique indice de suffisance \(J(s)\) valide simultanément pour les trois régions.
\end{theorem}

\begin{proof}
L'égalité à partir d'un certain rang de \eqref{act21:eq:df-separacion-celdas} rend non vide
l'ensemble qui définit \(j_c(s)\). Son minimum existe par le bon ordre de
\(\mathbb N\). À tout indice où l'égalité se produit, la borne
inférieure rationnelle fournit
\(\lfloor s\ell_c(j)\rfloor\leq\lfloor sv_c\rfloor\).
La borne supérieure, jointe à l'évitement d'une frontière de cellule, donne
\(\lfloor sv_c\rfloor\leq\lceil su_c(j)\rceil-1\).
Les deux extrémités entières coïncidant, l'entier intermédiaire coïncide avec
elles. Cela démontre \eqref{act21:eq:df-prefijo-correcto}. Pour l'affirmation
conjointe, il suffit de prendre le maximum des trois indices de suffisance obtenus
pour la même échelle. L'indice conjoint contrôle simultanément les trois
émissions, tandis que chaque région conserve ses propres intervalles.
\end{proof}

\begin{proposition}[Compatibilité exacte des préfixes]
\label{act21:prop:df-prefijos-compatibles}
Pour \(n,m\geq0\),
\begin{equation}
 \left\lfloor\frac{P_c(b,n+m)}{b^m}\right\rfloor=P_c(b,n).
 \label{act21:eq:df-truncamiento}
\end{equation}
En particulier, le bloc suivant
\(d_c(b,n)=P_c(b,n+1)\bmod b\) satisfait
\begin{equation}
 P_c(b,n+1)=bP_c(b,n)+d_c(b,n),\qquad 0\leq d_c(b,n)<b.
 \label{act21:eq:df-paso-posicional}
\end{equation}
\end{proposition}

\begin{proof}
Écrivons \(b^n v_c=N+\theta\), avec
\(N=P_c(b,n)\) et \(0\leq\theta<1\). Alors
\(P_c(b,n+m)=b^mN+\lfloor b^m\theta\rfloor\), et le second terme
appartient à \(\{0,\ldots,b^m-1\}\). La division euclidienne par \(b^m\)
récupère \(N\). La spécialisation \(m=1\) produit
\eqref{act21:eq:df-paso-posicional} et sa borne résiduelle.
\end{proof}

\subsection{Blocs régionaux de six trits à profondeur arbitraire}
\label{act21:subsec:df-bloques-arbitrarios}

% Fuente: deltas/integracion/JointRegionalFrontier.lean,
% publish_base729_eq_base3, generated_block_from_longer_prefix,
% generated_block_eq_finite, generated_signature_eq_finite.

L'égalité entière \(729=3^6\) détermine le regroupement de six positions ternaires en un bloc. Pour \(t\geq0\), nous définissons
\begin{equation}
 u_c(t)=P_c(729,t+1)\bmod729,
 \qquad
 B_{c,j}(t)=
 \left\lfloor\frac{u_c(t)}{3^{5-j}}\right\rfloor\bmod3,
 \quad 0\leq j<6.
 \label{act21:eq:df-bloque-y-banda}
\end{equation}
Ainsi, \(B_c(t)\) est un mot ordonné de six trits et
\begin{equation}
 u_c(t)=\sum_{j=0}^{5}B_{c,j}(t)3^{5-j}.
 \label{act21:eq:df-elevacion-banda}
\end{equation}
Le mot et son élévation entière sont deux représentations mutuellement récupérables du même bloc émis.

\begin{theorem}[Stabilité d'extraction à toute profondeur]
\label{act21:thm:df-bloque-estable}
On vérifie, pour tout \(n\),
\[
 P_c(729,n)=P_c(3,6n).
\]
Si \(t<n\), alors
\begin{equation}
 u_c(t)=
 \left\lfloor\frac{P_c(729,n)}{729^{n-t-1}}\right\rfloor\bmod729.
 \label{act21:eq:df-bloque-prefijo-largo}
\end{equation}
En conséquence, toute extension du préfixe préserve tous les
blocs précédemment émis et toutes les signatures calculées dessus.
\end{theorem}

\begin{proof}
La première égalité résulte de ce que les deux préfixes se séparent à la même
échelle, \(729^n=3^{6n}\), et du
théorème~\ref{act21:thm:df-emision-correcta}. Pour la seconde, nous appliquons
\eqref{act21:eq:df-truncamiento} aux profondeurs \(n\) et \(t+1\), puis
nous prenons le résidu modulo \(729\). L'extraction ternaire de
\eqref{act21:eq:df-bloque-y-banda} est déterministe. Toute fonction explicite
de ces dix-huit coordonnées conserve, par conséquent, le même résultat en
étendant le préfixe.
\end{proof}

La construction est quantifiée sur \(t\in\mathbb N\). Une table qui
matérialise cent blocs constitue une évaluation finie de cette
construction ; la formule \eqref{act21:eq:df-bloque-prefijo-largo} détermine
exactement sa compatibilité avec toute extension ultérieure.

\begin{example}[Premier bloc conjoint]
L'évaluation des lecteurs régionaux sélectionnés donne les mots
\[
 B_{\mathrm P}(0)=010211,\qquad
 B_{\mathrm E}(0)=201101,\qquad
 B_{\mathrm F}(0)=121200.
\]
Leur élévation par \eqref{act21:eq:df-elevacion-banda} produit
\[
 u_{\mathrm P}(0)=103,\qquad
 u_{\mathrm E}(0)=523,\qquad
 u_{\mathrm F}(0)=450.
\]
Par exemple,
\(201101_3=2\cdot243+27+9+1=523\).
Les six positions, y compris les zéros initiaux, sont conservées dans la
représentation en bande. Le zéro initial de \(010211\) possède une position
définie, bien que l'écriture décimale de \(103\) en soit dépourvue.
\end{example}

\subsection{Registres de transition et reconstruction des élévations linéaires}
\label{act21:subsec:df-registros-transicion}

% Fuentes: deltas/integracion/RegionalTransitionRecords.lean y
% GeneratedTransitionRecords.lean; lean/biblioteca/TPKLifts.lean.
% Los enlaces reconstruidos son R(0)->R(1) y R(2)->R(3).

Sur \(\mathbb F_3\), nous réunissons deux émissions consécutives de chaque région dans la matrice
\begin{equation}
 R(t)=
 \begin{pmatrix}
 B_{\mathrm P}(t)\\ B_{\mathrm P}(t+1)\\
 B_{\mathrm E}(t)\\ B_{\mathrm E}(t+1)\\
 B_{\mathrm F}(t)\\ B_{\mathrm F}(t+1)
 \end{pmatrix}.
 \label{act21:eq:df-registro-matricial}
\end{equation}
Cette définition fixe à la fois l'ordre régional et l'ordre temporel. Les
quatre premiers registres sont
\[
 X_0=R(0),\quad Y_0=R(1),\quad X_1=R(2),\quad Y_1=R(3).
\]
Ses valeurs, obtenues par extraction des préfixes régionaux, sont
\begin{align*}
 X_0&=\begin{pmatrix}
0&1&0&2&1&1\\0&1&2&2&2&2\\2&0&1&1&0&1\\
1&2&1&2&2&1\\1&2&1&2&0&0\\1&1&2&2&0&2
\end{pmatrix},&
 Y_0&=\begin{pmatrix}
0&1&2&2&2&2\\0&1&0&2&1&1\\1&2&1&2&2&1\\
1&0&2&0&1&1\\1&1&2&2&0&2\\1&2&1&0&2&0
\end{pmatrix},\\[1ex]
 X_1&=\begin{pmatrix}
0&1&0&2&1&1\\0&0&2&1&1&1\\1&0&2&0&1&1\\
0&1&2&2&2&2\\1&2&1&0&2&0\\0&1&0&2&1&0
\end{pmatrix},&
 Y_1&=\begin{pmatrix}
0&0&2&1&1&1\\1&1&0&2&2&1\\0&1&2&2&2&2\\
1&0&2&0&1&1\\0&1&0&2&1&0\\0&1&0&2&0&0
\end{pmatrix}.
\end{align*}

\begin{theorem}[Reconstruction unique des deux liens initiaux]
\label{act21:thm:df-elevaciones-reconstruidas}
Les équations \(X_0L_0=Y_0\) et \(X_1L_1=Y_1\) ont, chacune, une
unique solution dans \(M_6(\mathbb F_3)\). Ces solutions sont
\begin{align*}
 L_0&=\begin{pmatrix}
2&2&2&1&2&1\\2&1&2&2&1&1\\1&0&0&1&0&2\\
0&0&1&1&1&0\\2&2&2&0&2&1\\2&1&2&1&0&0
\end{pmatrix},&
 L_1&=\begin{pmatrix}
2&1&2&0&2&2\\1&2&1&1&2&0\\1&2&1&1&1&2\\
0&1&0&0&2&1\\2&0&2&1&0&1\\0&2&2&2&1&1
\end{pmatrix}.
\end{align*}
Les deux matrices sont inversibles.
\end{theorem}

\begin{proof}
Les matrices
\begin{align*}
 A_0&=\begin{pmatrix}
1&1&2&0&1&2\\0&0&1&0&0&1\\2&1&2&1&0&1\\
0&2&0&1&0&2\\2&1&1&0&2&2\\2&1&1&1&1&2
\end{pmatrix},&
 A_1&=\begin{pmatrix}
1&0&1&2&0&0\\0&0&1&1&2&2\\0&1&2&0&1&1\\
1&1&0&2&0&0\\1&1&2&1&1&2\\1&0&0&0&0&2
\end{pmatrix}
\end{align*}
satisfont \(A_iX_i=X_iA_i=I\), par multiplication dans \(\mathbb F_3\).
Par conséquent, les solutions sont imposées par \(L_i=A_iY_i\) ; la
multiplication produit les deux matrices montrées. Leurs inverses sont
\begin{align*}
 L_0^{-1}&=\begin{pmatrix}
1&2&0&2&0&2\\2&0&2&2&0&0\\2&1&2&0&2&0\\
1&0&0&0&2&0\\0&2&1&1&2&0\\2&2&2&2&2&2
\end{pmatrix},&
 L_1^{-1}&=\begin{pmatrix}
2&0&2&1&2&1\\2&1&0&0&2&0\\2&0&2&2&0&2\\
2&0&0&1&1&0\\1&1&1&1&1&0\\2&0&1&2&2&0
\end{pmatrix}.
\end{align*}
La vérification bilatérale \(L_iL_i^{-1}=L_i^{-1}L_i=I\) démontre la
récupération exacte de chaque mot transformé. Tous les produits
utilisent des résidus \(0,1,2\), avec leur interprétation en \(\mathbb F_3\).
\end{proof}

La lecture causale de ce résultat est précise : les émissions régionales produisent \(R(t)\) ; les registres \(R(0),\ldots,R(3)\) déterminent les deux élévations précédentes ; leurs inverses récupèrent les mots transformés.  
La continuité pour tout \(t\) appartient aux lecteurs coinductifs et au théorème ~\ref{act21:thm:df-bloque-estable}. La reconstruction matricielle démontrée ici a pour domaine les deux liens spécifiés dans son énoncé.

\subsection{Raffinement simultané dans les bases ternaire et décimale}
\label{act21:subsec:df-cilindros-mixtos}

% Fuente: deltas/integracion/RegionalCylinderTransition.lean.

Une émission ternaire de six trits et une émission décimale de trois chiffres
constituent deux lectures de bases différentes. Pour conserver leur relation, on
enregistre l'état
\[
 s=(n,k,N,D),
\]
où \(N\) est un préfixe de profondeur \(n\) en base \(729\), et \(D\) est
un préfixe de profondeur \(k\) en base \(1000\). Leurs cylindres
archimédiens sont
\begin{equation}
 I_{729}(s)=\left[\frac{N}{729^n},\frac{N+1}{729^n}\right),\qquad
 I_{1000}(s)=\left[\frac{D}{1000^k},\frac{D+1}{1000^k}\right).
 \label{act21:eq:df-cilindros}
\end{equation}
L'extrémité supérieure semi-ouverte attribue chaque valeur à une unique cellule d'une base et d'une profondeur données.

\begin{definition}[Défauts entiers de compatibilité]
On définit
\begin{align}
 A(s)&=N1000^k-D729^n,\label{act21:eq:df-defecto-inferior}\\
 B(s)&=(N+1)1000^k-D729^n=A(s)+1000^k.
 \label{act21:eq:df-defecto-superior}
\end{align}
L'état est compatible lorsque
\begin{equation}
 A(s)<729^n,\qquad B(s)>0.
 \label{act21:eq:df-compatibilidad}
\end{equation}
\end{definition}

\begin{lemma}[Critère entier d'intersection]
Les cylindres de \eqref{act21:eq:df-cilindros} ont une intersection non vide si et
seulement si \eqref{act21:eq:df-compatibilidad} est vérifiée.
\end{lemma}

\begin{proof}
Deux intervalles semi-ouverts \([a,b)\) et \([c,d)\), avec \(a<b\) et
\(c<d\), se coupent exactement lorsque \(a<d\) et \(c<b\). En appliquant cette
condition à \eqref{act21:eq:df-cilindros} et en multipliant par les dénominateurs
positifs, nous obtenons
\(N1000^k<(D+1)729^n\) et
\(D729^n<(N+1)1000^k\). Ce sont les deux inégalités de
\eqref{act21:eq:df-compatibilidad}.
\end{proof}

\begin{definition}[Mise à jour des profondeurs et des préfixes]
Pour \(0\leq u<729\) et \(0\leq d<1000\), nous définissons
\begin{align*}
 T_u(n,k,N,D)&=(n+1,k,729N+u,D),\\
 T_{u,d}(n,k,N,D)&=(n+1,k+1,729N+u,1000D+d).
\end{align*}
La première opération raffine uniquement le cylindre ternaire. La seconde  
raffine les deux lectures et conserve les deux profondeurs dans l'état.
\end{definition}

\begin{proposition}[Récurrences exactes de frontière]
\label{act21:prop:df-rec-frontera}
Les mises à jour précédentes satisfont
\begin{align}
 A(T_u s)&=729A(s)+u1000^k,\label{act21:eq:df-rec-ternaria}\\
 A(T_{u,d}s)&=729000A(s)+u1000^{k+1}-d729^{n+1}.
 \label{act21:eq:df-rec-conjunta}
\end{align}
\end{proposition}

\begin{proof}
Pour la première identité,
\[
 (729N+u)1000^k-D729^{n+1}
 =729(N1000^k-D729^n)+u1000^k.
\]
Pour la seconde,
\begin{align*}
 &(729N+u)1000^{k+1}-(1000D+d)729^{n+1}\\
 &\qquad=729000(N1000^k-D729^n)
      +u1000^{k+1}-d729^{n+1}.
\end{align*}
Les deux sont des identités entières antérieures à toute approximation réelle.
\end{proof}

\begin{theorem}[Raffinement des cylindres régionaux générés]
\label{act21:thm:df-refinamiento-generado}
Soit
\[
 s_c(n,k)=(n,k,P_c(729,n),P_c(1000,k)).
\]
Alors \(s_c(n,k)\) est compatible quels que soient \(n,k\). De plus,
\begin{align*}
 T_{u_c(n)}s_c(n,k)&=s_c(n+1,k),\\
 T_{u_c(n),d_c(1000,k)}s_c(n,k)&=s_c(n+1,k+1).
\end{align*}
La division euclidienne des préfixes raffinés récupère exactement les
préfixes précédents. Pour tout \(a,b\geq0\),
\[
 \frac{P_c(729,n+a)-P_c(729,n+a)\bmod729^a}{729^a}
 =P_c(729,n),
\]
et l'identité analogue est vérifiée en base \(1000\).
\end{theorem}

\begin{proof}
Le théorème~\ref{act21:thm:df-emision-correcta} place \(v_c\) simultanément dans
les deux intervalles \eqref{act21:eq:df-cilindros} ; le critère d'intersection donne
la compatibilité. Les identités de mise à jour sont
\eqref{act21:eq:df-paso-posicional} dans chaque base. La récupération s'obtient à partir de
\eqref{act21:eq:df-truncamiento}. Par conséquent, la récurrence du défaut et la
récupération du préfixe sont des conséquences d'une même mise à jour.
\end{proof}

\begin{example}[Premier couple de cylindres régionaux]
Pour \(n=k=1\), l'évaluation des trois régions donne
\[
\begin{array}{c|rr|rr}
 c&N&D&A&B\\\hline
 \mathrm P&2290&3141&211&1211\\
 \mathrm E&1981&2718&-422&578\\
 \mathrm F&1179&1618&-522&478
\end{array}
\]
Dans les trois lignes \(A<729\) et \(B>0\). La première colonne de préfixes
se décompose comme \(2290=3\cdot729+103\),
\(1981=2\cdot729+523\) et \(1179=729+450\) : réapparaissent exactement les
blocs de six trits déjà émis. Les signes différents de \(A\) expriment
la position relative des deux extrémités inférieures, sans altérer la
appartenance commune de la valeur régionale.
\end{example}

\subsection{Sélection du bloc suivant par le lecteur régional}

La compatibilité de deux intervalles exprime une relation entre lectures. La détermination du bloc suivant utilise en outre l'intervalle rationnel produit par la région. Cette opération peut se formuler sans éliminer l'histoire de profondeur.

\begin{definition}[Admissibilité de la prolongation régionale]
\(c,n\) étant fixés, soient \(s=729^{n+1}\) et \(j=j_c(s)\). Un résidu
\(u\in\{0,\ldots,728\}\) est admissible lorsque
\begin{equation}
 \lfloor s\ell_c(j)\rfloor
 \leq729P_c(729,n)+u
 \leq\lceil su_c(j)\rceil-1.
 \label{act21:eq:df-admisibilidad}
\end{equation}
\end{definition}

\begin{theorem}[Unicité de la prolongation compatible avec le lecteur]
\label{act21:thm:df-prolongacion-unica}
La condition \eqref{act21:eq:df-admisibilidad} est équivalente à \(u=u_c(n)\).
Par conséquent, il existe exactement une prolongation admissible de chaque région
à chaque profondeur.
\end{theorem}

\begin{proof}
À l'indice de séparation, les deux extrémités entières de
\eqref{act21:eq:df-admisibilidad} coïncident avec \(P_c(729,n+1)\).
La condition est ainsi équivalente à
\(729P_c(729,n)+u=P_c(729,n+1)\).
Par \eqref{act21:eq:df-paso-posicional}, le côté droit est
\(729P_c(729,n)+u_c(n)\). L'annulation du terme commun donne l'égalité des résidus. Réciproquement, cette égalité satisfait la condition.
\end{proof}

La dépendance effective de cette sélection reste complètement déterminée :
région sélectionnée, lecteur rationnel, profondeur, indice de séparation,
préfixe antérieur et résidu suivant. Le défaut \(A\) résume une relation
de frontière ; l'état \((n,k,N,D)\) et les intervalles régionaux conservent
les données employées par la mise à jour.

\subsection{Signature ternaire, retenue, incidence de Witt et horloge de lecture}
\label{act21:subsec:df-firma-conjunta}

% Fuentes: lean/regional/N69RegionalSignature.lean,
% deltas/integracion/JointRegionalFrontier.lean y JointCylinderSignature.lean.

Une profondeur \(t\) étant fixée, écrivons \(B_{c,j}=B_{c,j}(t)\). Pour chaque colonne \(j\), on calcule les entiers et les résidus suivants :
\begin{align}
 S_j&=B_{\mathrm P,j}+B_{\mathrm E,j}+B_{\mathrm F,j},
 &q_j&=S_j\bmod3,
 &c_j&=\left\lfloor S_j/3\right\rfloor,
 \label{act21:eq:df-firma-qc}\\
 a_j&=(B_{\mathrm P,j}+B_{\mathrm E,j}-B_{\mathrm F,j})\bmod3,
 &w_j&=\sum_c\mathbf1_{B_{c,j}\ne0},
 &\bar w_j&=w_j\bmod3.
 \label{act21:eq:df-firma-axial}
\end{align}
Le résidu de l'expression axiale est pris dans les entiers. Une implémentation
avec des naturels le calcule comme
\((B_{\mathrm P,j}+B_{\mathrm E,j}+3-B_{\mathrm F,j})\bmod3\), qui le conserve
exactement.

La matrice d'incidence ternaire utilisée par le lecteur est
\begin{equation}
 W=\begin{pmatrix}
0&1&1&1&1&1\\
1&0&1&1&2&2\\
1&1&0&2&1&2\\
1&1&2&0&2&1\\
1&2&1&2&0&1\\
1&2&2&1&1&0
\end{pmatrix}.
\label{act21:eq:df-matriz-witt}
\end{equation}
Pour chaque région, deux charges résiduelles sont conservées :
\begin{equation}
 v_c^{\mathrm{res}}=\sum_j B_{c,j}\bmod3,
 \qquad
 v_c^{\mathrm W}=\sum_j(B_cW)_j\bmod3.
 \label{act21:eq:df-cargas-residuales}
\end{equation}
La signature complète du lecteur est
\(\mathcal S(B)=(q,a,c,w,\bar w,v^{\mathrm{res}},v^{\mathrm W})\).
Sa structure conserve séparément le reste de colonne, la retenue, le
contraste axial, le poids de support et les deux lectures de charge.

\begin{theorem}[Récupération exacte et bornes de la signature]
\label{act21:thm:df-recuperacion-firma}
Pour chaque colonne et chaque région, on a
\begin{align}
 S_j&=q_j+3c_j,&0\leq q_j,a_j&<3,&0\leq c_j&\leq2,
 \label{act21:eq:df-recuperacion-total}\\
 B_{\mathrm F,j}&=2(q_j+3-a_j)\bmod3,&0\leq w_j&\leq3,
 &0\leq\bar w_j&<3.
 \label{act21:eq:df-recuperacion-eje}
\end{align}
De plus,
\begin{equation}
 v_c^{\mathrm W}
 =\left(\sum_j\bigl((B_cW)_j\bmod3\bigr)\right)\bmod3.
 \label{act21:eq:df-witt-reduccion}
\end{equation}
Les récupérations énoncées déterminent le total de chaque colonne et sa
coordonnée d'auto-échelle ; le registre régional conserve en outre
les trois bandes ordonnées.
\end{theorem}

\begin{proof}
L'identité \eqref{act21:eq:df-recuperacion-total} est la division euclidienne de
\(S_j\) par \(3\). Chaque terme de \(S_j\) est entre \(0\) et \(2\),
de telle sorte que \(0\leq S_j\leq6\) et \(c_j\leq2\). En \(\mathbb F_3\),
la soustraction des deux congruences qui définissent \(q_j\) et \(a_j\) donne
\(q_j-a_j=2B_{\mathrm F,j}\). Comme l'inverse de \(2\) est \(2\), on
obtient \eqref{act21:eq:df-recuperacion-eje} ; le terme \(3\) permet
d'écrire celle-ci avec des nombres naturels. Les bornes de \(w_j\) résultent de sommer trois
indicateurs. Enfin, réduire une somme entière modulo \(3\) équivaut
à réduire chaque terme et à réduire à nouveau la somme, ce qui prouve
\eqref{act21:eq:df-witt-reduccion}.
\end{proof}

\begin{example}[Signature du premier bloc conjoint]
Les trois bandes initiales produisent
\begin{align*}
 q&=(0,0,2,2,1,2),&a&=(1,2,0,1,1,2),\\
 c&=(1,1,0,1,0,0),&w&=(2,2,2,3,1,2),\\
 v^{\mathrm{res}}&=(2,2,0),&v^{\mathrm W}&=(2,1,1).
\end{align*}
La quatrième colonne a un total \(2+1+2=5\) ; sa signature enregistre \(q_3=2\) et \(c_3=1\), par conséquent \(q_3+3c_3=5\).
La récupération axiale donne \(2(2+3-1)\bmod3=2\), précisément la coordonnée d'auto-échelle.
Les images de Witt, réduites composante par composante, sont
\[
 (2,0,2,1,1,2),\qquad
 (0,0,0,2,0,2),\qquad
 (2,1,1,2,1,0).
\]
Sa somme résiduelle par ligne produit les trois charges \((2,1,1)\).
\end{example}

\begin{definition}[Horloge entière de lecture décimale]
L'indice de transition \(t\) et la profondeur décimale de représentation sont
reliés par
\begin{equation}
 \kappa(t)=\max\{k\in\mathbb N:1000^k\leq3^{6(t+1)}\}.
 \label{act21:eq:df-reloj}
\end{equation}
\end{definition}

\begin{proposition}[Caractérisation et monotonie de l'horloge]
\label{act21:prop:df-reloj}
L'entier \(\kappa(t)\) est déterminé de manière univoque par
\begin{equation}
 1000^{\kappa(t)}\leq729^{t+1}<1000^{\kappa(t)+1}.
 \label{act21:eq:df-reloj-cotas}
\end{equation}
La fonction \(\kappa\) est monotone et satisfait
\(\kappa(20)=\kappa(21)=20\).
\end{proposition}

\begin{proof}
Les puissances de \(1000\) sont strictement croissantes et non bornées ;
il existe donc un unique intervalle de puissances consécutives qui
contient \(729^{t+1}\). La monotonie de cette dernière expression démontre
celle de \(\kappa\). Les deux évaluations finales sont les comparaisons
entières
\[
 1000^{20}\leq729^{21}<1000^{21},\qquad
 1000^{20}\leq729^{22}<1000^{21}.
\]
Ces comparaisons s'effectuent avec des puissances entières exactes.
\end{proof}

L'égalité de profondeur décimale en deux transitions distinctes explique la nécessité de conserver les deux indices : l'horloge de lecture peut rester constante tandis que la bande ternaire avance et que sa signature est mise à jour. L'information chronologique se trouve dans le registre ordonné des transitions.

\subsection{Théorème de composition du bloc, du cylindre et de la signature}

\begin{theorem}[Compatibilité conjointe des lectures régionales]
\label{act21:thm:df-composicion-conjunta}
Pour chaque profondeur \(n\), il existe un unique triplet de blocs
\((u_{\mathrm P},u_{\mathrm E},u_{\mathrm F})\) qui satisfait
l'admissibilité régionale \eqref{act21:eq:df-admisibilidad} dans les trois régions.
Ce triplet est \((u_{\mathrm P}(n),u_{\mathrm E}(n),u_{\mathrm F}(n))\).
Simultanément :
\begin{enumerate}
\item chaque bloc met à jour son cylindre par les récurrences
\eqref{act21:eq:df-rec-ternaria} et \eqref{act21:eq:df-rec-conjunta} ;
\item ses six trits produisent la signature
\(\mathcal S(B_{\mathrm P}(n),B_{\mathrm E}(n),B_{\mathrm F}(n))\);
\item les totaux de colonne se retrouvent comme \(q+3c\), et les charges de Witt sont obtenues à partir des mêmes bandes par
\eqref{act21:eq:df-witt-reduccion} ;
\item tout préfixe de plus grande profondeur renvoie exactement ces
blocs, cylindres tronqués et signatures.
\end{enumerate}
\end{theorem}

\begin{proof}
Nous appliquons le théorème ~\ref{act21:thm:df-prolongacion-unica} à chacun des trois
indices régionaux. L'unicité coordonnée par coordonnée donne l'unicité
conjointe. Nous substituons ces mêmes résidus dans les opérations
\(T_u,T_{u,d}\), ce qui permet d'appliquer le
théorème ~\ref{act21:thm:df-refinamiento-generado}. L'application déterministe de
\eqref{act21:eq:df-bloque-y-banda} produit les bandes utilisées par
\eqref{act21:eq:df-firma-qc}--\eqref{act21:eq:df-cargas-residuales}. Les identités
de récupération sont celles du théorème ~\ref{act21:thm:df-recuperacion-firma}.
Finalement, le théorème ~\ref{act21:thm:df-bloque-estable} et le troncage de
chaque base démontrent la stabilité à plus grande profondeur. À chaque étape, on a conservé l'identité du bloc émis ; aucune des deux
lectures ne requiert de sélectionner un autre bloc.
\end{proof}

La conclusion est une compatibilité opérationnelle entre la contraction
cylindrique et la conservation de l'incidence. Les deux utilisent le même
registre généré, maintiennent les profondeurs et admettent la récupération
entière. Ce résultat est la forme explicite du lien
\[
 \text{bloc régional émis}
 \longrightarrow
 \begin{cases}
 \text{cylindre et frontière entière},\\
 \text{bande, signature, retenue et charge de Witt}.
 \end{cases}
\]
Le développement ultérieur de l'incidence et des champs emploie la deuxième lecture, conservant l'origine commune avec la première.

\subsection{Recensement des calendriers et compatibilité simultanée des biographies régionales}
Les lignes précédentes se regroupent, sans appliquer encore aucun lecteur décimal,
dans les trois biographies de blocs produites par les transitions :
\begin{align}
 B^{\rm cl}
  &=010211|012222|010211|002111|110221,\notag\\
 B^{\rm pr}
  &=201101|121221|102011|012222|102011,
 \label{act21:eq:biografias-tpk-tres-caracteres}\\
 B^{\rm au}
  &=121200|112202|121020|010210|010200.\notag
\end{align}
En effet, les paires consécutives des deux premiers tronçons sont les lignes de \(X_0,Y_0\), et celles des deux derniers sont les lignes de \(X_1,Y_1\). Ainsi, les biographies sont une autre écriture du registre de douze transitions, non une collection de préfixes conventionnels.

Pour vérifier le mot opératoire du calendrier, soit
\(c=c_1c_2c_3c_4\in\{0,1\}^4\) et, pour chaque régime
\(\chi\in\{{\rm cl},{\rm pr},{\rm au}\}\), définissons
\[
 u_0^\chi=b_0^\chi,\qquad
 u_j^\chi=u_{j-1}^\chi L_{c_j}\quad(1\le j\le4).
\]
Pour condenser le recensement, notons par
\[
 N_{\mathrm{ar}}(c):=
 \#\{(j,\chi):u_j^\chi=b_j^\chi,\ 1\le j\le4\},\qquad
 N_{\mathrm{bio}}(c):=
 \#\{\chi:(u_0^\chi,\ldots,u_4^\chi)=B^\chi\}.
\]
La multiplication exacte dans \(\FF_3\) des seize calendriers donne :
\[
\begin{array}{c|c|c}
 c&N_{\mathrm{ar}}(c)&N_{\mathrm{bio}}(c)\\ \hline
0000,0001&6&0\\
0010&9&0\\
0011&12&3\\
0100,0101,0110,0111&3&0\\
1000,1001,1010,1011&0&0\\
1100,1101,1110,1111&0&0
\end{array}
\tag{D}
\label{act21:eq:censo-dieciseis-calendarios}
\]
Le tableau épuise les \(2^4=16\) possibilités et chacune de ses entrées s'obtient
par quatre multiplications d'une ligne par l'une des deux matrices reproduites.
Par conséquent, \(0011\) est l'unique calendrier qui reconstruit
simultanément les trois biographies complètes. En notation opératorielle,
\begin{equation}
 (L_0,L_0,L_1,L_1)
 \label{act21:eq:calendario-0011-interno}
\end{equation}
ce n'est pas une prescription indépendante : c'est la seule lecture binaire compatible
avec les douze arêtes déjà produites par \(U_t\).

Le calendrier \eqref{act21:eq:calendario-0011-interno} produit quatre nouveaux
blocs de six composants. Leurs concaténations forment
\[
 w_6\longrightarrow w_{12}\longrightarrow w_{18}
 \longrightarrow w_{24}\longrightarrow w_{30}.
\]
Les sous-indices de \(w\) indiquent la longueur accumulée. La frontière \(R_{36}\) conserve la structure terminale et ouvre la relation de prolongement ; elle n'est pas renommée comme un dernier bloc périodique qui pourrait se répéter indéfiniment.

L'identité \(L=X^{-1}Y\) démontre l'unicité une fois les
transitions fixées ; elle ne remplace pas leur production par \(U_t\). Réciproquement,
\eqref{act21:eq:censo-dieciseis-calendarios} prouve dans l'article l'unicité
de l'ordre des quatre transports. La chaîne causale est ainsi close :
le recensement APP produit les régions initiales, le TRIT fixe le régime et
l'orientation, le TPK produit le registre des transitions, et ce registre
détermine \(L_0,L_1\) et \(0011\) avant toute réalisation archimédienne.


\clearpage
% Desarrollo reunido K31-07--K31-15. Inserción anterior a registro_k y alpha.
% Mantiene explícita la interfaz del repertorio terminal de ocho bloques.
% Propietarios y localizadores: metadata/K_ALPHA_PROPIETARIOS.md.
\section{Détermination régionale du descripteur dodécaphasique et sélection du registre de couplage}
\label{act21:chap:despliegue-selector-k}

La détermination du registre de couplage articule deux constructions
qu'il convient de déployer dans leur ordre effectif. La première produit les
antécédents du sélecteur à partir des représentations régionales du TPK :
bandes de six trits, signatures de colonne, calendrier de conversion,
descripteur orienté et tableau fonctionnel. La seconde opère sur ces
antécédents et sur le répertoire terminal de huit blocs, énumère leurs
ordonnancements admissibles et sélectionne un unique registre au moyen de deux
conditions simultanées : incidence exceptionnelle et hétérotypie de lecture.
La valeur du registre apparaît au terme de cette composition.

Le présent développement présuppose les représentations régionales obtenues
dans les chapitres précédents à partir d'APP--TRIT--TPK. Son propos est de rendre
explicite chaque opération intermédiaire qui conduit de ces représentations
au sélecteur fini. Sont ensuite incorporées la réalisation du même
registre par incidences signées, l'inversion de Hadamard et la
représentation de la constante de structure fine. Les trois constructions
conservent leurs domaines respectifs : sélection du registre, récupération
à partir de ses canaux et évaluation arithmétique de ses représentations.

\subsection{Bandes régionales, codification et signatures entières}

\subsubsection{Domaines et codification des bandes régionales}

On fixe l'ordre des canaux
\[
 \chi\in\{\pi,e,\varphi\},
 \qquad B_\chi(t)\in\mathbb F_3^6,
 \qquad t\in\mathbb N.
\]
Les symboles de canal identifient les régions déjà engendrées. Chaque
$B_\chi(t)$ conserve la place qu'il occupe dans sa suite de représentations ;
l'indice $t$ compte les blocs ternaires. L'évaluation naturelle d'une
bande $b=(b_0,\ldots,b_5)$ est
\begin{equation}
 \operatorname{enc}_6(b)=
 243b_0+81b_1+27b_2+9b_3+3b_4+b_5,
 \qquad 0\leq\operatorname{enc}_6(b)<729.
 \label{act21:eq:kdes-enc6}
\end{equation}
Son inverse, avec une largeur fixée, est
\[
 \operatorname{dig}_j(n)=
 \left\lfloor\frac{n}{3^{5-j}}\right\rfloor\bmod3,
 \qquad 0\leq j<6.
\]
Ainsi sont retenus les zéros initiaux : les mots $002001$ et $2001$
pourraient partager une écriture entière abrégée, mais le type du
premier contient six positions. Toutes les concaténations qui suivent
utilisent cette largeur.

Dans la fenêtre de cent blocs, la représentation de six cents trits du
canal $\chi$ détermine un entier $P_\chi^{(600)}$. L'extraction s'effectue
par
\begin{equation}
 b_\chi(t)=
 \left\lfloor\frac{P_\chi^{(600)}}{729^{99-t}}\right\rfloor\bmod729,
 \qquad 0\leq t<100,
 \qquad B_\chi(t)=\operatorname{enc}_6^{-1}(b_\chi(t)).
 \label{act21:eq:kdes-extraccion}
\end{equation}
L'entier de six cents trits est une représentation du producteur régional ;
la table de cent lignes est l'évaluation de
\eqref{act21:eq:kdes-extraccion}. Cette séparation permet de remplacer une table
historique par son producteur sans modifier le sélecteur qui l'utilise.

\subsubsection{Résidus, quotients, occupation et orientation}

Pour chaque colonne $j$ et temps $t$, on introduit
\begin{align}
 T_j(t)&=B_\pi(t)_j+B_e(t)_j+B_\varphi(t)_j,\\
 q_j(t)&=T_j(t)\bmod3,
 &c_j(t)&=\left\lfloor T_j(t)/3\right\rfloor,\\
 a_j(t)&=\bigl(B_\pi(t)_j+B_e(t)_j-B_\varphi(t)_j\bigr)\bmod3,\\
 w_j(t)&=\mathbf1_{B_\pi(t)_j\ne0}
       +\mathbf1_{B_e(t)_j\ne0}
       +\mathbf1_{B_\varphi(t)_j\ne0},
 &\bar w_j(t)&=w_j(t)\bmod3.
 \label{act21:eq:kdes-firmas}
\end{align}
La soustraction de la troisième expression s'effectue sur les entiers et se réduit
ensuite. Dans une implémentation sur les naturels, l'expression équivalente
est $(B_\pi+B_e+3-B_\varphi)\bmod3$ ; la soustraction tronquée altérerait le lecteur.

\begin{proposition}[Reconstruction de la somme de colonne]
Pour toutes les bandes admissibles,
\[
 T_j(t)=q_j(t)+3c_j(t),\qquad
 0\leq q_j(t)<3,\quad 0\leq c_j(t)\leq2,
 \quad0\leq w_j(t)\leq3.
\]
\end{proposition}
\begin{proof}
Chaque colonne additionne trois représentants de $\{0,1,2\}$, de sorte que
$0\leq T_j(t)\leq6$. La division euclidienne par trois fournit le reste
et le quotient indiqués. L'occupation compte trois conditions binaires ;
sa borne se déduit directement de cette définition.
\end{proof}

Cette identité conserve le quotient de la somme locale. La reconstruction
des trois bandes individuelles requiert les bandes elles-mêmes ou un lecteur
supplémentaire qui les distingue : une somme de colonne a moins de coordonnées
que le triplet qui la produit. Par conséquent, la ligne temporelle utilisée plus
tard conserve conjointement les trois bandes et les quatre champs
$q,a,c,\bar w$.

\subsection{Horloge de conversion et premier avancement nul}

\subsubsection{Capacité ternaire et profondeur décimale}

Une représentation de $t+1$ bandes contient $6(t+1)$ trits. Le calendrier
de conversion associé est défini par l'entier
\begin{equation}
 d_t=\max\{k\in\mathbb N:1000^k\leq3^{6(t+1)}\}
     =\max\{k\in\mathbb N:1000^k\leq729^{t+1}\}.
 \label{act21:eq:kdes-reloj}
\end{equation}
La définition compare des capacités entières. Le compteur $t$ et la profondeur $d_t$ enregistrent des grandeurs distinctes : une transition ajoute une bande de six trits ; l'horloge compte les puissances complètes de la base décimale mille contenues dans la capacité ternaire accumulée.

\begin{lemma}[Caractérisation entière de l'horloge]
Pour $t,k\in\mathbb N$ on a
\[
 d_t=k\quad\Longleftrightarrow\quad
 1000^k\leq729^{t+1}<1000^{k+1}.
\]
De plus, $d_t\leq d_{t+1}$.
\end{lemma}
\begin{proof}
Les puissances de mille forment une suite strictement croissante et non bornée. La puissance $729^{t+1}$ se situe entre deux termes consécutifs ; l'exposant inférieur est exactement le maximum de \eqref{act21:eq:kdes-reloj}. L'inégalité $729^{t+1}<729^{t+2}$ implique la monotonie du maximum.
\end{proof}

\begin{proposition}[Premier avancement nul]
Le plus petit temps qui satisfait $d_{t+1}=d_t$ est
\[
 t_*=20.
\]
En particulier,
\[
 d_0=0,\quad d_{19}=19,\quad d_{20}=20,
 \quad d_{21}=20,\quad d_{22}=21.
\]
\end{proposition}
\begin{proof}
Les comparaisons entières exactes
\[
 1000^{20}\leq729^{21},\qquad
 729^{22}<1000^{21},\qquad
 1000^{21}\leq729^{23}<1000^{22}
\]
donnent les valeurs aux temps vingt, vingt-et-un et vingt-deux. Pour
$0\leq t\leq20$, l'expression
$729^{t+1}/1000^t=729(729/1000)^t$ est décroissante et sa valeur en vingt
est supérieure ou égale à un. Par conséquent,
$1000^t\leq729^{t+1}<1000^{t+1}$ et $d_t=t$. Aucun temps inférieur à
vingt ne peut satisfaire $d_{t+1}=d_t$ ; en vingt, c'est le cas.
Toutes les comparaisons utilisées sont de puissances entières.
\end{proof}

Le nombre vingt est donc une sortie du critère minimal. La règle qui sélectionne l'instant est ``première avance nulle de l'horloge'' ; l'algorithme et sa preuve peuvent utiliser cette règle sans recevoir vingt comme paramètre externe.

\subsubsection{Retour nonadique et antécédent temporel}

La longueur du retour de phase détermine
\begin{equation}
 r=t_*-9=11.
 \label{act21:eq:kdes-retorno}
\end{equation}
Les requêtes du descripteur ont lieu en $t_*+1$, $r+1$ et $r$.
Cette distribution conserve la relation entre l'instant de pause, la
translation nonadique et la lecture de phase immédiatement associée. Les
temps onze, douze et vingt-et-un désignent des positions du calendrier
régional ; leur différence temporelle reste enregistrée même lorsque deux
temps produisent la même profondeur décimale.

\subsection{Construction orientée du descripteur de vingt-quatre trits}

\subsubsection{Défauts opposés de semi-bande}

On considère les défauts orientés de la coordinatisation régionale
\[
 \delta_+=(0,0,0,1,1,1),\qquad
 \delta_-=(0,0,0,2,2,2)\in\mathbb F_3^6.
\]
Son opposition se vérifie composante par composante :
\[
 \delta_++\delta_-=0\quad\text{en }\mathbb F_3^6.
\]
Le représentant $2$ réalise $-1$ dans la coordinatisation ternaire. La concaténation
des queues de trois positions, dans l'ordre négatif--positif, produit
le mot $222111$. La somme des défauts et la concaténation de leurs
queues sont des opérations distinctes : la première a pour codomaine
$\mathbb F_3^6$, tandis que la seconde conserve l'ordre de deux
segments de longueur trois.

\subsubsection{Charge duale et sélection de phase}

Dans la coordinatisation régionale, la transformation de six coordonnées utilise
\[
 A_{\rm reg}=
 \begin{pmatrix}
 0&1&1&1&1&1\\
 1&0&1&1&2&2\\
 1&1&0&2&1&2\\
 1&1&2&0&2&1\\
 1&2&1&2&0&1\\
 1&2&2&1&1&0
 \end{pmatrix}.
\]
Avec des vecteurs ligne, la charge duale de $b\in\mathbb F_3^6$ est la somme
des coordonnées de $bA_{\rm reg}$. Les sommes entières des lignes de
$A_{\rm reg}$ sont $(5,7,7,7,7,7)$ ; d'où résulte l'expression locale
\begin{equation}
 \eta^\vee(b)=2b_0+b_1+b_2+b_3+b_4+b_5\pmod3.
 \label{act21:eq:kdes-cargadual}
\end{equation}
La phase de lecture à l'instant $r$ est
$\theta_r=(r-1)\bmod3$. Pour chaque canal, on définit
\[
 \varepsilon_\chi(r)=
 \begin{cases}
 1,&\eta^\vee(B_\chi(r))=\theta_r,\\
 2,&\eta^\vee(B_\chi(r))\ne\theta_r.
 \end{cases}
\]
La représentation phase--charge est
\begin{equation}
 D(r)=\varepsilon_\pi(r)B_\pi(r)
      +\varepsilon_e(r)B_e(r)
      +\varepsilon_\varphi(r)B_\varphi(r)
      \quad\text{en }\mathbb F_3^6.
 \label{act21:eq:kdes-fasecarga}
\end{equation}
Le coefficient se détermine par une égalité de charges, avant d'évaluer le mot qui résultera de la combinaison.

\subsubsection{Composition du descripteur}

On désigne par $\Vert$ la concaténation de mots de largeur fixée.
La construction complète est
\begin{equation}
 \begin{split}
 \Omega_{24}={}&
 (B_e(t_*+1)+\delta_-)
 \Vert (\delta_-|_{3,4,5}\Vert\delta_+|_{3,4,5})\\
 &\Vert(B_\pi(r+1)+B_e(r+1))\Vert D(r).
 \end{split}
 \label{act21:eq:kdes-omega}
\end{equation}
Chaque terme de longueur six se calcule en $\mathbb F_3^6$ ; la
concaténation finale a une longueur $6+6+6+6=24$. Cette construction
sépare le descripteur $\Omega_{24}$ des étapes régionales
$w_{24}^{\chi}$ : les deux objets partagent la même longueur, mais leurs producteurs
et leurs places dans la composition sont différents.

\begin{proposition}[Évaluation régionale du descripteur]
Les représentations régionales et le critère de la première avancée nulle
déterminent
\[
 \Omega_{24}=020022\,222111\,211201\,021101,
 \qquad [\Omega_{24}]_3=66241910521.
\]
\end{proposition}
\begin{proof}
Les bandes nécessaires, conservant l'ordre $(\pi,e,\varphi)$, sont
\[
 \begin{array}{c|ccc}
 t&B_\pi(t)&B_e(t)&B_\varphi(t)\\\hline
 11&022212&110001&100021\\
 12&012012&202222&000120\\
 21&221022&020100&211021
 \end{array}
\]
et proviennent de l'extraction \eqref{act21:eq:kdes-extraccion}. Au temps
vingt-et-un,
$020100+000222=020022$ en $\mathbb F_3^6$. Les queues des défauts
apportent $222111$. Au temps douze,
$012012+202222=211201$.

Pour le temps onze, les charges duales sont respectivement $0,1,2$ ; la phase vaut $(11-1)\bmod3=1$. Les coefficients sont $(2,1,2)$ et
\[
 2(022212)+(110001)+2(100021)=021101
 \quad\text{en }\mathbb F_3^6.
\]
La concaténation de ces quatre résultats donne le mot annoncé.
Son évaluation entière s'obtient par Horner en base trois, ou par trois
concaténations successives en base $729$.
\end{proof}

Le tableau précédent joue un rôle vérificateur : chacune de ses entrées provient du producteur régional. La règle de construction \eqref{act21:eq:kdes-omega} est celle qui permet de reproduire le descripteur sans introduire en tant que donnée la chaîne finale de vingt-quatre trits.

\subsection{Panneau fonctionnel et répertoire terminal}

\subsubsection{Neuf positions régionales}

Le panneau fonctionnel conserve les trois premières bandes de chacun des trois canaux, dans l'ordre régional déclaré :
\begin{equation}
 \mathcal P_9=
 \bigl(b_\pi(0),b_\pi(1),b_\pi(2),
       b_e(0),b_e(1),b_e(2),
       b_\varphi(0),b_\varphi(1),b_\varphi(2)\bigr).
 \label{act21:eq:kdes-panel}
\end{equation}
Son évaluation est
\[
 \mathcal P_9=(103,161,103,523,457,301,450,398,438).
\]
Les neuf positions conservent l'origine régionale, bien que la valeur
103 apparaisse deux fois. Le panneau comporte neuf incidences ordonnées et
huit valeurs distinctes. L'énumération ultérieure utilise les incidences ;
la déduplication des candidats s'effectue à une étape postérieure.

\subsubsection{Répertoire non ordonné de huit blocs}

Le répertoire terminal qui intervient dans cette construction est
\begin{equation}
 \mathcal S_{\rm term}=
 \{002001,020111,200110,121012,
   010122,001001,221110,222110\}\subset\mathbb F_3^6.
 \label{act21:eq:kdes-sterm}
\end{equation}
Ses huit éléments sont distincts. La codification entière, ordonnée par
valeur uniquement pour fixer une représentation informatique, est
\[
 \{28,55,98,175,437,498,687,714\}.
\]
La provenance documentaire de ce répertoire est la segmentation terminale du mot de soixante-douze trits et la récupération par positions au moyen des incidences des lecteurs phase--charge et affine. Le sélecteur développé ci-après reçoit le répertoire sans son ordonnancement historique et met à l'épreuve ses $8!$ ordonnancements.

La formulation exacte de l'implémentation réunie maintient ici une
interface explicite : les producteurs régionaux remplacent matériellement
le descripteur, le panneau et les lignes temporelles ; le répertoire
\eqref{act21:eq:kdes-sterm} conserve son producteur documentaire antérieur. Cette
distinction identifie où intervient chaque antécédent. L'appartenance de
huit mots à un répertoire temporel de lecteurs, la sélection de ces
huit mots et la détermination de leur ordre terminal sont trois
énoncés différents. La preuve de sélection terminale démontre le
troisième pour le répertoire déclaré.

La récupération documentaire par positions peut s'exprimer de manière
exacte. Soit $\mathcal I$ le registre des incidences phase--charge et soit
$\mathcal J$ le registre de l'extension affine. Chaque incidence conserve
la position $\ell$ du mot de soixante-douze trits et le mot
émis $u$. Pour $5\leq\ell\leq12$, on forme
\[
 \mathcal S_\ell=
 \{u:(\ell,u)\in\mathcal I\cup\mathcal J\}.
\]
Le récupérateur vérifie que les huit positions apparaissent et que chaque
$\mathcal S_\ell$ est un singleton. Si $s_\ell$ est son élément unique,
la sortie livrée au sélecteur est
\[
 \mathcal S_{\rm term}=\{s_5,s_6,\ldots,s_{12}\}.
\]
La preuve de singleton par position préserve la cohérence des témoins historiques ; l'exportation ultérieure préserve les huit mots et supprime uniquement la mise en ordre historique. Le sélecteur reconstruit l'ordre par l'énumération complète et les conditions terminales, au lieu de le recevoir du registre d'origine.

% RESULTADO_RECUPERADO: K31-10/P10 del inventario REV05.
% Fuente de las incidencias:
% 16_CIERRE_GLOBAL_HMT_MD_2026-07-22/02_LECTURA_DODECAFASICA/continuacion_fases/
% verificar_obstruccion_y_dato_minimo_w72.py, lineas 45--60 y 141--233.
% RESULTADO_OBSTRUCCION_Y_DATO_MINIMO_W72.json, phase_charge_reader.
% Fuente de las bandas: N69_signatures_t0_t99.csv, tiempos impresos abajo.
% Consumidor: PAQUETE_CONTINUIDAD_K_UNIDAD_20260918_BASE_COMPARTIDA/python/
% selector_orbital_original.py, verify_historical_inputs, lineas 159--193.
% No se transfiere la restriccion del lector comparativo al generador completo.
\subsubsection{Incidences temporelles du répertoire terminal et récupération par position}
\label{act21:sec:repertorio-incidencias-completas}

La représentation du répertoire terminal utilise les mêmes blocs régionaux,
les mêmes retenues et le même calendrier qui produisent le descripteur initial
de vingt-quatre trits. Il convient d'expliciter séparément deux opérations de
cette composition. La première récupère huit mots à partir d'incidences
temporelles enregistrées. La seconde présente ces huit mots au sélecteur
dodécaphasique comme un ensemble non ordonné. Le sélecteur rétablit
l'ordre par ses conditions d'incidence, de cylindre et de clôture.
Cette distinction permet de conserver les témoins de provenance sans transformer
l'ordre d'une segmentation historique en donnée d'entrée du sélecteur.

Soit $B_t=(b_{t,0},b_{t,1},b_{t,2})\in(\mathbb F_3^6)^3$ la frontière
régionale déjà engendrée au temps $t$, avec $0\leq t<100$. Dans la coordinatisation
utilisée par le registre historique, la matrice d'incidence est
\[
 A_{\mathrm{hist}}=
 \begin{pmatrix}
 0&1&1&1&1&1\\
 1&0&1&1&2&2\\
 1&1&0&2&1&2\\
 1&1&2&0&2&1\\
 1&2&1&2&0&1\\
 1&2&2&1&1&0
 \end{pmatrix}.
\]
Les mots sont considérés comme des vecteurs ligne. Leurs deux charges et la phase
chronologique sont
\[
 r_{t,i}=\sum_{j=1}^{6}b_{t,i,j},\qquad
 d_{t,i}=\sum_{j=1}^{6}(b_{t,i}A_{\mathrm{hist}})_j,
 \qquad \vartheta_t=t-1\pmod 3.
\]
Toutes les opérations de ces trois expressions ont lieu dans
$\mathbb F_3$. L'indice historique identifie une coordinatisation précise de
l'incidence de Witt ; il maintient explicite la transformation de coordonnées
lorsqu'une autre présentation de cette même incidence est utilisée.

Pour $c\in\{r_t,d_t\}$, $\delta\in\mathbb F_3$ et
$\epsilon\in\{1,-1\}$, nous définissons
\[
 \eta_i(c,\delta,\epsilon)=
 \begin{cases}
  \epsilon,&c_i=\vartheta_t+\delta,\\
  -\epsilon,&c_i\ne\vartheta_t+\delta,
 \end{cases}
 \qquad
 L_{t,c,\delta,\epsilon}
   =\sum_{i=0}^{2}\eta_i(c,\delta,\epsilon)b_{t,i}.
\]
L'orientation $-1$ s'écrit $2$ lors de l'impression des coefficients
ternaires. Outre ces douze lecteurs comparatifs par frontière, le
registre dispose du lecteur affine
\[
 L^{\mathrm{af}}_t
    =\sum_{i=0}^{2}(d_{t,i}+\vartheta_t)b_{t,i}.
\]
L'addition de la phase dans le lecteur affine et la comparaison des charges dans
le lecteur précédent sont des opérations distinctes, toutes deux sur des données déjà
produites par la frontière TPK.

\begin{table}[htbp]
\centering\small
\begin{tabular}{rccc cc}
\toprule
$t$&$b_{t,0}$&$b_{t,1}$&$b_{t,2}$&$r_t$&$d_t$\\
\midrule
8 &200121&212212&110110&011&202\\
11&022212&110001&100021&001&012\\
30&210112&110202&211110&100&012\\
34&121022&000011&112110&220&021\\
37&120121&011202&112221&100&201\\
49&110212&200212&100122&110&201\\
58&111121&022121&122201&122&220\\
67&010012&001220&011120&122&122\\
75&220000&020111&100002&120&021\\
81&122201&020202&201122&202&001\\
90&022112&022110&121010&202&200\\
\bottomrule
\end{tabular}
\caption{Frontières temporelles qui réalisent les incidences imprimées du
répertoire terminal. Chaque entrée de six chiffres conserve l'ordre des
six coordonnées de la frontière régionale.}
\label{act21:tab:terminal-bandas-testigo}
\end{table}

Le tableau suivant déploie toutes les incidences comparatives du
registre sur les douze blocs de la segmentation considérée. Les
lettres $r$ et $d$ désignent respectivement la charge visible et la charge
duale. La colonne $j$ conserve l'indice de position du témoin historique ;
l'opération de livraison au sélecteur oublie cet indice après avoir formé
l'ensemble de mots.

\begin{table}[htbp]
\centering\small
\begin{tabular}{rrcrrc c}
\toprule
$j$&$t$&charge&$\delta$&$\epsilon$&$(\eta_0\eta_1\eta_2)$&$L_{t,c,\delta,\epsilon}$\\
\midrule
4&11&$d$&0& 1&212&021101\\
5&67&$d$&1&$-1$&211&002001\\
5&67&$d$&2& 1&211&002001\\
5&67&$r$&1&$-1$&211&002001\\
5&67&$r$&2& 1&211&002001\\
6&81&$d$&0& 1&222&020111\\
6&81&$r$&2& 1&222&020111\\
7&34&$r$&1&$-1$&111&200110\\
7&75&$d$&1&$-1$&211&200110\\
7&75&$r$&2&$-1$&211&200110\\
8&90&$r$&0& 1&121&121012\\
8&90&$r$&1&$-1$&121&121012\\
9&49&$d$&0&$-1$&121&010122\\
11&11&$d$&2&$-1$&211&221110\\
11&37&$d$&0&$-1$&121&221110\\
12&8&$d$&0&$-1$&111&222110\\
12&8&$r$&1&$-1$&111&222110\\
12&58&$d$&1&$-1$&111&222110\\
12&58&$r$&0&$-1$&111&222110\\
\bottomrule
\end{tabular}
\caption{Incidences phase--charge complètes sur la segmentation terminale  
aux cent frontières du registre. Les coïncidences de mots entre  
lignes conservent leurs temps, orientations et types de charge.}
\label{act21:tab:terminal-incidencias-completas}
\end{table}

Le dixième bloc dispose d'un témoin affine. Dans $t=30$ on a
$d_{30}=(0,1,2)$ et $\vartheta_{30}=2$, d'où
\[
 (d_{30,0}+\vartheta_{30},d_{30,1}+\vartheta_{30},
 d_{30,2}+\vartheta_{30})=(2,0,1)
\]
et, en utilisant les bandes du tableau ~\ref{act21:tab:terminal-bandas-testigo},
\[
 L^{\mathrm{af}}_{30}
 =2(210112)+0(110202)+(211110)
 =(001001)\quad\text{en }\mathbb F_3^6.
\]
L'égalité est coordonnée par coordonnée : avant de réduire modulo trois,  
le vecteur de la somme est $(6,3,1,3,3,4)$.

\begin{proposition}[Récupération intégrale du répertoire par incidences]
\label{act21:prop:repertorio-terminal-ocho-incidencias}
Considérons le registre des incidences comparatives du
tableau~\ref{act21:tab:terminal-incidencias-completas}, augmenté du témoin
affine $(j,t,L)=(10,30,001001)$. Pour chaque position $5\leq j\leq12$,
l'ensemble des mots de sortie de ses incidences est un singleton. L'
extraction de ces huit mots et l'oubli ultérieur de leur ordre
produisent exactement
\[
 \mathcal S_{\mathrm{term}}=
 \{002001,020111,200110,121012,010122,001001,221110,222110\}.
\]
Sa codification en tant que nombres entiers en base trois est
\[
 \{28,55,98,175,437,498,687,714\}.
\]
\end{proposition}
\begin{proof}
La multiplicité d'incidences par position $j=5,6,7,8,9,10,11,12$
est, respectivement,
\[
 (4,2,3,2,1,1,2,4).
\]
À chaque position, toutes les lignes impriment le même mot. Chaque ligne se vérifie en substituant ses trois bandes et les trois coefficients dans la formule de $L$ ; le dixième cas a été calculé explicitement au moyen de $L^{\mathrm{af}}_{30}$. Par exemple, en $t=67$ et avec les coefficients $(2,1,1)$, on obtient
\[
 2(010012)+(001220)+(011120)=(002001)\pmod3.
\]
Les huit résultats sont distincts en tant que mots de longueur six. Leur codification $\sum_{k=1}^{6}w_k3^{6-k}$ donne, dans l'ordre des positions, $(55,175,498,437,98,28,714,687)$. En oubliant l'ordre, on obtient l'ensemble des entiers énoncé. Ainsi, l'extraction préserve chaque mot complet, y compris ses positions nulles, et fournit huit éléments distincts au sélecteur.
\end{proof}

La preuve précédente est une récupération exacte à partir d'incidences
identifiées. Sa donnée d'entrée comprend le registre d'incidences
avec leurs positions et leurs témoins. La sélection orbitale ultérieure reçoit
uniquement $\mathcal S_{\mathrm{term}}$ et le descripteur régional initial ;
elle parcourt les $8!$ ordonnancements et détermine leur compatibilité avec les
frontières fonctionnelles et les incidences orientées de la roue
dodécaphasique. Les deux opérations sont ainsi contiguës, avec des domaines
explicites : provenance temporelle des huit mots et détermination
constructive de leur ordre.

Le tableau comparatif contient les dix-neuf incidences que sa formule produit sur la segmentation de douze blocs aux cent temps déclarés. Cette exhaustivité finie se vérifie en parcourant $100\cdot2\cdot3\cdot2=1200$ évaluations. Le lecteur affine étend la couverture des positions terminales de sept à huit. Ce décompte caractérise le répertoire local de lecteurs qui vient d'être défini ; les autres opérations du TPK conservent les domaines et les fonctions constructives établis dans leurs dérivations respectives.

\subsection{Lecteurs comparatifs, lecteur affine et incidence temporelle}

\subsubsection{Transformation de Witt du sélecteur}

Le sélecteur utilise la coordinatisation à six coordonnées
\[
 A_{\rm ter}=
 \begin{pmatrix}
 0&1&1&1&1&1\\
 1&0&1&2&2&1\\
 1&1&0&1&2&2\\
 1&2&1&0&1&2\\
 1&2&2&1&0&1\\
 1&1&2&2&1&0
 \end{pmatrix}.
\]
On définit $W(b)=bA_{\rm ter}$ dans $\mathbb F_3^6$, ainsi que les charges
\[
 \eta(b)=\sum_{j=0}^5 b_j\pmod3,
 \qquad\eta_{\rm ter}^\vee(b)=\eta(W(b)).
\]
La coordinatisation du descripteur et celle du sélecteur sont reliées par la
permutation $\sigma=(0,1,2,4,5,3)$, écrite comme liste d'images :
\begin{equation}
 (A_{\rm ter})_{ij}=(A_{\rm reg})_{\sigma(i),\sigma(j)}.
 \label{act21:eq:kdes-cambio-carta}
\end{equation}

\begin{lemma}[Compatibilité de la fonctionnelle de charge]
Pour toute bande $b$, on a
\[
 \eta_{\rm ter}^\vee(b)=\eta^\vee(b).
\]
\end{lemma}
\begin{proof}
La somme entière de chaque ligne correspondante des deux matrices coïncide :
elle vaut cinq pour la première ligne et sept pour les cinq autres. Par
distributivité,
\[
 \sum_j\sum_i b_i(A_{\rm ter})_{ij}
 =\sum_i b_i\sum_j(A_{\rm ter})_{ij}
 =\sum_i b_i\sum_j(A_{\rm reg})_{ij}.
\]
La réduction modulo trois donne l'égalité des charges. L'identité des matrices \eqref{act21:eq:kdes-cambio-carta} se vérifie par substitution de ses trente-six entrées et conserve, en outre, le changement de coordonnées.
\end{proof}

L'égalité de la fonctionnelle totale de charge permet de comparer ses
représentations. Les transformations vectorielles restent identifiées
par leurs matrices et par la permutation explicite ; la preuve précédente
conserve cette information de coordonnées.

\subsubsection{Douze lecteurs comparatifs par instant}

On définit la phase temporelle
\[
 \theta(t)=(t+2)\bmod3.
\]
Cette écriture réalise également $(t-1)\bmod3$ pour $t=0$, où la phase
vaut deux. Pour chaque type de charge
$\nu\in\{\eta,\eta_{\rm ter}^\vee\}$, déplacement
$o\in\{0,1,2\}$ et orientation $s\in\{1,2\}$, on prend
\[
 \epsilon_\chi^{\nu,o,s}(t)=s
 \begin{cases}
 1,&\nu(B_\chi(t))=\theta(t)+o\pmod3,\\
 2,&\nu(B_\chi(t))\ne\theta(t)+o\pmod3.
 \end{cases}
\]
Le mot comparatif correspondant est
\begin{equation}
 C_{\nu,o,s}(t)=
 \sum_{\chi\in\{\pi,e,\varphi\}}
 \epsilon_\chi^{\nu,o,s}(t)B_\chi(t)
 \quad\text{en }\mathbb F_3^6.
 \label{act21:eq:kdes-comparativo}
\end{equation}
Les paramètres engendrent douze représentations étiquetées par ligne. Deux
étiquettes peuvent produire le même mot ; les témoins temporels et
les étiquettes restent disponibles dans le calendrier complet.

\subsubsection{Lecteur affine de phase et de charge}

Le lecteur affine emploie des coefficients
\[
 a_\chi^{\rm af}(t)=\eta_{\rm ter}^\vee(B_\chi(t))+\theta(t)
 \pmod3,
\]
et produit
\begin{equation}
 F(t)=\sum_{\chi\in\{\pi,e,\varphi\}}
 a_\chi^{\rm af}(t)B_\chi(t)
 \quad\text{en }\mathbb F_3^6.
 \label{act21:eq:kdes-afin}
\end{equation}
La règle comparative distingue la coïncidence et la différence de charges ; la règle affine additionne charge et phase avant de combiner les bandes. C'est cette différence opératoire qu'utilise la sélection hétérotypique.

\subsubsection{Incidence mot--résidu}

Soit $R$ la rotation cyclique des six positions d'une bande. À la
ligne temporelle $t$, on forme le répertoire résiduel
\[
 \mathcal R(t)=\{R^jv:
      v\in\{q(t),a(t),c(t),\bar w(t)\},\ 0\leq j<6\}.
\]
Les relations d'incidence temporelle sont
\begin{align}
 \mathcal A_{\rm cmp}
 &=\{(C_{\nu,o,s}(t),u):
      0\leq t<100,\ u\in\mathcal R(t),\ \nu,o,s\text{ admissibles}\},\\
 \mathcal A_{\rm af}
 &=\{(F(t),u):0\leq t<100,\ u\in\mathcal R(t)\}.
 \label{act21:eq:kdes-atlas}
\end{align}
Chaque couple relie un mot lu à un résidu de la même ligne.
Le temps se quantifie de manière commune dans les deux coordonnées du couple.
Dans le calendrier employé, les temps $0,\ldots,99$ sont distincts ;
identifier une même ligne et identifier un même temps sont, par conséquent,
des opérations équivalentes.

\begin{proposition}[Compression existentielle du répertoire temporel]
La fonction
\[
 M(v,u)=\mathbf1_{(v,u)\in\mathcal A_{\rm cmp}}
        +2\mathbf1_{(v,u)\in\mathcal A_{\rm af}}
        \in\{0,1,2,3\}
\]
conserve exactement l'existence de chaque type de lecture pour le couple  
$(v,u)$.
\end{proposition}
\begin{proof}
Le premier bit vaut un exactement lorsqu'il existe un temps, un champ
résiduel, une rotation et une étiquette comparative qui produisent le couple.
Le second bit exprime la condition correspondant au lecteur affine.
L'union par disjonction des bits conserve chaque condition d'existence
lors du parcours de toutes les lignes. Un couple peut admettre les deux types et, dans ce
cas, le masque vaut trois.
\end{proof}

La fonction $M$ est un lecteur d'existence. Les temps témoins, leurs multiplicités et les étiquettes individuelles demeurent dans \eqref{act21:eq:kdes-atlas} et dans le calendrier qui les produit. L'implémentation finie emploie $M$ pour décider de l'hétérotypie, tandis que l'archive généalogique conserve les données de provenance plus riches.

\subsection{Énumération de cylindres et candidats décimaux}

\subsubsection{Concaténation de soixante-douze et soixante-dix-huit trits}

Pour un ordonnancement $\tau=(s_1,\ldots,s_8)$ de
$\mathcal S_{\rm term}$, on définit
\[
 P_{72}(\tau)=[\Omega_{24}\Vert s_1\Vert\cdots\Vert s_8]_3.
\]
De manière récursive, on part de $p_0=[\Omega_{24}]_3$ et on calcule
$p_i=729p_{i-1}+[s_i]_3$ pour $1\leq i\leq8$. Cette récurrence
conserve les zéros initiaux de chaque bloc par la multiplication
explicite par $729$. Pour chaque position $b$ du panneau fonctionnel,
\[
 P_{78}(\tau,b)=729P_{72}(\tau)+b.
\]
Les longueurs sont respectivement $24+8\cdot6=72$ et $72+6=78$.
Le cylindre ternaire correspondant est
\begin{equation}
 I_{\tau,b}=
 \left[\frac{P_{78}(\tau,b)}{3^{78}},
       \frac{P_{78}(\tau,b)+1}{3^{78}}\right).
 \label{act21:eq:kdes-cilindro}
\end{equation}

\subsubsection{Intersection exacte avec la grille de douze triades}

La lecture décimale utilise douze triades, c'est-à-dire la grille $N/10^{36}$. On utilise les abréviations $D=10^{36}$ et $Q=3^{78}$.

\begin{lemma}[Extrémités entières du cylindre]
Pour un préfixe entier $p$ de longueur soixante-dix-huit, les numérateurs
des points de la grille qui appartiennent à son cylindre sont
\begin{equation}
 \left\{N\in\mathbb Z:
 \left\lceil\frac{pD}{Q}\right\rceil
 \leq N<
 \left\lceil\frac{(p+1)D}{Q}\right\rceil\right\}.
 \label{act21:eq:kdes-rejilla}
\end{equation}
Chaque cylindre contient au plus un point de cette grille.
\end{lemma}
\begin{proof}
La condition $N/D\in[p/Q,(p+1)/Q)$ est équivalente à
$pD/Q\leq N<(p+1)D/Q$. Pour $N$ entier, la première inégalité
est équivalente à $\lceil pD/Q\rceil\leq N$ ; la seconde est équivalente à
$N<\lceil(p+1)D/Q\rceil$. La longueur de l'intervalle en coordonnée
$N$ est $D/Q<1$, car $10^{36}<3^{78}$. Par conséquent, il peut contenir,
au maximum, un entier.
\end{proof}

Les plafonds sont évalués par division entière :
$\lceil a/Q\rceil=\lfloor(a+Q-1)/Q\rfloor$ pour $a\geq0$.
Toute l'énumération utilise des entiers exacts ; le choix des points de frontière est fixé par l'extrémité droite semi-ouverte.

On conserve d'abord la liste indexée d'incidences
\[
 \mathcal C_{\rm ind}=
 \{(\tau,j,N):N/D\in I_{\tau,\mathcal P_9(j)}\},
 \qquad1\leq j\leq9,
\]
et l'on obtient ensuite l'ensemble des numérateurs
\[
 \mathcal C=\{N:\exists\tau,j\ (\tau,j,N)\in\mathcal C_{\rm ind}\}.
\]
Le passage à $\mathcal C$ élimine des répétitions d'un numérateur, conservées
dans $\mathcal C_{\rm ind}$ avec leur origine.

\begin{lemma}[Récupération de l'ordre à partir du point du cylindre]
Si $(\tau,j,N)\in\mathcal C_{\rm ind}$, le bloc numéro $i$ du
mot de soixante-dix-huit trits se récupère par
\[
 \left\lfloor\frac{N\,729^i}{D}\right\rfloor\bmod729,
 \qquad1\leq i\leq13.
\]
En particulier, les blocs cinquième à douzième récupèrent $\tau$.
\end{lemma}
\begin{proof}
L'appartenance au cylindre signifie
$p\leq QN/D<p+1$, où $p=P_{78}(\tau,\mathcal P_9(j))$.
Par conséquent, $\lfloor QN/D\rfloor=p$. L'extraction des blocs
précédents coïncide avec la division successive du préfixe $p$ par
des puissances de $729$ : tronquer une expansion positionnelle avant d'extraire
un bloc conserve ce bloc. La formule écrite réalise
exactement cette extraction. Les quatre premiers blocs forment
$\Omega_{24}$ et les huit suivants sont l'ordonnancement $\tau$.
\end{proof}

\begin{proposition}[Recensement exact des candidats]
Pour $\Omega_{24}$, $\mathcal S_{\rm term}$ et $\mathcal P_9$
définis précédemment,
\[
 \#\operatorname{Ord}(\mathcal S_{\rm term})=40320,
 \qquad\#\mathcal C_{\rm ind}=21902,
 \qquad\#\mathcal C=19446.
\]
\end{proposition}
\begin{proof}
Les huit blocs sont distincts, donc leurs ordonnements sont les
$8!=40320$ permutations. Pour chaque permutation et chacune des
neuf positions du panneau, on calcule $p=P_{78}(\tau,b)$, on applique
les extrêmes \eqref{act21:eq:kdes-rejilla} et on ajoute son unique entier lorsque
l'intervalle entier n'est pas vide. La somme des cardinaux de ces
intervalles est $21902$. La réunion de leurs entiers contient $19446$
éléments distincts. L'évaluation exhaustive de ces opérations
est certifiée par les théorèmes
\texttt{indexed\_cylinder\_count} et \texttt{candidate\_count} du
sélecteur fini. L'algorithme réalise exactement l'énumération
décrite : permutations, neuf incidences du panneau, limites par
division entière et élimination des doublons.
\end{proof}

Les nombres $21902$ et $19446$ comptent des objets différents. Le premier
compte des incidences indexées qui contiennent un point décimal ; le second
compte des numérateurs distincts. Le nombre total de requêtes de cylindre
est $40320\cdot9=362880$, y compris les cylindres dont l'intersection avec
la grille est vide.

\subsection{Incidence exceptionnelle et condition d'hétérotypie}

\subsubsection{Face exceptionnelle produite par le panneau}

On élève une bande à douze coordonnées par le biais de
\[
 \Gamma(b)=(b,W(b))\in\mathbb F_3^{12}.
\]
Le support d'un mot est l'ensemble des positions de ses coordonnées non nulles, avec les indices humains $1,\ldots,12$. Les premières bandes régionales du panneau produisent
\begin{equation}
 \mathcal B=\operatorname{supp}\Gamma(B_\pi(0))
          \cap\operatorname{supp}\Gamma(B_e(0))
          =\{4,6,9,10\}.
 \label{act21:eq:kdes-cara}
\end{equation}
L'opération utilise les mots régionaux $103$ et $523$ et la matrice de Witt du sélecteur. La face se calcule avant de consulter les coordonnées des candidats.

Pour $N\in\mathcal C$, ses douze triades sont
\begin{equation}
 K_i(N)=\left\lfloor\frac{N}{1000^{12-i}}\right\rfloor\bmod1000,
 \qquad1\leq i\leq12.
 \label{act21:eq:kdes-triadas}
\end{equation}
Le support de couronne est
\[
 \mathcal H(N)=\{i:K_i(N)\geq729\},
 \qquad
 \mathcal C_{\mathcal B}=\{N\in\mathcal C:\mathcal H(N)=\mathcal B\}.
\]
La frontière $729=3^6$ relie la bande de six trits avec la triade décimale. L'évaluation exhaustive de l'égalité des supports donne
\begin{equation}
 \#\mathcal C_{\mathcal B}=79.
 \label{act21:eq:kdes-79}
\end{equation}
Le support conserve les positions, tandis que le registre complet
conserve aussi les douze grandeurs. L'étape suivante utilise cette
information positionnelle conjointement avec des relations entre grandeurs.

\subsubsection{Mots ternaires et différences orientées du candidat}

La $i$-ième bande ternaire du point $N/D$ s'obtient par
\begin{equation}
 T_i(N)=\left\lfloor\frac{N\,729^i}{D}\right\rfloor\bmod729,
 \qquad1\leq i\leq13.
 \label{act21:eq:kdes-tercandidato}
\end{equation}
L'arête orientée $a\to b$ a un résidu
\begin{equation}
 R_{a,b}(N)=(K_a(N)-K_b(N))\bmod729.
 \label{act21:eq:kdes-resarista}
\end{equation}
La soustraction s'effectue en $\mathbb Z$ avant la réduction. Pour cette
arête, on consulte dans les relations \eqref{act21:eq:kdes-atlas} le couple
\[
 \bigl(\operatorname{enc}_6^{-1}(T_b(N)),
  \operatorname{enc}_6^{-1}(R_{a,b}(N))\bigr).
\]
C'est donc une consultation conjointe sur le mot cible et sur la différence orientée des triades.

On écrit $\mathrm{Cmp}(a,b;N)$ lorsque le couple appartient à $\mathcal A_{\rm cmp}$ et $\mathrm{Af}(a,b;N)$ lorsqu'il appartient à $\mathcal A_{\rm af}$.

\subsubsection{Détermination de l'orbite transversale}

Le descripteur de vingt-quatre trits détermine les trois premières triades décimales avant les filtrations terminales. En effet, son cylindre satisfait l'inclusion exacte
\[
 \left[\frac{66241910521}{3^{24}},
       \frac{66241910522}{3^{24}}\right)
 \subseteq
 \left[\frac{234543140}{10^9},
       \frac{234543141}{10^9}\right).
\]
Les deux inégalités aux extrémités se vérifient par des produits croisés
entiers. Tout candidat partage donc les trois positions de
référence $1,2,3$ et leurs représentations $(234,543,140)$.

La translation de pas quatre sur $\mathbb Z/12\mathbb Z$ possède les
quatre orbites
\[
 (1,5,9),\qquad(2,6,10),\qquad(3,7,11),\qquad(4,8,12).
\]
Les trois premières contiennent respectivement l'une des positions de référence ; la quatrième est la seule disjointe de $\{1,2,3\}$. La forêt orientée des pas trois et quatre est fixée au moyen des arêtes
\[
 \begin{array}{lll}
 1\to4,&1\to5,&2\to6,\\
 3\to7,&4\to8,&5\to9,\\
 6\to10,&7\to11,&8\to12.
 \end{array}
\]
Sa première arête a un pas de trois et les autres ont un pas de quatre. Les deux
arêtes internes de l'orbite distinguée sont précisément
$4\to8$ et $8\to12$. Ainsi est fixé le lieu de la condition de
lecture par la géométrie des positions et par le préfixe initial,
avant d'examiner les candidats de la face exceptionnelle.

\subsubsection{Condition terminale d'hétérotypie}

La condition terminale est
\begin{equation}
 \begin{split}
 \mathcal T(N)={}&
 [\mathrm{Cmp}(4,8;N)\wedge\mathrm{Af}(8,12;N)]\\
 &\vee[\mathrm{Af}(4,8;N)\wedge\mathrm{Cmp}(8,12;N)].
 \end{split}
 \label{act21:eq:kdes-heterotipia}
\end{equation}
Les deux affectations de types aux deux arêtes sont admises. La
sélection exige qu'il existe une réalisation de types distincts ; les
étiquettes individuelles et leurs temps restent disponibles dans les
témoins d'incidence.

\begin{theorem}[Sélection terminale du registre de couplage]
Sur le domaine de candidats produit par
$\Omega_{24}$, $\mathcal S_{\rm term}$, $\mathcal P_9$ et le calendrier
régional de cent lignes, existe un unique numérateur qui satisfait
simultanément l'incidence exceptionnelle et l'hétérotypie :
\[
 \{N\in\mathcal C_{\mathcal B}:\mathcal T(N)\}
 =\{234543140729659824621058914794146601\}.
\]
Son registre de douze triades est
\begin{equation}
 K=(234,543,140,729,659,824,621,058,914,794,146,601).
 \label{act21:eq:kdes-kfinal}
\end{equation}
\end{theorem}
\begin{proof}
Le recensement précédent produit les $79$ éléments de
$\mathcal C_{\mathcal B}$. Pour chacun, on extrait ses douze
coordonnées par \eqref{act21:eq:kdes-triadas} ; on calcule les bandes de
destination $T_8,T_{12}$ et les résidus $R_{4,8},R_{8,12}$ ; on consulte
les deux masques du répertoire temporel et on applique la disjonction
\eqref{act21:eq:kdes-heterotipia}. La liste résultante contient exactement
l'entier indiqué, selon l'évaluation exhaustive certifiée par
\texttt{terminal\_selection}. L'égalité de liste avec un singleton
implique existence et unicité dans le domaine déclaré. L'extraction
base mille de cet entier donne \eqref{act21:eq:kdes-kfinal}.

La procédure de sélection reçoit les producteurs et les règles énumérées. L'entier final intervient dans l'énoncé de l'égalité vérifiée, après avoir exécuté la sélection. Cette position logique permet de comparer le résultat sans l'utiliser comme critère de choix.
\end{proof}

\subsubsection{Témoin arithmétique d'appartenance à un cylindre}

Une des incidences du numérateur sélectionné utilise la mise en ordre
\[
 (002001,020111,200110,121012,010122,001001,221110,222110)
\]
et la frontière $438=121020_3$. Pour elle,
\[
 P_{72}=5283881584882673136514102926090957,
\]
\[
 P_{78}=3851949675379468716518781033120308091.
\]
La formule des extrémités entières produit
\begin{align*}
 \left\lceil\frac{P_{78}10^{36}}{3^{78}}\right\rceil
 &=234543140729659824621058914794146601,\\
 \left\lceil\frac{(P_{78}+1)10^{36}}{3^{78}}\right\rceil
 &=234543140729659824621058914794146602.
\end{align*}
L'unique entier de l'intervalle est donc $N_K$. Les treize bandes ternaires de son point décimal sont
\[
 \begin{array}{rrrr}
 020022&222111&211201&021101\\
 002001&020111&200110&121012\\
 010122&001001&221110&222110\\
 121020&&&
 \end{array}
\]
Ce calcul exhibe un témoin concret d'appartenance au domaine énuméré. L'unicité du registre procède de l'examen exhaustif de tous les candidats, outre ce témoin particulier.

\subsection{Composition régionale et continuité des réalisations}

\subsubsection{Composition des générateurs régionaux}

Soit $\operatorname{Sel}(\Omega,\mathcal S,\mathcal P,\mathcal N)$ le
sélecteur défini par l'énumération et les deux filtres précédents.
Les producteurs régionaux satisfont les égalités
\[
 \Omega_{\rm reg}=\Omega_{24},\qquad
 \mathcal P_{\rm reg}=\mathcal P_9,\qquad
 \mathcal N_{\rm reg}=\mathcal N_{100},
\]
où $\mathcal N_{\rm reg}$ se calcule par extraction de bandes et
lecture de signatures. Par congruence de fonctions,
\begin{equation}
 \begin{split}
 &\operatorname{Sel}(\Omega_{\rm reg},\mathcal S_{\rm term},
                    \mathcal P_{\rm reg},\mathcal N_{\rm reg})\\
 &\hspace{10mm}=
 \operatorname{Sel}(\Omega_{24},\mathcal S_{\rm term},
                    \mathcal P_9,\mathcal N_{100})=\{N_K\}.
 \end{split}
 \label{act21:eq:kdes-composicion}
\end{equation}
L'égalité remplace trois interfaces documentaires par leur construction
régionale effective. Le répertoire terminal conserve la provenance
spécifiée dans \eqref{act21:eq:kdes-sterm}.

Le registre se forme alors par les douze divisions euclidiennes de
\eqref{act21:eq:kdes-triadas}. Chaque coordonnée appartient à
$\{0,\ldots,999\}$ et la longueur est douze par construction. L'
implémentation qui prend le premier élément de la liste sélectionnée
est justifiée par l'égalité avec le singleton : la valeur par défaut
d'une opération informatique sur des listes vides n'intervient jamais dans
ce résultat.

\subsubsection{Relation avec les incidences signées et l'inversion}

La construction présente fournit une sélection prospective dans un
domaine fini explicite. Le développement ultérieur du registre
dodécaphasique produit ses canaux orientés à partir de l'état
terminal enrichi et démontre une inversion intégrale par
Hadamard. Cette inversion récupère un registre à partir de canaux
compatibles ; ici, on sélectionne un registre à partir de ses
antécédents régionaux et incidentiels. L'égalité de ses résultats
permet de composer les deux réalisations en maintenant l'origine de chaque
opération.

L'unicité de sélection et l'unicité d'inversion sont des propriétés distinctes. La première quantifie les candidats construits dans ce chapitre. La seconde quantifie les registres qui partagent une signature compatible. Leur conjonction justifie que le registre sélectionné puisse être transporté vers la réalisation signée et être récupéré à partir d'elle, sans convertir une récupération rétrospective en producteur du registre.

\subsubsection{Relation entre les recensements terminaux}

La suite
\[
 40320\ \longrightarrow\ 21902\ \longrightarrow\ 19446
 \ \longrightarrow\ 79\ \longrightarrow\ 1
\]
enregistre successivement des ordonnancements, des incidences indexées avec la grille, des numérateurs distincts, des candidats de la face exceptionnelle et des survivants hétérotypiques. La notation à flèches indique l'ordre des opérations ; le deuxième nombre compte une classe d'objets différente de la première.

Le recensement des catalogues régionaux
$196\cdot197\cdot196\longrightarrow331\longrightarrow2
\longrightarrow1$, développé à l'endroit correspondant, part d'une
autre population : blocs régionaux de trois catalogues et conditions
terminales de Gram, d'occupation, de distance et d'orientation. Les deux recensements
conservent leurs domaines. Leur convergence vers une même représentation
s'exprime par l'identification de leurs lecteurs et de leurs sorties ;
les chiffres intermédiaires décrivent des opérations différentes et restent
documentés séparément.

\subsubsection{Composition régionale de la constante de structure fine}

La sortie $K$ intervient ensuite dans la composition régionale
ordonnée et dans la retenue en base mille. Son emploi ultérieur conserve les
douze positions et la frontière de lecture. Le couplage avec la racine
analytique et les représentations à profondeur arbitraire utilise cette
même sortie, avec les coordonnées régionales déjà engendrées et les
règles de la coordinatisation analytique. Cette dépendance prépare le chapitre
consacré aux deux réalisations corrélées de la constante de
structure fine.

La portée de ce chapitre est précise : production du descripteur et de ses antécédents régionaux, détermination du domaine fini, sélection unique sur ce domaine et composition matérielle avec les producteurs. La récurrence analytique de tous les degrés et la compatibilité des retenues infinies sont développées ensuite avec leurs démonstrations propres.

\subsection{Certification exécutable et traçabilité des opérations}

Le développement formel conserve la correspondance entre les
définitions mathématiques et leurs réalisations exécutables. Les preuves
des signatures et de l'horloge se trouvent dans
\texttt{N69RegionalSignature.lean} ; l'extraction des bandes et la
reproduction des cent lignes, dans
\texttt{GeneratedN69Rows.lean} ; le critère minimal et la construction
de $\Omega_{24}$, dans \texttt{RegionalW24.lean}. La compatibilité des
deux matrices de Witt est réunie dans
\texttt{WittReaderCompatibility.lean}.

Les lecteurs temporels, les rotations, les résidus orientés et la
condition hétérotypique sont formalisés dans
\texttt{TerminalReaders.lean}. L'énumération et les recensements sont
formalisés dans \texttt{Terminal}\allowbreak\texttt{Orbital}\allowbreak\texttt{Selection.lean}. Le panneau
régional se relie par \texttt{Regional}\allowbreak\texttt{Panel}\allowbreak\texttt{Bridge.lean}, et la
composition \eqref{act21:eq:kdes-composicion} s'exprime en
\texttt{Selected}\allowbreak\texttt{From}\allowbreak%
\texttt{Regional}\allowbreak\texttt{Inputs.lean}.

Les vérifications finies d'énumération utilisent l'évaluation native de Lean ; leurs rapports enregistrent \texttt{Lean.ofReduceBool}, ainsi que les axiomes logiques habituels de la bibliothèque. La reconstruction régionale du descripteur utilise des preuves du noyau et une réduction ordinaire. Cette distinction décrit les mécanismes de vérification réellement employés. Les reçus de compilation certifient ses modules et dépendances déclarés ; les preuves mathématiques conservent les domaines, producteurs et hypothèses qui ont été explicités dans le corps du chapitre.

\subsection{Recensement complet des cylindres et sélection de la frontière terminale}
\label{act21:censo:terminal-completo}

La lecture terminale utilise conjointement deux informations : la compatibilité des cylindres régionaux et un profil d'incidence avec des canaux et des positions étiquetés. Dans cette section, on construit les catalogues complets, on énumère leurs triplets et on démontre l'effet de chaque condition du lecteur. Le point de départ est constitué des sorties régionales déjà produites dans la section~\ref{sec:generacion}, et non de valeurs introduites pour définir le transport APP--TRIT--TPK. La réduction à une matrice appartient à une lecture ultérieure de ces sorties et conserve l'ordre des trois canaux.

Le résultat est un théorème de recensement à structure de départ explicite.
Le profil de Gram, des poids, des distances et des marges est déclaré comme donnée du
lecteur terminal. L'exhaustivité de l'énumération démontre ce que sélectionne
ce profil sur les catalogues fixés ; elle ne démontre pas, à elle seule, que le
profil soit une loi universelle ni qu'il puisse être reconstruit à partir du seul triplet
qui le satisfait finalement. Cette séparation permet d'intégrer la preuve
finie complète sans inverser sa généalogie.

\subsubsection{Des sorties régionales aux catalogues ternaires}

Soient \(x_\pi,x_e,x_\varphi\in[0,1)\) les parties fractionnaires des
trois coordonnées régionales déjà construites. On utilise leurs cylindres
exprimés à mille positions décimales,
\[
 J_\chi^{(1000)}=
 \left[\frac{d_\chi}{10^{1000}},
       \frac{d_\chi+1}{10^{1000}}\right),
 \qquad \chi\in\{\pi,e,\varphi\}.
\]
Les entiers \(d_\chi\) sont des lectures finies de sortie. La longueur mille
est suffisante pour les entiers qui suivent, comme on le vérifie par les
deux bornes de chaque intervalle ; ce n'est pas une profondeur supposée infinie.
Définissons
\begin{equation}
 p=30,\quad r=1050,\quad L=p+r=1080,\quad k=171,
 \qquad D=3^L,\quad Q=1000^k,
 \label{act21:censo:parametros}
\end{equation}
et
\[
 P_\chi=\lfloor3^{30}x_\chi\rfloor,\qquad
 T_\chi=\lfloor Qx_\chi\rfloor.
\]
Ici, \(k\) compte des triades décimales et ne désigne pas le registre dodécaphasique \(K\). On a \(1000^{171}\leq3^{1080}<1000^{172}\).

\begin{lemma}[Extraction stable d'un entier]
\label{act21:censo:estabilidad}
Pour \(d,m,h\in\mathbb Z\), avec \(m,h>0\), la valeur de
\(\lfloor mx\rfloor\) est constante en \([d/h,(d+1)/h)\) si et seulement si
\begin{equation}
 \left\lfloor\frac{md}{h}\right\rfloor
 =\left\lfloor\frac{m(d+1)-1}{h}\right\rfloor.
 \label{act21:censo:prueba-estabilidad}
\end{equation}
Lorsqu'elle est satisfaite, chacun des deux membres donne cette valeur.
\end{lemma}
\begin{proof}
Le minimum de \(mx\) est atteint à l'extrémité gauche. Son supremum
est \(m(d+1)/h\) et n'est pas atteint. Le plus grand entier que peut prendre
\(\lfloor mx\rfloor\) est
\(\lceil m(d+1)/h\rceil-1\), égal au second membre car
le numérateur et le dénominateur sont des entiers. L'égalité des extrémités
de l'intervalle de valeurs prouve l'affirmation.
\end{proof}

La preuve de stabilité se vérifie sur les trois canaux pour
\(m=3^{30},Q,3^{1080}\). Le dernier choix sera utilisé uniquement
pour une comparaison ultérieure, pas pour sélectionner le profil. Les
premières trente positions récupérées sont précisément les cinq
étapes de six trits du transport régional :
\begin{center}\small
\begin{tabular}{@{}c l r@{}}\toprule
Canal&Préfixe de trente trits&\(P_\chi\)\\\midrule
\(\pi\)&\texttt{010211012222010211002111110221}&29152671743887\\
\(e\)&\texttt{201101121221102011012222102011}&147887858824447\\
\(\varphi\)&\texttt{121200112202121020010210010200}&127247717616687\\\bottomrule
\end{tabular}
\end{center}
La coïncidence conserve la prolongation déjà construite ; la lecture décimale
n'est pas utilisée pour remplacer rétroactivement ses matrices de transport.

Un canal étant fixé, une queue de \(r\) trits se représente par
\(0\leq R<3^r\). Son cylindre et le cylindre des \(k\) premières
triades sont
\begin{equation}
 I_{\chi,R}=\left[\frac{P_\chi3^r+R}{D},
                       \frac{P_\chi3^r+R+1}{D}\right),\qquad
 J_{\chi,T}=\left[\frac{T_\chi}{Q},\frac{T_\chi+1}{Q}\right).
 \label{act21:censo:dos-cilindros}
\end{equation}
Le catalogue inclut toutes les queues pour lesquelles ces deux cylindres
ont une intersection non vide.

\begin{proposition}[Catalogue par intersection exacte]
\label{act21:censo:catalogos}
Les queues compatibles sont les entiers de l'intervalle \([a_\chi,b_\chi]\),
avec
\begin{align}
 a_\chi&=\max\left(0,
       \left\lfloor\frac{T_\chi D}{Q}\right\rfloor-P_\chi3^r\right),
       \label{act21:censo:extremo-a}\\
 b_\chi&=\min\left(3^r-1,
       \left\lfloor\frac{(T_\chi+1)D-1}{Q}\right\rfloor-P_\chi3^r\right).
       \label{act21:censo:extremo-b}
\end{align}
Pour les cylindres régionaux fixés, les troncatures n'interviennent pas et l'on obtient
\begin{equation}
\begin{array}{c|r|r|r}
 \chi&a_\chi\bmod729&b_\chi\bmod729&b_\chi-a_\chi+1\\\hline
 \pi&67&262&196\\
 e&584&51&197\\
 \varphi&117&312&196
\end{array}
\label{act21:censo:tabla-catalogos}
\end{equation}
\end{proposition}
\begin{proof}
Écrivons \(N=P_\chi3^r+R\). Deux intervalles semi-ouverts de
\eqref{act21:censo:dos-cilindros} se coupent exactement lorsque
\[
 NQ<(T_\chi+1)D,\qquad (N+1)Q>T_\chi D.
\]
La deuxième inégalité est équivalente à  
\(N\geq\lfloor T_\chi D/Q\rfloor\) ; la première, à  
\(N\leq\lfloor((T_\chi+1)D-1)/Q\rfloor\). Soustraire  
\(P_\chi3^r\) et conserver \(0\leq R<3^r\) prouve les formules.  
L'évaluation par divisions entières des \(T_\chi\) extraits par  
\eqref{act21:censo:prueba-estabilidad} donne le tableau. Comme \(r\geq6\),  
\(P_\chi3^r\equiv0\pmod{729}\), de sorte que le reste de la première  
extrémité s'obtient également directement de  
\(\lfloor T_\chi D/Q\rfloor\bmod729\).
\end{proof}

Les six derniers trits d'une queue forment le mot
\(\nu_6(R\bmod729)\), où \(\nu_6\) est l'écriture ternaire
de longueur six, y compris ses zéros initiaux. D'après le tableau, chaque
catalogue possède moins de \(729\) éléments consécutifs, de sorte que
cette réduction y est injective. Ainsi, les trois catalogues de frontière
sont entièrement décrits, sans choisir les triplets finaux :
\begin{equation}
 \begin{split}
 \mathcal B_\pi&=\{\nu_6(67+i):0\leq i<196\},\\
 \mathcal B_e&=\{\nu_6((584+j)\bmod729):0\leq j<197\},\\
 \mathcal B_\varphi&=\{\nu_6(117+\ell):0\leq\ell<196\}.
 \end{split}
 \label{act21:censo:fronteras-explicitas}
\end{equation}
Exiger que l'extrémité gauche de \(I_{\chi,R}\) appartienne à
\(J_{\chi,T}\) serait une autre condition : elle perdrait le premier cylindre
qui coupe la frontière et donnerait \(195,196,195\). Le recensement utilise l'intersection des cylindres complets, conformément à sa définition.

\subsubsection{Profil déclaré et indicateurs de compatibilité}

Pour \(u,v\in\mathbb F_3^6\), soient
\(g(u,v)=\sum u_iv_i\in\mathbb F_3\),
\(n(u)=g(u,u)\), \(w(u)=\#\{i:u_i\ne0\}\) et
\(h(u,v)=\#\{i:u_i\ne v_i\}\). Les deux dernières fonctions prennent
des valeurs entières ; les deux premières, des valeurs dans \(\mathbb F_3\).
Le profil local du lecteur, en ordre \((\pi,e,\varphi)\), est
\begin{equation}
 \begin{split}
 (g_{\pi e},g_{\pi\varphi},g_{e\varphi})&=(2,2,0),\qquad
 (n_\pi,n_e,n_\varphi)=(0,0,0),\\
 (w_\pi,w_e,w_\varphi)&=(3,3,3),\qquad
 (h_{\pi e},h_{\pi\varphi},h_{e\varphi})=(2,5,3).
 \end{split}
 \label{act21:censo:perfil-local}
\end{equation}
Les mots sont ordonnés comme dans \eqref{act21:censo:fronteras-explicitas}.
Nous utiliserons \([E]\in\{0,1\}\) pour l'indicateur d'une condition \(E\).
Pour chaque paire de canaux \(\chi,\psi\), la matrice rectangulaire
\(A_{\chi\psi}\) a l'entrée \([g(u,v)=g_{\chi\psi}]\).
La matrice \(A^{h}_{\chi\psi}\) multiplie cette entrée par
\([h(u,v)=h_{\chi\psi}]\). Soient \(D_\chi^n\) la matrice diagonale
de \([n(u)=0]\), et \(D_\chi^{nw}\) celle de
\([n(u)=0][w(u)=3]\).

\begin{theorem}[Énumération exhaustive par sommes finies]
\label{act21:censo:teorema-conteo}
Le produit des trois catalogues contient \(7\,567\,952\) triplets.
Les conditions successives de \eqref{act21:censo:perfil-local} laissent
\begin{equation}
 7\,567\,952\longrightarrow275\,794\longrightarrow6235
 \longrightarrow3127\longrightarrow331.
 \label{act21:censo:cadena-local}
\end{equation}
Ces quatre comptages sont, respectivement,
\begin{align}
 N_g&=\operatorname{tr}(A_{\pi e}A_{e\varphi}A_{\varphi\pi}),
 \nonumber\\
 N_{gn}&=\operatorname{tr}(D_\pi^n A_{\pi e}
             D_e^n A_{e\varphi}D_\varphi^n A_{\varphi\pi}),
 \nonumber\\
 N_{gnw}&=\operatorname{tr}(D_\pi^{nw} A_{\pi e}
             D_e^{nw} A_{e\varphi}D_\varphi^{nw} A_{\varphi\pi}),
 \nonumber\\
 N_{gnwh}&=\operatorname{tr}(D_\pi^{nw} A^h_{\pi e}
             D_e^{nw} A^h_{e\varphi}D_\varphi^{nw} A^h_{\varphi\pi}).
 \label{act21:censo:trazas}
\end{align}
Tous les produits et traces de cette formule sont calculés dans les entiers,  
non dans le corps de trois éléments.
\end{theorem}
\begin{proof}
La taille du produit est \(196\cdot197\cdot196\). En développant
une trace de \eqref{act21:censo:trazas}, ses indices parcourent exactement un
triplet de mots. Le terme vaut un s'il satisfait toutes les conditions
représentées et zéro sinon. Chaque triplet est compté une fois.
Les matrices diagonales ajoutent les conditions individuelles ; les facteurs
rectangulaires ajoutent les conditions par paires. Cela démontre que les
traces sont précisément les cardinaux annoncés.

Pour expliciter son évaluation, une fois fixés les indices \(i,j\) des deux premiers canaux, on forme les ensembles d'indices \(\ell\) du troisième qui satisfont chaque condition par paires. On intersecte leurs indicateurs, on ajoute les indicateurs individuels pertinents et on additionne leur cardinal. La somme finale parcourt \(0\leq i<196\), \(0\leq j<197\). La substitution des mots de \eqref{act21:censo:fronteras-explicitas} et du profil \eqref{act21:censo:perfil-local} donne, dans cet ordre, \(275794,6235,3127,331\). Le supplément enregistre les quatre contributions de chacun des \(196\) indices \(i\) et les \(331\) triplets complets ; son programme évalue ces sommes avec des indicateurs binaires exacts et vérifie scalairement chaque survivant. Aucun échantillon ni aucune tolérance n'intervient.
\end{proof}

\subsubsection{Incidence étiquetée et charge d'orientation}

Pour un triplet survivant, soit \(B\in\mathbb F_3^{3\times6}\) sa
matrice de lignes ordonnées. Le second profil exige
\begin{equation}
 \operatorname{rank}_{\mathbb F_3}B=3,\qquad
 w_{\mathrm{col}}(B)=(1,1,2,2,3,0),\qquad
 q_{\partial}(B)=(1,2,2,1,2,0),
 \label{act21:censo:perfil-columnas}
\end{equation}
où \(w_{\mathrm{col}}\) compte les éléments non nuls de chaque
colonne et \(q_{\partial}\) additionne ses entrées modulo trois. Il s'agit d'un profil
de colonnes étiquetées ; les permuter sans transporter le profil définit un autre
lecteur. Notons \(\mathcal F\) l'ensemble des \(331\) triplets
du théorème précédent.

\begin{proposition}[Réduction du recensement au couple terminal]
\label{act21:censo:dos-matrices}
Les sommes d'indicateurs sur \(\mathcal F\), en exigeant successivement
le rang, les poids de colonne et les sommes de colonne, sont
\begin{equation}
 \sum_{B\in\mathcal F}[\operatorname{rank}B=3]=331,
 \qquad N_{\mathrm{rango,pesos}}=18,
 \qquad N_{\mathrm{rango,pesos,sumas}}=2.
 \label{act21:censo:conteo-columnas}
\end{equation}
Avec des indices de catalogue commençant à un, les deux triplets finaux sont
\((120,197,157)\) et \((129,197,148)\), et leurs matrices sont
\[
 B_0=\begin{pmatrix}0&2&0&2&2&0\\0&0&1&2&2&0\\1&0&1&0&1&0\end{pmatrix},
 \qquad
 B_1=\begin{pmatrix}0&2&1&0&2&0\\0&0&1&2&2&0\\1&0&0&2&1&0\end{pmatrix}.
\]
\end{proposition}
\begin{proof}
Dans chaque triplet de l'ensemble déjà énuméré, on calcule le rang par élimination en \(\mathbb F_3\), on compte les éléments non nuls des six colonnes et on les additionne modulo trois. Les produits des indicateurs de \eqref{act21:censo:perfil-columnas} donnent \eqref{act21:censo:conteo-columnas}.
La substitution des deux indices finaux en
\eqref{act21:censo:fronteras-explicitas} donne les matrices exposées. La liste
complète des survivants et les filtres exacts du supplément permettent
de reproduire l'élimination sans supposer ces deux matrices comme domaine.
\end{proof}

L'orientation \(\sigma=(1,1,-1)=(1,1,2)\in\mathbb F_3^3\)
sélectionne la droite
\(\langle\sigma\rangle=\{(0,0,0),(1,1,2),(2,2,1)\}\).
Les charges des lignes de \(B_0,B_1\) sont \((0,2,0)\) et \((2,2,1)\) ;
seule la seconde appartient à cette droite. On retrouve ainsi la dernière étape
du sélecteur de la section~\ref{act21:censo:dos-matrices}, désormais précédée de
la construction et de l'énumération de son domaine complet. La lecture stable
de \(\lfloor3^{1080}x_\chi\rfloor\), effectuée après les filtres,
donne les frontières \texttt{021020}, \texttt{001220}, \texttt{100210},
avec les mêmes rangs \((129,197,148)\). Cette comparaison ultérieure ne définit
pas le profil et ne participe pas aux sommes d'indicatrices.

\subsubsection{Information conservée et portée du résultat}

Le procédé conserve le préfixe de trente trits, la queue compatible, l'ordre régional, la position de chaque trit et l'orientation. La réduction modulo \(729\) conserve ici l'identité de chaque candidat car l'intervalle des queues a une longueur inférieure à \(729\) ; cette injectivité n'est pas affirmée pour une histoire arbitraire. Le recensement de \(331\) démontre également que le profil local isolé ne sélectionne pas une seule matrice : les marges étiquetées et la charge apportent des conditions effectives supplémentaires.

On a donc récupéré une énumération finie intégrale et son sélecteur sur des cylindres de sortie, avec un profil de lecture explicite. La construction du registre signé utilise en outre son lecteur de mémoire et ses canaux, comme cela est spécifié ci-dessous ; la matrice visible ne se confond pas avec ce registre par égalité du nombre de composants. Ce compte ne démontre pas non plus l'universalité du profil, ne génère pas à nouveau les coordonnées régionales et ne décide pas d'une affirmation sur alpha ou sur le continu complet.

\clearpage
% END NUCLEO_ADITIVO_20260921_SELECCION
\section{Registre dodécaphasique et constante rationnelle de phase}
\label{sec:k-registro}

La génération des coordonnées régionales n'épuise pas l'information de l'état APP--TRIT--TPK. Une lecture archimédienne détermine une valeur au moyen de cylindres compatibles ; une autre lecture conserve l'opposition des feuillets, l'orientation de la route et le registre des incidences traversées. Le registre dodécaphasique appartient à cette seconde lecture avant d'intervenir dans la clôture d'alpha. Il n'est donc pas défini en soustrayant une constante physique à une combinaison de valeurs préalablement connues. Il est reconstruit à partir d'une coordonnée signée du même état arithmétique avec mémoire de trajectoire qui produit les coordonnées régionales.

La distinction est nécessaire même lorsque deux objets possèdent douze
coordonnées. Un mot de douze blocs du catalogue, une liste de douze
événements orientés, un vecteur entier signé et un registre dodécaphasique en base mille sont des
objets distincts. Leur comparaison requiert les applications qui les
relient. L'égalité de longueur ne fournit pas ces applications, et la
perte des signes n'est pas une simple simplification de notation.

Ce chapitre maintient la délimitation causale fixée dans les chapitres précédents :
les neuf positions APP des axes horizontal et vertical produisent les
feuillets arithmétiques ; TRIT type le régime et l'orientation ; TPK sélectionne,
transporte, met à jour et prolonge, en conservant le résidu, le quotient, la retenue,
le feuillet, la route, la frontière et la mémoire. À partir de la structure discrète conjointe du
continu s'obtiennent les coordonnées utilisées ici. La chaîne
\(w_6\to w_{12}\to w_{18}\to w_{24}\to w_{30}\to R_{36}\to G_9\)
n'est pas remplacée par un calendrier fini. L'ordre
\(U_t\to\Gamma_9\to K_{\mathrm{ph}}\) n'est pas non plus inversé : la dernière application est une
lecture à mémoire réduite, et non le générateur du registre signé.

\input{sections/revision_k.tex}

\subsection{Extension cyclique et graduation de mémoire}
\label{subsec:k-calendario}

Le calendrier observable contient cent huit positions, organisées en douze fenêtres nonadiques. Avec des indices positifs, soient
\[
 E_m=\{9(m-1)+1,\ldots,9m\},\qquad 1\leq m\leq12.
\]
Le support ordonné \(E_1\sqcup\cdots\sqcup E_{12}\) conserve la position et l'
appartenance à une fenêtre. Le groupe \(C_{108}=\mathbb Z/108\mathbb Z\)
conserve, en revanche, la loi de translation cyclique. Pour le représentant
\(0\leq\widetilde\theta<108\), les coordonnées
\[
 \beta(\theta)=\left[\left\lfloor\frac{\widetilde\theta}{9}\right\rfloor\right]_{12},
 \qquad r(\theta)=1+(\widetilde\theta\bmod9)
\]
séparent le secteur dodécaphasique et la position nonadique intérieure. Aucune d'elles ne conserve à elle seule le nombre total de tours.

\begin{proposition}[La retenue du calendrier]
\label{prop:k-extension}
Les applications
\[
 \iota:C_{12}\longrightarrow C_{108},\quad[b]\longmapsto[9b],
 \qquad p:C_{108}\longrightarrow C_9,\quad[\theta]\longmapsto[\theta]_9
\]
forment une suite exacte non scindée
\[
 0\longrightarrow C_{12}\xrightarrow{\iota}C_{108}
 \xrightarrow{p}C_9\longrightarrow0.
\]
Pour la section ensembliste qui choisit les représentants \(0,\ldots,8\), le défaut d'additivité est la retenue
\[
 c(r,r')=\left[\left\lfloor\frac{\widetilde r+\widetilde r'}9\right\rfloor\right]_{12}.
\]
\end{proposition}

\begin{proof}
Le noyau de \(p\) est formé des multiples de neuf et coïncide avec
l'image de \(\iota\). La réduction est surjective et la multiplication
par neuf est injective sur le domaine indiqué. Si la suite était
scindée, le groupe central serait isomorphe à \(C_{12}\times C_9\), dont l'
exposant est \(\operatorname{mcm}(12,9)=36\) ; l'exposant de
\(C_{108}\) est 108. La contradiction prouve la non-scission. Enfin,
la division euclidienne de \(\widetilde r+\widetilde r'\) par neuf donne la
retenue écrite et vérifie l'identité de la section.
\end{proof}

Le résultat explique une différence qui réapparaîtra dans alpha : revenir à la
phase n'équivaut pas à revenir à l'état. Si \(q_{\mathrm{mem}}\) compte les
tours nonadiques, le calendrier ne conserve que
\(\beta=[q_{\mathrm{mem}}]_{12}\). Après cent huit pas,
\(\beta\) revient, mais \(q_{\mathrm{mem}}\) augmente de douze. Le
registre conserve en outre le cylindre, le livre d'incidences et la mémoire
verticale. La clôture d'une coordinatisation finie ne transforme pas cet historique en un
cycle sans mémoire.

\subsection{Le registre orienté et la coordinatisation intégrale \texorpdfstring{\(1+5+6\)}{1+5+6}}
\label{subsec:k-registro-integral}

Pour une histoire produite par TPK, soient \(\tau_m\) la sous-route sur
\(E_m\), \(\sigma_m\) son orientation, \(\kappa_m\) la retenue de
frontière et \(\Lambda_m\) le livre local des incidences. Le registre est
\[
 \mathcal R_{12}(x)=
 \bigl((E_m,\tau_m,\sigma_m,\kappa_m,\Lambda_m)\bigr)_{m=1}^{12}.
\]
Son domaine est l'espace des histoires arithmétiques avec mémoire de trajectoire, et non l'ensemble des
étiquettes de phase. La projection sur le calendrier oublie précisément certaines
des données dont le lecteur signé a besoin.

Une coordinatisation initiale du registre peut s’écrire au moyen d’événements
orientés. La construction de référence~(sous-section~\ref{subsec:k-registro-integral}) fixe l’ensemble des codes d’événement
\(\{2,4,6,8\}\) et l’orientation sectorielle
\[
 \Sigma_{12}=(+,+,+,-,-,-,+,+,+,-,-,-).
\]
Si \(d_t\) est le code de direction au pas \(t\), on obtient
\[
 e_i=\Sigma_{12}(i+1)
 \#\{t\in E_{i+1}:d_t\in\{2,4,6,8\}\},
 \qquad0\leq i<12.
\]
La formule maintient distincts le recensement non négatif des événements et
l'orientation. Elle n'identifie pas une direction cardinale à un signe TRIT. Le
vecteur \(e\), dont les composantes sont des recensements locaux signés, n'est pas non plus
encore le vecteur \(U\) qui apparaîtra plus bas.

Les positions opposées déterminent, pour \(0\leq i<6\),
\[
 p_i=e_i+e_{i+6},\qquad a_i=e_i-e_{i+6}.
\]
Avec
\[
 s_C=\sum_{i=0}^{5}p_i,\qquad M_i=p_i-p_5\quad(0\leq i<5),
\]
on définit la coordinatisation intégrale
\[
 \mathcal L_{12}(e)
 =(s_C,M_0,\ldots,M_4,a_0,\ldots,a_5).
\]
La décomposition \(1+5+6\) n'est pas une répartition arbitraire des coordonnées :
une coordonnée conserve la charge, cinq comparent les paires par rapport à une
paire de référence et six conservent l'opposition des demi-couronnes.

\begin{lemma}[Inversion de la coordinatisation intégrale]
\label{lem:k-carta-integral}
Un tuple entier \((s_C,M_0,\ldots,M_4,a_0,\ldots,a_5)\) appartient à l'image de \(\mathcal L_{12}\) si et seulement si
\[
 s_C-\sum_{i=0}^{4}M_i\equiv0\pmod6,
 \qquad p_i\equiv a_i\pmod2\quad(0\leq i<6),
\]
où
\[
 p_5=\frac{s_C-\sum_{i=0}^{4}M_i}{6},\qquad p_i=p_5+M_i\quad(i<5).
\]
Dans ce cas, sa préimage est unique et est donnée par
\[
 e_i=\frac{p_i+a_i}{2},\qquad e_{i+6}=\frac{p_i-a_i}{2}.
\]
\end{lemma}

\begin{proof}
La somme des cinq marges est \(s_C-6p_5\), d'où l'on obtient la
division par six. Une fois les \(p_i\) reconstruits, chaque paire d'équations
\(p_i=e_i+e_{i+6}\), \(a_i=e_i-e_{i+6}\) possède la solution indiquée.
Cette solution est entière exactement sous la congruence de parité. Toutes
les opérations s'inversent, de sorte qu'aucune liberté résiduelle ne subsiste.
\end{proof}

Cette réversibilité montre quelle information est conservée ; elle n'autorise pas à
supprimer le reste du livre. Le registre complet possède un domaine plus vaste
que cette coordinatisation des recensements. L'extraction des canaux qui alimentent
\(U\) utilise ce registre complet et ne se déduit pas de la liste
\((e_i)\) par une identification de noms.

% Propuesta separada; no modifica el bloque común del artículo I.
% Procedencia: artículo II, revisión 05, registro_incidencias.tex y
% registro_imagen_integral.tex; aclaraciones de índices de revisiones 06/07.
% Inserción prevista inmediatamente antes de subsec:k-terminal.
\subsection{Évaluation des incidences et détermination des coordonnées transversales}
\label{iii:r25:incidencias}

Le registre dodécaphasique conserve une histoire de visites et de traces
orientées. Ses coordonnées transversales s’obtiennent au moyen d’une
composition d’opérations : évaluation arithmétique des visites,
dénombrement par fenêtres, normalisation décimale et différences cycliques.
La reconstruction du registre à partir de ces différences constitue
une autre opération, dont l’existence et l’unicité seront démontrées ensuite.
Cette séparation permet de préciser tant les données utilisées par
l’extracteur que l’information qu’il transmet à la lecture harmonique.

Soit $\mathscr L$ un registre d’incidences produit par la sélection,
le transport et la mise à jour d’une histoire APP--TRIT--TPK.
Chaque visite conserve son chemin, sa position APP, son feuillet, son instant,
sa multiplicité et sa classe de frontière ; chaque trace conserve ses membres,
son orientation, sa multiplicité et sa longueur dans l’unité déclarée par son
lecteur. Les fenêtres sont
\[
 E_m=\{9(m-1)+1,\ldots,9m\},\qquad 1\leq m\leq12.
\]
L’évaluation du feuillet de somme ou de produit détermine l’entier
$n(v)$ de chaque visite et sa division positive
\[
 n(v)=r(v)+9q(v),\qquad 1\leq r(v)\leq9.
\]
Le résidu se calcule sur cette évaluation. La classe de frontière d’une
visite de résidu $9$ est conservée au moyen de
$\epsilon_\partial(v)=1$ dans la couronne et de $-1$ à l’intérieur.

La présentation d’incidence du registre utilise les trois dénombrements
\begin{align}
 A_m&=\sum_{v\in E_m}w(v)
       \mathbf1_{\{1,3,5,7\}}(r(v)),\label{iii:r25:recuento-a}\\
 C_m&=\sum_{j\geq0}
       \bigl(\nu^-_{120,m}(j)-\nu^+_{120,m}(j)\bigr),
       \label{iii:r25:recuento-c}\\
 V_m&=\sum_{v\in E_m}w(v)
       \mathbf1_{\{9\}}(r(v))\epsilon_\partial(v).
       \label{iii:r25:recuento-v}
\end{align}
Ici, $w(v)$ est la multiplicité d’incidence et
$\nu^\pm_{120,m}(j)$ compte les traces orientées du canal indiqué
de longueur réduite $120(1+3j)$, selon l’énumération du registre.
La famille de traces et leur multiplicité sont des données de l’évaluation,
avec la règle de production qui leur correspond. La somme entière nécessite
un support fini ou une construction qui détermine son sens.
La longueur réduite d’une trace et le nombre d’instants de la
fenêtre sont des grandeurs typées différentes ; l’unité du lecteur
spécifie leur relation. Les symboles $A_m,C_m,V_m$ désignent ici
des dénombrements entiers, distincts de la paire angulaire et de l’opérateur constitutif.

Soit $[x]_{10}\in\{0,\ldots,9\}$ le représentant euclidien.
La normalisation conserve également le quotient dans
\[
 x=[x]_{10}+10\lfloor x/10\rfloor.
\]
On forme le bloc
\begin{equation}
 z_m=100[A_m]_{10}+10[C_m]_{10}+[V_m]_{10},
 \qquad 0\leq z_m\leq999.
 \label{iii:r25:bloque}
\end{equation}
Avec des indices cycliques modulo $12$, soient
\[
 (D_a z)_m=z_m-z_{m+a},\qquad a\in\{3,4\},
 \qquad Q(z)=\sum_{m=1}^{12}z_m.
\]
La composition d’incidence s’exprime par
\begin{equation}
 \mathcal E_{\rm cnt}(\mathscr L)
 =\mathcal D(z):=(D_3z,D_4z,Q(z))\in\mathbb Z^{25}.
 \label{iii:r25:composicion}
\end{equation}
Ses arguments sont les visites et les traces avec leurs règles de dénombrement.
L’évaluation ne reçoit ni $K$, ni $U$, ni $\alpha$, ni les coordonnées
transversales attendues. Pour l’identifier à
$\mathcal E_{108}^{90,120}$, on conserve en outre la composition avec
la famille et les incidences effectives de l’état terminal.
L’inversibilité de $\mathcal D$ détermine une reconstruction ultérieure ;
à elle seule, elle ne sélectionne pas cette famille d’incidences.

\begin{proposition}[Information résiduelle déterminée par l’extracteur]
\label{iii:r25:residuos}
Soient $N,N'\in\mathbb Z^{12\times3}$ deux registres de dénombrements,
ordonnés par $(A,C,V)$. Pour les opérations
\eqref{iii:r25:bloque}--\eqref{iii:r25:composicion},
\[
 \mathcal E_{\rm cnt}(N)=\mathcal E_{\rm cnt}(N')
 \quad\Longleftrightarrow\quad N-N'\in10\mathbb Z^{36}.
\]
Les $36$ résidus sectoriels déterminent exactement la sortie de
$25$ coordonnées. Les quotients et les provenances constituent
une information supplémentaire du registre.
\end{proposition}
\begin{proof}
L’égalité modulo $10$ produit les mêmes blocs et leurs mêmes
coordonnées transversales. Réciproquement, l’égalité des
coordonnées donne $D_3(z-z')=D_4(z-z')=0$. Les pas $3$ et $4$
engendrent le cycle de longueur $12$, donc $z-z'$ est constant.
L’égalité de $Q$ annule cette constante. Le développement de chaque entier
de $0$ à $999$ en trois chiffres décimaux est unique ; on retrouve ainsi
les trois résidus de chaque fenêtre.
\end{proof}

La proposition est valable sur tout le domaine entier indiqué. Elle ne
postule pas que les $36$ classes soient indépendantes sur le sous-ensemble
des registres admissibles produit par TPK.

\subsection{Image entière et inversion des coordonnées transversales}
\label{iii:r25:imagen-seccion}

Pour exprimer l’inversion, on réindexe les blocs par
$i\in\mathbb Z/12\mathbb Z$, de $0$ à $11$. Soient
$(Sz)_i=z_{i+1}$, $D_3=I-S^3$ et $D_4=I-S^4$.
Étant donné un triplet entier $(b,c,q)$, définissons
\begin{equation}
 w_i=c_i-b_{i+1},\qquad
 P_0=0,\quad P_i=\sum_{j=0}^{i-1}w_j\ (1\leq i\leq12),
 \qquad T(w)=\sum_{i=0}^{11}P_i.
 \label{iii:r25:incrementos}
\end{equation}

\begin{theorem}[Caractérisation de l’image entière]
\label{iii:r25:imagen}
Un triplet $(b,c,q)\in\mathbb Z^{25}$ appartient à
$\operatorname{im}\mathcal D$ si et seulement si
\begin{align}
 b_i&=w_i+w_{i+1}+w_{i+2},\qquad 0\leq i<12,
       \label{iii:r25:relaciones}\\
 \sum_{i=0}^{11}w_i&=0,\label{iii:r25:cierre}\\
 q+T(w)&\equiv0\pmod {12}.\label{iii:r25:congruencia}
\end{align}
Son unique préimage entière est
\begin{equation}
 z_i=\frac{q+T(w)}{12}-P_i,\qquad0\leq i<12.
 \label{iii:r25:inversa}
\end{equation}
Les douze relations \eqref{iii:r25:relaciones} et la relation
\eqref{iii:r25:cierre} sont linéairement indépendantes sur
$\mathbb Q$.
\end{theorem}
\begin{proof}
Pour une sortie de $\mathcal D$,
\[
 c_i-b_{i+1}
 =(z_i-z_{i+4})-(z_{i+1}-z_{i+4})=z_i-z_{i+1}.
\]
Trois accroissements consécutifs ont pour somme $b_i$ et leur somme cyclique
totale est nulle. De plus, $P_i=z_0-z_i$, d’où
$q=12z_0-T(w)$. On obtient les conditions et la formule inverse.
Réciproquement, la congruence rend entières les coordonnées de
\eqref{iii:r25:inversa} et la clôture permet de les prolonger
cycliquement. Leurs différences adjacentes sont $w_i$, donc $D_3z=b$
et
\[
 (D_4z)_i=w_i+w_{i+1}+w_{i+2}+w_{i+3}
 =w_i+b_{i+1}=c_i.
\]
La somme donne $Q(z)=q$. La formule détermine également l’unicité.
Enfin, le changement entier inversible
$(b,c)\mapsto(b,w=c-Sb)$ transforme les douze premières relations
en $b=(I+S+S^2)w$ : chacune isole une coordonnée différente de
$b$. La clôture $\sum_iw_i=0$ ajoute une relation indépendante.
\end{proof}

L’image rationnelle a pour dimension $12$. La condition entière ajoute
la congruence d’ordre $12$ de \eqref{iii:r25:congruencia}.
Les onze accroissements $w_0,\ldots,w_{10}$ et la charge $q$ sont
des coordonnées suffisantes, avec $w_{11}=-\sum_{i=0}^{10}w_i$.
Les $25$ coordonnées conservent ces $12$ degrés de liberté
au moyen de $13$ relations linéaires. Si le triplet provient en outre de
\eqref{iii:r25:bloque}, l’inverse appartient à
$\{0,\ldots,999\}^{12}$ et retrouve ses trois résidus par fenêtre.

\subsubsection{Concordance avec la transformation harmonique}

La transformation $\mathcal H_{12}$ agit au moyen de $H_4$ sur les
orbites $(r,r+3,r+6,r+9)$, $r=0,1,2$, où
\[
 H_4=\begin{pmatrix}
 1&1&1&1\\1&1&-1&-1\\1&-1&1&-1\\1&-1&-1&1
 \end{pmatrix},\qquad H_4^2=4I_4.
\]
Pour un triplet compatible, soient
\[
 B_r=\sum_{j=0}^3c_{r+3j},\qquad
 a_0=\frac{q+2B_0+B_1}{3},\quad
 a_1=a_0-B_0,\quad a_2=a_1-B_1.
\]
Le lecteur harmonique a pour blocs
\[
 (\Pi_H(b,c,q))^{(r)}
 =(a_r,b_{r+3}-b_{r+9},b_r+b_{r+6},b_r-b_{r+6}).
\]

\begin{proposition}[Concordance des reconstructions]
\label{iii:r25:hadamard}
Sur $\operatorname{im}\mathcal D$, on a
\[
 \Pi_H\mathcal D=\mathcal H_{12},\qquad
 \mathcal D^{-1}=\tfrac14\mathcal H_{12}\Pi_H.
\]
L’image de $\Pi_H$ restreinte à ce domaine est
\[
 \mathcal L_H=
 \{u\in\mathbb Z^{12}:\mathcal H_{12}u\equiv0\pmod4\}.
\]
\end{proposition}
\begin{proof}
Pour $z=\mathcal D^{-1}(b,c,q)$, soient
$h_r=\sum_{j=0}^3z_{r+3j}$. Le décalage de quatre positions
envoie chaque orbite sur la suivante ; par conséquent, $B_r=h_r-h_{r+1}$.
L’identité $q=h_0+h_1+h_2$ donne $a_r=h_r$.
Sur une orbite, écrivons $(x_0,x_1,x_2,x_3)$ pour ses coordonnées
et $d_j=x_j-x_{j+1}$. On vérifie
\[
 \begin{aligned}
 d_1-d_3&=x_0+x_1-x_2-x_3,\\
 d_0+d_2&=x_0-x_1+x_2-x_3,\\
 d_0-d_2&=x_0-x_1-x_2+x_3.
 \end{aligned}
\]
Ce sont les trois lignes restantes de $H_4x$. L’inversion et la
caractérisation entière découlent de $\mathcal H_{12}^2=4I$.
\end{proof}

La condition de congruence harmonique contrôle la sortie de $\Pi_H$ ;
l’appartenance des canaux à $\operatorname{im}\mathcal D$
nécessite en outre \eqref{iii:r25:relaciones}--\eqref{iii:r25:congruencia}.
Un contrôle négatif exact consiste à prendre $b=0$, $q=0$ et
$c=e_0-e_3$. Les trois sommes orbitales de $c$ s’annulent, de sorte
que $\Pi_H(b,c,q)=0$, mais les relations d’image ne sont pas satisfaites.
La lecture harmonique ne remplace pas le contrôle de compatibilité des
canaux.

\subsection{Conjugaison, covariance et mémoire du registre d’incidence}
\label{iii:r25:transporte}

Soit $\rho(m)=13-m$. Si la conjugaison des incidences inverse
l’ordre des fenêtres, conserve les classes arithmétique et de frontière
des visites et échange les canaux des traces, alors
\[
 (A_m^\dagger,C_m^\dagger,V_m^\dagger)
 =(A_{\rho(m)},-C_{\rho(m)},V_{\rho(m)}).
\]
Avec $c_m=[C_m]_{10}$, la normalisation décimale donne
\begin{equation}
 z_m^\dagger=z_{\rho(m)}
 +100\mathbf1_{\{c_{\rho(m)}\ne0\}}-20c_{\rho(m)}.
 \label{iii:r25:conjugacion}
\end{equation}
En effet, $[-C]_{10}=0$ si $[C]_{10}=0$ et vaut
$10-[C]_{10}$ dans les autres cas. La négation du canal
s’applique avant de reconstruire le bloc décimal.

L’application à une involution concrète conserve également les
extrémités des fenêtres. Pour un chemin
$\gamma=(x_0,\ldots,x_n)$ lu avant l’avancée,
\begin{equation}
 \sum_{k=0}^{n-1}f(x_{n-k})
 =\sum_{k=0}^{n-1}f(x_k)+f(x_n)-f(x_0).
 \label{iii:r25:extremos}
\end{equation}
Les deux sommes contiennent les mêmes termes intérieurs ; la différence
est l’extrémité finale moins l’initiale. Sur les chemins fermés, cette
correction s’annule. Dans les fenêtres ouvertes, elle est incorporée avant de
réduire modulo $10$.

\begin{proposition}[Conséquence d’une réduction périodique]
\label{iii:r25:periodo}
Si l’observable d’une famille fixe de chemins satisfait
$o(t+54)=o(t)$, si ses multiplicités sont conservées sous ce retour
et si le même lecteur local agit dans des fenêtres de neuf instants,
alors
\[
 (A,C,V)_{m+6}=(A,C,V)_m,\qquad z_{m+6}=z_m.
\]
\end{proposition}
\begin{proof}
La translation $t\mapsto t+54$ envoie chaque fenêtre bijectivement
sur celle située six positions plus loin et conserve chaque incidence
comptée et sa multiplicité. L’égalité se transmet aux dénombrements
et à leurs normalisations décimales.
\end{proof}

Une présentation à douze blocs distincts conserve donc
une information qui ne se factorise pas par cette observable de période $54$ :
sélection des incidences, orientation, extrémités, multiplicité ou
mémoire. La proposition délimite l’effet de cette réduction ; elle
n’attribue pas une période $54$ à l’histoire complète.

\subsubsection{Détermination sur les orbites et accumulation d’événements}

Fixer uniquement $Q(z)=q$ laisse une classe affine du réseau
\[
 \mathcal Z_0=\{v\in\mathbb Z^{12}:\sum_i v_i=0\}
 =\bigoplus_{i=0}^{10}\mathbb Z(e_i-e_{11}).
\]
Par exemple, $100\mathbf1+e_0$ et $100\mathbf1+e_1$ ont pour
charge $1201$ et satisfont les contrôles d’image au moyen de leurs
canaux respectifs. Ce sont des témoins des équations de compatibilité,
sans attribution d’appartenance aux histoires terminales HMT.

\begin{proposition}[Covariance cyclique]
\label{iii:r25:covariancia}
Soit $\rho$ la translation des fenêtres d’un registre $R$, de
période minimale $p\mid12$, et soit $S$ la translation de ses
coordonnées. Un lecteur covariant sur cette orbite est déterminé
par $z\in\operatorname{Fix}(S^p)$. La condition $Q(z)=q$
admet des solutions entières si et seulement si $12/p$ divise $q$ ;
lorsqu’elles existent, elles forment une classe affine de rang $p-1$.
\end{proposition}
\begin{proof}
La règle $F(\rho^jR)=S^jz$ est bien définie exactement si
$S^pz=z$. Ces vecteurs répètent $12/p$ fois un bloc de
longueur $p$, et leur charge est $\frac{12}{p}\sum_{i=0}^{p-1}z_i$.
La formule prouve la divisibilité et le rang. La covariance des
canaux découle de $D_aS=SD_a$ et $QS=Q$.
\end{proof}

Dans la coordinatisation harmonique, l’action correspondante est
$\widehat S_H=\mathcal H_{12}S\mathcal H_{12}/4$.
Si l’origine est marquée, l’action transporte également cette marque.
La période de l’orbite indiquée ici ne s’identifie pas à celle de
l’histoire enrichie. Pour une lecture initiale $F_{108}$ déjà
définie, la naturalité sous troncature s’exprime au moyen de
$F_N=F_{108}\circ\tau_{N,108}$ ; cette identité conserve
l’évaluation initiale sans sélectionner sa règle de production.

\begin{proposition}[Composition additive des contributions d’incidence]
\label{iii:r25:eventos}
Supposons que l’état initial et chaque événement $a_t$ produisent,
au moyen de règles préalablement fixées sur le registre, des contributions
$(\delta w_t,\delta q_t)$ telles que
\[
 \sum_i\delta w_{t,i}=0,
 \qquad \delta q_t+T(\delta w_t)\equiv0\pmod {12}.
\]
Avec des conditions initiales compatibles $(w^{(0)},q^{(0)})$,
la mise à jour par somme détermine un unique registre à chaque étape.
Son accroissement est
\[
 \delta z_{t,i}=
 \frac{\delta q_t+T(\delta w_t)}{12}-P_i(\delta w_t).
\]
La construction respecte la concaténation des segments et conserve
la lecture d’un préfixe sous les prolongements ultérieurs.
\end{proposition}
\begin{proof}
Les conditions de compatibilité se conservent par somme et
l’inverse \eqref{iii:r25:inversa} est additif en $(w,q)$.
La récurrence fixe chaque état. La somme sur des segments concaténés
est la somme de leurs contributions ; un préfixe reçoit les mêmes
termes avant et après le prolongement.
\end{proof}

Cette proposition donne un critère suffisant pour l’accumulation
additive, non une exigence de linéarité de tout lecteur TPK.
Lorsque les événements se transportent par des actions non triviales, leurs
contributions s’expriment dans la coordinatisation terminale au moyen du cocycle
affine correspondant. Les règles d’événement et l’état initial
conservent dans les deux cas leur provenance d’incidence.

Le développement précédent détermine le passage
\[
 \mathscr L\longmapsto(A,C,V)\longmapsto z
 \longmapsto(D_3z,D_4z,Q(z))\longmapsto\mathcal H_{12}z
\]
avec les conditions de chaque application. Son domaine est distinct du
registre régional de l’appendice~\ref{iii:nucleo-finito} : les
équations $X_aL_a=Y_a$ y déterminent les relèvements à partir de
transitions ternaires. Ici, les incidences des visites et des traces
déterminent des blocs entiers et leurs canaux. L’unicité d’une
reconstruction ne remplace pas la production du registre qu’utilise
l’autre. Les deux conservent leurs primitives, opérations et preuves dans
la chaîne APP--TRIT--TPK.


\subsection{Deux lectures de l'état terminal}
\label{subsec:k-terminal}

Désignons par \(x_{\mathrm{term}}\) l'état produit par la composition de sélection, de transport et de mise à jour jusqu'à la coupe terminale. L'écriture opératorielle est
\[
 x_{\mathrm{term}}=
 (\mathcal U_{108}\circ\cdots\circ\mathcal U_1)(x_{\mathrm{seed}}),
 \qquad \mathcal U_t=\mathsf{Upd}_t\circ\mathsf{Tra}_t\circ\mathsf{Sel}_t.
\]
Ses deux lectures pertinentes sont
\[
 x_{\mathrm{term}}
 \xrightarrow{(\pi_{\mathrm{term}},R_{\mathrm{sgn}})}
 (B_{\mathrm{term}},U).
\]
La matrice \(B_{\mathrm{term}}\) retient la lecture ternaire visible ; le vecteur \(U\) conserve une autre coordonnée de la même histoire. On n'insère pas entre eux une flèche \(B_{\mathrm{term}}\to U\).

Le sélecteur terminal de la construction de référence~(sous-section~\ref{subsec:k-terminal}) opère sur des catalogues
de tailles 196, 197 et 196. Les signatures locales réduisent leurs
\(196\cdot197\cdot196=7\,567\,952\) triplets à 331, et la marge des
colonnes \(q_{\partial}=(1,2,2,1,2,0)\) laisse deux matrices :
\[
 B_0=\begin{pmatrix}0&2&0&2&2&0\\0&0&1&2&2&0\\1&0&1&0&1&0\end{pmatrix},
 \qquad
 B_1=\begin{pmatrix}0&2&1&0&2&0\\0&0&1&2&2&0\\1&0&0&2&1&0\end{pmatrix}.
\]
L'orientation des lignes est \(\sigma=(1,1,-1)=(1,1,2)\) dans
\(\mathbb F_3^3\). Les charges des lignes des candidates sont
\((0,2,0)\) et \((2,2,1)\). Seule la seconde appartient à
\(\langle\sigma\rangle=\{(0,0,0),(1,1,2),(2,2,1)\}\) ; par conséquent,
\(B_{\mathrm{term}}=B_1\). Cette vérification prouve la dernière étape du
sélecteur à partir du recensement déclaré. Elle ne transforme pas la marge des colonnes
en une conséquence de la charge des lignes, et ne remplace pas l'énumération
précédente par l'inspection de ces deux matrices.

Le lecteur signé possède la composition typée
\begin{equation}
 R_{\mathrm{sgn}}
 =\Pi_H\circ\mathcal E_{108}^{90,120}\circ\mathcal R_{12}:
 \mathcal X_{\mathrm{TPK}}^{\mathrm{term,mem}}
 \longrightarrow\mathbb Z^{12}.
 \label{eq:k-lector-firmado}
\end{equation}
Le registre des canaux détermine
\[
 \mathcal E_{108}^{90,120}(\mathcal R_{12}(x_{\mathrm{term}}))
 =(b^{90},b^{120},Q_{\mathrm{TPK}}),
\]
avec les coordonnées
\begin{align*}
 b^{90}={}&(-495,-116,-684,108,601,-90,\\
            &\hspace{19mm}-173,-88,313,560,-397,461),\\
 b^{120}={}&(-425,-281,-481,671,-255,30,\\
             &\hspace{19mm}475,-543,680,251,6,-128),\\
 Q_{\mathrm{TPK}}={}&6263.
\end{align*}
Ces quantités sont utilisées comme coordonnées du registre préalable, non
comme paramètres ajustés au moyen d’alpha. La reconstruction démontrée
dans les paragraphes suivants commence exactement à cette sortie typée.
La trace matérielle de l’extracteur antérieur et la portée de sa transcription
autonome sont conservées dans la note technique de provenance ; l’arithmétique
ultérieure n’est pas présentée comme un substitut à cette trace.

\subsection{Du registre des canaux au vecteur signé}
\label{subsec:k-pi-h}

Le lecteur harmonique \(\Pi_H\) peut s'écrire sans utiliser ni \(K\) ni
alpha. Pour éviter de confondre ses sommes orbitales avec les blocs d'alpha,
nous les noterons \(s_1,s_2,s_3\). Soient
\[
 B_r=\sum_{j=0}^{3}b^{120}_{r+3j},\qquad
 d_{r,j}=b^{90}_{r+3j},\qquad1\leq r\leq3,
\]
et
\begin{equation}
 s_1=\frac{Q_{\mathrm{TPK}}+2B_1+B_2}{3},\qquad
 s_2=s_1-B_1,\qquad s_3=s_2-B_2.
 \label{eq:k-sumas-armonicas}
\end{equation}
La division par trois appartient au lecteur des trois orbites. Son caractère entier est vérifié sur les sorties compatibles du registre ; il n'est pas supposé pour un élément arbitraire de \(\mathbb Z^{25}\).

Les trois blocs signés sont
\begin{equation}
 U^{(r)}=(s_r,d_{r,1}-d_{r,3},d_{r,0}+d_{r,2},d_{r,0}-d_{r,2}).
 \label{eq:k-bloques-firmados}
\end{equation}
La substitution entière donne
\[
 (B_1,B_2,B_3)=(972,-1073,101),\qquad
 (s_1,s_2,s_3)=(2378,1406,2479),
\]
et, par conséquent,
\begin{align*}
 U^{(1)}&=(2378,-452,-668,-322),\\
 U^{(2)}&=(1406,998,-204,-28),\\
 U^{(3)}&=(2479,-551,-371,-997).
\end{align*}
En entrelaçant les colonnes, on obtient
\begin{equation}
 U=(2378,1406,2479,-452,998,-551,
       -668,-204,-371,-322,-28,-997).
 \label{eq:k-u-firmado}
\end{equation}
La présence de valeurs négatives et de composantes supérieures à 999
montre que \(U\) n'est pas un mot de chiffres en base mille. Ce n'est pas non plus
\(U_6\Vert U_6\), concaténation du catalogue de période six. La différence
concerne le domaine et l'information conservée, et non seulement les valeurs.

L'ordre constructif jusqu'ici est
\[
 x_{\mathrm{term}}\longrightarrow\mathcal R_{12}
 \longrightarrow(b^{90},b^{120},Q_{\mathrm{TPK}})
 \longrightarrow U.
\]
L'inversion qui suit produit le registre dodécaphasique. Le fait que nous puissions ensuite récupérer les canaux à partir du registre dodécaphasique constitue un contrôle de réversibilité ; il n'inverse pas cet ordre causal.

\subsection{Hadamard par orbites et reconstruction entière}
\label{subsec:k-hadamard}

Dans le cycle à douze positions, le pas trois détermine trois orbites :
\[
 (1,4,7,10),\qquad(2,5,8,11),\qquad(3,6,9,12).
\]
La séparation de quatre secteurs de chaque orbite est celle qui, une fois
l'origine et l'orientation fixées, admet la lecture angulaire de quatre-vingt-dix degrés. La
transformation n'agit pas sur trois quadruplets contigus de l'ordre naturel.

Soit
\[
 H_4=\begin{pmatrix}
 1&1&1&1\\1&1&-1&-1\\1&-1&1&-1\\1&-1&-1&1
 \end{pmatrix}.
\]
Si \(P_{\mathrm{orb}}\) réordonne selon les trois orbites indiquées, nous définissons
\[
 \mathcal H_{12}=P_{\mathrm{orb}}^{-1}
 (H_4\oplus H_4\oplus H_4)P_{\mathrm{orb}}.
\]

\begin{lemma}[Inversion de Hadamard et intégralité]
\label{lem:k-hadamard}
On a
\[
 H_4^2=4I_4,\qquad H_4^{-1}=\frac14H_4,\qquad\det H_4=-16.
\]
Pour \(u\in\mathbb Z^4\), la préimage \(H_4^{-1}u\) est entière si et
seulement si \(H_4u\equiv0\pmod4\).
\end{lemma}

\begin{proof}
Les lignes sont orthogonales et le carré de leur norme vaut quatre ; la matrice est symétrique. Cela donne \(H_4^2=4I_4\) et l'inverse. Une élimination entière, ou le développement du déterminant, donne le signe négatif et la valeur \(-16\). La condition de caractère entier est exactement la divisibilité par quatre des quatre coordonnées de \(H_4u\).
\end{proof}

\begin{definition}[Transformation entière dodécaphasique]
Le sous-réseau des registres signés et sa coordinatisation entière sont
\[
 \mathcal L_H=\{u\in\ZZ^{12}:\mathcal H_{12}u\equiv0\pmod4\},
 \qquad \mathscr K(u)=\tfrac14\mathcal H_{12}u.
\]
L'application \(\mathscr K:\mathcal L_H\to\ZZ^{12}\) est un isomorphisme
de groupes abéliens, d'inverse \(k\mapsto\mathcal H_{12}k\).
Le registre canonique \(K\) est l'image du registre signé généré
par les opérations précédentes, avec une origine de phase et une orientation fixées.
\end{definition}

\begin{proof}
La congruence définit exactement l'intégralité de \(\mathscr K(u)\).
L'identité \(\mathcal H_{12}^2=4I\) démontre les deux compositions
inverses. En particulier, tout \(k\in\ZZ^{12}\) a pour préimage
\(\mathcal H_{12}k\in\mathcal L_H\).
\end{proof}

\begin{proposition}[Registre entier canonique]
\label{prop:k-sello}
L'inversion \(K^{(r)}=H_4U^{(r)}/4\) des blocs de
\eqref{eq:k-bloques-firmados} produit
\begin{align*}
 K^{(1)}&=(234,729,621,794),\\
 K^{(2)}&=(543,659,58,146),\\
 K^{(3)}&=(140,824,914,601).
\end{align*}
Dans l'ordre naturel, le registre dodécaphasique est
\begin{equation}
 K=(234,543,140,729,659,824,621,058,914,794,146,601).
 \label{eq:k-vector}
\end{equation}
C'est l'unique préimage rationnelle de \(U\), et elle appartient à \(\{0,\ldots,999\}^{12}\).
\end{proposition}

\begin{proof}
Les trois produits avant division sont
\begin{align*}
 H_4U^{(1)}&=(936,2916,2484,3176),\\
 H_4U^{(2)}&=(2172,2636,232,584),\\
 H_4U^{(3)}&=(560,3296,3656,2404).
\end{align*}
Toutes les coordonnées sont divisibles par quatre. Diviser et défaire la
permutation des orbites donne \eqref{eq:k-vector}. L'unicité procède de l'
inversibilité rationnelle, et non d'une recherche parmi des candidats proches d'
une valeur physique.
\end{proof}

Le caractère entier peut également se formuler en termes de réseaux. Les diviseurs déterminantaux de \(H_4\) sont \(1,2,4,16\) : les entrées ont pour plus grand commun diviseur un ; les mineurs d'ordre deux sont pairs et l'un d'eux est de module deux ; ceux d'ordre trois sont des multiples de quatre et l'un d'eux est de module quatre. On en déduit
\[
 \operatorname{SNF}(H_4)=\operatorname{diag}(1,2,2,4),
\]
et pour la transformation globale,
\[
 \operatorname{coker}\mathcal H_{12}
 \cong(\mathbb Z/2\mathbb Z)^6\oplus(\mathbb Z/4\mathbb Z)^3.
\]
L'indice de l'image est \(16^3=4096\). Ainsi, tout vecteur signé
entier n'admet pas nécessairement une préimage entière ; le vecteur obtenu satisfait les
congruences précises exigées par l'inversion. Le facteur \(1/4\) est
imposé par la matrice et n'a pas le statut d'un coefficient libre.

\subsection{Différences cycliques et composante uniforme}
\label{subsec:k-transversal}

Avec des indices cycliques modulo douze, soient
\[
 (D_3z)_m=z_m-z_{m+3},\qquad(D_4z)_m=z_m-z_{m+4},
 \qquad Q(z)=\sum_{m=1}^{12}z_m.
\]
Le premier opérateur compare des secteurs séparés par trois positions ; le
second, par quatre. Les noms de canal 90 et 120 se réfèrent ici à
ces séparations dans le cycle orienté, et non à une identification avec un
opérateur physique d'un autre domaine.

\begin{proposition}[Complétude du système transversal]
\label{prop:k-transversal}
L'opérateur
\[
 \mathcal D=(D_3,D_4,Q):\mathbb Z^{12}\longrightarrow\mathbb Z^{25}
\]
a un rang égal à douze et des facteurs invariants non nuls
\[
 \operatorname{diag}(1,\ldots,1,12),
\]
avec onze unités. En particulier, leurs sorties déterminent un unique registre dodécaphasique.
\end{proposition}

\begin{proof}
Si \(D_3z=D_4z=0\), le vecteur est constant le long des deux pas.
Comme \(3\) et \(4\) engendrent le cycle modulo douze,
\(\ker D_3\cap\ker D_4=\langle\mathbf1\rangle\). La charge totale
élimine cette dernière liberté parce que \(Q(\mathbf1)=12\).

Pour l'affirmation entière, les lignes de \((D_3,D_4)\) sont des incidences orientées d'un graphe connexe. Un arbre couvrant fournit onze facteurs invariants unitaires. Tout mineur d'ordre douze doit contenir la ligne \(Q\). En ajoutant les onze premières colonnes à la dernière, celle-ci s'annule dans les lignes d'incidence et vaut douze dans la ligne \(Q\). Tous les mineurs maximaux sont divisibles par douze et l'arbre en produit un de module exactement douze. Le dernier facteur est donc douze.
\end{proof}

La vérification sur \eqref{eq:k-vector} donne
\[
 D_3K=b^{90},\qquad D_4K=b^{120},\qquad Q(K)=6263.
\]
Les formules \eqref{eq:k-sumas-armonicas} et
\eqref{eq:k-bloques-firmados} reconstruisent à leur tour \(U\) à partir de ces
différences. On ferme ainsi un triangle de représentations exactes du
même objet : registre transversal, vecteur signé et registre dodécaphasique. L'égalité
des résultats ne fait pas de leurs codomaines un seul et même domaine ; elle démontre que
les applications déclarées restituent l'information pertinente.

\subsection{Deux lectures rationnelles : \texorpdfstring{\(\kappa_{36}\)}{kappa36} et \texorpdfstring{\(\kappa_{\mathrm{per}}\)}{kappa périodique}}
\label{subsec:k-kappas}

Une fois \(K\) construit, la coordinatisation en base mille permet deux opérations
distinctes. La première lit un seul tour ; la seconde prolonge périodiquement le registre
dodécaphasique. Soit \(B=1000\) et définissons son entier de concaténation
\begin{equation}
 N_K=\sum_{m=1}^{12}K_mB^{12-m}
 =234543140729659824621058914794146601.
 \label{eq:k-numerador}
\end{equation}
Les zéros initiaux d'un bloc, comme celui de 058, font partie de son écriture
à trois chiffres. Ils ne changent pas sa valeur entière, mais permettent de concaténer
sans ambiguïté les douze blocs.

\begin{definition}[Lectures finie et périodique du registre dodécaphasique]
\label{def:k-kappas}
La lecture d'un tour est
\[
 \kappa_{36}=\sum_{m=1}^{12}K_mB^{-m}=\frac{N_K}{B^{12}}.
\]
Le prolongement périodique est
\[
 K_j^{\mathrm{per}}=K_{1+((j-1)\bmod12)},\qquad
 \kappa_{\mathrm{per}}=\sum_{j\geq1}K_j^{\mathrm{per}}B^{-j}.
\]
\end{definition}

\begin{proposition}[Valeur et période exactes du registre dodécaphasique prolongé]
\label{prop:k-periodo}
On a
\begin{equation}
 \kappa_{\mathrm{per}}=\frac{N_K}{B^{12}-1},\qquad
 \kappa_{\mathrm{per}}-\kappa_{36}
 =\frac{N_K}{B^{12}(B^{12}-1)}
 =B^{-12}\kappa_{\mathrm{per}}.
 \label{eq:k-diferencia-kappas}
\end{equation}
Le mot infini du registre dodécaphasique a pour période minimale douze en base mille et pour période minimale trente-six en base dix.
\end{proposition}

\begin{proof}
La somme du premier tour est \(N_KB^{-12}\) ; chaque nouveau tour
multiplie sa contribution par \(B^{-12}\). La série géométrique produit
\(N_K/(B^{12}-1)\), et la soustraction donne la seconde identité.

Les douze composantes de \(K\) sont distinctes. Une période propre du cycle
identifierait deux d'entre elles ; la période minimale des blocs est donc
douze. Le développement décimal est canonique, car le registre dodécaphasique n'a pas de queue
de blocs 999. Si sa période décimale minimale était \(d\), elle diviserait
36. Un regroupement par trois produirait une période de blocs
\(d/\gcd(d,3)\). Aucun diviseur propre de 36 ne donne douze par cette
opération. La période décimale minimale est donc 36.
\end{proof}

Les deux lectures sont rationnelles, mais ne sont pas le même rationnel. La première se termine après douze blocs ; la seconde répète le tour. En outre, \(0<\kappa_{\mathrm{per}}-\kappa_{36}<B^{-12}\). Leur proximité ne permet pas de substituer l'une à l'autre dans une identité exacte. En particulier, la clôture finie d'alpha utilise \(\kappa_{36}\), tandis que la somme de la section prolongée utilise \(\kappa_{\mathrm{per}}\).

\input{sections/k_reversibilidad.tex}

\subsection{Périodicité et conoyau de l'opérateur de transport}
\label{subsec:k-clase-acarreo}

La périodicité que nous venons de démontrer appartient au registre dodécaphasique. Elle n'implique
pas la périodicité des coordonnées régionales, des retenues de la clôture
ni d'alpha. Il est utile de préciser cette séparation au moyen d'un second
quotient, purement entier.

Soit \(S\) le décalage cyclique de douze coordonnées, d'orientation fixée, et soit \(b\geq2\) une base entière. L'opérateur de retenue périodique est \(\partial_b=S-bI\). Si une identité périodique s'écrit
\[
 A=P+E-\Phi-K+\partial_bC,
\]
alors les transformations
\[
 K\longmapsto K+\partial_bh,\qquad C\longmapsto C+h
\]
conservent son membre de droite. Cette invariance compare les représentants d'une
classe de retenue ; elle n'autorise pas à modifier le registre dodécaphasique déjà
sélectionné par le registre.

\begin{proposition}[Quotient de la retenue périodique]
\label{prop:k-cociente-acarreo}
Avec l'orientation précédente,
\[
 \operatorname{coker}(S-bI)\cong\mathbb Z/(b^{12}-1)\mathbb Z.
\]
\end{proposition}

\begin{proof}
Le module cyclique se présente sous la forme
\(\mathbb Z[t]/(t^{12}-1)\). Imposer \(S=bI\) équivaut à imposer
\(t=b\) et laisse la relation \(b^{12}-1=0\). L'évaluation des
coefficients en \(b\), avec l'ordre de concaténation choisi, représente
la classe résultante.
\end{proof}

Le dénominateur \(B^{12}-1\) de \(\kappa_{\mathrm{per}}\) et ce quotient proviennent de la même longueur cyclique, mais leurs objets restent distincts : l'un est une évaluation rationnelle et l'autre un quotient d'un module de retenues. La clôture finie d'alpha sera une chaîne orientée à frontière droite. Son prolongement transportera une frontière distante compatible, et non une retenue identifiée par décret entre la première et la treizième position.

Le chapitre suivant est ainsi préparé. Le registre dodécaphasique a été fixé avant
de résoudre la récurrence d'alpha ; sa reconstruction par Hadamard est
entière et unique ; ses différences transversales restituent le registre ; et
ses deux évaluations rationnelles possèdent des domaines et des usages distincts. Aucun
de ces faits n'utilise alpha comme entrée, et aucun ne décide encore du
type arithmétique de la sortie de la clôture complète.

\clearpage\section{Dérivation de la constante de structure fine}
\label{sec:alpha}

La constante de structure fine reçoit dans ce développement une détermination arithmétique qui réunit les trois coordonnées régionales et le registre intégral dodécaphasique. L'opération combine ses blocs avec orientation fixée et transporte les quotients de la normalisation positionnelle. Le paramètre résultant exprime la compatibilité de cette différence orientée avec la prolongation des cylindres. Sa valeur et les données d'incidence utilisées dans sa construction resteront liées pendant le développement exceptionnel.

L'ordre interne est important. Les coordonnées \(\pi\), \(e\) et \(\varphi\), déjà construites dans la section précédente, sont des coordonnées
précédentes du même état arithmétique avec mémoire de trajectoire, et peuvent agir légitimement comme
entrées d'un lecteur ultérieur. Elles ne sont pas des constantes conventionnelles
injectées pour sélectionner alpha. La génération corrélée des quatre coordonnées partage une origine et
une compatibilité coinductive, sans que cela impose d'évaluer ses quatre
composantes en un seul pas.

Les deux voies expriment des fonctions mathématiques différentes du même paramètre.
La première réalise une clôture entière par division euclidienne en base
\(1000\). La seconde étudie l’équilibre analytique des vacances au moyen
de la résolvante \(E_{21/22}\) et de ses prolongements gradués. Une voie
détermine des blocs et des quotients ; l’autre détermine une racine par une
équation et un domaine d’unicité. Leur identification est formulée ensuite
dans les cylindres communs à toute profondeur. L’exposé distingue
ces trois résultats : construction positionnelle, détermination analytique
et compatibilité des deux.

\subsection{Première voie : détermination positionnelle par clôture entière}
\subsubsection{Domaine et orientation de la clôture entière}
\label{subsec:alpha-dominio}

Soient \(P_j,E_j,\Phi_j\in\{0,\ldots,999\}\) les blocs canoniques des
trois coordonnées régionales, et prolongeons le registre dodécaphasique par
\[
 K_j=K_{1+((j-1)\bmod12)}.
\]
L'orientation du canal est \((1,1,-1)\). La contribution du registre dodécaphasique se soustrait selon la clôture dodécaphasique déjà typée. Pour \(B=1000\), l'opération locale distingue le bilan signé
\[
 Z_j:=P_j+E_j-\Phi_j-K_j
\]
de la somme qui incorpore le quotient entrant. La normalisation est
\begin{equation}
 s_j=Z_j+c_{j+1},\qquad
 A_j=s_j-B\left\lfloor\frac{s_j}{B}\right\rfloor,
 \qquad c_j=\left\lfloor\frac{s_j}{B}\right\rfloor.
 \label{eq:alpha-recurrencia}
\end{equation}
Ainsi, \(A_j\in\{0,\ldots,B-1\}\), même lorsque \(s_j\) est
négatif. La retenue qui entre depuis la droite appartient au domaine de
l'opération ; elle ne peut pas être omise car le tableau s'imprime de gauche
à droite.

\begin{definition}[Clôture de profondeur finie]
\label{def:alpha-finita}
Pour une profondeur \(N\) et une frontière entière \(t\), l'application
\(\mathcal C_N\) reçoit
\((P_{1:N},E_{1:N},\Phi_{1:N},K_{1:N},c_{N+1}=t)\) et applique
\eqref{eq:alpha-recurrencia} dans l'ordre \(N,N-1,\ldots,1\). Sa sortie
est la paire de mots
\[
 (A_{1:N},c_{1:N+1}).
\]
\end{definition}

La division euclidienne ne possède pas un choix de signe résiduel : le reste se
fixe en \(\{0,\ldots,999\}\) et le quotient reste déterminé. Par
exemple, \(-431=569-1000\) et \(-1200=800-2\cdot1000\). Les
retenues négatives de la clôture ne sont pas des anomalies ; elles sont l'information
entière nécessaire pour maintenir l'égalité orientée.

\subsubsection{Fenêtre dodécaphasique et conditions de retenue}
\label{subsec:alpha-ventana}

Les douze premières triades obtenues sont
\begin{align*}
 P={}&(141,592,653,589,793,238,462,643,383,279,502,884),\\
 E={}&(718,281,828,459,045,235,360,287,471,352,662,497),\\
 \Phi={}&(618,033,988,749,894,848,204,586,834,365,638,117).
\end{align*}
Dans le canal de fermeture, les cinq régions qui déterminent la coordonnée \(P\) conservent leurs antécédents distincts. Partager cette lecture ne les identifie pas en tant que régions ou en tant qu'histoires.

Avec frontière \(c_{13}=0\), la récurrence produit le tableau suivant.
Chaque entrée est montrée, la retenue entrante, le total avant normalisation,
le bloc obtenu et la retenue sortante. Le tableau est un calcul entier
reproductible, non une liste de blocs sélectionnés par leur ressemblance avec
un objectif.

\begin{table}[htbp]
\centering
\small
\setlength{\tabcolsep}{3.5pt}
\begin{tabular}{r|rrrr|rr|rr}
\hline
\(j\)&\(P_j\)&\(E_j\)&\(\Phi_j\)&\(K_j\)&\(c_{j+1}\)&\(s_j\)&\(A_j\)&\(c_j\)\\
\hline
1&141&718&618&234&0&7&007&0\\
2&592&281&033&543&0&297&297&0\\
3&653&828&988&140&\(-1\)&352&352&0\\
4&589&459&749&729&\(-1\)&\(-431\)&569&\(-1\)\\
5&793&045&894&659&\(-2\)&\(-717\)&283&\(-1\)\\
6&238&235&848&824&\(-1\)&\(-1200\)&800&\(-2\)\\
7&462&360&204&621&0&\(-3\)&997&\(-1\)\\
8&643&287&586&058&\(-1\)&285&285&0\\
9&383&471&834&914&\(-1\)&\(-895\)&105&\(-1\)\\
10&279&352&365&794&0&\(-528\)&472&\(-1\)\\
11&502&662&638&146&0&380&380&0\\
12&884&497&117&601&0&663&663&0\\
\hline
\end{tabular}
\caption{Clôture initiale avec frontière droite nulle. La dépendance se propage de la dernière ligne à la première.}
\label{tab:alpha-ventana}
\end{table}

Le mot des retenues est
\[
 (c_1,\ldots,c_{13})=(0,0,0,-1,-1,-2,-1,0,-1,-1,0,0,0),
\]
et la lecture de la fenêtre est le rationnel
\begin{equation}
 \alpha_{12}^{(0)}
 =\frac{7297352569283800997285105472380663}{10^{36}}
 =0.007297352569283800997285105472380663.
 \label{eq:alpha-ventana-cero}
\end{equation}
L'exposant rappelle la frontière choisie. Cette fenêtre est exacte dans son
domaine fini ; on ne l'identifie pas d'emblée à la troncature de la
section infinie. La différence entre les deux lectures deviendra visible lors du
transport de la retenue de la queue lointaine.

\subsubsection{Télescopage et réversibilité dans une fibre de frontière}
\label{subsec:alpha-telescopia}

La forme locale de l'égalité est
\begin{equation}
 P_j+E_j-\Phi_j-K_j-A_j=Bc_j-c_{j+1}.
 \label{eq:alpha-identidad-local}
\end{equation}
Le terme de retenue est un lien entre des positions contiguës. Si on le supprime coordonnée par coordonnée, on obtient une égalité fausse. Seule une évaluation qui conserve ses extrémités permet de le télescoper.

\begin{proposition}[Identité entière de profondeur arbitraire]
\label{prop:alpha-telescopia}
Pour \(\iota_N(z)=\sum_{j=1}^{N}z_jB^{N-j}\), toute exécution de
\(\mathcal C_N\) satisfait
\begin{equation}
 \iota_N(P)+\iota_N(E)-\iota_N(\Phi)-\iota_N(K)-\iota_N(A)
 =B^Nc_1-c_{N+1}.
 \label{eq:alpha-telescopia}
\end{equation}
\end{proposition}

\begin{proof}
Multiplier \eqref{eq:alpha-identidad-local} par \(B^{N-j}\) et ajouter  
produit du côté droit
\[
 \sum_{j=1}^{N}(Bc_j-c_{j+1})B^{N-j}=B^Nc_1-c_{N+1}.
\]
Chaque retenue intérieure apparaît deux fois avec un signe opposé ; seuls les deux extrêmes survivent.
\end{proof}

Pour le tableau ~\ref{tab:alpha-ventana}, les deux extrémités sont nulles. Si
\(p_{12},e_{12},f_{12}\) sont les trois sommes fractionnaires tronquées,
\[
 \alpha_{12}^{(0)}=p_{12}+e_{12}-f_{12}-\kappa_{36}.
\]
La constante du registre dodécaphasique qui apparaît est \(\kappa_{36}\), non
\(\kappa_{\mathrm{per}}\). Le domaine de cette identité contient
des troncatures et une frontière finie ; sa forme n'autorise pas à effacer les deux
conditions en passant à la section complète.

L'inversion locale est aussi exacte :
\[
 K_j=P_j+E_j-\Phi_j+c_{j+1}-A_j-Bc_j.
\]
Étant donnés \(P,E,\Phi\), les blocs \(A\) et les retenues complètes,  
on restaure le registre dodécaphasique. La concaténation restaure en outre
\[
 \iota_N(K)=\iota_N(P)+\iota_N(E)-\iota_N(\Phi)-\iota_N(A)
             -B^Nc_1+c_{N+1}.
\]
Cette réversibilité est une propriété de la fibre avec données et frontières fixes. Elle ne doit pas devenir une nouvelle recette causale qui commence par alpha et fabrique ensuite le registre. La direction de construction reste encore registre signé, registre dodécaphasique, clôture et sortie.

\subsubsection{Prolongation et normalisation canonique de la section complète}
\label{subsec:alpha-completa}

La prolongation conserve le registre dodécaphasique périodique et reçoit les suites
régionales complètes. Soient
\[
 p=\sum_{j\geq1}P_jB^{-j},\qquad
 e_0=\sum_{j\geq1}E_jB^{-j},\qquad
 f=\sum_{j\geq1}\Phi_jB^{-j}.
\]
Chaque série est une réalisation du caractère interne correspondant.  
Les parties entières déjà obtenues sont 3, 2 et 1, de sorte que  
\(p=\pi-3\), \(e_0=e-2\) et  
\(f=\varphi-1\).

\begin{proposition}[Existence et unicité de la normalisation complète]
\label{prop:alpha-normalizacion}
La série
\begin{equation}
 y=p+e_0-f-\kappa_{\mathrm{per}}
   =\sum_{j\geq1}(P_j+E_j-\Phi_j-K_j)B^{-j}
 \label{eq:alpha-serie}
\end{equation}
converge absolument. Dans le domaine des blocs obtenus,
\(0<y<10^{-2}\). Il existe une unique expansion canonique \((A_j)\) de
\(y\), sans queue finalement constante 999, et une unique suite entière
bornée \((c_j)\) qui satisfait \eqref{eq:alpha-identidad-local} avec
\(c_1=0\). Ces retenues appartiennent à l'ensemble invariant
\(\mathcal C=\{-2,-1,0,1,2\}\).
\end{proposition}

\begin{proof}
Les coefficients de la série sont bornés par \(2(B-1)\) en module,  
de sorte qu'elle est dominée par une série géométrique. Si  
\(Z_j=P_j+E_j-\Phi_j-K_j\) et  
\(R_j(Z)=\sum_{r\geq0}Z_{j+r}B^{-r-1}\), chaque queue est la différence  
de deux sommes de deux queues canoniques et appartient à \((-2,2)\).  
Le premier coefficient est  
\(Z_1=141+718-618-234=7\) ; par conséquent  
\(y=(7+R_2(Z))/1000\in(0.005,0.009)\).

Soit \((A_j)\) l'expansion canonique de \(y\). Définissons
\[
 c_j=R_j(Z)-R_j(A).
\]
En multipliant l'égalité des deux séries par \(B^{j-1}\) et en séparant
ses premiers \(j-1\) termes on obtient
\[
 c_j=\sum_{r=1}^{j-1}(A_r-Z_r)B^{j-1-r}\in\mathbb Z,
\]
et \(c_1=0\). De plus, \(R_j(A)\in[0,1)\), donc
\(-3<c_j<2\), ce qui inclut toutes les retenues dans \(\mathcal C\).
Les identités de déplacement des queues donnent
\(Bc_j=Z_j-A_j+c_{j+1}\), qui est la récurrence requise.

Réciproquement, toute solution bornée à sortie canonique et
\(c_1=0\) se télescope à la limite en la même série \(y\). Le développement
canonique est unique et, une fois fixé, la formule des queues détermine
toutes les retenues. Enfin, si l'on applique une exécution finie avec
\(c_{j+1}\in\mathcal C\), on a \(-2000\leq s_j\leq2000\),
de sorte que son quotient euclidien appartient de nouveau à \(\mathcal C\).
\end{proof}

La proposition explicite la normalisation du lecteur ultérieur : elle ne remplace pas la génération des trois suites par une série conventionnelle introduite comme donnée. Une fois ces suites produites par l'état coinductif, leur combinaison et la frontière de la normalisation sont déterminées sans consulter alpha.

La procédure finie conservée dans le corpus propage les cinq
frontières de \(\mathcal C\) et conserve uniquement le préfixe commun à
toutes. Si deux profondeurs ont coalescé avant une frontière,
la récurrence déterministe de droite à gauche donne les mêmes blocs
à sa gauche. C'est la compatibilité avec la troncature. On ne
choisit pas une frontière parce qu'elle donne des chiffres favorables, et on ne déduit pas la
compatibilité à toute profondeur de l'observation d'une longue exécution.

La comparaison initiale montre pourquoi cette précaution importe. Dans
la prolongation stabilisée on obtient \(c_{13}^{\infty}=-1\), tandis
que la fenêtre isolée imposait \(c_{13}=0\). Ainsi,
\[
 884+497-117-601-1=662
\]
dans le dernier bloc du premier tour. Les onze premières triades ne
changent pas, mais la prolongation commence
\[
 (007,297,352,569,283,800,997,285,105,472,380,662,999,\ldots).
\]
Une apparition du bloc 999 ne viole pas la convention canonique : ce
qui est exclu est une queue qui soit finalement constante 999. La fenêtre
à frontière nulle se termine en 663 ; le troncage de la section complète
se termine en 662. Les deux affirmations décrivent des opérations différentes
et exactes. Elles ne doivent pas être fusionnées en une seule notation de projection.

\subsubsection{Définition opérationnelle et équation additive complète}
\label{subsec:alpha-definicion}

\begin{definition}[Coordonnée alpha de la clôture]
\label{def:alpha-operativa}
La voie entière détermine la coordonnée
\begin{equation}
 \alpha_A=\sum_{j\geq1}A_jB^{-j}
 =p+e_0-f-\kappa_{\mathrm{per}},
 \label{eq:alpha-via-a}
\end{equation}
où \((A_j,c_j)\) est la normalisation canonique compatible de
\eqref{eq:alpha-recurrencia}. Dans la réalisation scalaire des
coordonnées HMT,
\begin{equation}
 \boxed{\alpha_A=
 \pi+e-\varphi
 -4-\kappa_{\mathrm{per}}.}
 \label{eq:alpha-identidad-completa}
\end{equation}
\end{definition}

L'entier quatre n'est pas calibré : il est
\(3+2-1\), combinaison des parties entières des trois coordonnées.
La base mille provient de la coordinatisation cubique et la période douze du registre dodécaphasique. Les
signes sont ceux du canal orienté et de la clôture, et la frontière distante est
celle qu'impose la section compatible. Telles sont les origines effectives des
coefficients de l'équation.

À ce niveau, on peut décrire alpha comme la coordonnée scalaire du défaut de fermeture entre les trois réalisations régionales et le retour rationnel dodécaphasique, calculé avec une retenue compatible. Le terme \emph{défaut} désigne une différence algébrique déterminée, et non une perte arbitraire ni une imperfection du système. Cette description ne supprime pas la généalogie : sans les régions, le registre signé, l'orientation et la loi de retenue, l'égalité scalaire ne serait pas une exposition suffisante de la construction.

Il faut également maintenir séparés trois objets que la notation d'alpha peut accompagner dans des lecteurs ultérieurs : \(A\) est un mot de blocs ; le support négatif de pré-retenue s'obtient à partir du signe des totaux avant normalisation ; et un vecteur radial exceptionnel appartient à un codomaine d'incidence. Aucun d'eux n'est interchangeable avec la valeur réelle \(\alpha_A\).

\input{sections/revision_alpha_notacion.tex}
\input{sections/revision_alpha_frontera.tex}

\input{sections/alpha_dependencias.tex}

\subsection{Seconde voie : détermination analytique au moyen des vacances}
\subsubsection{Invariants du transport et résolvante graduée}
\label{subsec:alpha-via-b}

La voie analytique reçoit des invariants de l'état orienté, et non la sortie
\(\alpha_A\). Sa troncature d'ordre six est
\begin{equation}
 P_6(x)=2\pi x-\frac74x^2
 +\frac{x^3}{2\pi}
 +\frac{x^4}{20}-\frac{2x^5}{21}-\frac{x^6}{46},
 \label{eq:alpha-p6}
\end{equation}
et la condition de vacance utilise
\[
 \delta=\log_{10}\frac{10}{9}.
\]
Le rapport \(10/9\) compare la coordinatisation décimale aux neuf positions
APP de la construction nonadique. Son logarithme appartient à cette lecture
ultérieure de compactification ; ce n'est pas une valeur cible d'alpha.

La généalogie des coefficients doit être conservée complète, mais avec sa
force logique exacte. Les développements récents réunissent une signature marquée
de l'état. La matrice de synchronisation et le recensement associé sont
\[
 D_{\mathrm{syn}}=\begin{pmatrix}85&112\\41&52\end{pmatrix},
 \qquad N_9=43.
\]
De son déterminant on extrait
\[
 a=-\frac{\det D_{\mathrm{syn}}}{N_9}
   =\frac{172}{43}=4.
\]
L'horloge induite sur les trous courts a un spectre secondaire \(\{6,7\}\) ; son extrémité supérieure donne \(b=7\). La cardinalité tritique \(m=3\) détermine
\[
 t_*=mb=21,\qquad k_*=mb-1=20.
\]
La lecture de vacance dans la coordinatisation en base mille donne
\[
 v_-=\lfloor1000\delta\rfloor=45,\qquad
 v_+=\lceil1000\delta\rceil=46.
\]
Ces entiers sont des antécédents typés de l'évaluation de la troncature. Ils ne sont obtenus ni à partir de sa racine ni à partir d'une approximation métrologique de celle-ci.

La branche inférieure coïncide en outre avec le déterminant du bloc APP
\[
 B_c=\begin{pmatrix}7&2\\2&7\end{pmatrix},\qquad\det B_c=45.
\]
Le rang transversal du registre du chapitre précédent donne
\(r_{\mathrm{dod}}=11\), et la demi-couronne de phase fournit
\(f_{\mathrm{ph}}=8\binom62=120\). En particulier,
\(45r_{\mathrm{dod}}=495\).

Avec \(\omega=2\pi\), la signature
\[
 \mathcal I_9=(\omega,m,a,b,k_*,t_*,v_-,v_+,f_{\mathrm{ph}},r_{\mathrm{dod}})
\]
alimente le lecteur gradué
\begin{align}
 \mathfrak E_9(\mathcal I_9)(x)
 ={}&\omega x-\frac ba x^2+\omega^{-1}x^3
 +\frac{x^4}{k_*}-\frac{2x^5}{t_*}-\frac{x^6}{v_+}\notag\\
 &-\frac{x^7}{f_{\mathrm{ph}}}+\frac{x^8}{v_-}
 +\frac{2x^9}{v_-r_{\mathrm{dod}}}.
 \label{eq:alpha-firma-coeficientes}
\end{align}
La substitution restitue les coefficients de \(P_6\) et de sa
prolongation d'ordre neuf. Le couple \(45/46\) traverse la séparation
entre les deux : la branche supérieure occupe le degré six et la branche inférieure réapparaît
aux degrés huit et neuf.

La factorisation de \eqref{eq:alpha-firma-coeficientes} se formule  
pour un état qui conserve l'origine, l'ordre des canaux et l'orientation.  
Dans ce domaine marqué, chaque coefficient reste fixé par  
\((\omega,m,a,b,k_*,t_*,v_-,v_+,f_{\mathrm{ph}},r_{\mathrm{dod}})\).  
Deux lecteurs qui préservent la signature et appliquent la même évaluation  
graduée coïncident terme à terme jusqu'au degré neuf. Cette  
unicité ne s'étend pas aux lecteurs qui oublient ces données.

\subsubsection{Résolvante nilpotente aux degrés \(7\), \(8\) et \(9\)}
\label{subsec:alpha-resolvente}

La première prolongation admet une construction opératoire finie
spécialement explicite. Sur l'espace à base \(e_7,e_8,e_9\),
soit
\[
 Ne_7=-\frac83e_8,\qquad Ne_8=\frac2{11}e_9,\qquad Ne_9=0,
 \qquad v_7=-\frac1{120}e_7.
\]
Le symbole \(N\) désigne ici un opérateur nilpotent, et non une profondeur de troncature. Son domaine et sa graduation sont fixés avant d'évaluer le polynôme en un nombre.

\begin{lemma}[Prolongement nilpotent exact]
\label{lem:alpha-resolvente}
On vérifie \(N^3=0\) et
\[
 (I-N)^{-1}v_7
 =-\frac{e_7}{120}+\frac{e_8}{45}+\frac{2e_9}{495}.
\]
L'évaluation \(e_n\mapsto x^n\) produit
\begin{equation}
 T_{7:9}(x)=-\frac{x^7}{120}+\frac{x^8}{45}+\frac{2x^9}{495}.
 \label{eq:alpha-t79}
\end{equation}
\end{lemma}

\begin{proof}
L'opérateur augmente le degré et s'annule sur le dernier vecteur de la base,
de sorte que son cube est nul. L'inverse est la résolvante finie
\(I+N+N^2\). De plus,
\[
 Nv_7=\frac1{45}e_8,\qquad N^2v_7=\frac2{495}e_9.
\]
Additionner les trois termes et évaluer donne la formule.
\end{proof}

Ici, les dénominateurs 120, 45 et 495 ne sont pas des ajustements décimaux : ils proviennent de la demi-couronne, du déterminant APP et de son produit par le rang dodécaphasique. Les transferts \(-8/3\) et \(2/11\), ainsi que le signe et le degré du germe, appartiennent au lecteur typé. Conserver seulement la symétrie du cycle et remplacer cet opérateur par un autre opérateur symétrique ne conserve pas nécessairement les coefficients.

\subsubsection{Racine de multiplicité un de la troncature d'ordre neuf}
\label{subsec:alpha-jet}

Définissons
\[
 P_9=P_6+T_{7:9},\qquad E_{21/22}^{(9)}(x)=P_9(x)-\delta,
 \qquad I_\alpha=(0,10^{-2}).
\]
Cet opérateur fini a un résultat propre d'existence et d'unicité.

Comme vérification ultérieure des caractères déjà construits, l'isolement rationnel peut se reproduire sans une table décimale. Soient
\[
 A_n(t)=\sum_{j=0}^{n}\frac{(-1)^jt^{2j+1}}{2j+1},\qquad
 L_n(t)=2\sum_{j=0}^{n}\frac{t^{2j+1}}{2j+1}\quad(0<t<1).
\]
Le critère alternant et la somme géométrique de la queue donnent
\begin{align}
 A_{2n+1}(t)&<\arctan t<A_{2n}(t),
 \label{eq:alpha-cotas-arctan}\\
 L_n(t)&<\log\frac{1+t}{1-t}
 <L_n(t)+\frac{2t^{2n+3}}{(2n+3)(1-t^2)}.
 \label{eq:alpha-cotas-log}
\end{align}
Nous appliquons \eqref{eq:alpha-cotas-arctan} à l'identité de Machin
\(\pi=16\arctan(1/5)-4\arctan(1/239)\), et
\eqref{eq:alpha-cotas-log} à
\[
 \log(10/9)=\log\frac{1+1/19}{1-1/19},\qquad
 \log 10=\log(10/9)+2\log\frac{1+1/2}{1-1/2}.
\]
Écrivons
\[
 \delta:=\frac{\log(10/9)}{\log 10}.
\]
Avec les ordres de troncature \(80,20,40,120\), appliqués
respectivement aux quatre séries précédentes, on obtient des intervalles
rationnels pour \(\pi\) et \(\delta\), tous deux de largeur inférieure à
\(10^{-70}\). Ces bornes vérifient une réalisation ultérieure ; elles ne
sélectionnent pas l'état APP--TRIT--TPK ni ne reçoivent une valeur d'alpha.

\begin{proposition}[Racine de multiplicité un de la troncature]
\label{prop:alpha-raiz-jet}
L'opérateur \(E_{21/22}^{(9)}\) possède une unique racine
\(\xi_9\) dans \(I_\alpha\), et cette racine est simple. Elle possède en outre l'isolement rationnel
\begin{equation}
 \frac{729735256928380099}{10^{20}}
 <\xi_9<\frac{729735256928380100}{10^{20}}.
 \label{eq:alpha-intervalo-jet}
\end{equation}
\end{proposition}

\begin{proof}
Les bornes des coordonnées internes
\(3.14<\pi<3.15\) et
\(0.045<\delta<0.046\) sont suffisantes pour l'existence. En zéro,
\(E_{21/22}^{(9)}(0)=-\delta<0\). En \(10^{-2}\), la contribution
linéaire moins les termes négatifs est supérieure à \(0.0626\), donc la
valeur de l'opérateur est positive.

Pour \(0\leq x\leq10^{-2}\), en éliminant les termes positifs de
la dérivée on obtient
\[
 (E_{21/22}^{(9)})'(x)
 \geq2\pi-\frac72\cdot10^{-2}
 -\frac{10}{21}\cdot10^{-8}-\frac6{46}\cdot10^{-10}
 -\frac7{120}\cdot10^{-12}>6.24.
\]
La continuité donne l'existence ; la borne sur la dérivée donne l'unicité et la simplicité. Pour l'extrémité gauche \(l\) et l'extrémité droite \(u\) de \eqref{eq:alpha-intervalo-jet}, la substitution rationnelle des bornes \eqref{eq:alpha-cotas-arctan}--\eqref{eq:alpha-cotas-log} dans le polynôme donne, avec arrondi dirigé,
\[
 -4.6\,10^{-20}<E_{21/22}^{(9)}(l)<-4.5\,10^{-20}<0,
 \qquad
 0<1.6\,10^{-20}<E_{21/22}^{(9)}(u)<1.8\,10^{-20}.
\]
La monotonie déjà prouvée place l'unique racine entre les deux extrémités.
L'intervalle se vérifie après avoir construit l'opérateur et n'intervient pas dans la sélection de ses coefficients.
\end{proof}

Le résultat concerne \(\xi_9\). Il n'autorise pas à écrire \(\xi_9=\alpha_A\) comme une égalité exacte de nombres complets. La racine d'une troncature peut partager un préfixe avec la racine de l'opérateur prolongé et différer à une position ultérieure. La simplicité permet de contrôler le déplacement de la racine sous perturbations ; elle n'annule pas ces perturbations.

\subsection{Critère d'identification des lectures compatibles}
\label{subsec:alpha-limite-dos-vias}

L'identification de deux lectures complètes exige de conserver leurs cylindres
et les applications de restriction. Pour une collection graduée compatible,
\[
 E_{21/22}^{(6)}\longleftarrow E_{21/22}^{(9)}
 \longleftarrow E_{21/22}^{(12)}\longleftarrow\cdots,
 \qquad
 E_{21/22}=\varprojlim_m E_{21/22}^{(6+3m)},
\]
la limite désigne le lecteur avec ses transports, et non une queue sélectionnée
librement. Le critère suivant sépare la conséquence de la
compatibilité de la construction effective qui l'établit ensuite.

\begin{proposition}[Identification des limites compatibles]
\label{prop:alpha-dos-limites}
Si les valeurs des voies entière et analytique appartiennent aux mêmes cylindres \(C_N^\alpha\) pour tout \(N\), et
\(\operatorname{diam}C_N^\alpha\to0\), alors
\[
 \alpha_A=\alpha_B=\alpha.
\]
\end{proposition}

\begin{proof}
Pour tout \(N\), la distance entre les deux valeurs est bornée
par \(\operatorname{diam}C_N^\alpha\). Faire tendre \(N\) vers l'infini
donne une distance nulle. L'existence des sections compatibles assure
qu'il ne s'agit pas de l'intersection vide d'un système d'approximations.
\end{proof}

La différence peut s'exprimer algébriquement. Soit \(F=\mathbb Q(\pi,\delta)\) et soit \(j_9:F[[x]]\to F[x]/(x^{10})\) la troncature. Alors
\[
 \ker j_9=x^{10}F[[x]].
\]
Tous les opérateurs
\[
 E_h(x)=E_{21/22}^{(9)}(x)+x^{10}h(x)
\]
ont la même troncature d'ordre neuf. Cependant, en prenant
\(h=0\) et \(h=1/729\),
\[
 E_{1/729}(\xi_9)=\frac{\xi_9^{10}}{729}>0,
\]
de sorte qu'ils ne partagent pas cette racine. L'observation ne permet pas de choisir une queue étrangère au transport HMT ; elle démontre précisément pourquoi la queue générée et sa compatibilité doivent être conservées. Une série formelle arbitraire ne suffit pas non plus pour parler de racines analytiques sans son domaine de convergence ou sans la réalisation du système de cylindres.

\input{sections/alpha_publicacion_completa.tex}

\clearpage\section{Prolongement exceptionnel : de l'incidence dodécaphasique à Moonshine}
\label{exc:seccion}

L'expression numérique n'épuise pas l'état qui l'a générée. Dans la construction précédente, APP conserve les deux feuillets et la division avec reste et quotient ; le TRIT détermine le régime local et son orientation ; le TPK compose sélection, transport, retenue, mémoire et retour. La lecture archimédienne contracte des cylindres compatibles. La lecture que nous étudions maintenant conserve en revanche l'incidence des positions et leur transport entre repères. Toutes deux partent du même état arithmétique avec mémoire de trajectoire. On ne reconstruit pas une géométrie exceptionnelle à partir des décimales finales de $\pi$, $\varphi$, $e$ ou $\alpha$.

Ce point permet de préciser le sens du prolongement
$\alpha/K\longrightarrow\text{incidence exceptionnelle}$. La flèche
n'affirme pas qu'un nombre réel isolé détermine un code, un réseau ou un
groupe. Elle transporte le registre dodécaphasique et les données signées qui sont
intervenues dans la détermination de $\alpha$, avec les mots et les
cadres issus d'APP--TRIT--TPK. En particulier, le registre ordonné
$K$ peut être retrouvé à partir de sa constante rationnelle $\kappa$ avec les
conventions de base, de longueur et d'origine déjà fixées. La projection sur un
support ne conserve en revanche que les incidences sélectionnées.
Pour sa part, $K_{\rm ph}$
désigne une autre lecture, à mémoire réduite, postérieure à l'holonomie
$\Gamma_9$.

L'exposé sépare trois opérations mathématiques : construire les objets
finis et leurs applications ; les prolonger en un réseau et en une algèbre d'opérateurs
de vertex ; reconnaître ensuite les objets obtenus au moyen des théorèmes
classiques pertinents. Les noms Paley, Witt, Mathieu, Golay, Leech et
Monstre désignent ces reconnaissances et les réalisations explicites
qui les rendent vérifiables. Ce ne sont pas des sélecteurs rétrospectifs des
germes ni des coefficients de la génération.

\input{sections/revision_incidencia_argumento.tex}



\subsection{Compatibilité des réalisations archimédienne et d'incidence}
\label{exc:doble-lectura}

Soit $X_n^{\rm cal}$ l'espace des préfixes admissibles de longueur $n$ de l'état arithmétique avec mémoire de trajectoire, avec les troncatures $p_{n,m}:X_n^{\rm cal}\to X_m^{\rm cal}$ pour $m\leq n$. L'exposant rappelle que le repère initial déclaré et son transport sont conservés, et non qu'un paramètre est ajusté à une valeur expérimentale. Nous écrivons $(w_j,f_j)$ pour le mot visible de six trits et le repère au pas $j$. Les repères satisfont une loi causale
\[
 f_{j+1}=g_j(x)f_j,
\]
où $g_j(x)$ dépend du préfixe déjà construit. En conséquence, une
prolongation n'altère pas les cadres antérieurs.

La chaîne
\[
 w_6\longrightarrow w_{12}\longrightarrow w_{18}
 \longrightarrow w_{24}\longrightarrow w_{30}
 \longrightarrow R_{36}\longrightarrow G_9
\]
est conservé dans cette lecture. Le retour de la phase ne réinitialise ni le
cadre ni la mémoire de l'état. De même, on maintient distincts le
recensement de $468$ émissions décimales, les $243$ mots visibles de six
trits et l'espace ambiant $\mathbb F_3^6$, de $729$ mots. L'espace ambiant sera
utilisé pour construire un code linéaire ; il ne se substitue pas au domaine
admissible des émissions.

La lecture archimédienne prend, pour un bloc $b$ de six trits, son
adresse entière $\nu(b)\in\{0,\ldots,728\}$. Un préfixe détermine
\[
 q_n=m+\sum_{j=0}^{n-1}\frac{\nu(b_j)}{729^{j+1}},
 \qquad I_n=[q_n,q_n+729^{-n}].
\]
La compatibilité des préfixes donne des cylindres emboîtés dont le diamètre tend
vers zéro. Leur classe de Cauchy appartient d'abord à la complétion interne
$\mathbb R_{\rm HMT}$ ; son identification ultérieure à la droite réelle
usuelle n'intervient pas dans le choix des blocs.

Pour la lecture d'incidence, nous introduisons, sur le même préfixe, l'application de graphe
\begin{equation}
 \gamma_A:\mathbb F_3^6\longrightarrow\mathbb F_3^{12},
 \qquad w\longmapsto(w,wA).
 \label{exc:gamma}
\end{equation}
Les mots s'écrivent comme des vecteurs lignes. Si $A^f=fAf^{-1}$, le changement de repère satisfait l'identité concrète
\begin{equation}
 \gamma_{A^f}(wf^{-1})
   =\gamma_A(w)(f^{-1}\oplus f^{-1}).
 \label{exc:covariancia}
\end{equation}
Le codomaine calibré conserve $f$ avec le mot codé.
Nous définissons alors
\begin{equation}
 Q_n(x)=\bigl((f_j,\gamma_{f_jA_Wf_j^{-1}}(w_j))\bigr)_{0\leq j<n}.
 \label{exc:qn}
\end{equation}

\begin{proposition}[Compatibilité et prolongement de la lecture d'incidence]
\label{exc:limite-inverso}
Si $r_{n,m}$ tronque les suites codées, on a
\[
 r_{n,m}Q_n=Q_m p_{n,m}.
\]
Les $Q_n$ déterminent donc une unique application $Q_\infty$ entre les limites inverses respectives. Dans chaque repère, l'application conserve exactement le mot visible, puisque $\operatorname{pr}_1\gamma_A=\operatorname{id}$.
\end{proposition}

\begin{proof}
Deux préfixes compatibles possèdent les mêmes mots et les mêmes
opérateurs de transport jusqu'à l'indice $m-1$. Ils possèdent, par récurrence,
les mêmes cadres à ces indices. Les $m$ premières composantes de
\eqref{exc:qn} coïncident, ce qui prouve le carré. La propriété
universelle de la limite inverse donne $Q_\infty$. L'affirmation de récupération
est la projection sur les six premières coordonnées de chaque graphe.
\end{proof}

Cette injectivité locale a un domaine précis. Elle ne démontre pas que
$Q_\infty$ récupère toute la mémoire profonde de l'état TPK lorsque seule
la suite des mots visibles et des repères a été retenue. Elle démontre encore moins
l'injectivité après le passage des mots aux supports.
En outre, un changement général de base dans $\operatorname{GL}_6(\mathbb F_3)$
ne conserve pas à lui seul le poids de Hamming ordinaire. Dans une coordinatisation
générale, on transporte aussi l'incidence et la métrique de la coordinatisation
initiale. Les calculs de supports qui suivent sont effectués dans la coordinatisation
de Paley déclarée ; les réétiquetages monomiaux conservent directement leur
interprétation en coordonnées.

\subsection{La réalisation finie et le code ternaire}
\label{exc:paley-codigo}

L'opérateur elliptique du TRIT admet la réalisation finie suivante dans
la coordinatisation à six positions. Celles-ci coordinatisent les six unités
$\{1,2,4,5,7,8\}$ modulo neuf, conservées par le transport modulaire.
Sur $\mathbb P^1(\mathbb F_5)$, une fois
l'ordre $\infty,0,1,2,3,4$ déclaré, soit $C$ la matrice de diagonale
nulle, avec $C_{\infty,a}=C_{a,\infty}=1$ et
$C_{a,b}=\chi(b-a)$ aux cinq positions finies. Ici,
$\chi(1)=\chi(4)=1$, $\chi(2)=\chi(3)=-1$ et $\chi(0)=0$.
Il s'agit d'une réalisation en coordonnées de la structure quadratique déjà
transportée ; l'identification des six positions à la droite
projective est une donnée de coordinatisation, et non une nouvelle entrée numérique dans le
générateur.

Les identités de la somme quadratique donnent $C^2=5I_6$. En réduisant modulo
trois, nous obtenons
\begin{equation}
 A_W=
 \begin{pmatrix}
 0&1&1&1&1&1\\
 1&0&1&2&2&1\\
 1&1&0&1&2&2\\
 1&2&1&0&1&2\\
 1&2&2&1&0&1\\
 1&1&2&2&1&0
 \end{pmatrix},
 \qquad A_W^2=2I_6=-I_6.
 \label{exc:aw}
\end{equation}
L'égalité quadratique réalise le régime elliptique ; elle n'identifie pas le
corps $\mathbb F_3$ à une réalisation archimédienne de l'unité
imaginaire. Ce sont des réalisations distinctes d'une relation opératoire
commune, avec leurs domaines respectifs.

\Needspace{8\baselineskip}
\begin{definition}[Code de la coordinatisation dodécaphasique]
\label{exc:codigo}
Le code ternaire associé à \eqref{exc:aw} est
\[
 C_W=\{(w,wA_W):w\in\mathbb F_3^6\}
       \subset\mathbb F_3^{12}.
\]
Son ensemble d'hexades est
\[
 \mathcal H=
 \{\operatorname{supp}(c):c\in C_W,
                           \operatorname{wt}(c)=6\}.
\]
\end{definition}

\begin{proposition}[Autodualité, distance et système de Witt]
\label{exc:codigo-witt}
Le code $C_W$ est autodual de paramètres $[12,6,6]_3$. Son énumérateur
des poids est
\begin{equation}
 \sum_{c\in C_W}z^{\operatorname{wt}(c)}
   =1+264z^6+440z^9+24z^{12}.
 \label{exc:enumerador}
\end{equation}
L'ensemble $\mathcal H$ possède $132$ éléments et constitue un système
$S(5,6,12)$ : chaque ensemble de cinq positions est contenu dans une
unique hexade.
\end{proposition}

\begin{proof}
La matrice génératrice $[I_6\mid A_W]$ est de rang six et, par symétrie de $A_W$, satisfait
\[
 [I_6\mid A_W][I_6\mid A_W]^{\mathsf T}
   =I_6+A_W^2=0.
\]
Le code est auto-orthogonal, de dimension égale à la moitié de sa longueur, et donc
autodual. L'évaluation finie des $3^6$ mots donne
\eqref{exc:enumerador} ; il s'agit d'une énumération sur la matrice explicite,
sans valeurs réelles ni cibles externes. En particulier, la distance
minimale est six.

Si deux mots de poids six ont le même support, parmi leurs six quotients de coordonnées non nulles, l'un des signes $1,-1$ apparaît au moins trois fois. En soustrayant le multiple correspondant, on obtiendrait un mot non nul de poids au plus trois, ce qui est impossible. Un support correspond donc exactement à la paire $\{c,-c\}$ et il y a $264/2=132$ supports. Si deux supports distincts partageaient cinq positions, le même argument annulerait au moins trois coordonnées d'une union de taille sept, produisant un poids au plus quatre. Un quintuplet appartient donc à au plus une hexade. Enfin,
\[
 132\binom65=792=\binom{12}5,
\]
de sorte qu'il appartient à une et une seule.
\end{proof}

La reconnaissance ultérieure est le code de Golay ternaire étendu et
le petit système de Witt \cite{conwaysloane}. La preuve précédente explique quel objet est
reconnu. Elle ne remplace pas la généalogie APP--TRIT--TPK par une recherche
parmi les codes connus.

\subsection{Drapeau dodécaphasique associé à \texorpdfstring{$K$ et $\alpha$}{K et alpha}}
\label{exc:bandera}

\input{sections/incidencia_registro.tex}

Dans le cadre hérité du registre dodécaphasique, les données sont
\begin{align}
 K={}&(234,543,140,729,659,824,621,\mathtt{058},914,794,146,601),
 \label{exc:k}\\
 w_\pi={}&010211, &w_e={}&201101,\nonumber\\
 w_\varphi={}&121200, &w_{\rm APP}={}&022020.\nonumber
\end{align}
En appliquant $\gamma_{A_W}$, on obtient les incidences suivantes :
\begin{align*}
 P_\pi&=\{2,4,5,6,7,8,9,10,11\},\\
 H_e&=\{1,3,4,6,9,10\},\\
 H_\varphi&=\{1,2,3,4,7,9\},\\
 H_{\rm APP}&=\{2,3,5,10,11,12\}.
\end{align*}
Le premier est le support d'un mot de poids neuf ; les trois autres
sont des hexades. Le registre fournit, par sa lecture de bloc haut,
\begin{equation}
 B_K=\{i:K_i\geq729\}=\{4,6,9,10\}=H_e\cap P_\pi.
 \label{exc:bk}
\end{equation}
L'entier $729$ indique ici une frontière du format de blocs
déclaré. Il n'est pas identifié par son seul cardinal à l'espace ambiant
ternaire utilisé dans \eqref{exc:codigo}. Dans ce registre précis, tous
les seuils entiers entre $660$ et $729$ produisent le même support :
le support ne restitue ni le seuil ni les douze blocs.

La voie arithmétique de $\alpha$ conserve, avant la retenue, le vecteur
signé
\begin{equation}
 Z_0=(7,297,353,-430,-715,-1199,-3,286,-894,-528,380,663).
 \label{exc:precarry}
\end{equation}
Son support négatif est
\begin{equation}
 H^-_\alpha=\{4,5,6,7,9,10\}.
 \label{exc:halpha}
\end{equation}
Ces signes ne sont pas extraits du développement décimal du scalaire
$\alpha$ ; ils sont conservés dans l'état qui précède son expression numérique.

\input{sections/incidencia_normalizacion.tex}

\begin{proposition}[Coïncidence des deux sélecteurs du drapeau]
\label{exc:selector-bandera}
Dans $\mathcal H$, le support \eqref{exc:halpha} est l'unique hexade $H$ qui satisfait
\[
 B_K\subset H\subset P_\pi,
 \qquad
 (|H\cap H_e|,|H\cap H_\varphi|,|H\cap H_{\rm APP}|)=(4,3,2).
\]
Ainsi, la lecture signée et la lecture d'incidence sélectionnent le même
drapeau $B_K\subset H^-_\alpha$.
\end{proposition}

\begin{proof}
Les quatre hexades qui contiennent $B_K$ s'obtiennent en ajoutant respectivement les paires
\[
 \{1,3\},\qquad\{2,11\},\qquad\{5,7\},\qquad\{8,12\}.
\]
Seules les deuxième et troisième paires sont contenues dans $P_\pi$. Pour la
deuxième, l'intersection avec $H_\varphi$ a une taille de trois et avec
$H_{\rm APP}$ une taille de trois ; elle ne satisfait pas la dernière condition. Pour la
troisième, les trois intersections ont les tailles $4,3,2$. Cette dernière
produit précisément le support négatif de \eqref{exc:precarry}.
\end{proof}

Il convient de conserver tant le drapeau que son cadre marqué :
\[
 \mathcal I_*=(B_K\subset H^-_\alpha),\qquad
 \widetilde{\mathcal I}_*
   =(\mathcal I_*;P_\pi,H_e,H_\varphi,H_{\rm APP}).
\]
Les opérations d'union, d'intersection et de différence sont équivariantes sous un réétiquetage simultané du plan combinatoire. C'est la classe marquée, et non les numéros des positions pris isolément, qui est l'objet transporté.

\begin{lemma}[Origine sélectionnée par l'incidence marquée]
\label{exc:origen}
Pour le cadre précédent,
\[
 o_*=(H_e\cap H_\varphi)
           \setminus(H^-_\alpha\cup H_{\rm APP})=\{1\}.
\]
La règle est équivariante sous tout automorphisme appliqué simultanément
aux données de $\widetilde{\mathcal I}_*$.
\end{lemma}

\begin{proof}
La première intersection est $\{1,3,4,9\}$. La réunion que l'on soustrait
contient $3,4,9$ et ne contient pas $1$. L'équivariance découle de ce qu'une
bijection commute avec les opérations ensemblistes.
\end{proof}

Cette origine sélectionnera une direction du réseau voisin lorsque la métrique $A_2$ et la chambre positive seront déclarées. L'incidence seule ne contient pas cette métrique. Le passage au réseau conservera explicitement les deux données.

\input{sections/incidencia_bandera.tex}

\subsection{Entrelacement isométrique des représentations d'incidence}
\label{exc:entrelazador}

L'incidence définit une application linéaire dont l'action peut être vérifiée
position par position :
\begin{equation}
 I:\mathbb R^{12}\longrightarrow\mathbb R^{\mathcal H},
 \qquad (Iv)(H)=\sum_{p\in H}v_p.
 \label{exc:incidencia}
\end{equation}
Soit $V_{11}=\mathbf1^\perp\subset\mathbb R^{12}$. Les deux réalisations
euclidiennes sont ici des langages ultérieurs pour le même plan fini.

\begin{proposition}[Isométrie de l'écran dodécaphasique]
\label{exc:isometria}
On a
\[
 I^{\mathsf T}I=36I_{12}+30J_{12},
\]
où $J_{12}$ est la matrice dont les entrées sont un. En conséquence,
\[
 U_W=\frac16 I\big|_{V_{11}}:
 V_{11}\xrightarrow{\ \sim\ }E_{\rm inc}:=I(V_{11})
\]
est une isométrie équivariante pour les actions par permutation des
automorphismes du système.
\end{proposition}

\begin{proof}
Chaque point appartient à
$r=\binom{11}{4}/\binom54=66$ hexades, et chaque paire à
$\lambda_2=\binom{10}{3}/\binom43=30$. Ce sont, respectivement, les
entrées diagonales et non diagonales de $I^{\mathsf T}I$. Sur
$\mathbf1^\perp$, le terme $J_{12}$ s'annule et il reste $36I$.
Enfin, si $g$ est un automorphisme, $p\in H$ équivaut à
$g(p)\in g(H)$ ; d'où $I\rho_{12}(g)=\rho_{\mathcal H}(g)I$.
\end{proof}

L'information d'incidence possède également une lecture spectrale concrète. Définissons
\[
 G_{H,H'}=\binom{|H\cap H'|}{4}.
\]

\Needspace{8\baselineskip}
\begin{lemma}[Action du noyau d'intersection]
\label{exc:espectro}
Si $J_{132,12}$ est la matrice rectangulaire dont tous les coefficients valent un,
\[
 GI=30I+15J_{132,12}.
\]
Ainsi, $G$ agit par le scalaire $30$ sur $E_{\rm inc}$.
\end{lemma}

\begin{proof}
L'entrée $(H,p)$ de $GI$ compte les couples $(T,H')$ tels que
$T\subset H\cap H'$, $|T|=4$ et $p\in H'$. Si $p\in H$, il y a dix
quadruplets $T$ contenant $p$, chacun contenu dans quatre hexades,
et cinq qui ne le contiennent pas, chacun ayant une unique extension du
quintuplet $T\cup\{p\}$. Le total est $10\cdot4+5=45$.
Si $p\notin H$, les quinze quadruplets de $H$ donnent chacun une
extension : le total est $15$. D'où l'identité. Le terme rectangulaire
s'annule sur $V_{11}$.
\end{proof}

Il existe bien ici un opérateur d'entrelacement matériel : son domaine, son codomaine, ses entrées et sa normalisation sont déterminés. Par exemple, tout projecteur orthogonal $P$ déjà construit sur $V_{11}$ se transporte en $Q=U_WPU_W^*$ sur $E_{\rm inc}$, avec $U_WP=QU_W$. Cette affirmation n'identifie pas par leur seule trace tous les opérateurs qui agissent dans ces espaces.

\subsection{L'étoile de Witt et la génération de \texorpdfstring{$M_{12}$}{M12}}
\label{exc:estrella}

\begin{definition}[Étoile sur un quadruplet]
\label{exc:def-estrella}
Pour $B\in\binom{[12]}4$, l'étoile de Witt est la collection des
quatre hexades qui contiennent $B$. Ses pétales sont les paires $H\setminus B$.
La permutation $\sigma_B$ fixe $B$ et échange les deux points de chaque
pétale.
\end{definition}

L'existence de quatre hexades se déduit de
\[
 \lambda_4=\frac{\binom81}{\binom21}=4.
\]
Deux pétales ne peuvent pas partager un point, car deux hexades contiendraient alors le même quintuplet. Ce sont donc quatre paires disjointes qui partitionnent $[12]\setminus B$. Cette description n'est ni une étoile à cinq branches ni une métaphore géométrique : c'est une incidence quatre-à-un du plan combinatoire.

\begin{figure}[htbp]
\centering
\begin{tikzpicture}[x=1cm,y=1cm, every node/.style={font=\small}]
 \node[draw,rounded corners,minimum width=3.8cm,minimum height=.7cm]
  (b) at (0,0) {$B_K=\{4,6,9,10\}$};
 \node[draw,minimum width=2.7cm,minimum height=.65cm]
  (p1) at (-4.6,-1.7) {$\{1,3\}$};
 \node[draw,minimum width=2.7cm,minimum height=.65cm]
  (p2) at (-1.55,-1.7) {$\{2,11\}$};
 \node[draw,very thick,minimum width=2.7cm,minimum height=.65cm]
  (p3) at (1.55,-1.7) {$\{5,7\}$};
 \node[draw,minimum width=2.7cm,minimum height=.65cm]
  (p4) at (4.6,-1.7) {$\{8,12\}$};
 \draw (b) -- (p1); \draw (b) -- (p2);
 \draw[very thick] (b) -- (p3); \draw (b) -- (p4);
 \node at (-4.6,-2.35) {$H_e$};
 \node at (-1.55,-2.35) {$B_K\cup\{2,11\}$};
 \node at (1.55,-2.35) {$H^-_\alpha$};
 \node at (4.6,-2.35) {$B_K\cup\{8,12\}$};
 \node at (0,-3.05)
  {$\sigma_{B_K}=(1\ 3)(2\ 11)(5\ 7)(8\ 12)$};
\end{tikzpicture}
\caption{L'étoile concrète du drapeau dodécaphasique. Chaque ligne
indique l'union du quadruplet central avec la paire à son extrémité ; la
ligne épaisse signale l'hexade également sélectionnée par les signes
antérieurs à la retenue de $\alpha$. Les quatre paires épuisent les huit
positions extérieures.}
\label{exc:fig-estrella}
\end{figure}

\begin{proposition}[La collection d'étoiles]
\label{exc:m12}
Les $\binom{12}4=495$ permutations $\sigma_B$ préservent $\mathcal H$.
Le groupe qu'elles engendrent est d'ordre $95040$ et est le groupe de Mathieu
$M_{12}$ dans son action sur le petit système de Witt.
\end{proposition}

\begin{proof}
L’affirmation de préservation est une vérification finie sur les
$132$ hexades de \eqref{exc:codigo} : pour chaque quadruplet, on construit
les quatre pétales et l’on vérifie
$\{\sigma_B(H):H\in\mathcal H\}=\mathcal H$.
La clôture par composition de ces permutations produit $95040$
éléments. Dans l’énumération lexicographique, il suffit d’ajouter successivement
six étoiles qui ne sont pas encore contenues dans le groupe ; les ordres intermédiaires
sont $2,6,18,360,7920,95040$, et les $495$ étoiles appartiennent au groupe
final. Le procédé et les six générateurs sont explicités dans
la note technique de provenance. Enfin, on applique la reconnaissance
du groupe d’automorphismes de $S(5,6,12)$ comme $M_{12}$, dont l’ordre est
$12\cdot11\cdot10\cdot9\cdot8=95040$.
\end{proof}

Une étoile individuelle engendre un groupe cyclique d'ordre deux. Elle
n'engendre pas à elle seule $M_{12}$, et encore moins le Monstre. L'
existence de $495$ générateurs d'un groupe de permutations ne prouve pas non plus que
chacun soit le transport d'un lacet spécifique du TPK. Cette affirmation
requiert l'application de lacets correspondante ; le résultat établi ici
est l'action d'incidence finie.

\begin{remark}[L'égalité des traces n'identifie pas les opérateurs]
Sur $V_{11}$, une étoile a pour spectre $+1$ avec multiplicité sept
et $-1$ avec multiplicité quatre. Sa trace normalisée est $3/11$.
Un projecteur de rang trois sur le même espace possède également une trace
normalisée $3/11$, mais un spectre $1^3,0^8$. Ils ne sont pas conjugués : leurs
polynômes minimaux et leurs spectres sont différents. L'opérateur d'entrelacement
\eqref{exc:incidencia} transporte des opérateurs définis ; il ne transforme pas
cette coïncidence scalaire en une équivalence d'opérateurs.
\end{remark}

\subsection{Le résidu binaire et les cinq régions de \texorpdfstring{$\pi$}{pi}}
\label{exc:cinco-regiones}

La reconnaissance suivante utilise le code de Golay binaire
étendu $G_{24}$, de paramètres $[24,12,8]_2$. Ses poids sont
$0,8,12,16,24$, et les supports de poids huit forment le plan
$S(5,8,24)$. Nous fixons une dodécade $D$, support d'un mot de poids
douze. Cette donnée et l'isomorphisme d'incidence qui sera déclaré
ci-après appartiennent à la coordinatisation de réalisation ; ils ne sont pas utilisés pour
engendrer ni sélectionner les décimales de $\pi$.

\begin{proposition}[Résidu d'une dodécade et élévation d'hexades]
\label{exc:residuo-binario}
L'ensemble
\[
 \mathcal H(D)=
 \{O\cap D: O\text{ est une octade},\ |O\cap D|=6\}
\]
est un système $S(5,6,12)$ sur $D$. Chacune de ses hexades a une unique élévation en une octade. Pour tout quadruplet $B\subset D$, les cinq octades qui le contiennent se répartissent en quatre élévations résiduelles et une unique octade $O_B^\perp$ avec $O_B^\perp\cap D=B$.
\end{proposition}

\begin{proof}
Un quintuplet $A\subset D$ appartient à une unique octade $O$. Si
$j=|O\cap D|$, le mot binaire somme de $O$ et $D$ a pour poids
$20-2j$. Parmi les poids permis, et avec $0\leq j\leq8$, seuls
$j=2,4,6$ sont possibles. Comme $A\subset O\cap D$, on a $j\geq5$
et donc $j=6$. Cela donne l'existence et l'unicité de l'hexade
résiduelle, ainsi que l'unicité de son élévation en choisissant cinq de ses
points. D'autre part,
\[
 \lambda_4(S(5,8,24))=\frac{\binom{20}1}{\binom41}=5,
 \qquad
 \lambda_4(S(5,6,12))=4.
\]
Les quatre hexades résiduelles sur $B$ donnent quatre octades distinctes. La cinquième ne peut pas couper $D$ en six points, car elle serait une autre hexade résiduelle ; puisqu'elle contient $B$, son intersection est de taille quatre et coïncide avec $B$.
\end{proof}

Une coordinatisation $\ell:[12]\to D$ qui préserve les hexades identifie
$\mathcal H$ à $\mathcal H(D)$. Son existence est la reconnaissance
du plan déjà construit comme le petit plan de Witt ; un choix
précis de $\ell$ fixe le réétiquetage. Le relèvement
\[
 \iota_D:\mathcal H\longrightarrow
       \{\text{octades de }G_{24}\},
 \qquad H\longmapsto O\text{ tel que }O\cap D=\ell(H),
\]
est injective parce que l'intersection avec $D$ redonne $\ell(H)$. Cette injectivité a pour domaine les hexades, et non le registre $K$.

La relation avec les cinq régions de $\pi$ conserve des données plus riches
que le nombre cinq. Les émissions et leurs signatures sont :
\begin{center}
\small
\begin{tabular}{@{}crll@{}}
\toprule
Région & Multiplicité & Émission $U_6$ & Signature \\
\midrule
$\perp$ & $432$ & $(501,614,498,169,272,272)$ & $555555$ \\
$++_1$ & $144$ & $(810,923,870,169,272,272)$ & $339555$ \\
$++_2$ & $144$ & $(810,983,810,169,272,272)$ & $393555$ \\
$++_3$ & $144$ & $(870,923,810,169,272,272)$ & $933555$ \\
$--$ & $144$ & $(870,923,810,418,674,521)$ & $933717$ \\
\bottomrule
\end{tabular}
\end{center}
Dans chaque bloc décimal de trois chiffres, la signature utilise la différence
modulaire des deux premiers chiffres. Les cinq émissions ont le
même mot tritique visible $w_\pi=010211$, mais pas la même
généalogie régionale. Les trois régions positives conservent la queue
$169|272|272$ ; la région négative la remplace par $418|674|521$ ;
la région transversale a une signature constante et une multiplicité supérieure.

L'extension marquée de la coordinatisation régionale envoie la région transversale
sur $O_{\ell(B_K)}^\perp$, la région négative sur $\iota_D(H^-_\alpha)$ et les
trois régions positives sur les trois autres relèvements de l'étoile. Ainsi se
réalise la partition $4+1$ des cinq octades sur $\ell(B_K)$.
La correspondance individuelle entre les trois noms positifs et
les trois pétales restants exige un ordre de coordinatisation ; l'incidence ne la détermine
pas sans cette donnée. L'énoncé précis est donc une
correspondance marquée des régions, compatible avec leurs rôles et avec
le drapeau, et non une bijection canonique déduite du cardinal cinq.

L'identité $432+4\cdot144=1008$ conserve les multiplicités
régionales. On a aussi $1008=36\cdot28$, et une octade possède
$\binom82=28$ paires. Ces égalités sont des contrôles arithmétiques du
catalogue, et non des substituts à l'application $\iota_D$ ni à la coordinatisation régionale.

\input{sections/estrella_pentadica_recuperada.tex}

\subsection{Du code ternaire au réseau de Niemeier}
\label{exc:niemeier}

Le passage à la dimension vingt-quatre ne nécessite pas de parcourir d'abord le
code binaire. Il s'effectue directement à partir de $C_W$ par recollement des
discriminants. On distingue ainsi deux prolongations compatibles :
l'une conserve le résidu $W_{12}\subset W_{24}$ ; l'autre construit le
réseau sur lequel sera réalisée l'algèbre d'opérateurs de vertex.

Soit $A_2=\mathbb Z a_1\oplus\mathbb Z a_2$ de matrice de Gram
\[
 \begin{pmatrix}2&-1\\-1&2\end{pmatrix},
 \qquad \rho=a_1+a_2,\qquad \rho^2=2.
\]
Nous représentons les trois classes de $A_2^*/A_2$ par
\[
 \varpi_0=0,\qquad
 \varpi_1=\frac{2a_1+a_2}{3},\qquad
 \varpi_2=\frac{a_1+2a_2}{3}.
\]
L'application $j\mapsto\varpi_j+A_2$ identifie additivement
$\mathbb F_3$ à ce discriminant. Pour $c\in C_W$, nous écrivons
$\phi(c)=(\varpi_{c_1},\ldots,\varpi_{c_{12}})$ et définissons
\begin{equation}
 N(C_W)=\bigcup_{c\in C_W}\bigl(A_2^{12}+\phi(c)\bigr).
 \label{exc:niemeier-def}
\end{equation}

\begin{theorem}[Pegado ternario]
\label{exc:niemeier-teorema}
L'objet \eqref{exc:niemeier-def} est un réseau positif, pair et
unimodulaire de rang $24$. Son système de racines est exactement
$A_2^{12}$ ; il est donc reconnu comme le réseau de Niemeier ayant
ce système de racines.
\end{theorem}

\begin{proof}
La linéarité du code rend l'union de classes stable par addition. L'auto-orthogonalité annule l'accouplement discriminant : celui-ci est un multiple de $c\cdot d/3$ modulo $\mathbb Z$. Le réseau est donc entier. La norme d'une classe de recollement est $2\operatorname{wt}(c)/3$ modulo $2\mathbb Z$. Les poids de \eqref{exc:enumerador} sont des multiples de trois, de sorte qu'il est pair.

L'indice de $A_2^{12}$ dans $N(C_W)$ est $|C_W|=3^6$. En conséquence,
\[
 \det N(C_W)=\frac{3^{12}}{(3^6)^2}=1.
\]
La positivité procède de la somme orthogonale des douze métriques $A_2$.
Dans une classe non nulle de $A_2^*/A_2$, la norme minimale est $2/3$ ;
tout mot non nul possède au moins six positions non nulles.
Par conséquent, un vecteur d'une classe de recollement non triviale a une norme
au moins égale à $6\cdot2/3=4$. Aucune racine de norme deux n'apparaît hors de
$A_2^{12}$. À l'intérieur de celui-ci, les racines sont précisément celles de ses
douze composantes.
\end{proof}

\subsection{Sélection par incidence du voisin unimodulaire sans racines}
\label{exc:leech}

\input{sections/incidencia_reticulo.tex}

L'étiquette $o_*=1$ du drapeau se réalise maintenant comme l'une des
douze composantes $A_2$. Cette opération déclare la métrique précédente et
la chambre radiale positive
\[
 v(a)=\sum_{i=1}^{12}a_i\rho_i,
 \qquad a_i=1+3b_i,\quad b_i\in\mathbb Z_{\geq0}.
\]
L'exigence de parité du voisin d'indice trois est $v(a)^2\equiv0\pmod{18}$. Puisque
\[
 v(a)^2=2\sum_i(1+3b_i)^2,
\]
cette congruence équivaut à $\sum_i b_i\equiv1\pmod3$. La plus petite
valeur de la norme dans la chambre s'obtient avec un unique $b_i=1$ et
tous les autres nuls ; elle vaut $54$. L'origine marquée sélectionne
\begin{equation}
 v_\alpha=4\rho_1+\rho_2+\cdots+\rho_{12},
 \qquad v_\alpha^2=54.
 \label{exc:vector}
\end{equation}
Le mot ``plus petit'' se rapporte à cette chambre positive. Il n'affirme pas la minimalité dans toute la classe résiduelle : le vecteur $(a_1-2a_2,\rho,\ldots,\rho)$ représente la même classe modulo trois et a pour norme $14+11\cdot2=36$. La restriction de chambre fait partie des données de la construction, et n'est pas une conclusion que l'on puisse omettre.

Soit $N=N(C_W)$. Nous définissons
\begin{equation}
 N_{\alpha,0}=\{x\in N:\langle x,v_\alpha\rangle\equiv0\pmod3\},
 \qquad
 N_\alpha=N_{\alpha,0}+\mathbb Z\frac{v_\alpha}{3}.
 \label{exc:vecino}
\end{equation}

\begin{theorem}[Voisin de Leech déterminé par le drapeau marqué]
\label{exc:teorema-leech}
Avec la métrique et la chambre déclarées, $N_\alpha$ est un réseau
positif, pair et unimodulaire de rang $24$, sans vecteurs de norme deux.
D'après le théorème de classification correspondant, il est isométrique au
réseau de Leech $\Lambda_{24}$.
\end{theorem}

\begin{proof}
Le caractère $x\mapsto\langle x,v_\alpha\rangle\bmod3$ est non nul :
une racine simple d'une composante s'apparie avec $a_i\rho_i$ en
$a_i\equiv1\pmod3$. Par conséquent, $[N:N_{\alpha,0}]=3$ et
$\det N_{\alpha,0}=9$. De plus,
$v_\alpha/3\in N_{\alpha,0}^*$ et $v_\alpha\in N_{\alpha,0}$.
Si $v_\alpha/3$ appartenait à $N$, l'intégralité de $N$ rendrait
tous les appariements avec $v_\alpha$ divisibles par trois,
ce qui contredirait le caractère non nul. La classe de $v_\alpha/3$ est
donc d'ordre trois. On en déduit
\[
 [N_\alpha:N_{\alpha,0}]=3,
 \qquad \det N_\alpha=9/3^2=1.
\]
L'extension est entière, et pour $x\in N_{\alpha,0}$,
\[
 \left(x+m\frac{v_\alpha}{3}\right)^2
   =x^2+2m\frac{\langle x,v_\alpha\rangle}{3}+6m^2
   \in2\mathbb Z.
\]
Elle est donc paire.

Il reste à exclure les racines, ce qui est l'étape décisive. Les racines de $N$ appartiennent à $A_2^{12}$. Leurs accouplements avec $v_\alpha$ sont congrus à $\pm1$ ou à $\pm2$ modulo trois ; aucune n'appartient à $N_{\alpha,0}$. Pour les deux autres classes du voisin, chaque composante appartient à l'une des classes locales
\[
 A_2+\varpi_j\pm\rho/3,\qquad j=0,1,2.
\]
Comme $4\rho/3\equiv\rho/3\pmod{A_2}$, cette description vaut
également dans la composante distinguée. En écrivant un vecteur local
sous la forme $(u a_1+v a_2)/3$, sa norme est
$2(u^2-uv+v^2)/9$. Les six classes précédentes ont des résidus non
nuls et un minimum $2/9$ : les représentants de forme quadratique minimale
sont $(\pm1,0),(0,\pm1),(1,1),(-1,-1)$ dans les résidus appropriés.
Un vecteur d'une classe nouvelle a donc une norme au moins égale à
$12\cdot2/9=8/3>2$. Aucune des trois classes ne contient de racines.
La classification des réseaux pairs unimodulaires positifs de
rang vingt-quatre identifie l'unique cas sans racines à
$\Lambda_{24}$.
\end{proof}

L'argument ne reconnaît pas Leech par une coïncidence de dimension.
Il construit le recollement, prouve l'intégralité, la parité et le déterminant, et
élimine les racines au moyen d'une borne locale. Ce n'est qu'ensuite qu'il applique le
théorème de classification. La dépendance à l'égard de $\alpha$ est conservée dans
le drapeau qui sélectionne $o_*$ et dans le vecteur \eqref{exc:vector} ;
on ne transforme pas la valeur scalaire de la constante en une longueur du
réseau.

\subsection{L'algèbre d'opérateurs de vertex et le Monstre}
\label{exc:voa}

Soit $\Lambda=N_\alpha$, déjà reconnu comme le réseau de Leech. L'étape suivante utilise son réseau entier, et non un développement décimal. La construction de l'algèbre d'opérateurs de vertex de réseau s'écrit
\[
 V_\Lambda=M(1)\otimes\mathbb C_\varepsilon[\Lambda].
\]
Ici, $M(1)$ est l'espace des oscillateurs de Heisenberg en rang vingt-quatre ; $\mathbb C_\varepsilon[\Lambda]$ est l'algèbre de groupe tordue par le cocycle de l'extension centrale du réseau. Le cocycle et le relèvement des isométries font partie de la construction : une somme vectorielle nue ne détermine pas les produits de vertex.

La négation $\lambda\mapsto-\lambda$ admet le relèvement standard
$\theta$. Avec son module tordu et le choix de parité compatible,
la construction de Frenkel--Lepowsky--Meurman produit
\begin{equation}
 V^\natural_{\rm HMT}
   =V_\Lambda^+\oplus(V_\Lambda^T)^+.
 \label{exc:flm}
\end{equation}
Le signe $+$ désigne la partie paire de l'action relevée. L'expression
inclut le produit de vertex de l'orbifold ; elle ne se limite pas à une somme
d'espaces gradués.

\begin{theorem}[Prolongement jusqu'au module Moonshine]
\label{exc:moonshine}
L'algèbre \eqref{exc:flm}, construite sur le voisin
\eqref{exc:vecino}, est une réalisation de l'algèbre Moonshine
$V^\natural$. Elle a une charge centrale $24$, une composante de poids un nulle
et un groupe d'automorphismes isomorphe au Monstre $\mathbb M$. Son caractère
gradué est
\begin{equation}
 \operatorname{Tr}_{V^\natural_{\rm HMT}}q^{L_0-1}
   =J(\tau)=j(\tau)-744
   =q^{-1}+196884q+21493760q^2+\cdots,
 \qquad q=e^{2\pi i\tau}.
 \label{exc:j}
\end{equation}
\end{theorem}

\begin{proof}
Le théorème \ref{exc:teorema-leech} vérifie les hypothèses sur le réseau :
positivité, parité, unimodularité, rang vingt-quatre et absence de
racines. On applique alors la construction et le théorème de
Frenkel--Lepowsky--Meurman pour le réseau de Leech. Ceux-ci identifient
l'orbifold à $V^\natural$, établissent son caractère et son groupe d'
automorphismes. La composition du drapeau, du voisin et du foncteur
de construction de VOA donne la réalisation indiquée ici. Le groupe ne se déduit
pas du coefficient $196884$.
\end{proof}

Le théorème utilisé est une reconnaissance classique ultérieure dont les hypothèses sont vérifiées, et non une nouvelle démonstration dans cet article de la construction FLM \cite{flm}. En outre, la notation $q=e^{2\pi i\tau}$ est la variable conventionnelle de Fourier employée pour reconnaître le caractère. Elle n'introduit ni $\pi$ ni $e$ comme entrées du générateur coinductif qui les a déjà déterminés.

Une partie du mécanisme de trace peut être observée avant l'orbifold.
Seul le vecteur nul du réseau est fixé par la négation ; chaque
oscillateur contribue avec un signe alterné. D'où
\[
 \operatorname{Tr}_{V_\Lambda}\theta q^{L_0-1}
   =t_2(\tau)
   :=q^{-1}\prod_{n\geq1}(1+q^n)^{-24}.
\]
Cette trace dans l'algèbre de réseau ne doit pas être confondue avec le
caractère sans insertion de $V^\natural$. Elle n'identifie pas non plus à elle
seule une classe de conjugaison du Monstre : l'espace, l'insertion
et le relèvement doivent être spécifiés avant de comparer les fonctions.

\input{sections/incidencia_automorfismos.tex}

\subsection{Traces graduées de Moonshine et représentation du Monstre}
\label{exc:alcance}

Pour $g\in\operatorname{Aut}(V^\natural_{\rm HMT})$, la trace avec insertion est
\[
 T_g(\tau)=
 \operatorname{Tr}_{V^\natural_{\rm HMT}}gq^{L_0-1}.
\]
Le théorème du Moonshine monstrueux identifie ces séries aux
Hauptmoduln des groupes de genre zéro correspondants
\cite[théorème~1.1]{borcherds}.
Le caractère $J$, le groupe $\mathbb M$ et la collection $g\mapsto T_g$
sont des réalisations de l'algèbre déjà construite. Ils ne forment pas une suite
causale dans laquelle on devinerait d'abord $J$ pour en déduire ensuite le Monstre.

L'unité possède ici une localisation algébrique précise. Le sous-espace de poids zéro est la droite du vide, \(V^\natural_0=\CC\mathbf1\), et le coefficient de \(q^{-1}\) dans \(J\) est donc \(1\). En poids un, la dimension est zéro, d'où l'absence du terme constant dans \(J=j-744\). Un morphisme d'algèbres de vertex conserve le vecteur du vide. Cette propriété fournit le sens technique de la conservation de l'unité à ce niveau ; l'unité des observables cylindriques est conservée par la précomposition démontrée précédemment. La relation entre les deux niveaux doit suivre les applications de construction avec leurs unités.

En poids deux apparaît l'égalité
\[
 196884=1+196883=196560+324=54(5\cdot729+1).
\]
La première séparation correspond à la droite conforme triviale et à la
composante irréductible de dimension $196883$ de la représentation
du Monstre. Les autres expressions ont des lectures en termes de réseaux ou
d'arithmétique. Leur égalité n'identifie pas automatiquement leurs termes
à des régions, des mots ou des routes du TPK. Pour ce faire, il faudrait
présenter les applications entre ces espaces et prouver qu'elles respectent les
actions pertinentes.

En particulier, une isométrie $N_\alpha\simeq\Lambda_{24}$ induit
une réalisation de \eqref{exc:flm}, et un choix d'isomorphisme
avec $V^\natural$ transporte son groupe d'automorphismes. Cela n'a pas
construit une action du Monstre sur les chiffres de
$\pi$, sur les blocs de $K$ ni sur toutes les histoires TPK.
Une telle affirmation supplémentaire requerrait une représentation
$\rho_{\rm TPK}$ et une application $F$ avec un carré explicite
\[
 F\rho_{\rm TPK}(g)=\rho_{V^\natural}(g)F,
\]
en incluant le domaine, le codomaine et l'information conservée. L'absence
de cette étape dans l'affirmation présente ne retire ni l'opérateur d'entrelacement
d'incidence $U_W$, ni le drapeau ni le voisin déjà construits ; elle
délimite seulement une promotion différente, celle de l'action sur l'état source.

La chaîne établie dans cette section peut se résumer, en gardant visibles ses données de transport, comme suit
\begin{align*}
 &\text{APP--TRIT--TPK et état arithmétique avec mémoire de trajectoire}
   \longrightarrow (w_j,f_j;K,Z_0)\\
 &\longrightarrow (C_W,\mathcal H;B_K\subset H^-_\alpha,
                   P_\pi,H_e,H_\varphi,H_{\rm APP})\\
 &\longrightarrow \bigl(N(C_W),o_*,\text{métrique et chambre}\bigr)
   \longrightarrow N_\alpha
   \longrightarrow V^\natural_{\rm HMT}.
\end{align*}
Au niveau de l'incidence, les étoiles engendrent $M_{12}$ et la
coordinatisation résiduelle réalise la structure régionale $4+1$ dans $W_{24}$.
Au niveau des réseaux et des VOA, la construction donne Leech et le module
Moonshine. Ce sont des prolongements d'une même source marquée, avec
des objets et des applications différents. La continuité causale n'oblige pas à
identifier leurs codomaines ni à effacer les choix explicites
de repère, de chambre et de relèvement.

Le résultat de la section est donc une dérivation réunie :
les preuves finies et réticulaires rendent effectif le passage depuis
l’incidence dodécaphasique jusqu’aux hypothèses des théorèmes de
reconnaissance. Ce résultat se rapporte aux constructions et
aux hypothèses spécifiées ; il n’ajoute pas de validation expérimentale
à la comparaison des constantes.

\endgroup

% Dependencias literales del artículo II, revisión 05.
% Compilar desde la raíz III. Los rótulos técnicos se distinguen sin editar fuentes.
% Las tres entradas de BIBITEMS_ADICIONALES_II.tex van dentro de la bibliografía final.
\begingroup
\makeatletter
\def\input@path{{base_articulo_II/}{base_articulo_I/}}
\makeatother
\providecommand{\HMT}{\mathrm{HMT}}
\providecommand{\Tr}{\operatorname{Tr}}
\providecommand{\dd}{\mathrm d}
\providecommand{\R}{\mathbb R}
\providecommand{\Z}{\mathbb Z}
\definecolor{ink}{HTML}{183D4A}
\definecolor{accent}{HTML}{8C3D2E}
\ifcsname apunte\endcsname\else
\newenvironment{apunte}[1]{\par\Needspace{6\baselineskip}\begin{quote}\small\raggedright\noindent\textbf{Note de révision: #1.} }{\end{quote}}
\fi
\NewCommandCopy\HMTIILabelOriginal\label
\NewCommandCopy\HMTIIRefOriginal\ref
\makeatletter
\let\HMTIIMathLabelOriginal\label@in@display
\def\label@in@display#1{\HMTIIMathLabelOriginal{hmtII:#1}}
\makeatother
\expandafter\def\csname HMTIILabel@eq:cycle\endcsname{1}
\expandafter\def\csname HMTIILabel@eq:det\endcsname{1}
\expandafter\def\csname HMTIILabel@eq:el-cola-regional\endcsname{1}
\expandafter\def\csname HMTIILabel@eq:el-cpi\endcsname{1}
\expandafter\def\csname HMTIILabel@eq:el-da\endcsname{1}
\expandafter\def\csname HMTIILabel@eq:el-decada\endcsname{1}
\expandafter\def\csname HMTIILabel@eq:el-eta\endcsname{1}
\expandafter\def\csname HMTIILabel@eq:el-h4\endcsname{1}
\expandafter\def\csname HMTIILabel@eq:el-h5\endcsname{1}
\expandafter\def\csname HMTIILabel@eq:el-lector-determinantal\endcsname{1}
\expandafter\def\csname HMTIILabel@eq:el-ratio\endcsname{1}
\expandafter\def\csname HMTIILabel@eq:el-secciones-accion\endcsname{1}
\expandafter\def\csname HMTIILabel@eq:eta\endcsname{1}
\expandafter\def\csname HMTIILabel@eq:events\endcsname{1}
\expandafter\def\csname HMTIILabel@eq:gamma\endcsname{1}
\expandafter\def\csname HMTIILabel@eq:ii-accion-cartas\endcsname{1}
\expandafter\def\csname HMTIILabel@eq:ii-barbero-cartas\endcsname{1}
\expandafter\def\csname HMTIILabel@eq:ii-canales\endcsname{1}
\expandafter\def\csname HMTIILabel@eq:ii-elipse-area-fase\endcsname{1}
\expandafter\def\csname HMTIILabel@eq:ii-elipse-bohr\endcsname{1}
\expandafter\def\csname HMTIILabel@eq:ii-elipse-fases\endcsname{1}
\expandafter\def\csname HMTIILabel@eq:ii-elipse-focal\endcsname{1}
\expandafter\def\csname HMTIILabel@eq:ii-elipse-impedancia\endcsname{1}
\expandafter\def\csname HMTIILabel@eq:ii-elipse-lc\endcsname{1}
\expandafter\def\csname HMTIILabel@eq:ii-elipse-radio-polar\endcsname{1}
\expandafter\def\csname HMTIILabel@eq:ii-elipse-reciproca\endcsname{1}
\expandafter\def\csname HMTIILabel@eq:ii-par-angular\endcsname{1}
\expandafter\def\csname HMTIILabel@eq:minus12\endcsname{1}
\expandafter\def\csname HMTIILabel@eq:revision-h5-coeficientes\endcsname{1}
\expandafter\def\csname HMTIILabel@eq:revision-renovacion\endcsname{1}
\expandafter\def\csname HMTIILabel@fig:ii-elipse-calculada\endcsname{1}
\expandafter\def\csname HMTIILabel@intsuccpass:eq:c12-fundamental-chain\endcsname{1}
\expandafter\def\csname HMTIILabel@intsuccpass:eq:dodecaphase-forces-12minus1\endcsname{1}
\expandafter\def\csname HMTIILabel@mem:balance\endcsname{1}
\expandafter\def\csname HMTIILabel@mem:identidad-finita\endcsname{1}
\expandafter\def\csname HMTIILabel@prop:ii-desplazamiento-angular\endcsname{1}
\expandafter\def\csname HMTIILabel@prop:ii-elipse-area\endcsname{1}
\expandafter\def\csname HMTIILabel@prop:ii-elipse-bohr\endcsname{1}
\expandafter\def\csname HMTIILabel@prop:ii-elipse-focos\endcsname{1}
\expandafter\def\csname HMTIILabel@prop:ii-elipse-lc\endcsname{1}
\expandafter\def\csname HMTIILabel@prop:ii-inversion-canales-angulares\endcsname{1}
\expandafter\def\csname HMTIILabel@sec:action\endcsname{1}
\expandafter\def\csname HMTIILabel@sec:barbero\endcsname{1}
\expandafter\def\csname HMTIILabel@sec:ellipse\endcsname{1}
\expandafter\def\csname HMTIILabel@sec:eta\endcsname{1}
\expandafter\def\csname HMTIILabel@sec:ii-angulos\endcsname{1}
\expandafter\def\csname HMTIILabel@sec:ii-elipse-bohr\endcsname{1}
\expandafter\def\csname HMTIILabel@sec:ii-elipse-fase-area\endcsname{1}
\expandafter\def\csname HMTIILabel@sec:ii-elipse-geometria\endcsname{1}
\expandafter\def\csname HMTIILabel@sec:ii-elipse-lc\endcsname{1}
\expandafter\def\csname HMTIILabel@sec:revision-planck\endcsname{1}
\expandafter\def\csname HMTIILabel@sec:revision-planck-contraste\endcsname{1}
\expandafter\def\csname HMTIILabel@sec:sello\endcsname{1}
\expandafter\def\csname HMTIILabel@subsec:ii-memoria-resolvente\endcsname{1}
\renewcommand{\label}[1]{\HMTIILabelOriginal{hmtII:#1}}
\RenewDocumentCommand{\ref}{s m}{\ifcsname HMTIILabel@#2\endcsname\IfBooleanTF{#1}{\HMTIIRefOriginal*{hmtII:#2}}{\HMTIIRefOriginal{hmtII:#2}}\else\IfBooleanTF{#1}{\HMTIIRefOriginal*{#2}}{\HMTIIRefOriginal{#2}}\fi}
\clearpage\section{Mise à jour dodécaphasique et mémoire exprimée}
\label{sec:eta}
Soit $d_t$ le code entier de direction conservé par une trace TPK. Ce code et la direction cardinale du curseur sont des coordonnées distinctes. Les douze fenêtres $E_m=\{9m+1,\ldots,9m+9\}$, $0\le m<12$, divisent les 108 pas du registre. L'orientation est
\[
\Sigma_{12}=(+,+,+,-,-,-,+,+,+,-,-,-).
\]
L'événement d'une fenêtre est le résultat
\begin{equation}
a_m=\Sigma_{12}(m)
\sum_{t\in E_m}\mathbf1_{\{2,4,6,8\}}(d_t).
\label{eq:events}
\end{equation}
La notation $a_m$ évite de confondre cet événement avec le nombre d’Euler.
Le registre source conserve en outre l’orientation TRIT, la signature,
la retenue et le segment d’histoire correspondant (voir \ref{sec:k-registro}).

\begin{proposition}[Accroissement produit par le transport]
Soit $\mathbf e_m$ la base de $\R^{12}$ et
$S\mathbf e_m=\mathbf e_{m+1\pmod{12}}$. Si
$z_{m+1}=z_m+a_m\mathbf e_m$ et $y_m=S^{-m}z_m$, alors
\begin{equation}
y_{m+1}=S^{-1}y_m+\eta_m,
\qquad
\eta_m=a_mS^{-1}\mathbf e_0.
\label{eq:eta}
\end{equation}
Après un tour complet,
\begin{equation}
y_{12}=y_0+\sum_{m=0}^{11}a_m\mathbf e_m.
\label{eq:cycle}
\end{equation}
\end{proposition}
\begin{proof}
Multiplier la mise à jour de $z_m$ par $S^{-(m+1)}$ donne
\[
y_{m+1}=S^{-1}y_m+a_mS^{-(m+1)}\mathbf e_m
=S^{-1}y_m+a_mS^{-1}\mathbf e_0.
\]
En itérant, le coefficient de l’événement $a_m$ est
$S^{-(12-m)}\mathbf e_0=\mathbf e_m$. Comme $S^{12}=I$,
on obtient \eqref{eq:cycle}.
\end{proof}

Par conséquent, $\eta_m$ n’est pas une force libre choisie pour imposer une
conservation : c’est une représentation de l’événement TPK. L’égalité de phase
après un tour n’annule pas les événements accumulés. Ce lecteur
conserve les dénombrements, tandis que le registre complet conserve en outre leur
ordre et leur provenance. On distingue $\eta_m$, accroissement vectoriel, de
$\eta_{\rm ret}$, mantisse scalaire d’action, et de $\eta_{\rm el}$,
paramètre de forme qui sera introduit dans la section~\ref{sec:ellipse}.

\subsection{Bilan quadratique}
On définit
\[
\nabla_3=I-S^3,\quad \nabla_4=I-S^4,\quad
B=\nabla_3^{\mathsf T}\nabla_3+\nabla_4^{\mathsf T}\nabla_4,\qquad
\mathcal M(y)=\tfrac12 y^{\mathsf T}By.
\]
Les lecteurs orientés du registre signé satisfont
$(D_jz)_m=z_m-z_{m+j}$ ; avec la convention
$S\mathbf e_m=\mathbf e_{m+1}$, on a
$D_j=I-S^{-j}=\nabla_j^{\mathsf T}$. Les deux sens produisent la même
forme quadratique $B$, car les puissances de $S$ commutent.
On a $S^{\mathsf T}BS=B$ et $B_{00}=4$. En substituant
\eqref{eq:eta}, on déduit exactement
\begin{equation}
\mathcal M(y_{m+1})-\mathcal M(y_m)
=a_m(By_m)_0+2a_m^2.
\label{mem:balance}
\end{equation}
De même, la charge $q(y)=\sum_jy_j$ satisfait
$q(y_{m+1})-q(y_m)=a_m$.
La preuve est le développement du polynôme quadratique après retrait de la
permutation $S^{-1}$. Aucune valeur métrologique n’entre dans ce calcul.
La forme $\mathcal M$ mesure ce registre ; son interprétation comme une
énergie physique déterminée nécessiterait le lecteur dimensionnel correspondant.

\subsection{Résolvantes du transport et conservation de la mémoire}
\label{subsec:ii-memoria-resolvente}

L’histoire prolongée détermine les fenêtres
\(I_m=\{9m+1,\ldots,9m+9\}\), pour \(m\geq0\).
Avec la numérotation de cette section, \(I_m=E_m\) pendant le premier cycle.
Soit \(\sigma_m\in\{-1,1\}\) l’orientation conservée par cette histoire ;
pour \(0\leq m<12\), elle coïncide avec \(\Sigma_{12}(m)\). On définit
\[
 a_m=\sigma_m\sum_{t\in I_m}\mathbf1_{\{2,4,6,8\}}(d_t),
 \qquad |a_m|\leq9.
\]
Le registre et le repère mobile se prolongent par
\[
 z_{m+1}=z_m+a_m\mathbf e_{m\bmod12},\qquad y_m=S^{-m}z_m.
\]
L’état initial \(z_0\) conserve l’histoire antérieure ; il vaut zéro pour
un préfixe initialement vide. Dans cette partie, nous écrivons
\(R=S^{-1}\) pour le transport du repère. La mise à jour
\eqref{eq:eta} est maintenue dans toute fenêtre. L’orientation et les événements
ultérieurs s’obtiennent à partir de l’histoire prolongée : le retour du repère
ne présuppose pas la périodicité de la suite d’événements.

\begin{proposition}[Identité de troncature avec état terminal]
Pour $N\geq1$, soient
\[
Y_N(u)=\sum_{m=0}^N y_mu^m,\qquad
H_{\eta,N}(u)=\sum_{m=0}^{N-1}\eta_mu^m.
\]
On a l’identité polynomiale exacte
\begin{equation}
(I-uR)Y_N(u)=y_0+uH_{\eta,N}(u)-u^{N+1}Ry_N.
\label{mem:identidad-finita}
\end{equation}
\end{proposition}
\begin{proof}
Aux degrés $1\leq m\leq N$, le coefficient du membre de gauche est
$y_m-Ry_{m-1}=\eta_{m-1}$, d’après \eqref{eq:eta}.
Ses coefficients aux degrés zéro et $N+1$ sont $y_0$ et $-Ry_N$,
respectivement. La comparaison des coefficients prouve l’égalité.
\end{proof}

Le terme terminal conserve la mémoire transportée lorsqu’on coupe une histoire.
L’identité permet de comparer une évaluation finie à son prolongement
en maintenant explicite l’état qui reste à la frontière.


\subsubsection{Correspondance avec le développement commun de mémoire}

Le changement de numérotation des fenêtres ne change pas le registre. Dans le
premier cycle, les conventions des deux développements sont reliées par
\[
 E^{(I)}_{m+1}=I_m=E^{(II)}_m,\qquad 0\leq m<12.
\]
Les signes et les événements de chaque fenêtre coïncident après ce
changement d’indice. La convention du transport coïncide également :
\(S\mathbf e_j=\mathbf e_{j+1\bmod12}\) et \(R=S^{-1}\).
Par conséquent, la mise à jour \eqref{eq:eta} est la même identité que
l’équation~\textup{(\HMTIIRefOriginal{mem:actualizacion})} du noyau
commun, avec le même état préalable conservé.

La proposition~\HMTIIRefOriginal{mem:prop-serie} démontre l’identité
de la série génératrice, sa convergence pour \(|u|<1\) et la borne uniforme
de queue~\textup{(\HMTIIRefOriginal{mem:cota})}. L’identité finie
\eqref{mem:identidad-finita} ajoute ici le terme de frontière
\(-u^{N+1}Ry_N\) : elle permet de couper l’histoire sans remplacer son état
terminal par zéro. Le paramètre \(u\) compte, dans les deux cas, des fenêtres
de neuf transitions.

L’élimination de composantes de mémoire conserve l’opérateur, la source
et le lecteur selon la proposition~\HMTIIRefOriginal{mem:prop-schur} et
les équations~\textup{(\HMTIIRefOriginal{mem:schur-fuente})} et
\textup{(\HMTIIRefOriginal{mem:schur-lector})}. L’exemple du cycle
de douze fenêtres~\textup{(\HMTIIRefOriginal{mem:ejemplo-ciclo})}
distingue la compression de la résolvante de la résolvante de la compression.
La composition de trois segments et la transitivité de la réduction
conservent leurs preuves avec les identités
\textup{(\HMTIIRefOriginal{mem:cociclo-tres})} et
\textup{(\HMTIIRefOriginal{mem:schur-transitivo})}.
Ces renvois appartiennent au développement intégral inclus dans ce même
article. Le bilan \eqref{mem:balance} et l’état terminal de
\eqref{mem:identidad-finita} conservent, dans leur présente application,
les accroissements et la frontière du registre ; ils ne déterminent pas à eux seuls
de coefficients analytiques supplémentaires.



\input{sections/ii_sello_y_lectura.tex}

\clearpage\section{Décomposition cyclique et représentations déterminantales}
Les projecteurs cycliques
\[
P_3=\tfrac14(I+S^3+S^6+S^9),\qquad
P_4=\tfrac13(I+S^4+S^8)
\]
ont des images de dimensions trois et quatre. Si $E=P_3-P_4$, alors
\begin{equation}
\frac{\Tr E}{12}=-\frac1{12}.
\label{eq:minus12}
\end{equation}
Le scalaire provient d'une différence de rangs normalisée, non d'une
somme ordinaire d'entiers positifs.

\begin{theorem}[Représentation déterminantale du lecteur de canal]
Soient $\mathcal H_3=\operatorname{Fix}(S^3)$ et
$\mathcal H_4=\operatorname{Fix}(S^4)$. Pour $0<s<1$,
\begin{equation}
F(s)=\frac{\det_{\mathcal H_3}(I-sS)}{\det_{\mathcal H_4}(I-sS)}
=\frac{1-s^3}{1-s^4}
=\frac{1+s+s^2}{1+s+s^2+s^3}.
\label{eq:det}
\end{equation}
Son prolongement continu en $s=1$ vaut $3/4$.
\end{theorem}
\begin{proof}
La restriction de $S$ à chaque image est un cycle à trois ou quatre
positions, respectivement. Leurs polynômes caractéristiques donnent
$\det(I-sS)=1-s^r$ pour $r=3,4$. Après simplification du facteur commun $1-s$,
on obtient la dernière expression, régulière en $s=1$.
\end{proof}

La substitution $s=q^{30}$ exprime
$F(q^{30})=(1-q^{90})/(1-q^{120})$.
La même paire d’espaces qui produit \eqref{eq:minus12}
produit le quotient de canal. Cette égalité n’identifie pas
une matrice à douze coordonnées à tout l’état TPK.

\subsection{Information orientée et secteur complexe}
Dans le secteur antipodal $V_-=\operatorname{Ran}((I-S^6)/2)$,
l'opérateur $J=S^3$ satisfait $J^2=-I$ et $J^{\mathsf T}=-J$.
Avec $Q=P_4(I-S^6)/2$, on a
\[
\operatorname{rank}Q=2,\qquad E|_{V_-}=-Q,\qquad B|_{V_-}=5I-3Q.
\]
Sur $\operatorname{Ran}Q$, $S^2=-I$ et le déterminant est $1+s^2$.
Ainsi,
\[
F(s)=\frac{1+s+s^2}{1+s}\,\frac1{1+s^2}.
\]
L'inversion d'orientation change le signe de la forme
$\omega(u,v)=\langle Ju,v\rangle$, mais conserve $B$ et $F$.
L'évaluation scalaire ne distingue pas tous les transports orientés.

La trace transportée produit en outre
\[
a_n^{\rm tr}=\Tr(S^nE)=3\mathbf1_{3\mid n}-4\mathbf1_{4\mid n},
\]
et, pour $|s|<1$,
\[
-\log F(s)=\sum_{n\ge1}\frac{a_n^{\rm tr}}n\,s^n.
\]
Cette famille de moments relie le présent bloc aux lecteurs
spectraux compilés. Son développement arithmétique complet, y compris
Catalan, aura sa propre résidence ; on ne présuppose pas ici qu’une
identité de fonctions décide à elle seule de sa nature arithmétique.

\clearpage\section{Constantes de Planck et structure déterminantale de l'action}
\label{sec:action}
\subsection{Échelle déterminantale d'action et valuation décimale}

L'échelle absolue et le correcteur relatif proviennent d'objets
distincts. La première utilise le bloc local du registre dodécaphasique
de pas trois,
\begin{equation}
 H_4=\begin{pmatrix}
 1&1&1&1\\1&1&-1&-1\\1&-1&1&-1\\1&-1&-1&1
 \end{pmatrix},\qquad
 G_H=\mathbb Z^4/H_4\mathbb Z^4.
 \label{eq:el-h4}
\end{equation}
L'opération utilise un quotient local. Les trois blocs du registre
complet ne sont pas remplacés par une tensorisation qui multiplierait à
nouveau le nombre de feuillets de ce lecteur.

\begin{lemma}[Nombre de feuillets du quotient local]
La forme normale de Smith de \(H_4\) est
\(\operatorname{diag}(1,2,2,4)\). Par conséquent,
\[
 G_H\simeq\mathbb Z/2\mathbb Z\oplus\mathbb Z/2\mathbb Z
 \oplus\mathbb Z/4\mathbb Z,\qquad |G_H|=16.
\]
\end{lemma}
\begin{proof}
Le plus grand commun diviseur des entrées vaut un. Les mineurs d'ordre deux
sont pairs et l'un d'eux vaut \(-2\), de sorte que leur plus grand commun diviseur
vaut deux. Comme \(H_4H_4^{\mathsf T}=4I\) et \(\det H_4=-16\), les
cofacteurs valent \(\pm4\) ; le troisième diviseur déterminantal vaut quatre.
Le quatrième vaut \(16\). Les quotients successifs des diviseurs
\(1,2,4,16\) sont \(1,2,2,4\). Leur produit donne l'ordre du quotient.
\end{proof}

\begin{definition}[Lecteur déterminantal d'action]
La norme multiplicative de la section scalaire \(\alpha\), constante
sur les feuillets de \(G_H\), et une application de l'auto-échelle
\(\varphi\) produisent
\begin{equation}
 \Sigma_H=\varphi\,\mathrm N_{G_H}(\alpha),\qquad
 \mathrm N_{G_H}(\alpha)=\prod_{g\in G_H}\alpha=\alpha^{16}.
 \label{eq:el-lector-determinantal}
\end{equation}
La coordinatisation décimale associée détermine
\(\operatorname{ord}_{10}(\Sigma_H)=\lfloor\log_{10}\Sigma_H\rfloor\).
\end{definition}

La puissance seize provient ainsi du nombre de feuillets du quotient, une
fois ce lecteur local fixé. Le facteur \(\varphi\) agit une fois sur
la base de la droite d'action. Le calcul de Smith, à lui seul, ne
sélectionnerait pas n'importe quelle fonctionnelle possible sur le quotient : la norme
multiplicative et l'auto-échelle font partie de l'opération déclarée.

\begin{theorem}[Décade de la section d'action]
Le lecteur \eqref{eq:el-lector-determinantal}, appliqué aux coordonnées
HMT déjà engendrées, satisfait
\begin{equation}
 10^{-34}<\Sigma_H<10^{-33},\qquad
 \operatorname{ord}_{10}(\Sigma_H)=-34.
 \label{eq:el-decada}
\end{equation}
\end{theorem}
\begin{proof}
Les intervalles rationnels suivants, beaucoup plus grossiers que les cylindres
disponibles, suffisent :
\[
 \frac{1618}{1000}<\varphi<\frac{1619}{1000},\qquad
 \frac{7297}{10^6}<\alpha<\frac{7298}{10^6}.
\]
La fonction \((x,a)\mapsto xa^{16}\) est croissante en ses deux arguments
positifs. Les deux inégalités entières
\[
 1618\cdot7297^{16}>10^{65},\qquad
 1619\cdot7298^{16}<10^{66}
\]
impliquent
\[
 10^{-34}
 <\frac{1618}{1000}\left(\frac{7297}{10^6}\right)^{16}
 <\Sigma_H
 <\frac{1619}{1000}\left(\frac{7298}{10^6}\right)^{16}
 <10^{-33}.
\]
La définition de l'ordre décimal conclut la preuve. Ni une
unité SI ni une masse de référence n'intervient dans ces inégalités.
\end{proof}

La section normalisée commence par
\(\Sigma_H=1.0462438381\ldots\times10^{-34}\). Sa mantisse ne
s'identifie pas à la coordonnée d'action avant ou après le retour.
Le lecteur déterminantal fixe l'ordre décimal ; les deux coordonnées de
retour s'obtiennent au moyen des lecteurs suivants.


\subsection{Sections d'action et quotient adimensionnel de retour}

Soit \(\mathcal L_S\) une droite orientée d'action et \(\mathcal U_S\)
sa base positive. Le caractère d'ordre cinq et le terme régional de
retour ont pour expressions
\begin{align}
 H_5={}&\frac12\sqrt{\frac{1000\alpha}{\varphi}}-\alpha
   +\frac9{16}\alpha^2-\frac59\alpha^3
   +\frac7{48}\alpha^4-\frac1{54}\alpha^5,
   \label{eq:el-h5}\\
 D_A={}&\exp\!\left(-\frac{100\pi\alpha}{9}\right),
   \label{eq:el-da}\\
 \mathcal C_\pi(\alpha)={}&169\alpha^6+
       \frac{D_A\alpha^7}{1-D_A\alpha},
   \label{eq:el-cpi}\\
 \eta_{\mathrm{ret}}={}&H_5-\frac{90}{\pi}\mathcal C_\pi(\alpha).
   \label{eq:el-eta}
\end{align}
Ce sont des lectures ultérieures de la fermeture dodécaphasique et de son
incidence régionale. L'entier \(169\) est conservé comme le sélecteur
régional déjà construit, et les coefficients de \(H_5\) comme ceux de son
caractère local ; ils ne sont pas remplacés par une interpolation de la constante de
Planck. Les formules explicitent exactement les lectures nécessaires
à l'étape électronique. Il n'est pas nécessaire d'introduire ici d'autres
fonctionnelles de la théorie des masses.

Pour \(0<\alpha<1\), on a \(0<D_A<1\), de sorte que
\(D_A\alpha<1\). Le terme rationnel de \eqref{eq:el-cpi} conserve la
somme complète de la queue géométrique
\begin{equation}
 \frac{D_A\alpha^7}{1-D_A\alpha}
 =\sum_{n=0}^{\infty}D_A^{n+1}\alpha^{n+7}.
 \label{eq:el-cola-regional}
\end{equation}
Il ne s'agit pas d'une troncature que l'on pourrait modifier silencieusement en changeant la
précision. Après le terme d'indice \(N\), la queue restante est
\(D_A^{N+2}\alpha^{N+8}/(1-D_A\alpha)\), ce qui permet de contrôler
l'évaluation sans recourir à une grandeur cible.

\begin{definition}[Sections d'action et lecteur de retour]
Les sections avant et après le retour sont
\begin{equation}
 \hbar_{\mathrm{pre}}=H_5\,10^{-34}\mathcal U_S,
 \qquad
 \hbar_{\mathrm{ret}}=\eta_{\mathrm{ret}}\,10^{-34}\mathcal U_S.
 \label{eq:el-secciones-accion}
\end{equation}
Lorsque \(\eta_{\mathrm{ret}}>0\), on définit
\begin{equation}
 \delta_{\mathrm{ret}}=H_5-\eta_{\mathrm{ret}},\qquad
 R_{\mathrm{act}}=\frac{\hbar_{\mathrm{pre}}}{\hbar_{\mathrm{ret}}}
 =\frac{H_5}{\eta_{\mathrm{ret}}}.
 \label{eq:el-ratio}
\end{equation}
\end{definition}

\begin{proposition}[Positivité et invariance du rapport]
Dans les intervalles utilisés dans la preuve de \eqref{eq:el-decada},
\(0<\eta_{\mathrm{ret}}<H_5\), et donc
\(R_{\mathrm{act}}>1\). Le rapport ne change pas lorsqu'on remplace la base
\(\mathcal U_S\) par toute autre base positive commune.
\end{proposition}
\begin{proof}
Les termes de \(\mathcal C_\pi\) sont positifs. Pour obtenir une borne
élémentaire, il suffit d'utiliser \(\alpha<0.008\), \(\varphi<1.619\),
\(\alpha>0.007\), \(D_A<1\) et \(\pi>3\). La racine dans
\eqref{eq:el-h5} est supérieure à un, tandis que la somme des termes
négatifs est inférieure à \(0.008001\) ; en particulier, \(H_5>0.99\).
De même,
\[
 0<\frac{90}{\pi}\mathcal C_\pi
 <30\left(169(0.008)^6+
           \frac{(0.008)^7}{1-0.008}\right)<10^{-6}.
\]
Ainsi, \(\eta_{\mathrm{ret}}>0.989999>0\) et
\(\eta_{\mathrm{ret}}<H_5\). Enfin, le changement de base divise
les deux coordonnées d'action par le même facteur positif, qui se simplifie dans
leur quotient.
\end{proof}

L'évaluation ultérieure fournit
\begin{align*}
 H_5&=1.0545718176461544696\ldots,\\
 \eta_{\mathrm{ret}}&=1.0545718169150390611\ldots,\\
 R_{\mathrm{act}}&=1.0000000006932817630\ldots.
\end{align*}
La différence et le quotient sont deux lectures des mêmes sections,
l'une additive et l'autre multiplicative. Ils ne représentent pas deux constantes de
Planck indépendantes. Le passage de l'action angulaire à l'action par tour
complet s'exprime, dans chaque section, par
\(h=2\pi\hbar\). La coordinatisation ultérieure
\(\mathcal U_S\mapsto\mathrm{J\,s}\) assigne une unité, mais ne
rétroagit pas sur l'ordre décimal déjà démontré.

Le retour régional peut aussi s'écrire dans une coordinatisation angulaire.
Cette reparamétrisation n'est pas une dépendance supplémentaire indispensable :
\eqref{eq:el-h5}--\eqref{eq:el-eta} fournissent directement
\(H_5\), \(\eta_{\mathrm{ret}}\) et \(R_{\mathrm{act}}\). Si l'on
introduit l'angle et le contre-angle, tous deux doivent être transportés dans la même
coordinatisation ; on ne mélange pas degrés et radians au sein d'une différence.

\input{sections/revision_planck.tex}

\begingroup\fontsize{10}{12.1}\selectfont
\setlength{\parskip}{1.5pt}\setlength{\abovedisplayskip}{4pt}\setlength{\belowdisplayskip}{4pt}
\subsection{Coordonnées angulaires, normalisation du tour et unités d'action}
\label{sec:ii-angulos}

La grandeur d’action et la coordonnée angulaire appartiennent à des espaces
distincts. La section \(\hbar_{\rm ret}\) occupe la droite d’action
\(\mathcal L_S\) ; \(\eta_{\rm ret}\) est son coefficient dans la base
déclarée. La représentation angulaire agit sur la paire adimensionnelle
\((\alpha,\eta_{\rm ret})\), après sa construction :
\begin{equation}
 \mathcal P(\alpha,\eta_{\rm ret})
 =\bigl([1000\alpha]_{360},[2(\eta_{\rm ret}+\alpha)]_{360}\bigr).
 \label{eq:ii-par-angular}
\end{equation}
La complétion nonadique détermine le facteur 1000 ; la clôture
bidirectionnelle détermine le facteur deux. La coordinatisation circulaire de 360
unités rend visibles les partitions
\(12\cdot30=4\cdot90=3\cdot120\). Sur la branche positive de
l’article, les représentants sont
\[
 A_{\deg}=1000\alpha,\qquad
 C^*_{\deg}=2(\eta_{\rm ret}+\alpha).
\]

\begin{proposition}[Décalage de la représentation en base mille]
\label{prop:ii-desplazamiento-angular}
Soit \(\alpha=\sum_{n\geq1}b_n1000^{-n}\), avec
\(b_n\in\{0,\ldots,999\}\), le développement compatible déjà engendré.
La coordonnée relevée \(A_{\deg}=1000\alpha\) satisfait
\[
 A_{\deg}=b_1+\sum_{n\geq1}b_{n+1}1000^{-n}.
\]
La représentation conserve la suite des blocs à partir du deuxième ;
le bloc initial passe à la partie entière.
\end{proposition}
\begin{proof}
La série converge absolument. Multiplier chaque terme par 1000
et séparer le terme initial fournit l’identité. L’opération
agit sur le développement construit, avec ses blocs et sa mémoire ;
la représentation circulaire ultérieure prend sa classe modulo 360.
\end{proof}

\begin{definition}[Coordinatisations de la coordonnée angulaire]
La coordinatisation en degrés exprime un tour complet au moyen de 360 unités.
La coordinatisation en radians l’exprime au moyen de \(2\pi\). Le changement de
coordonnées et ses inverses sont
\[
 x=\frac{\pi}{180}A_{\deg},\qquad
 y=\frac{\pi}{180}C^*_{\deg},\qquad
 A_{\deg}=\frac{180}{\pi}x,\qquad
 C^*_{\deg}=\frac{180}{\pi}y.
\]
Sur un parcours dont le nombre de tours est conservé, on utilise les
coordonnées relevées ; le quotient circulaire exprime sa classe modulo
le tour correspondant.
\end{definition}

La définition caractérise les degrés et les radians par leur relation exacte
avec le tour. Le registre retient en outre l’orientation, les tours et la
mémoire de transport. La représentation de sa classe angulaire conserve
la phase, tandis que le relèvement distingue des parcours ayant la même
phase finale. La lecture physique d’action nécessite de préciser laquelle de
ces coordonnées intervient dans l’intégrale.

\begin{proposition}[Covariance du lecteur d'action]
Pour \(h_a=2\pi\hbar_a\), l'action d'une coordonnée relevée
satisfait
\begin{equation}
 \mathcal S_a(\theta)
 =\hbar_a\theta_{\rm rad}
 =h_a\frac{\theta_{\deg}}{360}.
 \label{eq:ii-accion-cartas}
\end{equation}
Un tour positif exprime \(h_a\).
\end{proposition}
\begin{proof}
Substituer \(\theta_{\rm rad}=\pi\theta_{\deg}/180\) donne
\(\hbar_a\pi\theta_{\deg}/180
=(2\pi\hbar_a)\theta_{\deg}/360\). Pour un tour,
\(\theta_{\deg}=360\), d'où \(h_a\).
\end{proof}

La périodicité d'une représentation est déclarée séparément. Si la
représentation transportée satisfait \(U(2\pi)=-I\), sa restitution
opératoire se produit en \(4\pi\). Cette propriété du relèvement
coexiste avec la normalisation par tour de
\eqref{eq:ii-accion-cartas}. Chacune conserve son objet : grandeur
intégrée, phase circulaire et état représenté.

\subsection{Covariance des canaux et de la fonctionnelle d'incidence}

Les canaux s'écrivent de manière équivalente
\begin{equation}
 q_\pm=\exp[-(x\pm y)]
 =\exp\!\left[-\frac{\pi}{180}(A_{\deg}\pm C^*_{\deg})\right].
\label{eq:ii-canales}
\end{equation}
\begin{proposition}[Récupération des coordonnées paire et impaire]
\label{prop:ii-inversion-canales-angulares}
Dans la coordinatisation réelle relevée, les canaux positifs déterminent
de manière univoque les deux coordonnées :
\[
 A_{\deg}=-\frac{90}{\pi}\log(q_+q_-),\qquad
 C^*_{\deg}=\frac{90}{\pi}\log\frac{q_-}{q_+}.
\]
L'échange \((q_+,q_-)\mapsto(q_-,q_+)\) conserve
\(A_{\deg}\) et inverse \(C^*_{\deg}\).
\end{proposition}
\begin{proof}
La somme et la différence des logarithmes de
\eqref{eq:ii-canales} sont respectivement
\(-\pi A_{\deg}/90\) et \(-\pi C^*_{\deg}/90\).
Le logarithme réel est univoque sur les canaux positifs.
Le produit est symétrique et le quotient s'inverse lors de leur échange.
\end{proof}
Dans le registre à douze positions, l'involution antipodale
\(r\mapsto r+6\pmod {12}\) forme six paires. La distribution
uniforme de la coordonnée impaire sur ces paires est
\(C^*_{\deg}/6\) : sommer ses six contributions restitue
\(C^*_{\deg}\). Cette coordonnée sectorielle se distingue de
la constante de Planck \(h\), qui appartient à la droite d'action.

La conversion est appliquée une seule fois à chaque argument. La grandeur
\(\alpha\) demeure adimensionnelle et la grandeur \(\hbar\)
appartient toujours à la droite d'action ; leurs coordonnées antérieures
interviennent dans l'application \(\mathcal P\).

La fonctionnelle de Barbero--Immirzi contient, outre les canaux, une
composante bilinéaire dans les représentants angulaires. Son transport est
\begin{equation}
 \frac{\pi}{180}A_{\deg}C^*_{\deg}
 =\frac{180}{\pi}xy.
 \label{eq:ii-barbero-cartas}
\end{equation}
L'identité s'obtient en remplaçant les deux représentants en degrés
par leurs expressions en radians. La formule complète en
coordonnées analytiques est donc
\[
 \gamma_{\HMT}
 =\frac{180}{\pi}xy
 +12\bigl[S_{90}(e^{-x+y})-S_{90}(e^{-x-y})\bigr]
 -\bigl[S_{120}(e^{-x+y})-S_{120}(e^{-x-y})\bigr].
\]
L’équivalence des deux coordinatisations transporte la composante bilinéaire
et les fonctions exponentielles simultanément. Cette transformation
explicite le rôle précis du degré dans la construction et permet
de reproduire l’évaluation sans ambiguïté d’unités angulaires.
\par\endgroup

\clearpage% Construcción dodecafásica del funcional de retorno.
\section{Chaîne fondamentale dodécaphasique et fonctionnelle hexadique--octadique}

L’incidence active se construit avant toute réalisation gravitationnelle.
Soit
\[
 H_{\mathrm{act}}=\{1,2,4,5,7,8\},
 \qquad
 \mathcal P_2(H_{\mathrm{act}})
 =\{A\subset H_{\mathrm{act}}:|A|=2\}.
\]
Alors
\[
 |H_{\mathrm{act}}|=6,
 \qquad
 |\mathcal P_2(H_{\mathrm{act}})|=15,
\]
et les incidences antérieure et postérieure à l’extension ont pour cardinaux
\[
 \mathcal F_{90}=H_{\mathrm{act}}\times
 \mathcal P_2(H_{\mathrm{act}}),
 \qquad |\mathcal F_{90}|=90,
\]
\[
 \mathcal F_{120}=\{1,\ldots,8\}\times
 \mathcal P_2(H_{\mathrm{act}}),
 \qquad |\mathcal F_{120}|=120.
\]
Ces deux cardinaux sont des sorties de la règle de phase ; ce ne sont ni des angles ni
des paramètres reçus de la réalisation physique.

Notons \(C_{12}=\mathbb Z/12\mathbb Z\) le registre temporel du TPK et
\(\partial_{\rm ret}\) son unique couture de retour par tour. Dans le
module libre des événements orientés, on considère la chaîne fondamentale
\begin{equation}
 c_{12}^{+}
 =\sum_{j\in C_{12}}[j,\mathcal F_{90}]
  -[\partial_{\rm ret},\mathcal F_{120}].
 \label{intsuccpass:eq:c12-fundamental-chain}
\end{equation}
La somme contient exactement les douze secteurs de la roue ; le second
terme est unique parce qu’un tour possède une seule couture de clôture. Son signe
est le signe de frontière. L’involution bidirectionnelle envoie
\(c_{12}^{+}\) sur \(c_{12}^{-}=-c_{12}^{+}\) : elle change l’orientation de la
chaîne, mais n’altère ni ses multiplicités absolues ni ses types
d’incidence.

Pour \(m\in\{90,120\}\), soit
\[
 S_m(q)=\frac{q^m}{1-q^{3m}}.
\]
Le lecteur de retour \(\mathscr R\) est défini sur les générateurs de
\eqref{intsuccpass:eq:c12-fundamental-chain} par
\[
 \mathscr R([j,\mathcal F_{90}])=S_{90},
 \qquad
 \mathscr R([\partial_{\rm ret},\mathcal F_{120}])=S_{120}.
\]
La première égalité est constante sur \(C_{12}\) par covariance sous la
rotation de phase ; la seconde a un domaine distinct et n’agit que sur la
couture étendue. Par linéarité,
\begin{equation}
 \boxed{\mathscr R(c_{12}^{+})=12S_{90}-S_{120}.}
 \label{intsuccpass:eq:dodecaphase-forces-12minus1}
\end{equation}
Ainsi, les coefficients \(12:-1\) sont l’image de la chaîne fondamentale : douze
cellules de phase et une frontière orientée. Ils ne sont sélectionnés ni au moyen de la valeur
conventionnelle du paramètre de Barbero--Immirzi, ni au moyen d’une comparaison
décimale, ni en imposant au préalable le coefficient de pôle qui sera calculé
ci-après.

\begin{theorem}[Coefficient de pôle de la chaîne dodécaphasique]
La fonctionnelle \(F(q)=\mathscr R(c_{12}^{+})(q)
=12S_{90}(q)-S_{120}(q)\), déterminée par les douze secteurs
et leur couture orientée, a pour coefficient \(1/24\) dans
\((1-q)^{-1}\). Son résidu complexe en \(q=1\) est \(-1/24\).
Le coefficient de pôle est une sortie de la fonctionnelle déjà construite.
\end{theorem}

\begin{proof}
Pour une fonctionnelle ayant un pôle simple en \(q=1\), nous écrivons
\[
 p_1(F):=\lim_{q\uparrow1}(1-q)F(q).
\]
L’identité
\[
 1-q^{3m}=(1-q)\sum_{j=0}^{3m-1}q^j
\]
implique
\[
 p_1(S_m)
 =\lim_{q\uparrow1}\frac{q^m}{\sum_{j=0}^{3m-1}q^j}
 =\frac1{3m}.
\]
Par linéarité,
\[
 p_1\bigl(\mathscr R(c_{12}^{+})\bigr)
 =\frac{12}{270}-\frac1{360}
 =\frac{48-3}{1080}=\frac1{24}.
\]
Ces égalités s’effectuent dans \(\mathbb Q\). Le résidu complexe
est le coefficient de \((q-1)^{-1}\), de sorte que
\[
 \operatorname*{Res}_{q=1}\mathscr R(c_{12}^{+})(q)
 =-p_1\bigl(\mathscr R(c_{12}^{+})\bigr)=-\frac1{24}.
\]
\end{proof}

L’équation diophantienne
\[
 \frac a{270}-\frac b{360}=\frac1{24}
 \quad\Longleftrightarrow\quad
 4a-3b=45
\]
demeure un certificat arithmétique ultérieur. À elle seule, elle admet la
famille \((a,b)=(12+3t,1+4t)\) ; la chaîne
\eqref{intsuccpass:eq:c12-fundamental-chain} élimine cette ambiguïté avant le calcul,
car changer \((12,1)\) altère la multiplicité de l’orbite ou de sa couture
et, par conséquent, ne représente plus un tour fondamental du TPK.

Avec les canaux conjugués \(q_\pm\) déjà engendrés par la branche angulaire, on
obtient enfin
\[
 \mathscr R(c_{12}^{+})(q_-)-
 \mathscr R(c_{12}^{+})(q_+)
 =12\Delta S_{90}-\Delta S_{120}.
\]
C’est l’entrée interne de la fonctionnelle hexadique--octadique. Son évaluation
comme paramètre de Barbero--Immirzi et son insertion dans l’opérateur d’aire sont
des représentations ultérieures : elles ne participent à la sélection
ni de \(c_{12}^{+}\), ni de \(90/120\), ni de \(12:-1\), ni de \(1/24\).

\paragraph{Falsificateurs.}
La construction est réfutée si le recensement de phase ne donne pas douze secteurs, si
un tour contient plus d’une couture, si les incidences n’ont pas pour
cardinaux \(90\) et \(120\), si l’involution n’inverse pas la chaîne complète,
ou si le coefficient exact de \((1-q)^{-1}\) de son image diffère de \(1/24\).

\clearpage\section{Incidence hexadique--octadique et fonctionnelle de Barbero--Immirzi}
\label{sec:barbero}
Soit $H\hookrightarrow O$ une inclusion marquée, $|H|=6$, $|O|=8$.
Les objets d'incidence
\[
\mathcal F_6=H\times\binom H2,\qquad
\mathcal F_8=O\times\binom H2
\]
ont pour cardinaux $90$ et $120$. Leur différence normalisée donne
\[
\frac{|\mathcal F_6|-|\mathcal F_8|}{24\binom62}=-\frac1{12}.
\]
C'est une réalisation combinatoire du même scalaire.
Une égalité numérique avec \eqref{eq:minus12} n'efface pas
la distinction entre les espaces sur lesquels il a été calculé.

\begin{figure}[htbp]
\centering
\begin{tikzpicture}[>=Latex,every node/.style={font=\small}]
\draw[rounded corners,draw=ink] (-1.8,-1.35) rectangle (1.8,1.35);
\foreach \a in {0,60,...,300}{
\fill[ink] (\a:0.83) circle(2pt);}
\node at (0,-1.7) {$H:\ 6$ puntos};
\draw[->,thick,ink] (2.1,0)--(3.6,0);
\node[above] at (2.85,0.1) {inclusion marquée};
\draw[rounded corners,draw=ink] (4,-1.35) rectangle (7.6,1.35);
\foreach \a in {0,60,...,300}{
\fill[ink] ($(5.8,0)+(\a:0.83)$) circle(2pt);}
\fill[accent] (4.5,0) circle(2pt);
\fill[accent] (7.1,0) circle(2pt);
\node at (5.8,-1.7) {$O:\ 8$ puntos};
\node at (0,-2.35) {$6\binom62=90$};
\node at (5.8,-2.35) {$8\binom62=120$};
\end{tikzpicture}
\caption{Incidence marquée de six à huit points. Le dessin représente
les ensembles et l'inclusion ; il ne prétend pas transformer les faces ou sommets
d'un cube en blocs de Steiner sans leur application d'incidence.}
\end{figure}

\subsection{Les poids précèdent le coefficient du pôle}
Dans le module libre des événements orientés, la source fixe la chaîne
fondamentale du tour (voir \ref{intsuccpass:eq:dodecaphase-forces-12minus1})
\[
c_{12}^{+}=
\sum_{j\in\Z/12\Z}[j,\mathcal F_6]
-[\partial_{\rm ret},\mathcal F_8].
\]
La multiplicité douze correspond aux secteurs et la multiplicité un
au raccordement orienté. L'involution change l'orientation de la chaîne
complète. Soit
\[
S_m(q)=\frac{q^m}{1-q^{3m}},\qquad 0<q<1.
\]
Le lecteur linéaire attribue $S_{90}$ à chaque secteur et $S_{120}$ au raccordement ;
par conséquent,
\[
\mathscr R(c_{12}^{+})=12S_{90}-S_{120}.
\]

\begin{theorem}[Coefficient du pôle et évaluation conjuguée]
L'image de la chaîne fondamentale a pour coefficient $1/24$
devant $(1-q)^{-1}$ et pour résidu $-1/24$ devant $(q-1)^{-1}$.
Son évaluation dans les canaux préalablement engendrés détermine
\begin{equation}
\gamma_{\HMT}=
\frac{\pi AC^*}{180}
+12\bigl[S_{90}(q_-)-S_{90}(q_+)\bigr]
-\bigl[S_{120}(q_-)-S_{120}(q_+)\bigr].
\label{eq:gamma}
\end{equation}
\end{theorem}
\begin{proof}
La factorisation $1-q^{3m}=(1-q)\sum_{j=0}^{3m-1}q^j$ implique
$\lim_{q\uparrow1}(1-q)S_m(q)=1/(3m)$. Par linéarité,
\[
\frac{12}{270}-\frac1{360}=\frac1{24}.
\]
La coordonnée $q-1$ inverse le signe du résidu.
L'évaluation conjuguée et la composante angulaire de la source
produisent \eqref{eq:gamma}.
\end{proof}

Le certificat diophantien $4a-3b=45$ admet les paires
$(12+3t,1+4t)$, $t\ge0$ ; il certifie la minimalité de $(12,1)$,
mais ne remplace pas l’origine des multiplicités dans la chaîne.
L’évaluation conservée dans le document propriétaire est
\[
\gamma_{\HMT}
=0.27400767269445927474725526077398041940\ldots.
\]
Ce nombre est une sortie de la fonctionnelle source ; ce n'est pas une valeur externe
utilisée pour sélectionner $A$, $C^*$ ou leurs coefficients.

\subsection{Parité et sensibilité structurelle}
Inverser $C^*$ échange $q_+$ et $q_-$ et change le signe
de la composante angulaire. Ainsi,
\[
\gamma_{\HMT}(A,-C^*)=-\gamma_{\HMT}(A,C^*).
\]
La fonctionnelle conserve l'orientation, tandis que d'autres lecteurs,
comme $q_+q_-$, sont pairs. Cette différence précède son interprétation physique.

\subsection{Dépendance structurelle de la fonctionnelle et normalisation aréale}

La construction précédente situe le paramètre de Barbero--Immirzi à la
rencontre de trois déterminations : incidence marquée, orientation
dodécaphasique et évaluation angulaire. L'incidence apporte les espaces
d'événements ; l'orientation fixe la chaîne sur ces espaces ;
l'évaluation transporte la chaîne vers une quantité adimensionnelle. Cette
décomposition permet de conserver l'origine de chaque coefficient lors de
l'utilisation ultérieure de la fonctionnelle dans la géométrie d'aire.

L'inclusion \(\mathcal F_6\hookrightarrow\mathcal F_8\) maintient
la paire marquée de l'hexade. Son complément a pour cardinal
\(2\binom62=30\). Cet incrément réapparaît comme déterminant de
l'application qui reconstruit les occupations du bord à partir de l'aire et
de l'observable pondéré. L'égalité de ces deux entiers possède ici
une explication concrète : les poids sectoriels diffèrent précisément
du nombre d'événements ajoutés par l'inclusion marquée.

\begin{proposition}[Incidence et récupérabilité des occupations]
Soient \(n_6,n_8\) les occupations agrégées de deux secteurs et soient
\(r_6=|\mathcal F_6|\), \(r_8=|\mathcal F_8|\). L'application
\[
 (n_6,n_8)\longmapsto(n_6+n_8,r_6n_6+r_8n_8)
\]
est inversible sur son image lorsque l'inclusion ajoute au moins un
événement. Pour l'inclusion construite, son déterminant vaut 30.
\end{proposition}
\begin{proof}
La matrice de l'application est
\(\left(\begin{smallmatrix}1&1\\r_6&r_8\end{smallmatrix}\right)\).
Son déterminant est \(r_8-r_6=|\mathcal F_8\setminus\mathcal F_6|\).
L'inclusion considérée ajoute \(2\binom62=30\) événements. L'inversion
linéaire restitue les deux occupations et conserve les contraintes entières
propres à son image.
\end{proof}

Le résultat restitue des occupations agrégées. L'information concernant
chaque liaison et l'ordre des événements demeure dans l'état complet
et dans ses lecteurs supplémentaires. Cette hiérarchie donne une formulation
précise de la différence entre valeur d'aire et généalogie de frontière :
l'aire conserve une somme ; le couple d'observables conserve la
décomposition de cette somme par secteurs ; le registre conserve en outre
sa distribution et son histoire.

Le paramètre \(\gamma_{\HMT}\) transporte cette incidence à l'échelle
aréale. Son interprétation structurelle est celle d'un coefficient orienté
de la fonctionnelle de retour hexadique--octadique. La réalisation gravitationnelle
le reconnaît ensuite dans l'action de Holst. On distingue ainsi sa
définition intrinsèque, sa valeur et sa fonction dans l'équation physique.

\clearpage\section{Ellipse d’action et réciprocité modulaire}
\label{sec:ellipse}
Les sorties $A,C^*,\hbar>0$ étant fixées dans la coordinatisation
$|C^*/A|<1$, nous définissons
\[
\eta_{\rm el}=\operatorname{artanh}(C^*/A),\quad
a_S=\sqrt{2\hbar}\,e^{\eta_{\rm el}/2},\quad
b_S=\sqrt{2\hbar}\,e^{-\eta_{\rm el}/2}.
\]
Les coordonnées canoniques équilibrées ont pour unité la racine d’action ;
on n’additionne pas ici la position et la quantité de mouvement sans coordinatisation dimensionnelle.
On a
\[
a_Sb_S=2\hbar,\qquad
\tau_{\rm el}=i\,\frac{b_S}{a_S},\qquad
R_{\rm el}=\sqrt{\frac{b_S}{a_S}}
=\left(\frac{A-C^*}{A+C^*}\right)^{1/4}.
\]
La racine positive détermine la composition $\tau_{\rm el}=iR_{\rm el}^2$,
et permet de retrouver le rapport angulaire :
\[
\frac{C^*}{A}=\frac{1-R_{\rm el}^4}{1+R_{\rm el}^4}.
\]
L’inversion de forme envoie $R_{\rm el}$ sur $R_{\rm el}^{-1}$ et
$\tau_{\rm el}$ sur $-1/\tau_{\rm el}$ (voir \ref{sec:ii-elipse-geometria}).

\begin{proposition}[Énergie coïncidente, action discernable]
Soit
\[
D(\eta)=\begin{pmatrix}e^\eta&0\\0&e^{-\eta}\end{pmatrix},
\quad
S_2=\begin{pmatrix}0&1\\1&0\end{pmatrix},
\quad
J_2=\begin{pmatrix}0&-1\\1&0\end{pmatrix}.
\]
Les deux applications satisfont $T^{\mathsf T}D(-\eta)T=D(\eta)$.
Cependant, pour $\omega=\dd Q\wedge\dd P$,
\[
S_2^*\omega=-\omega,\qquad J_2^*\omega=\omega.
\]
Si $\mathcal I(\gamma)=\oint_\gamma P\,\dd Q$ sur une courbe fermée,
\[
\mathcal I(S_2\gamma)=-\mathcal I(\gamma),\qquad
\mathcal I(J_2\gamma)=\mathcal I(\gamma).
\]
\end{proposition}
\begin{proof}
La conjugaison échange les entrées diagonales de $D$.
Les déterminants de $S_2,J_2$ sont $-1,1$, respectivement, ce qui
prouve la transformation de $\omega$. Pour $\lambda=P\,\dd Q$,
\[
S_2^*\lambda=\dd(PQ)-\lambda,\qquad
J_2^*\lambda=\lambda-\dd(PQ).
\]
Intégrer la différentielle exacte sur la courbe fermée complète la preuve.
\end{proof}

L’énergie quadratique et le module ne suffisent pas à distinguer les deux
transports. L’orientation de l’action les distingue, même
dans la forme circulaire. Le marquage géométrique fait ainsi partie de la lecture
récupérable, ce n’est pas un ornement ajouté à la valeur.

% Adición separada: no sustituye 07_elipse.tex.
% Procedencia: integral 2249, 52.4.1--2, 52.6.2--3, 53.28 y 90.4.6.
% Recuperación de resultados y explicitación geométrica; certificado local
% verificar_elipse_II.py. No usa un radio físico como entrada generadora.

\subsection{Demi-axes, foyers et rayons intérieurs}
\label{sec:ii-elipse-geometria}

La figure s’obtient à partir des sorties angulaires et d’action des sections
précédentes. APP conserve les feuillets ainsi que leurs résidus et quotients ; TRIT fixe
l’orientation ; TPK transporte la frontière, la retenue et la mémoire jusqu’aux
représentations du même état enrichi et de la structure discrète
conjointe du continu. On reçoit ici la paire \((A,C^*)\) et la section
positive \(\hbar\) après cette construction. La géométrie explicite
l’information que conservent ses lecteurs, sans utiliser une longueur observée
pour sélectionner l’état source (voir \ref{sec:ii-elipse-geometria}).

Soit \(\xi=C^*/A\in(-1,1)\), avec \(A\ne0\) et les deux angles dans la
même unité. Écrivons \(\eta=\operatorname{artanh}\xi\) et
\[
 a_S=\sqrt{2\hbar}\,e^{\eta/2},\qquad
 b_S=\sqrt{2\hbar}\,e^{-\eta/2},\qquad
 \gamma(\theta)=(a_S\cos\theta,b_S\sin\theta).
\]
Les demi-axes sont étiquetés par leurs coordonnées : si \(\eta<0\),
l’axe vertical est le plus grand. Tous deux appartiennent à la droite de racine
d’action de la coordinatisation équilibrée, non à la droite des longueurs spatiales.

\begin{proposition}[Géométrie focale et intérieur radial]
\label{prop:ii-elipse-focos}
Soient \(a=\max(a_S,b_S)\), \(b=\min(a_S,b_S)\). La demi-distance
focale et l'excentricité sont
\begin{equation}
 f_S=\sqrt{a^2-b^2}
 =\sqrt{4\hbar\sinh|\eta|},\qquad
 e_f=\frac{f_S}{a}=\sqrt{1-e^{-2|\eta|}}.
 \label{eq:ii-elipse-focal}
\end{equation}
Pour \(\eta>0\), les foyers sont \((\pm f_S,0)\) ; pour
\(\eta<0\), ils sont \((0,\pm f_S)\). En \(\eta=0\), ils coïncident avec
le centre. Le rayon de frontière dans la direction polaire \(\phi\) est
\begin{equation}
 \rho_{\partial}(\phi)
 =\left(\frac{\cos^2\phi}{a_S^2}
        +\frac{\sin^2\phi}{b_S^2}\right)^{-1/2}.
 \label{eq:ii-elipse-radio-polar}
\end{equation}
Le segment intérieur dans cette direction est constitué des points
\(r(\cos\phi,\sin\phi)\), avec \(0\le r<\rho_{\partial}(\phi)\).
En particulier, \(b\le\rho_{\partial}(\phi)\le a\).
\end{proposition}

\begin{proof}
Les carrés du grand et du petit demi-axe sont
\(2\hbar e^{|\eta|}\) et \(2\hbar e^{-|\eta|}\). Leur différence
donne \eqref{eq:ii-elipse-focal}. Pour vérifier le foyer, supposons
\(\eta>0\). En \(X=\gamma(\theta)\), les distances à
\(F_+=(f_S,0)\) et \(F_-=(-f_S,0)\) satisfont
\[
 |X-F_+|^2=(a_S-f_S\cos\theta)^2,\qquad
 |X-F_-|^2=(a_S+f_S\cos\theta)^2.
\]
Comme \(a_S>f_S\), leurs racines positives ont pour somme \(2a_S\).
Pour \(\eta<0\), échanger les coordonnées donne la même preuve.
Dans le cas circulaire, \(f_S=0\). Enfin, substituer
\((Q,P)=r(\cos\phi,\sin\phi)\) dans
\(Q^2/a_S^2+P^2/b_S^2=1\) donne
\eqref{eq:ii-elipse-radio-polar}. Le coefficient de \(r^2\) est
compris entre \(a^{-2}\) et \(b^{-2}\), ce qui démontre les bornes et la
caractérisation de l'intérieur.
\end{proof}

\subsection{Phase paramétrique, angle polaire et aire d'action}
\label{sec:ii-elipse-fase-area}

\begin{proposition}[Loi aréale et deux coordonnées angulaires]
\label{prop:ii-elipse-area}
La phase \(\theta\) de \(\gamma\) et son angle polaire continu
\(\phi\), relevés à partir de zéro, satisfont
\begin{equation}
 \phi(\theta)=\arg(a_S\cos\theta+i b_S\sin\theta),\qquad
 \frac{d\phi}{d\theta}
 =\frac{a_Sb_S}{a_S^2\cos^2\theta+b_S^2\sin^2\theta}>0.
 \label{eq:ii-elipse-fases}
\end{equation}
Dans une coordinatisation où les deux tangentes sont définies,
\(\tan\phi=(b_S/a_S)\tan\theta\). L’aire orientée balayée est
\begin{equation}
 \mathcal A(\theta)=\frac12\int_0^\theta(QP'-PQ')\,dt
 =\hbar\theta,\qquad
 \mathcal A(1)=\hbar,\qquad \mathcal A(2\pi)=h.
 \label{eq:ii-elipse-area-fase}
\end{equation}
\end{proposition}

\begin{proof}
La dérivée de l'argument d'un vecteur non nul est
\((QP'-PQ')/(Q^2+P^2)\). Ici, le numérateur est
\(a_Sb_S=2\hbar\), ce qui démontre la dérivée et la monotonie.
Le quotient \(P/Q\) donne la relation des tangentes ; le relèvement
continu fixe les quadrants et le nombre de tours. L'intégrale aréale
vaut \(a_Sb_S\theta/2=\hbar\theta\), et
\(h=2\pi\hbar\) donne la dernière identité.
\end{proof}

La phase paramétrique ne coïncide pas en général avec l’angle polaire. Un radian
de phase balaie \(\hbar\), mais ne s’identifie pas pour autant à un arc
circulaire ni à un secteur d’angle polaire unitaire. En coordonnées polaires,
la même loi s’écrit
\(d\mathcal A=\rho_{\partial}(\phi)^2d\phi/2\).
Avec l’orientation de \(\gamma\) utilisée ici,
\(\oint P\,dQ=-h\) : l’aire positive est \(h\), tandis que
l’intégrale de cette forme de degré un conserve son signe. Cela concorde avec la
distinction entre les transports symplectique et antisymplectique de la
section précédente.

\begin{figure}[htbp]
 \centering
 \input{figures/elipse_accion_II.tex}
 \caption{Demi-axes, foyers et réciprocité de l'ellipse d'action dans les
 coordonnées normalisées \(u=Q/\sqrt{2\hbar}\),
 \(v=P/\sqrt{2\hbar}\). Le secteur marqué possède la phase paramétrique
 \(\theta=\pi/6\), distincte de son angle polaire \(\phi\) ; son aire
 d'action est \(\hbar\pi/6\). La figure utilise le cas algébrique
 \(\xi=1/3\) pour rendre les relations lisibles : elle n'attribue pas cette valeur
 à la sortie angulaire HMT. Les foyers sont calculés à partir des demi-axes.
 À droite, l'échange des axes conserve l'aire totale \(h\)
 et inverse le rayon modulaire adimensionnel.}
 \label{fig:ii-elipse-calculada}
\end{figure}

Pour le cas de contrôle de la figure, \(\xi=1/3\), les coordonnées
normalisées sont exactement
\[
 \eta=\tfrac12\log2,\quad
 \frac{a_S}{\sqrt{2\hbar}}=2^{1/4}=1.1892071150\ldots,\quad
 \frac{b_S}{\sqrt{2\hbar}}=
 \frac{f_S}{\sqrt{2\hbar}}=R_{\rm el}=2^{-1/4}=0.8408964153\ldots.
\]
Ce sont des nombres de contrôle de non-régression du dessin, non des valeurs attribuées au générateur
angulaire. Le certificat local vérifie également la forme circulaire et les
orientations \(\xi=\pm1/3,\pm3/5\).

\subsection{Réciprocité de forme et distance intérieure}

La représentation rectangulaire retrouve
\begin{equation}
 R_{\rm el}=\sqrt{b_S/a_S}=e^{-\eta/2},\qquad
 R_{\rm el}^4=\frac{1-\xi}{1+\xi},\qquad
 \xi=\frac{1-R_{\rm el}^4}{1+R_{\rm el}^4}.
 \label{eq:ii-elipse-reciproca}
\end{equation}
La preuve consiste à substituer
\(2\eta=\log((1+\xi)/(1-\xi))\). L'échange des axes
envoie \(\eta\) sur \(-\eta\), \(R_{\rm el}\) sur
\(R_{\rm el}^{-1}\) et \(\tau=iR_{\rm el}^2\) sur \(-1/\tau\).
Le produit \(a_Sb_S\) est conservé. Le rayon \(R_{\rm el}\) est adimensionnel ;
le rayon polaire \(\rho_{\partial}\) a pour unité la racine carrée d'action.

L'intérieur admet également la distance projective de Hilbert. Si
\(z_-,x,y,z_+\) sont ordonnés sur une corde de l'ellipse, on définit
\[
 d_H(x,y)=\frac12\log
 \frac{|y-z_-|\,|x-z_+|}{|x-z_-|\,|y-z_+|}.
\]
C'est une distance adimensionnelle : les unités se simplifient dans le birapport.
La normalisation \((Q,P)\mapsto(Q/a_S,P/b_S)\) envoie l'ellipse
sur le disque unité et préserve ce rapport, car les longueurs sur une
même droite sont multipliées par un facteur commun. Elle transporte ainsi la métrique
vers le modèle de Beltrami--Klein. Ni cette distance intérieure ni la distance
entre les foyers ne s'identifie à une longueur atomique.

\subsection{Rayon de Bohr : conséquence de l'échelle d'action}
\label{sec:ii-elipse-bohr}

Dans la réalisation électrodynamique, l'échelle de Bohr s'exprime par
\(a_{\rm B}=\hbar/(m_ec\alpha)\) \cite{codata2022}. Dans la chaîne
du présent article, cette lecture succède aux constructions de
l'action, du vide, de la structure fine et de la masse électronique. On reçoit ici
\(\hbar,m_e,c,\alpha>0\) dans une même réalisation dimensionnelle.
\begin{proposition}[Expression aréale de l'échelle de Bohr]
\label{prop:ii-elipse-bohr}
Dans cette réalisation,
\begin{equation}
 a_{\rm B}=\frac{a_Sb_S}{2m_ec\alpha}
 =\frac{\mathcal A(2\pi)}{2\pi m_ec\alpha}
 =\frac{\bar\lambda_C(m_e)}{\alpha},\qquad
 \bar\lambda_C(m_e)=\frac{\hbar}{m_ec}.
 \label{eq:ii-elipse-bohr}
\end{equation}
\end{proposition}
\begin{proof}
Les identités \(a_Sb_S=2\hbar\) et
\(\mathcal A(2\pi)=2\pi\hbar\), substituées dans l'expression
de Bohr, donnent les deux premiers membres. La définition de
\(\bar\lambda_C(m_e)\) donne le troisième.
\end{proof}

Le numérateur a la dimension d'une action et \(m_ec\) celle d'une quantité de mouvement ; le
quotient est une longueur. Cette composition utilise le produit des
demi-axes et les trois autres grandeurs : elle n'affirme pas que Bohr soit un demi-axe,
un foyer, une distance de Hilbert ou le rayon modulaire. Elle n'assimile pas non plus
la masse électronique au quantum circulaire d'ordre 108. La notation
\(a_{\rm B}\) évite de confondre cette longueur avec la normalisation
\(a_0\) de l'opérateur d'aire.

\subsection{Réalisation inductive--capacitive et continuité constitutive}
\label{sec:ii-elipse-lc}

Le pont électrodynamique de l’ellipse conserve une coordinatisation supplémentaire :
une horloge \(\omega>0\) et une impédance de référence \(Z_*>0\).
La construction retrouvée dans le corpus définit
\begin{equation}
 L_\eta=\frac{Z_*}{\omega}e^{-\eta},\qquad
 C_\eta=\frac{1}{Z_*\omega}e^\eta,
 \qquad Z_\eta=\sqrt{L_\eta/C_\eta}.
 \label{eq:ii-elipse-lc}
\end{equation}
Ici, \(C_\eta\) est la capacité électrique, distincte du contre-angle \(C^*\).
\begin{proposition}[Produit dynamique et quotient constitutif]
\label{prop:ii-elipse-lc}
Dans la coordinatisation précédente,
\begin{equation}
 L_\eta C_\eta=\omega^{-2},\qquad
 \frac{Z_\eta}{Z_*}=e^{-\eta}
 =\frac{b_S}{a_S}=R_{\rm el}^{2}.
 \label{eq:ii-elipse-impedancia}
\end{equation}
Avec \(Q=\sqrt{Z_*}\,q\), \(P=\Phi/\sqrt{Z_*}\), l'énergie
\(H_{LC}=q^2/(2C_\eta)+\Phi^2/(2L_\eta)\) prend la forme
\[
 H_{LC}=\frac\omega2(e^{-\eta}Q^2+e^\eta P^2).
\]
Sur \(\gamma(\theta)\), ses composantes électrique et magnétique sont
\(U_E=\hbar\omega\cos^2\theta\) et
\(U_B=\hbar\omega\sin^2\theta\) ; leur somme vaut \(\hbar\omega\).
\end{proposition}
\begin{proof}
Dans le produit de \eqref{eq:ii-elipse-lc}, les exponentielles et
\(Z_*\) se simplifient ; son quotient vaut \(Z_*^2e^{-2\eta}\). Prendre la racine
positive et utiliser \eqref{eq:ii-elipse-reciproca} donne
\eqref{eq:ii-elipse-impedancia}. Substituer \(q=Q/\sqrt{Z_*}\)
et \(\Phi=\sqrt{Z_*}P\) dans \(H_{LC}\) donne sa forme équilibrée.
Enfin, les demi-axes de \(\gamma\) simplifient les facteurs
\(e^{\mp\eta}\), et \(\cos^2\theta+\sin^2\theta=1\).
\end{proof}

Ce pont distingue la forme, l'horloge et le transducteur. Le choix
du sens dynamique doit s'accorder avec la convention symplectique : pour
les équations \(\dot Q=\partial_PH\),
\(\dot P=-\partial_QH\), la paramétrisation précédente satisfait
\(\dot\theta=-\omega\). Les identités énergétiques ne dépendent pas
de l'inversion de ce parcours.

L’article ultérieur sur le vide prolongera les canaux déjà produits
\(q_\pm=e^{-(A_{\rm rad}\pm C^*_{\rm rad})}\) et l’incidence
\(90/120\) au moyen de son lecteur résolvant. Ses réponses
\(r_\pm=(1-q_\pm^{90})/(1-q_\pm^{120})\) expriment
\(\widehat Z=r_-/r_+\) et \(\widehat c=(r_+r_-)^{-1}\).
Ce sont des réponses spectrales, non des rayons géométriques. L’identité
\eqref{eq:ii-elipse-impedancia} fixe ici la réalisation LC ;
l’identifier à l’impédance constitutive du vide conserve l’application
entre les deux lecteurs et leurs coordinatisations. On n’égalise pas leurs valeurs parce qu’ils partagent
l’antécédent angulaire.

\subsection{Confrontation de la constante de Planck réduite}
\label{sec:revision-planck-contraste}

Le Système international en vigueur depuis le 20 mai 2019 fixe
exactement \(h=6.62607015\times10^{-34}\,\mathrm{J\,s}\)
\cite{bipm2018}. Sa constante réduite est donc
\[
 \hbar_{\mathrm{SI}}=\frac{6.62607015}{2\pi}\,10^{-34}\,
 \mathrm{J\,s}
 =1.0545718176461563912\ldots\times10^{-34}\,\mathrm{J\,s}.
\]
Les points de suspension correspondent au développement d'une expression
exacte, et non à une incertitude expérimentale de \(h\) dans le SI actuel.
Cette valeur de référence est utilisée ici après l'évaluation des
lecteurs de HMT ; sa fonction est comparative.

Les deux coordonnées du retour de l'article sont
\begin{align*}
 H_5&=1.0545718176461544696\ldots,\\
 \eta_{\mathrm{ret}}&=1.0545718169150390611\ldots.
\end{align*}
Dans l'identification de \(\mathcal U_S\) à \(\mathrm{J\,s}\), la
différence relative de la section antérieure par rapport à
\(\hbar_{\mathrm{SI}}\) est d'environ \(-1.82\times10^{-15}\).
Celle de la section postérieure est d'environ \(-6.93\times10^{-10}\).
La seconde différence contient l'effet du retour utilisé par le
correcteur électronique ; elle n'équivaut pas à l'erreur d'évaluation de la
première section.

La précision de ces comparaisons permet de reconnaître la proximité
quantitative de la construction de \(H_5\) et, simultanément, de distinguer
proximité et identité exacte. Dans le présent manuscrit, les expressions
présentées ne fournissent pas une égalité exacte avec \(h/(2\pi)\).
Le résultat démontré sur \(\Sigma_H\) est sa décade \(-34\) ;
le résultat analytique sur \(H_5\) est l'évaluation du caractère
déclaré. L'identification physique est examinée dans cette coordinatisation et avec les
différences explicites que nous venons de donner. La valeur SI n'intervient
ni dans la définition de la norme déterminantale ni dans la récurrence régionale.

\endgroup

% Ampliación propia de III; no modifica el núcleo común copiado de I.
% Residencia propuesta: antes de Retorno de hoja, incidencia areal--volumétrica.
% Procedencia y controles: VACANCIAS_PREFACTOR_PROCEDENCIA.md, en esta carpeta.
\section{Polarisation discrète et mode tritique des projections arithmétiques}
\label{iii:sec:vacancias-prefactor}

Le raffinement et la comparaison entre les feuillets d’APP fournissent deux
observations complémentaires. La première suit l’écart accumulé
entre le nombre de transitions et la capacité acquise ; la seconde conserve
le mode du sélecteur tritique dans l’appariement des fibres additives et
multiplicatives. La polarisation discrète et le préfacteur électronique
s’obtiennent à partir de ces opérations respectives. Leur articulation avec le vide
requiert de maintenir tant l’information temporelle que les espaces sur
lesquels agissent les opérateurs.

\subsection{Discrépance du raffinement et convention des extrémités}

Les cardinaux du bloc ternaire et de sa représentation décimale fixent
\[
 \rho=\log_{1000}729=\log_{10}9,\qquad
 \delta=1-\rho=\log_{10}(10/9).
\]
On a $0<\delta<1$. En outre, $\delta$ est irrationnel : une égalité $\delta=p/q$, avec $q>0$, impliquerait $10^{q-p}=9^q$, incompatible avec la décomposition en facteurs premiers. Ce décalage intervient déjà dans le codage du raffinement. Sa lecture centrée est explicitée ici, en conservant l'origine de phase et la convention des extrémités.

\begin{definition}[Codages inférieur et supérieur]
Pour $\theta\in\mathbb R$ et $n\in\mathbb Z$, on définit
\begin{align}
 z_n^{\downarrow}(\theta)
 &=\lfloor(n+1)\delta+\theta\rfloor-\lfloor n\delta+\theta\rfloor,
 \label{iii:vp:lower}\\
 z_n^{\uparrow}(\theta)
 &=\lceil(n+1)\delta+\theta\rceil-\lceil n\delta+\theta\rceil.
 \label{iii:vp:upper}
\end{align}
Les deux prennent leurs valeurs dans $\{0,1\}$. La flèche distingue exclusivement la convention des extrémités des intervalles ; l'orientation tritique conserve sa définition précédente.
\end{definition}

Le codage inférieur coïncide avec celui utilisé dans le développement
du raffinement. Le codage supérieur permet de retrouver la formulation de la
polarisation avec la fonction plafond. Les deux expressions sont reliées par
un terme de frontière exactement déterminé.

\begin{proposition}[Changement de convention et terme de frontière]
Soit $b_n(\theta)=\mathbf1_{\mathbb Z}(n\delta+\theta)$. Alors
\begin{equation}
 z_n^{\uparrow}(\theta)-z_n^{\downarrow}(\theta)
 =b_n(\theta)-b_{n+1}(\theta).
 \label{iii:vp:endpoint}
\end{equation}
En particulier, pour $N\geq0$,
\[
 \sum_{n=0}^{N-1}(z_n^{\uparrow}-z_n^{\downarrow})
 =b_0-b_N.
\]
\end{proposition}
\begin{proof}
Pour tout réel $x$, $\lceil x\rceil=\lfloor x\rfloor+1-
\mathbf1_{\mathbb Z}(x)$. Appliquer cette identité aux deux extrémités
de chaque accroissement donne \eqref{iii:vp:endpoint} ; leur somme est télescopique.
L’irrationalité de $\delta$ implique que, $\theta$ étant fixée, au plus
un indice entier satisfait $n\delta+\theta\in\mathbb Z$.
\end{proof}

Avec $\theta=0$, le premier symbole supérieur est $1$ et l'inférieur est $0$ ; pour $n\geq1$, ils coïncident. Le terme initial appartient au registre et est conservé lors du changement de convention. Les densités et les bornes de discrépance sont transportées avec ce terme, sans modifier la mise à jour du TPK ni recommencer l'histoire à chaque retour de phase.

\subsection{Polarisation bornée et bilan temporel exact}

\begin{definition}[Polarisation discrète centrée]
La polarisation du codage supérieur, relativement à la densité
$\delta$, est
\begin{equation}
 P_N^{\uparrow}(\theta)
 =\sum_{n=0}^{N-1}\bigl(z_n^{\uparrow}(\theta)-\delta\bigr),
 \qquad P_0^{\uparrow}=0.
 \label{iii:vp:polarization}
\end{equation}
Son accroissement temporel est
\begin{equation}
 J_N^{\uparrow}(\theta)
 =P_{N+1}^{\uparrow}(\theta)-P_N^{\uparrow}(\theta)
 =z_N^{\uparrow}(\theta)-\delta.
 \label{iii:vp:current}
\end{equation}
\end{definition}

Le centrage sépare la croissance uniforme de densité $\delta$ de
l’information de phase. Ainsi, $P_N^{\uparrow}$ mesure un déséquilibre accumulé,
non le nombre total de vacances. L’accroissement $J_N^{\uparrow}$ conserve
l’impulsion locale et le fond par rapport auquel elle est évaluée.

\begin{theorem}[Borne uniforme et continuité discrète temporelle]
Pour $N\geq0$ et toute phase $\theta$,
\begin{equation}
 P_N^{\uparrow}(\theta)
 =\lceil N\delta+\theta\rceil-\lceil\theta\rceil-N\delta,
 \qquad |P_N^{\uparrow}(\theta)|<1.
 \label{iii:vp:bound}
\end{equation}
Pour tous entiers $0\leq a<b$,
\begin{equation}
 \sum_{n=a}^{b-1}J_n^{\uparrow}(\theta)
 =P_b^{\uparrow}(\theta)-P_a^{\uparrow}(\theta),
 \qquad
 \left|\sum_{n=a}^{b-1}J_n^{\uparrow}(\theta)\right|<1.
 \label{iii:vp:continuity}
\end{equation}
La densité de la suite des vacances est $\delta$.
\end{theorem}
\begin{proof}
La somme des accroissements de plafond de \eqref{iii:vp:upper} produit la première identité. En écrivant $r(x)=\lceil x\rceil-x\in[0,1)$, on obtient $P_N^{\uparrow}=r(N\delta+\theta)-r(\theta)$ ; sa valeur absolue est inférieure à un. Sur l'intervalle $[a,b)$, on obtient de même $r(b\delta+\theta)-r(a\delta+\theta)$, ce qui prouve la seconde borne, plus précise que la somme de deux bornes individuelles. Enfin, $N^{-1}\sum_{n=0}^{N-1}z_n^{\uparrow}
=\delta+N^{-1}P_N^{\uparrow}\to\delta$.
\end{proof}

Le codage inférieur admet les mêmes démonstrations en remplaçant
$r(x)$ par $\lfloor x\rfloor-x\in(-1,0]$. En particulier,
$P_N^{\uparrow}-P_N^{\downarrow}=b_0-b_N$. L’écart borné
utilisé précédemment dans les lectures de Dirichlet et la polarisation
présente conservent donc une relation explicite d’origine et d’extrémités.
L’apériodicité n’introduit pas non plus de dérive séculaire de ce déséquilibre :
l’histoire peut se prolonger indéfiniment tandis que l’écart
centré reste uniformément borné.

\subsection{Réalisation dans les droites de charge et de temps}

Pour réaliser dimensionnellement le bilan, on spécifie une droite de charge
$\mathcal L_Q$ de section non nulle $u_Q$ et une suite croissante
d’instants $t_N$ dans une droite temporelle orientée. Avec
$\Delta t_N=t_{N+1}-t_N>0$, la charge centrée et le courant moyen de
chaque intervalle sont définis par
\begin{equation}
 Q_N=P_N^{\uparrow}u_Q,\qquad
 I_N=\frac{Q_{N+1}-Q_N}{\Delta t_N}
     =\frac{u_Q}{\Delta t_N}\bigl(z_N^{\uparrow}-\delta\bigr).
 \label{iii:vp:dimensional}
\end{equation}
Ainsi, $Q_N\in\mathcal L_Q$ et $I_N\in\mathcal L_Q\otimes\mathcal L_T^{-1}$. La section $u_Q$ et les intervalles temporels sont des données de cette réalisation ; la construction arithmétique précédente a déjà déterminé $P_N^{\uparrow}$ et $J_N^{\uparrow}$.

\begin{proposition}[Conservation du bilan sous réalisation dimensionnelle]
Pour $a<b$,
\begin{equation}
 Q_b-Q_a=\sum_{n=a}^{b-1}\Delta t_n I_n,
 \qquad
 \left|\frac{Q_b-Q_a}{u_Q}\right|<1.
 \label{iii:vp:dimensional-balance}
\end{equation}
Si tous les intervalles ont longueur $t_0$, les deux courants possibles
sont $-\delta u_Q/t_0$ et $(1-\delta)u_Q/t_0$.
\end{proposition}
\begin{proof}
Multiplier chaque identité de \eqref{iii:vp:dimensional} par
$\Delta t_n$ et sommer fait télescoper les charges. La borne est
\eqref{iii:vp:continuity} ; les deux valeurs locales s’obtiennent à partir de
$z_n^{\uparrow}\in\{0,1\}$.
\end{proof}

La continuité démontrée est temporelle et se rapporte à la charge centrée. Une réalisation spatiale incorpore en outre la localisation des charges, l'incidence entre cellules et l'identification des flux de frontière. Dans une description constitutive par canaux, le couplage est spécifié au moyen d'une application d'excitation $\iota:\mathcal L_Q\to\mathcal H_{\rm ch}\otimes\mathcal L_Q$ et d'une application d'observation $\ell:\mathcal H_{\rm ch}\otimes\mathcal L_Q\to\mathcal L_Q$, où $\mathcal H_{\rm ch}$ est une réalisation réelle de l'espace des canaux. Si les deux sont linéaires et indépendantes de l'indice temporel, la réponse d'un opérateur $\mathcal V$ fixe est transportée par
\[
 Q_N^{\rm resp}=\ell(\mathcal V\otimes I_{\mathcal L_Q})\iota(Q_N),
 \qquad
 I_N^{\rm resp}=\frac{Q_{N+1}^{\rm resp}-Q_N^{\rm resp}}{\Delta t_N}.
\]
La linéarité conserve le bilan télescopique. Cette composition
explicite les données d’une identification constitutive : la séquence
scalaire détermine l’impulsion, tandis que $\iota$, $\mathcal V$ et $\ell$
déterminent son excitation, son transport et son observation par canaux. Les
identités \eqref{iii:vp:polarization}--\eqref{iii:vp:dimensional-balance}
établissent la première construction indépendamment de ce choix de
réalisation. L’opérateur du vide est construit dans la
section~\ref{iii:sec:constitutiva} sur ses canaux et ses projecteurs propres.

\subsection{Appariement des fibres additives et multiplicatives}

La seconde observation provient de la réalisation tridimensionnelle d’APP,
$\mathcal A_3=D_9^3$, avec une mesure uniforme et
$D_9=\{1,\ldots,9\}$. On conserve les observables
\[
 \Sigma(x)=\rho_9(x_1+x_2+x_3),\qquad
 \Pi(x)=\rho_9(x_1x_2x_3),
\]
où $\rho_9(n)=1+[(n-1)\bmod9]$ pour $n>0$. Dans
$\ell^2(\mathcal A_3)$, soient $P_{\Sigma,3}$ et $P_{\Pi,3}$ les
espérances conditionnelles sur leurs fibres. Par exemple,
\[
 (P_{\Sigma,3}f)(x)=
 \frac1{|\Sigma^{-1}(\Sigma(x))|}
 \sum_{y:\,\Sigma(y)=\Sigma(x)}f(y).
\]
Ce sont des projections orthogonales définies sur le même support. La matrice
d’incidences $N_{ab}=|\Sigma^{-1}(a)\cap\Pi^{-1}(b)|$ résulte de
l’énumération des $729$ triplets, en additionnant et en multipliant leurs coordonnées avant
d’appliquer $\rho_9$ :
\begin{equation}
 N=9\begin{pmatrix}
 0&1&1&0&1&1&0&1&4\\
 1&0&1&1&0&1&1&0&4\\
 1&0&2&0&1&2&0&0&3\\
 0&1&1&0&1&1&0&1&4\\
 1&0&1&1&0&1&1&0&4\\
 0&0&2&1&0&2&0&1&3\\
 0&1&1&0&1&1&0&1&4\\
 1&0&1&1&0&1&1&0&4\\
 0&1&2&0&0&2&1&0&3
 \end{pmatrix}.
 \label{iii:vp:joint-matrix}
\end{equation}
Les lignes ont pour somme $81$ et les sommes des colonnes sont
\[
 (m_1,\ldots,m_9)=(36,36,108,36,36,108,36,36,297).
\]
L’appariement des indicatrices normalisées des fibres
est donc $C_{ab}=N_{ab}/\sqrt{81m_b}$. La normalisation recueille les
multiplicités produites par APP, non un choix de valeurs singulières.

\begin{proposition}[Mode singulier du sélecteur tritique]
Dans la coordinatisation canonique, le vecteur du sélecteur est
$m_{\rm ph}=(1,-1,0)^3$. On a
\begin{equation}
 Cm_{\rm ph}=C^{\mathsf T}m_{\rm ph}=-\tfrac12m_{\rm ph}.
 \label{iii:vp:trit-mode}
\end{equation}
Ainsi, la valeur singulière correspondante est $s=1/2$.
\end{proposition}
\begin{proof}
Dans les six colonnes où $m_{\rm ph}$ est non nul, $m_b=36$ et
$C_{ab}=N_{ab}/54$. Multiplier les lignes de
\eqref{iii:vp:joint-matrix} par $(1,-1,0)^3$ produit
$-m_{\rm ph}/2$. Pour le produit transposé, les sommes signées
des colonnes $3,6,9$ sont nulles ; les six autres donnent les mêmes
valeurs $-m_{{\rm ph},b}/2$. Cela prouve les deux identités et, par
composition, $CC^{\mathsf T}m_{\rm ph}=m_{\rm ph}/4$.
\end{proof}

\subsection{Norme du commutateur et préfacteur électronique}

\begin{theorem}[Norme exacte d’incompatibilité arithmétique]
Les deux projections satisfont
\begin{equation}
 \bigl\|[P_{\Sigma,3},P_{\Pi,3}]\bigr\|=\frac{\sqrt3}{4}.
 \label{iii:vp:prefactor}
\end{equation}
Le maximum est atteint dans le plan principal correspondant au mode
tritique de \eqref{iii:vp:trit-mode}.
\end{theorem}
\begin{proof}
La matrice rationnelle $M=CC^{\mathsf T}$ a pour entrées $M_{ac}=\sum_bN_{ab}N_{cb}/(81m_b)$. Son spectre complet se vérifie au moyen des sous-espaces indépendants suivants : les vecteurs $(1,\ldots,1)$, $(1,-1,0)^3$ et $(1,1,-2)^3$ ont pour valeurs propres $1$, $1/4$ et $7/132$, respectivement ; le sous-espace des vecteurs à support dans $\{3,6,9\}$ dont la somme est nulle a pour dimension $2$ et pour valeur propre $1/18$ ; les vecteurs de somme nulle à support dans $\{1,4,7\}$ ou dans $\{2,5,8\}$ forment un noyau de dimension $4$. La substitution dans la matrice vérifie ces actions et les dimensions ont pour somme $9$. Par conséquent,
\[
 \operatorname{spec}M=
 \{1,1/4,7/132,1/18,1/18,0,0,0,0\}.
\]
Pour deux projections orthogonales, une valeur singulière $s\in(0,1)$ de
l’appariement détermine un plan principal dans lequel leurs matrices sont
\[
 P=\begin{pmatrix}1&0\\0&0\end{pmatrix},\qquad
 Q=\begin{pmatrix}s^2&s\sqrt{1-s^2}\\
 s\sqrt{1-s^2}&1-s^2\end{pmatrix}.
\]
Cette forme s’obtient en prenant un vecteur unitaire $u$ de la première
image avec $PQu=s^2u$ et en le complétant par
$(Qu-s^2u)/(s\sqrt{1-s^2})$. Leur commutateur est
$s\sqrt{1-s^2}\bigl(\begin{smallmatrix}0&1\\-1&0\end{smallmatrix}\bigr)$.
Les secteurs de valeur singulière $0$ ou $1$ contribuent par zéro.
Par conséquent, la norme totale est le maximum de
$\sqrt{\lambda(1-\lambda)}$ sur le spectre de $M$. Ce maximum
est $\sqrt{3/16}$, atteint en $\lambda=1/4$.
\end{proof}

La norme est le préfacteur qui intervient dans la composition électronique.
Son obtention conserve le support tridimensionnel, les fibres et le mode
du sélecteur ; l’évaluation scalaire ne remplace pas cette structure. En
particulier, connaître $\sqrt3/4$ permet de connaître la grandeur maximale de
l’incompatibilité, tandis que les projecteurs conservent les directions
dans lesquelles elle agit.

\clearpage
\subsection{Défaut quadratique et réponse des secteurs \(90/120\)}

Le mode tritique détermine simultanément
\begin{equation}
 1-s^2=\frac34=\frac{90}{120},\qquad s=\frac12.
 \label{iii:vp:defect-ratio}
\end{equation}
La première expression est le défaut quadratique du plan principal d'APP. La dernière conserve les longueurs des deux secteurs du calendrier. La réponse constitutive utilise ces secteurs sur un autre objet : un canal contractant $q\in(0,1)$. Son évaluation est
\[
 r(q)=\frac{1-q^{90}}{1-q^{120}}.
\]
La relation entre le quotient des secteurs et cette réponse peut être
précisée sans identifier leurs opérateurs.

\begin{proposition}[Limite du quotient de réponse]
Pour $0<q<1$,
\begin{equation}
 \frac34<r(q)<1,\qquad
 \lim_{q\uparrow1}r(q)=\frac34=1-s^2.
 \label{iii:vp:response-limit}
\end{equation}
\end{proposition}
\begin{proof}
Avec $u=q^{30}\in(0,1)$, la factorisation de $1-u^3$ et
$1-u^4$ donne
$r(q)=(1+u+u^2)/(1+u+u^2+u^3)$. La limite en $u=1$ est
$3/4$. De plus,
$4(1+u+u^2)-3(1+u+u^2+u^3)
=(1-u)(3u^2+2u+1)>0$, et le dénominateur dépasse le numérateur
de $u^3>0$. On obtient les deux inégalités.
\end{proof}

La coïncidence exacte concerne le défaut du mode singulier et la limite
de la réponse des secteurs. À l’intérieur contractif, la réponse
dépend de $q$ et est strictement supérieure à $3/4$. Le développement
constitutif conserve ses deux canaux $q_\pm$ et ses projecteurs
$P_\pm$ ; les projecteurs $P_{\Sigma,3}$ et $P_{\Pi,3}$ conservent,
pour leur part, les fibres arithmétiques qui ont déterminé le préfacteur.
Une équivalence opératorielle entre les deux réalisations exige une application
entre leurs espaces qui entrelace les actions correspondantes.
Les équations \eqref{iii:vp:trit-mode}, \eqref{iii:vp:prefactor} et
\eqref{iii:vp:response-limit} précisent l’information partagée et
maintiennent visibles les données de chaque réalisation.

\section[Retour de feuillet, incidence aréale--volumétrique et orientation constitutive]{Retour de feuillet, incidence aréale--volumétrique\\et orientation constitutive}
\label{iii:sec:hojas}

La structure des feuillets intervient avant l’évaluation des grandeurs
électromagnétiques. Sa fonction consiste à distinguer la fermeture d’une projection,
la restitution d’une orientation et l’évolution du registre transporté.
Ces opérations agissent sur des objets liés, mais de types différents : une
trajectoire projetée peut s’être refermée tandis que son relèvement reste
sur le feuillet opposé ; l’orientation peut avoir été restituée tandis que la mémoire
conserve les deux tours effectués. La construction constitutive doit recevoir
cette information avec la structure qui la détermine.

\subsection{Retour projeté et restitution du revêtement orienté}

Soit $\ell$ une boucle de retour d’une extension survivante et soit
$\varepsilon:\langle\ell\rangle\to\{+1,-1\}$ le caractère de son
relèvement, avec $\varepsilon(\ell)=-1$. Le fibré d’intervalles associé
a pour présentation
\begin{equation}
 \mathcal M=([0,1]\times[-1,1])/{(0,u)\sim(1,-u)}.
 \label{iii:eq:mobius}
\end{equation}
La donnée décisive est le signe du transport sur la boucle spécifiée.
L’identification des extrémités réalise géométriquement ce caractère.

\begin{proposition}[Deux ordres de clôture]
Un tour de la boucle ferme la base et inverse l’orientation de la fibre. Deux
tours restituent l’orientation. Dans la paramétrisation angulaire du lecteur de
bord, ces retours correspondent respectivement à la clôture aréale $2\pi$ et
à la clôture volumétrique $4\pi$ du revêtement orienté.
\end{proposition}
\begin{proof}
La relation d’équivalence de \eqref{iii:eq:mobius} identifie la fibre terminale
à la fibre initiale au moyen de $u\mapsto-u$. La composition de deux transports est
$u\mapsto-(-u)=u$. Paramétrer un tour de la base par $2\pi$ attribue $4\pi$
à son carré. L’assertion concerne ce relèvement et cette lecture.
\end{proof}

La mémoire du TPK conserve en outre la trajectoire parcourue et les mises à jour du registre. Le carré du caractère de signe est l'identité, mais n'impose pas l'annulation de ces mises à jour. La représentation spinorielle, lorsqu'elle sera utilisée, devra transporter cette donnée au moyen de sa propre application.

\subsection{Descente de l’incidence depuis le calendrier tritique}

Dans le bloc primitif $(+1,-1,0)$, le TPK met à jour le mode arithmétique au moyen de
$+1:m\mapsto\Sigma$, $-1:m\mapsto\Pi$ et $0:m\mapsto m$. Son orbite cyclique
est $(\Sigma,\Pi,\Pi)$. Le lecteur de feuillets attribue $\Sigma$ au feuillet aréal et
$\Pi$ au feuillet volumétrique. Dans l’ordre des bases $(e_{\rm ar},e_{\rm vol})$,
l’incidence des arêtes est représentée par
\begin{equation}
 F_{\rm av}=\begin{pmatrix}0&1\\1&1\end{pmatrix}.
 \label{iii:eq:fav}
\end{equation}
Les trois paires consécutives sont $\Sigma\to\Pi$, $\Pi\to\Pi$ et $\Pi\to\Sigma$, chacune de multiplicité un. Les calendriers de longueurs $9$, $27$ et $108$ contiennent $3$, $9$ et $36$ copies du bloc primitif. Leurs matrices d'incidence sont donc ces multiples de $F_{\rm av}$.

Cette construction conserve la différence entre fréquence chronologique et mesure
de l’espace des continuations. La première assigne des poids $(1/3,2/3)$ au cycle
des modes. La seconde se définit sur l’enveloppe symbolique minimale de mémoire
un qui admet exactement les couples précédents. Sa matrice d’adjacence est
$F_{\rm av}$ ; les dénombrements de trajectoires appartiennent à cette enveloppe.

\Needspace{10\baselineskip}
\begin{proposition}[Poids de retour des deux feuillets]
Soient $F_0=0$, $F_1=1$ et $F_{n+1}=F_n+F_{n-1}$. Pour $n\geq1$,
\begin{equation}
 F_{\rm av}^{n}=
 \begin{pmatrix}F_{n-1}&F_n\\F_n&F_{n+1}\end{pmatrix}.
 \label{iii:eq:fav-powers}
\end{equation}
La mesure de Parry de cette enveloppe attribue aux feuillets des poids proportionnels à $(1,\varphi^2)$, et le rapport des décomptes diagonaux converge vers $\varphi^{-2}$.
\end{proposition}
\begin{proof}
La formule des puissances se prouve par induction en utilisant la récurrence.
Le carré de $F_{\rm av}$ est strictement positif. Son vecteur positif de
Perron est $(1,\varphi)$, de valeur propre $\varphi$, car
$\varphi^2=\varphi+1$. La matrice étant symétrique, les poids de Parry sont
proportionnels aux carrés des composantes de ce vecteur. L’autre racine
du polynôme caractéristique est $-\varphi^{-1}$ ; son module est strictement
inférieur à $\varphi$, ce qui donne la limite des quotients diagonaux.
\end{proof}

L’identification spectrale de cette incidence est postérieure à sa construction
à partir du calendrier. La coordonnée régionale $\varphi$ conserve sa généalogie
coinductive antérieure : la reconnaissance de la valeur propre ne sélectionne pas le calendrier.

\subsection{Marquage de la frontière et rapport surfacique--volumétrique}

Soit $\Pi_{\rm HMT}(X_\infty)$ la valeur du lecteur régional de clôture déjà construit. Avec les projecteurs diagonaux de feuillet, on définit
\begin{equation}
 D_\pi=\Pi_{\rm HMT}(X_\infty)P_{\rm ar}+P_{\rm vol},\qquad
 \Theta_n=
 \frac{\langle e_{\rm ar},D_\pi F_{\rm av}^{n}e_{\rm ar}\rangle}
 {\langle e_{\rm vol},F_{\rm av}^{n}e_{\rm vol}\rangle}.
 \label{iii:eq:marked-return}
\end{equation}
La normalisation volumétrique est unitaire ; la clôture régionale marque une seule fois
le retour aréal. En appliquant \eqref{iii:eq:fav-powers},
\begin{equation}
 \Theta_n=\Pi_{\rm HMT}(X_\infty)\frac{F_{n-1}}{F_{n+1}}
 \longrightarrow\frac{\pi}{\varphi^2}.
 \label{iii:eq:av-limit}
\end{equation}
Le rapport de retour est déterminé par l'incidence, la mesure et le marquage. Le quotient $2\pi/4\pi$ décrit, en revanche, les deux ordres de clôture du revêtement. Les deux relations participent à la lecture des feuillets sans se substituer l'une à l'autre.

\subsection{Transport de l’orientation aux canaux constitutifs}

Dans une représentation à deux canaux, soit $\mathcal R$ une involution
autoadjointe et $\mathcal S$ un opérateur unitaire d’échange qui satisfait
$\mathcal S\mathcal R\mathcal S^{-1}=-\mathcal R$. Les coordonnées angulaires
déjà construites s’écrivent $x=A_{\rm rad}$ et $y=C^*_{\rm rad}$, avec $x>|y|$.
Alors
\begin{equation}
 B=xI+y\mathcal R,\qquad T=e^{-B},\qquad
 \mathcal ST(y)\mathcal S^{-1}=T(-y).
 \label{iii:eq:channel-conjugation}
\end{equation}
En effet, la conjugaison transforme $B(y)$ en $B(-y)$ et commute avec sa série
exponentielle convergente. Les valeurs propres $q_\pm=e^{-x\mp y}$ sont échangées.
Pour fixer l’ordre des réponses, on utilise la chambre $x>y>0$, où
$0<q_+<q_-<1$. L’application $(x,y)\mapsto(x,-y)$ la transporte dans la
chambre opposée $x>-y>0$, où $0<q_-<q_+<1$. Toutes deux appartiennent au domaine
contractif $x>|y|$. À la frontière $y=0$, les deux canaux coïncident.
La conjugaison échange les chambres, elle ne conserve pas l’ordre strict de
leurs étiquettes ; elle agit sur les projecteurs au moyen de
$\mathcal SP_\pm\mathcal S^{-1}=P_\mp$, avec
$P_\pm=(I\pm\mathcal R)/2$.
Cette application concrète transmet l’orientation à l’opérateur constitutif
développé ci-après.

\section{Opérateur constitutif et reconstruction des canaux du vide}
\label{iii:sec:constitutiva}

Les grandeurs constitutives s'obtiennent au moyen d'une opération sur le transport angulaire. L'information de départ comprend les canaux, leurs projecteurs et l'orientation qui les distingue. Cette organisation permet de suivre séparément la construction de l'opérateur, l'évaluation de ses fonctions spectrales et sa réalisation dimensionnelle ultérieure.

\subsection{Réponse positive sur deux feuillets}

Dans la chambre orientée $x>y>0$ de la section~\ref{iii:sec:hojas}, soient
$P_\pm=(I\pm\mathcal R)/2$ et
$T=q_+P_++q_-P_-$, avec $0<q_+<q_-<1$. Les incidences construites par le
calendrier fixent les secteurs $90$ et $120$. La réponse constitutive est
\begin{equation}
 \mathcal V=(I-T^{90})(I-T^{120})^{-1}.
 \label{iii:eq:vacuum-op}
\end{equation}
Cette définition reçoit le transport provenant du TPK et conserve les
attributions de feuillet. La règle de réponse et son domaine font explicitement partie
de la construction.

Le quotient provient de la comparaison des deux incidences du calendrier sur un même transport, et non de l'ajustement séparé de quatre grandeurs. Si $S=T^{30}$ et $\Sigma_j(S)=I+S+\cdots+S^{j-1}$, la factorisation de $I-S^j$ donne
\begin{equation}
 \mathcal V\,\Sigma_4(S)=\Sigma_3(S),\qquad
 \mathcal V=\Sigma_3(S)\Sigma_4(S)^{-1}.
 \label{iii:eq:completion-response}
\end{equation}
Les facteurs $\Sigma_3$ et $\Sigma_4$ expriment respectivement les
secteurs $90=3\cdot30$ et $120=4\cdot30$. Comme le spectre de $S$ est
contenu dans $(0,1)$, $\Sigma_4(S)$ est inversible : l’équation de
complétion détermine une unique réponse pour le transport fixé.
Cette unicité concerne la règle constitutive déclarée et ses
incidences. L’identification électromagnétique ajoute l’attribution
orientée des feuillets et les droites dimensionnelles précisées plus
bas.

\begin{proposition}[Décomposition et positivité]
La réponse est bien définie et
\begin{equation}
 \mathcal V=r_+P_++r_-P_-,\qquad
 r_\pm=\frac{1-q_\pm^{90}}{1-q_\pm^{120}}.
 \label{iii:eq:rchannels}
\end{equation}
Ses valeurs propres satisfont $3/4<r_-<r_+<1$ dans cette chambre.
\end{proposition}
\begin{proof}
Chaque coefficient $1-q_\pm^{120}$ est positif. L'inverse de $I-T^{120}$ est la somme des projecteurs divisés par ces coefficients, ce qui prouve \eqref{iii:eq:rchannels}. Avec $s=q^{30}$, on obtient
\begin{equation}
 f(s)=\frac{1+s+s^2}{1+s+s^2+s^3},\qquad
 f'(s)=-\frac{s^2(3+2s+s^2)}{(1+s+s^2+s^3)^2}<0.
 \label{iii:eq:monotone}
\end{equation}
Les limites en $0$ et en $1$ sont $1$ et $3/4$, respectivement. La monotonie
et l’ordre des canaux concluent la preuve.
\end{proof}

\subsection{Quatre coordonnées d’une réponse commune}

On définit les évaluations normalisées
\begin{equation}
 \widehat\varepsilon=r_+^2,\qquad
 \widehat\mu=r_-^2,\qquad
 \widehat Z=\frac{r_-}{r_+},\qquad
 \widehat c=\frac1{r_+r_-}.
 \label{iii:eq:four}
\end{equation}
Le premier couple conserve les réponses quadratiques assignées à chaque feuillet.
Le produit recueille le mode commun ; le quotient distingue leur orientation relative.
On vérifie exactement
\begin{equation}
 \widehat\mu\widehat\varepsilon=\widehat c^{-2},\qquad
 \frac{\widehat\mu}{\widehat\varepsilon}=\widehat Z^2,
 \qquad
 r_+=\frac1{\sqrt{\widehat Z\widehat c}},\quad
 r_-=\sqrt{\frac{\widehat Z}{\widehat c}}.
 \label{iii:eq:four-identities}
\end{equation}
Les racines positives fixent l'inversion. Sous la conjugaison des feuillets, lue par rapport aux marques $P_\pm$ fixes, $\widehat\varepsilon$ et $\widehat\mu$ sont échangés, $\widehat Z$ est inversé et $\widehat c$ est conservé. La réponse est alors transportée vers la chambre $x>-y>0$ : les formules restent valables, avec les ordres $q_-<q_+$ et $r_+<r_-$ inversés. Par conséquent, le mode commun et l'orientation sont des données récupérables au moyen d'observables complémentaires.

\begin{proposition}[Détermination des évaluations quadratiques]
Soit $W=r_+P_++r_-P_-$ une réponse positive sur les deux feuillets.
Étant fixées les lectures de produit et de quotient
\[
 \widehat c^{-1}=\det W=r_+r_-,
 \qquad \widehat Z=r_-/r_+,
\]
il existe une unique paire positive $(\widehat\varepsilon,\widehat\mu)$
qui vérifie
$\widehat\mu\widehat\varepsilon=\widehat c^{-2}$ et
$\widehat\mu/\widehat\varepsilon=\widehat Z^2$.
C’est précisément $(r_+^2,r_-^2)$.
\end{proposition}
\begin{proof}
La positivité permet de prendre les racines sans ambiguïté :
\[
 \widehat\mu=\frac{\widehat Z}{\widehat c}
 =\frac{r_-}{r_+}r_+r_-=r_-^2,\qquad
 \widehat\varepsilon=\frac1{\widehat Z\widehat c}
 =\frac{r_+}{r_-}r_+r_-=r_+^2.
\]
Les expressions vérifient les deux relations et en sont les seules racines
positives.
\end{proof}

Ainsi, les carrés ne constituent pas deux choix indépendants ajoutés
aux canaux : ils sont déterminés par les lectures multiplicative et
orientée, une fois fixées les relations constitutives. Cette proposition
établit l’unicité au sein de ce système de règles. L’interprétation
des règles comme réponse électromagnétique conserve son contenu
supplémentaire de représentation physique.

La lecture déterminantale du transport et la réponse constitutive
doivent également conserver leur ordre. Pour $T=e^{-xI-y\mathcal R}$, on a
$\det T=e^{-2x}$, tandis que
\[
 \det\mathcal V=f(q_+^{30})f(q_-^{30})=\widehat c^{-1}.
\]
Les deux proviennent du même opérateur à deux feuillets, mais évaluent des fonctions différentes de son spectre. En général, appliquer d'abord le déterminant puis une fonction rationnelle n'équivaut pas à prendre le déterminant de cette fonction de l'opérateur. Maintenir les deux valeurs propres jusqu'à l'achèvement de la réponse conserve l'information d'orientation qu'un déterminant isolé ne distingue plus.

\subsection{Inversion cubique et récupération angulaire}

\begin{theorem}[Inversion constitutive dans le domaine contractif]
Pour chaque $r\in(3/4,1)$, il existe une unique $s\in(0,1)$ telle que $f(s)=r$.
Cette $s$ est l’unique racine dans cet intervalle de
\begin{equation}
 r s^3+(r-1)(1+s+s^2)=0.
 \label{iii:eq:cubic}
\end{equation}
La réponse étiquetée reconstruit $s_\pm$, $q_\pm=s_\pm^{1/30}$ et
\begin{equation}
 x=-\frac12\log(q_+q_-),\qquad
 y=\frac12\log\frac{q_-}{q_+}.
 \label{iii:eq:recover-angular}
\end{equation}
\end{theorem}
\begin{proof}
La dérivée et les limites de \eqref{iii:eq:monotone} prouvent que $f$ est une
bijection décroissante entre les intervalles indiqués. Multiplier $f(s)=r$
par son dénominateur produit \eqref{iii:eq:cubic}. La racine positive d’ordre
$30$ retrouve le canal. Enfin, la somme et la différence de
$\log q_\pm=-x\mp y$ donnent \eqref{iii:eq:recover-angular}.
\end{proof}

La reconstruction conserve les étiquettes de feuillet. Avec l’opérateur complet
$\mathcal V$, on conserve aussi les projecteurs et on applique la fonction
inverse par calcul spectral. Récupérer deux coordonnées angulaires est une
opération définie sur cette représentation ; la trace historique complète
du TPK conserve des données supplémentaires.

L’intervalle $r\in(3/4,1)$ et la fonction inverse de ce théorème appartiennent
à la normalisation interne de \eqref{iii:eq:four}. Si l’on dispose de
coordonnées exprimées dans une autre coordinatisation d’unités, on retrouve d’abord
cette normalisation au moyen de la transformation de la
section~\ref{iii:sec:evaluacion} ; on applique ensuite l’inversion cubique.

% Recuperación expositiva de INTERRELACIONES_ESTRUCTURALES_VACIO.md, §§6--8.
% Insertar después de la inversión cúbica y antes de las coordenadas logarítmicas.
\subsection{Composition transportée des réponses constitutives}
\label{amp-comp-sec}

Le transport angulaire et les incidences $90$ et $120$ proviennent de la
même structure APP--TRIT--TPK : les feuillets conservent les évaluations et
leurs retenues ; le régime tritique conserve l’orientation ; le calendrier
TPK transporte la phase et la mémoire jusqu’aux canaux angulaires. La réponse
constitutive applique à ces canaux la fonction $f$ de
\eqref{iii:eq:monotone}. Son inversion permet d’exprimer, dans les coordonnées
de réponse, la composition des transports du même repère. On conserve
ainsi l’ordre constructif : transport, composition et évaluation spectrale.
L’opération obtenue agit sur cette représentation de l’état enrichi ;
son inverse retrouve les canaux, tandis que l’histoire complète de mémoire
demeure dans la structure de provenance.

Soient $s=q^{30}$ et $\psi=f^{-1}$. L’exposant $30$ est le plus grand commun
diviseur des incidences $90$ et $120$. L’inversion cubique détermine
$\psi:(3/4,1)\longrightarrow(0,1)$. En adjoignant la réponse de frontière
$f(1)=3/4$, nous étendons la bijection à
\[
 \psi:D\longrightarrow(0,1],\qquad D=[3/4,1).
\]
Sur cette frontière, on utilise la fonction rationnelle $f$ ; le quotient
$(1-q^{90})/(1-q^{120})$ y a une singularité éliminable.

\begin{proposition}[Monoïde des réponses et coordonnée additive]
\label{amp-comp-monoid}
L’opération
\begin{equation}
 r\odot t=f\bigl(\psi(r)\psi(t)\bigr),\qquad r,t\in D,
 \label{amp-comp-law}
\end{equation}
fait de $D$ un monoïde commutatif et simplifiable d’élément neutre $3/4$.
L’application
\begin{equation}
 \ell(r)=-\frac1{30}\log\psi(r)
 \label{amp-comp-character}
\end{equation}
est un isomorphisme de ce monoïde avec $(\mathbb R_{\geq0},+)$.
L’intervalle $(3/4,1)$ est stable par $\odot$ ; le neutre appartient à
l’extension de frontière. Le seul élément inversible dans $D$
est le neutre.
\end{proposition}
\begin{proof}
Le produit de deux éléments de $(0,1]$ appartient au même intervalle et
\[
 \psi(r\odot t)=\psi(r)\psi(t).
\]
Pour trois réponses, les deux ordres d’association ont pour image
$\psi(r)\psi(t)\psi(u)$ ; l’injectivité de $\psi$ prouve
l’associativité. La commutativité provient du produit scalaire. Si
$r\odot t=r\odot u$, la positivité de $\psi(r)$ permet de simplifier et
d’obtenir $\psi(t)=\psi(u)$, donc $t=u$. Le paramètre $1$ correspond à
$3/4$ et fournit l’identité. Le logarithme transforme le produit en
somme, et son image dans $(0,1]$ prouve que $\ell$ est une bijection sur
$\mathbb R_{\geq0}$. En particulier, pour la réponse d’un canal $q$,
$\ell(r)=-\log q$. Si les deux paramètres sont strictement inférieurs à
$1$, leur produit l’est aussi. Enfin, deux nombres de $(0,1]$
ont pour produit $1$ uniquement s’ils sont tous deux $1$.
\end{proof}

\begin{proposition}[Simplification admissible et prolongement inverse]
\label{amp-comp-inverse}
Pour $r,u\in D$, l’équation $r\odot t=u$ a une solution $t\in D$
exactement lorsque $u\geq r$. Cette solution est unique et donnée par
\begin{equation}
 t=f\!\left(\frac{\psi(u)}{\psi(r)}\right).
 \label{amp-comp-cancellation}
\end{equation}
La même fonction $f$ est une bijection de $(0,\infty)$ sur $(0,1)$.
Avec cette extension de $\psi$, la loi~\eqref{amp-comp-law} transporte le
groupe multiplicatif des paramètres positifs sur l’intervalle $(0,1)$.
L’inversion du groupe satisfait
\begin{equation}
 r^{\ominus}=\psi(r)r,\qquad
 r\odot r^{\ominus}=\frac34.
 \label{amp-comp-group-inverse}
\end{equation}
\end{proposition}
\begin{proof}
L’équation équivaut à
$\psi(t)=\psi(u)/\psi(r)$. Ce quotient appartient à $(0,1]$
exactement lorsque $\psi(u)\leq\psi(r)$, ce qui, par monotonie décroissante,
équivaut à $u\geq r$. La bijection prouve l’existence et l’unicité.
La dérivée de \eqref{iii:eq:monotone} est négative pour tout $s>0$ ;
les limites de $f$ en $0$ et à l’infini sont $1$ et $0$, respectivement.
Cela établit la bijection étendue. En multipliant le numérateur et le
dénominateur par $s^3$, on obtient
\[
 f(s^{-1})=\frac{s^3+s^2+s}{s^3+s^2+s+1}=s f(s).
\]
L’inverse multiplicatif de $s=\psi(r)$ a donc pour réponse
$\psi(r)r$. Sa composition avec $r$ correspond au paramètre $1$.
\end{proof}

Le prolongement à $s>1$, ou de manière équivalente à $q>1$, est l’extension
algébrique du lecteur. Le secteur physique contractant du chapitre conserve
son domaine déclaré. La coordonnée $\ell$ se prolonge en un isomorphisme
avec $(\mathbb R,+)$ et change de signe sous l’inversion du groupe.

\subsubsection{Réalisation opératorielle dans un repère de feuillets commun}

\begin{proposition}[Composition avec projecteurs marqués]
\label{amp-comp-operator}
Soient $\mathcal R=\mathcal R^*$, $\mathcal R^2=I$ et
$P_\pm=(I\pm\mathcal R)/2$. Pour $j=a,b$, soient
\[
 T_j=q_{j,+}P_++q_{j,-}P_-,\qquad
 \mathcal V_j=f(T_j^{30}),\qquad 0<q_{j,\pm}<1.
\]
La composition $T_{a\circ b}=T_aT_b$ satisfait
\begin{equation}
 \mathcal V_{a\circ b}
 =f\bigl(\psi(\mathcal V_a)\psi(\mathcal V_b)\bigr).
 \label{amp-comp-operator-law}
\end{equation}
Ses deux réponses sont $r_{a,\pm}\odot r_{b,\pm}$. Si
$T_j=\exp[-(x_jI+y_j\mathcal R)]$ avec $x_j>|y_j|$, alors
\begin{equation}
 T_aT_b=
 \exp[-((x_a+x_b)I+(y_a+y_b)\mathcal R)],
 \label{amp-comp-angular-addition}
\end{equation}
et $x_a+x_b>|y_a+y_b|$.
\end{proposition}
\begin{proof}
L’orthogonalité $P_+P_-=0$ donne
\[
 T_aT_b=q_{a,+}q_{b,+}P_++q_{a,-}q_{b,-}P_-.
\]
En particulier, $T_a$ et $T_b$ commutent. Le calcul fonctionnel dans ces
projecteurs fournit
$\psi(\mathcal V_j)=T_j^{30}$ et
$(T_aT_b)^{30}=T_a^{30}T_b^{30}$ ; appliquer $f$ prouve
\eqref{amp-comp-operator-law}. Pour toute fonction $g$ définie sur les
deux valeurs propres, on a, plus explicitement,
\[
 P_\pm g(T_j)=g(q_{j,\pm})P_\pm.
\]
Comme $q_{j,\pm}=\exp[-(x_j\pm y_j)]$, leurs produits prouvent
\eqref{amp-comp-angular-addition}. L’inégalité triangulaire conclut
$x_a+x_b>|y_a|+|y_b|\geq|y_a+y_b|$.
\end{proof}

Les coordonnées retrouvées peuvent s’écrire directement comme
\[
 x=\frac{\ell(r_+)+\ell(r_-)}2,\qquad
 y=\frac{\ell(r_+)-\ell(r_-)}2.
\]
La composition additionne ces paires. L’échange simultané de feuillets
permute les deux canaux de chaque transport, agit comme
$(x,y)\mapsto(x,-y)$ et est un automorphisme de la loi composante par
composante. L’inversion des deux transports dans l’extension algébrique
agit comme $(x,y)\mapsto(-x,-y)$. Les deux involutions commutent et
conservent des fonctions distinctes : l’une échange les marques d’orientation ;
l’autre inverse le transport.

Pour représenter l’échange par conjugaison, on utilise un
opérateur $L$ satisfaisant $L\mathcal R L^{-1}=-\mathcal R$. Dans la
représentation
\[
 \mathcal R=\begin{pmatrix}0&1\\1&0\end{pmatrix},\qquad
 L=\begin{pmatrix}1&0\\0&-1\end{pmatrix},
\]
la multiplication directe donne $LP_\pm L^{-1}=P_\mp$ et, par conséquent,
$LT_jL^{-1}=q_{j,-}P_++q_{j,+}P_-$. La conjugaison par
$\mathcal R$ conserve ses propres projecteurs. Une trace qui réunit les
canaux ne remplace pas non plus les projecteurs dans cette identité fonctionnelle :
l’information de leurs étiquettes intervient dans la composition.

L’hypothèse de repère commun est substantielle. Par exemple, les opérateurs
positifs définis
\[
 A_0=\begin{pmatrix}1/2&0\\0&1/3\end{pmatrix},\qquad
 B_0=\begin{pmatrix}1/2&1/6\\1/6&1/2\end{pmatrix}
\]
ont des valeurs propres positives, mais
\[
 A_0B_0=\begin{pmatrix}1/4&1/12\\1/18&1/6\end{pmatrix}
\]
n’est pas autoadjoint. La proposition spécifie la composition des
transports dans les projecteurs communs ; son domaine ne comprend pas un mélange
arbitraire de milieux matériels. La réalisation d’une concaténation de
chemins HMT conserve en outre ses données de mémoire : la fermeture algébrique du
cône de paramètres décrit la représentation et ne sélectionne pas à elle seule
de nouveaux chemins primaires.

\subsubsection{Produit de canaux et composition de transports}

La vitesse normalisée d’un état reste déterminée par
\[
 \widehat c=(r_+r_-)^{-1}.
\]
Ici, le produit ordinaire réunit les deux réponses de ce même état.
La loi $\odot$ évalue une autre opération : la composition de deux transports
dans chaque feuillet. Le calcul exact suivant distingue leurs arguments :
\begin{equation}
 f(1/2)=\frac{14}{15},\qquad
 \frac{14}{15}\odot\frac{14}{15}=f(1/4)=\frac{84}{85}
 \ne\frac{196}{225}=\left(\frac{14}{15}\right)^2.
 \label{amp-comp-rational-example}
\end{equation}
Les fractions de cet exemple sont des paramètres algébriques de vérification.
L’expression de $\widehat c$ et les identités constitutives précédentes
restent inchangées. Le contenu ajouté est une loi explicite pour
transporter la composition et retrouver sa coordonnée angulaire additive.


\subsection{Coordonnées logarithmiques et raffinement}

Soient $\mathcal E=\log(r_+r_-)$ et $\mathcal B=\log(r_-/r_+)$. Alors
\begin{equation}
 \log\mathcal V=\tfrac12\mathcal E I-\tfrac12\mathcal B\mathcal R.
 \label{iii:eq:logvac}
\end{equation}
Cette décomposition exhibe la composante invariante et la composante orientée.
Dans l’amplification $\mathcal V_m=\mathcal V\otimes I_{9^m}$, les
espérances partielles normalisées vérifient
$E_m(\mathcal V_{m+1})=\mathcal V_m$. Pour tout polynôme $p$,
$p(\mathcal V_m)=p(\mathcal V)\otimes I_{9^m}$ ; d’où
$E_m(p(\mathcal V_{m+1}))=p(\mathcal V_m)$. La continuité uniforme
étend l’identité aux fonctions continues du spectre. Les projecteurs,
les puissances et le logarithme sont ainsi conservés dans la tour indiquée.

\subsection{Types dimensionnels des quatre évaluations}

Avec les droites d’action, de charge et de vitesse, munies de sections
$u_S,u_Q,u_C$, la réalisation dimensionnelle s’écrit
\begin{align}
 Z_{\rm vac}&=\widehat Z u_Su_Q^{-2},&
 c_{\rm vac}&=\widehat c u_C,\\
 \mu_{\rm vac}&=\widehat\mu u_Su_C^{-1}u_Q^{-2},&
 \varepsilon_{\rm vac}&=\widehat\varepsilon u_S^{-1}u_C^{-1}u_Q^2.
 \label{iii:eq:dimensions}
\end{align}
En multipliant et en divisant ces expressions, on préserve
$\mu_{\rm vac}\varepsilon_{\rm vac}=c_{\rm vac}^{-2}$ et
$\mu_{\rm vac}/\varepsilon_{\rm vac}=Z_{\rm vac}^2$.
La coordinatisation SI évalue ensuite les sections. La comparaison physique nécessite
d’exposer leur construction et leur normalisation avec l’évaluation ; les quatre
coordonnées normalisées de \eqref{iii:eq:four} conservent leur type adimensionnel.

% INSERCIÓN FOCAL AUTORIZADA: no introduce un capítulo ni el concepto de radión.
% Incluir dentro de la sección constitutiva, después de la inversión cúbica
% y de la declaración de las cartas dimensionales. Etiquetas fuera de grupos
% de prefijado automático. Conserva íntegramente el bloque geométrico anterior.
\subsection{Reconstruction de la forme d’action à partir de la réponse constitutive}
\label{iii:subsec:puente-forma-LC}

La réponse constitutive et la réalisation inductive--capacitive reçoivent
le même couple angulaire construit à partir d’APP--TRIT--TPK. La première évalue
des fonctions rationnelles des deux canaux ; la seconde utilise le rapport de
leurs exposants pour déterminer la forme de l’ellipse d’action.
L’inversion cubique permet de composer les deux lectures sans identifier leurs
quotients numériques.

Soient $x=A_{\rm rad}$ et $y=C^*_{\rm rad}$, avec $x>|y|$, les coordonnées
relevées déjà construites, et conservons leurs étiquettes de feuillet. Nous écrivons
\begin{equation}
 q_\pm=e^{-(x\pm y)},\qquad s_\pm=q_\pm^{30},\qquad
 r_\pm=f(s_\pm),\qquad
 f(s)=\frac{1+s+s^2}{1+s+s^2+s^3}.
 \label{iii:eq:forma-canales}
\end{equation}
Le domaine contractant donne $s_\pm\in(0,1)$ et $r_\pm\in(3/4,1)$. Dans la chambre $x>y>0$, on a $s_+<s_-$ et $r_+>r_-$ ; l'échange des feuillets transporte cette chambre vers l'orientation opposée.

L’image des deux coordonnées constitutives normalisées est le domaine
\begin{equation}
 \mathcal U=\left\{(z,v)\in(0,\infty)^2:
   (zv)^{-1/2}\in(3/4,1),\quad
   (z/v)^{1/2}\in(3/4,1)\right\}.
 \label{iii:eq:forma-dominio}
\end{equation}
Ici, $(z,v)=(\widehat Z,\widehat c)$. Les restrictions expriment l'appartenance des canaux reconstruits au domaine du théorème d'inversion ; elles ne constituent pas un choix ultérieur de leurs valeurs.

\Needspace{23\baselineskip}
\begin{proposition}[Composition constitutive de la forme et de l’impédance LC]
\label{iii:prop:forma-LC}
Pour $(\widehat Z,\widehat c)\in\mathcal U$, soient $s_\pm\in(0,1)$ les
racines uniques de
\begin{equation}
 r_\pm s_\pm^3+(r_\pm-1)(1+s_\pm+s_\pm^2)=0,
 \qquad
 r_+=\frac1{\sqrt{\widehat Z\widehat c}},\quad
 r_-=\sqrt{\frac{\widehat Z}{\widehat c}}.
 \label{iii:eq:forma-inversion}
\end{equation}
Le rapport angulaire, l’anisotropie de l’ellipse d’action et l’impédance
LC relative sont reconstruits exactement par
\begin{align}
 \frac{y}{x}
 &=\frac{\log s_+-\log s_-}{\log s_++\log s_-},
 \label{iii:eq:forma-razon}\\
 \eta_{\rm el}
 &=\frac12\log\frac{\log s_+}{\log s_-},
 \label{iii:eq:forma-eta}\\
 \frac{Z_{\rm LC}}{Z_*}
 &=e^{-\eta_{\rm el}}
  =\sqrt{\frac{\log s_-}{\log s_+}}.
 \label{iii:eq:forma-LC}
\end{align}
Dans la dernière identité, $Z_{\rm LC}=Z_{\eta_{\rm el}}$ est l’impédance
de la coordinatisation LC de \eqref{hmtII:eq:ii-elipse-lc}.
\end{proposition}
\begin{proof}
L'inversion cubique de \eqref{iii:eq:cubic} détermine $s_\pm$ de manière univoque, car $f$ est une bijection décroissante de $(0,1)$ sur $(3/4,1)$. Ses racines positives d'ordre $30$ récupèrent $q_\pm$. Les identités
\[
 \log s_+=-30(x+y),\qquad \log s_-=-30(x-y)
\]
prouvent \eqref{iii:eq:forma-razon} par différence et somme. Les deux logarithmes
sont négatifs ; leur quotient est positif et leur somme n’est pas nulle. Par conséquent,
\[
 \frac12\log\frac{\log s_+}{\log s_-}
   =\frac12\log\frac{x+y}{x-y}
   =\operatorname{artanh}(y/x)=\eta_{\rm el}.
\]
La relation $Z_{\eta_{\rm el}}/Z_*=e^{-\eta_{\rm el}}$, démontrée dans
\eqref{hmtII:eq:ii-elipse-impedancia}, donne \eqref{iii:eq:forma-LC}.
\end{proof}

La composition conserve l’information nécessaire à chaque reconstruction.
Avec $\hbar$, on retrouve
$a_S=\sqrt{2\hbar}\,e^{\eta_{\rm el}/2}$ et
$b_S=\sqrt{2\hbar}\,e^{-\eta_{\rm el}/2}$ ; avec la coordinatisation
$(\omega,Z_*)$, on retrouve l’inductance et la capacité déjà définies.
L’action, l’horloge et la référence dimensionnelle accompagnent la forme comme
données typées. L’inversion de deux coordonnées adimensionnelles ne
reconstruit ni leurs unités ni toute la mémoire de l’état.

La normalisation intervient avant de résoudre les équations cubiques.
Si l’on reçoit $Z_{\rm vac}$ et $c_{\rm vac}$ dans une coordinatisation dimensionnelle,
les relations \eqref{iii:eq:dimensions} déterminent d’abord
\[
 \widehat Z=\frac{Z_{\rm vac}}{u_Su_Q^{-2}},\qquad
 \widehat c=\frac{c_{\rm vac}}{u_C}.
\]
Ce sont les coordonnées auxquelles s'applique \eqref{iii:eq:forma-inversion}. La référence $Z_*$ appartient à la réalisation LC et est expressément conservée dans \eqref{iii:eq:forma-LC}.

Les deux quotients sont reliés au moyen des canaux, mais évaluent
des fonctions différentes :
\begin{equation}
 \widehat Z=\frac{f(s_-)}{f(s_+)},\qquad
 \frac{Z_{\rm LC}}{Z_*}
       =\sqrt{\frac{\log s_-}{\log s_+}}.
 \label{iii:eq:forma-dos-cocientes}
\end{equation}
La reconstruction précédente fournit l'application entre les deux lectures. En particulier, échanger les feuillets inverse les deux quotients, remplace $\eta_{\rm el}$ par $-\eta_{\rm el}$ et conserve $\widehat c$. Dans la limite symétrique $y=0$, les canaux coïncident et les deux quotients valent $1$. En dehors de ce cas, l'égalité d'origine n'impose pas l'égalité des deux fonctions. La forme d'action est récupérable à partir de la réponse constitutive en conservant l'orientation et la normalisation.

% RESULTADO_RECUPERADO: funtor dimensional integral, eq:cZ-internos,
% eq:longitud-ruta; composición reunida en Artículo V REV02, sección 69.
% Incorporación autónoma a III: mismo transporte, mismas hojas y misma carta.
\subsection{Vitesse constitutive et réalisation cinématique de l’horloge}
\label{iii:sec:enlace-velocidad-constitutiva}

La réponse constitutive et le lecteur cinématique réalisent une même
section de vitesse. Leur identification conserve l’ordre de construction :
les feuillets d’APP, leur orientation TRIT et le transport angulaire exprimé
par le TPK précèdent les lectures de vitesse et le choix des unités.
Dans la représentation à deux feuillets de la
section~\ref{iii:sec:hojas}, avec $x>|y|$ et les projecteurs de rang un
$P_\pm$, le transport est
\[
 T=e^{-xI-y\mathcal R}=q_+P_++q_-P_-,
 \qquad q_\pm=e^{-x\mp y},\qquad 0<q_\pm<1.
\]
Ses blocs temporels de longueurs $90$ et $120$ produisent la réponse
$\mathcal V$ de \eqref{iii:eq:vacuum-op}. Les coordonnées $x,y$ et les
canaux sont des sorties du développement angulaire préalable ; cette composition
applique ses lecteurs à ce transport déjà construit.

Le lecteur logarithmique symétrique du transport temporel est
\begin{equation}
 \mathcal E=
 \log\!\bigl[(1-q_+^{90})(1-q_-^{90})\bigr]
 -\log\!\bigl[(1-q_+^{120})(1-q_-^{120})\bigr].
 \label{iii:eq:velocidad-lector-simetrico}
\end{equation}
Sur la droite dimensionnelle de vitesse, avec une section positive $u_C$, le
lecteur cinématique exprime $c_{\mathrm{int}}=\exp(-\mathcal E)u_C$.
La réponse constitutive exprime
$c_{\mathrm{vac}}=\widehat c u_C$ au moyen de \eqref{iii:eq:four}.

\begin{proposition}[Identité constitutive et cinématique]
\label{iii:prop:identidad-velocidad}
Dans la même coordinatisation dimensionnelle, on vérifie
\begin{equation}
 c_{\mathrm{int}}=c_{\mathrm{vac}}=\widehat c\,u_C,
 \qquad \widehat c=(r_+r_-)^{-1}=\exp(-\mathcal E).
 \label{iii:eq:identidad-velocidad}
\end{equation}
L'identification conserve l'échange des feuillets.
\end{proposition}
\begin{proof}
Tous les facteurs de \eqref{iii:eq:velocidad-lector-simetrico} sont
positifs. La loi du logarithme et \eqref{iii:eq:rchannels} donnent
\[
 \mathcal E=\log(r_+r_-)=\log\det\mathcal V,
 \qquad e^{-\mathcal E}=(\det\mathcal V)^{-1}.
\]
Le déterminant correspond à la réalisation à deux feuillets de rang un.
Les deux définitions expriment donc la même section. L’échange
de feuillets permute $r_+$ et $r_-$ et conserve leur produit. Le déterminant
$\det T=e^{-2x}$ correspond au transport élémentaire, tandis que
$\det\mathcal V=r_+r_-$ correspond à sa réponse temporelle ;
\eqref{iii:eq:identidad-velocidad} utilise le second.
\end{proof}

\subsubsection{Longueur de trajectoire et durée élémentaire}

Le repère dimensionnel positif $(u_S,u_C,u_T,u_Q)$ induit la base de longueur $u_L=u_Cu_T$. Pour une trajectoire enrichie $\gamma$ de $n=|\gamma|$ transitions et à canal constant, les lecteurs additifs sont
\begin{equation}
 T_{\mathrm{int}}(\gamma)=nu_T,
 \qquad L_{\mathrm{int}}(\gamma)=n\widehat c\,u_L.
 \label{iii:eq:velocidad-longitud-reloj}
\end{equation}
La concaténation additionne le nombre de transitions et leurs réalisations.
Les routes enrichies conservent en outre leur orientation et leur mémoire.

\begin{corollary}[Normalisation de la longueur élémentaire]
Pour $n>0$, $L_{\mathrm{int}}(\gamma)/T_{\mathrm{int}}(\gamma)
=c_{\mathrm{vac}}$. Dans la normalisation $t_0=u_T$,
\begin{equation}
 \ell_0=c_{\mathrm{int}}t_0=\widehat c\,u_L.
 \label{iii:eq:longitud-elemental-constitutiva}
\end{equation}
Dans une coordinatisation temporelle $t_0=\tau_0u_T$, avec $\tau_0>0$, on obtient
$\ell_0=\tau_0\widehat c\,u_L$.
\end{corollary}
\begin{proof}
Diviser les deux expressions de \eqref{iii:eq:velocidad-longitud-reloj} simplifie $n$ et utilise $u_L/u_T=u_C$. Multiplier la section de vitesse par $t_0$ prouve les deux formules élémentaires.
\end{proof}

Lorsque le canal dépend de l’arête, le lecteur conserve la forme additive
$L_{\mathrm{int}}(\gamma)=\sum_{e\in\gamma}\widehat c(e)u_L$.
Son quotient par $nu_T$ est la moyenne des vitesses de ces arêtes.
L’expression à canal constant est la spécialisation homogène de
cette lecture.

\subsubsection{Changement de coordinatisation et normalisation adaptée}

\begin{proposition}[Covariance de l’identification]
Soient $u_C'=d u_C$ et $u_T'=\eta u_T$, avec $d,\eta>0$. Les coordonnées des mêmes sections satisfont
\begin{equation}
 \widehat c'=d^{-1}\widehat c,
 \qquad \tau_0'=\eta^{-1}\tau_0,
 \qquad u_L'=d\eta u_L.
 \label{iii:eq:velocidad-cambio-carta}
\end{equation}
Si une autre représentation écrit $c'=\widehat c' u_C'$ avec $u_C'=d u_C$,
l’égalité des sections équivaut exactement à $d\widehat c'=\widehat c$.
\end{proposition}
\begin{proof}
L'égalité de chaque section dans les deux bases donne $\widehat c'u_C'=\widehat c u_C$ et $\tau_0'u_T'=\tau_0u_T$. En particulier,
\[
 \tau_0'\widehat c'u_L'
 =\frac{\tau_0}{\eta}\frac{\widehat c}{d}\,d\eta u_L
 =\tau_0\widehat c u_L=\ell_0.
\]
La dernière équivalence provient de la même substitution pour la vitesse.
\end{proof}

Le choix adapté $d=\widehat c$ donne $u_C'=c_{\mathrm{vac}}$ et
$\widehat c'=1$ : il exprime la section construite dans une base adaptée.
Le coefficient de la coordinatisation interne reste $(r_+r_-)^{-1}$, et c’est
la normalisation utilisée par l’inversion cubique. En maintenant les
bases d’action et de charge, \eqref{iii:eq:dimensions} produit
\[
 \widehat\mu'=d\widehat\mu,\qquad
 \widehat\varepsilon'=d\widehat\varepsilon,\qquad
 \widehat\mu'\widehat\varepsilon'=(\widehat c')^{-2}.
\]
Dans la coordinatisation adaptée, on a $\widehat\mu'=r_-/r_+$ et
$\widehat\varepsilon'=r_+/r_-$, par substitution de
$\widehat\mu=r_-^2$ et $\widehat\varepsilon=r_+^2$.

\subsubsection{Conservation du lecteur sous amplification}

Pour le raffinement $\mathcal V_m=\mathcal V\otimes I_{9^m}$, la trace normalisée $\tau_m=\operatorname{Tr}/(2\cdot9^m)$ satisfait
\begin{equation}
 \exp[-2\tau_m(\log\mathcal V_m)]
 =\exp[-\log r_+-\log r_-]=\widehat c.
 \label{iii:eq:velocidad-traza-normalizada}
\end{equation}
En effet, chaque valeur propre apparaît $9^m$ fois ; cette multiplicité
s’annule avec le dénominateur de la trace normalisée. Le déterminant
ordinaire amplifié est $(r_+r_-)^{9^m}$. La normalisation à deux feuillets,
exprimée de manière équivalente par \eqref{iii:eq:velocidad-traza-normalizada},
conserve le lecteur constitutif pendant l’amplification. Cette
réalisation conserve le transport enrichi dont elle est issue.

% Recuperación autocontenida: sector angular de árbol y reconstrucción Higgs--vacío.
% Propietarios y alcance: higgs_procedencia.md. Incorporar después del operador constitutivo.
\section{Reconstruction conjointe de la jauge, de Higgs et de la réponse constitutive}
\label{amp-higgs-seccion}

La réponse du vide et le secteur électrofaible admettent des lectures complémentaires
du transport angulaire déjà construit. La première conserve les deux canaux
marqués ; le déterminant conserve leur contraction commune ; la composante angulaire
de Higgs conserve en outre l’orientation qui disparaît dans ce déterminant.
La composition de ces lectures permet de reconstruire leur antécédent spectral
et d’exprimer l’impédance au moyen de l’autocouplage scalaire et des couplages
de jauge, dans une même réalisation à l’arbre.

\subsection{Provenance angulaire et signatures des chemins}

On conserve la chaîne APP--TRIT--TPK : le prolongement survivant détermine
le chemin, celui-ci produit son registre orienté et le lecteur du registre exprime
la signature angulaire. Les coordonnées $\pi=\pi_{\rm HMT}$ et
$\alpha=\alpha_{\rm HMT}$ sont des sorties préalables de la même génération ;
$A$ et $C^*$ sont leurs représentations angulaires ultérieures. Dans ce bloc, on
utilise des radians,
\begin{equation}
 x=\frac{\pi A_{\deg}}{180},\qquad
 y=\frac{\pi C^*_{\deg}}{180},\qquad
 A_{\deg}=1000\alpha,\qquad 0<y<x.
 \label{amp-higgs-angulos}
\end{equation}
Par conséquent, $x=50\pi\alpha/9$. Les canaux et les marques de feuillet sont ceux
de la section~\ref{iii:sec:hojas} :
\begin{equation}
 \begin{aligned}
 q_+&=e^{-x-y},& q_-&=e^{-x+y},& D&=q_+q_-=e^{-2x},\\
 P_\pm&=\frac{I\pm\mathcal R}{2},&
 T&=q_+P_++q_-P_-.
 \end{aligned}
 \label{amp-higgs-transporte}
\end{equation}
Ici, $\mathcal R$ est l’involution autoadjointe marquée et chaque projecteur
est de rang un dans la réalisation bicouche de base.

Le lecteur angulaire de durée--perdurance est
\begin{equation}
 \chi_{\rm dur}(n,s)=\exp\!\left(nx+\frac{s}{6}y\right).
 \label{amp-higgs-caracter}
\end{equation}
Pour conserver la provenance finie des coefficients qui seront utilisés,
nous réunissons les empreintes nominales des trois chemins du calendrier de longueur
$108$. Les colonnes de déplacements cardinaux se lisent comme $(N,E,S,O)$ ;
la paire suivante conserve le feuillet et le résidu exprimés par ce chemin.
Le tableau conserve, pour chaque chemin identifié, l’empreinte et la signature
angulaire exprimées par la chaîne précédente.
\begin{center}
\small
\begin{tabular}{c c c c c}
\hline
Secteur & Chemin et axe & $(N,E,S,O)$ & Feuillet/résidu & $(n,s)$\\
\hline
$W$ & $15/37$, E--O & $(23,44,11,30)$ & $(41,8)$ & $(94,-1)$\\
$Z$ & $20/23$, E--O & $(33,45,10,20)$ & $(57,7)$ & $(95,-1)$\\
$H$ & $24/29$, N--S & $(40,34,22,12)$ & $(47,10)$ & $(97,9)$\\
\hline
\end{tabular}
\end{center}
Les trois empreintes typées conservent l’identifiant de chemin, l’orientation
et le registre de provenance des signatures employées. L’évaluation angulaire
commence après l’extraction à partir du registre enrichi ; le tableau
réunit ses sorties pour composer les caractères dans cette réalisation.
Les marques $W_\pm$ qui suivent désignent des coordonnées d’enroulement,
distinctes du nom de l’espèce $W$ :
\begin{equation}
 \begin{aligned}
 W_\pm&=6n\pm s, &
 n&=\frac{W_++W_-}{12},& s&=\frac{W_+-W_-}{2},\\
 (n,s)_W&=(94,-1), & (W_+,W_-)_W&=(563,565),\\
 (n,s)_Z&=(95,-1), & (W_+,W_-)_Z&=(569,571),\\
 (n,s)_H&=(97,9), & (W_+,W_-)_H&=(591,573).
 \end{aligned}
 \label{amp-higgs-firmas}
\end{equation}
La matrice de la première application est
$\bigl(\begin{smallmatrix}6&1\\6&-1\end{smallmatrix}\bigr)$, de déterminant
$-12$ et d’image $\{(u,v)\in\mathbb Z^2:u+v\in12\mathbb Z\}$.
En effet, les formules inverses précédentes sont entières précisément sur
cette image. Le caractère conserve les deux orientations au moyen de
\begin{equation}
 \chi_{\rm dur}(n,s)
 =\left(q_+^{-W_+}q_-^{-W_-}\right)^{1/12}.
 \label{amp-higgs-caracter-canales}
\end{equation}
La preuve s’obtient en prenant les logarithmes : le numérateur est
$W_+(x+y)+W_-(x-y)=12nx+2sy$.
Ainsi sont identifiés la racine douzième et le diviseur $6$ du lecteur.

\begin{proposition}[Quotients angulaires déterminés par les signatures]
Dans une section d’échelle commune, avec les facteurs de raffinement partagés
qui permettent leur annulation dans les quotients, on a
\begin{equation}
 \frac{m_W^{\angle}}{m_Z^{\angle}}=e^{-x},\qquad
 \frac{m_H^{\angle}}{m_W^{\angle}}
 =e^{3x+5y/3},\qquad
 \left(\frac{m_H^{\angle}}{m_W^{\angle}}\right)^2
 =e^{6x+10y/3}.
 \label{amp-higgs-cocientes}
\end{equation}
\end{proposition}
\begin{proof}
Diviser les caractères de $W$ et $Z$ soustrait leurs signatures : $(94,-1)-(95,-1)=(-1,0)$.
Pour $H/W$, la différence est $(97,9)-(94,-1)=(3,10)$.
L’équation~\eqref{amp-higgs-caracter} donne $3x+(10/6)y$ ; le carré
multiplie cet exposant par deux. L’échelle commune et les raffinements
partagés s’annulent avant ces opérations.
\end{proof}
Ces quotients conservent l’identifiant de la coordinatisation angulaire. L’application
à des masses ayant des raffinements sectoriels distincts conserve leurs facteurs
supplémentaires ; la condition d’annulation fait partie de l’énoncé.

\subsection{Réalisation rationalisée à l’arbre}

La lecture électrofaible de cette coordinatisation définit
$c_W^2=D$, $s_W^2=1-D$ et $e_g^2=4\pi\alpha$. La racine positive fixe
$g=e_g/s_W$ et $g'=e_g/c_W$. On obtient
\begin{equation}
 g^2=\frac{4\pi\alpha}{1-D},\qquad
 g'^2=\frac{4\pi\alpha}{D}.
 \label{amp-higgs-calibre}
\end{equation}
La quantité $e_g$ est un couplage adimensionnel dans la normalisation
rationalisée, distinct d’une section dimensionnelle de charge.
La convention du potentiel scalaire est
\begin{equation}
 \begin{gathered}
 V(H)=-\mu_H^2H^\dagger H+\lambda_H(H^\dagger H)^2,\\
 m_H^2=2\lambda_Hv^2,\qquad m_W^2=\frac{g^2v^2}{4}.
 \end{gathered}
 \label{amp-higgs-potencial}
\end{equation}
La même $v$ intervient dans les deux relations. En les divisant, son échelle
s’annule et $m_H^2/m_W^2=8\lambda_H/g^2$. Avec
\eqref{amp-higgs-cocientes}, on obtient
\begin{equation}
 \boxed{\lambda_H=\frac{g^2}{8}\exp\!\left(6x+\frac{10}{3}y\right).}
 \label{amp-higgs-lambda}
\end{equation}
Le facteur $8$ provient des deux normalisations quadratiques de
\eqref{amp-higgs-potencial} ; les coefficients $6$ et $10/3$ proviennent des
signatures. La distinction conserve l’opération qui produit chaque facteur.
La composition reçoit les sorties HMT préalables et utilise cette réalisation physique
à l’arbre à échelle commune ; un changement de normalisation du champ, d’échelle ou
de schéma se transporte également dans ses relations de masses et de couplages.

\subsection{Inverse conjoint et domaine exact}

\begin{theorem}[Reconstruction angulaire au moyen de la jauge et de Higgs]
\label{amp-higgs-inversa-teorema}
Pour $\pi>0$ fixé, l’application de
\eqref{amp-higgs-transporte}, \eqref{amp-higgs-calibre} et
\eqref{amp-higgs-lambda}, sur $\alpha>0$, $x>y>0$, est injective et a
pour inverse
\begin{equation}
 \begin{aligned}
 D&=\frac{g^2}{g^2+g'^2},&
 x&=-\frac12\log D,&
 \alpha&=\frac{g^2g'^2}{4\pi(g^2+g'^2)},\\
 y&=\frac3{10}\left[\log\!\left(\frac{8\lambda_H}{g^2}\right)-6x\right].
 \end{aligned}
 \label{amp-higgs-inversa}
\end{equation}
Son image exacte est formée des triplets $g,g',\lambda_H>0$ tels que,
avec $D$ et $x$ reconstruits comme ci-dessus,
\begin{equation}
 \frac{g^2}{8}e^{6x}<\lambda_H<\frac{g^2}{8}e^{28x/3}.
 \label{amp-higgs-dominio}
\end{equation}
\end{theorem}
\begin{proof}
Les deux équations de jauge donnent
$g^2/(g^2+g'^2)=D$ et
$g^2g'^2/(g^2+g'^2)=4\pi\alpha$. Le logarithme réel retrouve $x$ et
l’équation scalaire retrouve $y$. Les conditions $g,g'>0$ impliquent
$0<D<1$ et $x>0$. L’inégalité~\eqref{amp-higgs-dominio} équivaut,
par le logarithme strictement croissant, à
$0<\log(8\lambda_H/g^2)-6x<10x/3$, c’est-à-dire $0<y<x$.
Réciproquement, les formules reconstruites produisent
$4\pi\alpha/(1-D)=g^2$, $4\pi\alpha/D=g'^2$ et
$g^2e^{6x+10y/3}/8=\lambda_H$. Les racines positives retrouvent les
couplages originaux. Cela prouve la nécessité, la suffisance et l’unicité.
\end{proof}

L’inversion précédente décrit exactement cette application entre lecteurs.
La section engendrée par HMT conserve en outre la relation de
\eqref{amp-higgs-angulos}, qui s’exprime comme
\begin{equation}
 -\frac12\log\!\left(\frac{g^2}{g^2+g'^2}\right)
 =\frac{25}{18}\frac{g^2g'^2}{g^2+g'^2}.
 \label{amp-higgs-compatibilidad-alpha}
\end{equation}
Elle s’obtient en substituant l’inverse de $\alpha$ dans $x=50\pi\alpha/9$.
Le contre-angle conserve également son équation de retour antérieure.
Par conséquent, l’image du lecteur et la section produite par le générateur
sont des objets distincts : les compatibilités préalables sélectionnent la section
au sein du domaine reconstructible. L’opération inverse agit après
la génération et permet de vérifier ses sorties conservées.

\subsection{Reconstruction Higgs--impédance}

Pour exprimer le contenu orienté sans altérer les canaux,
nous définissons
\begin{equation}
 \rho_{\rm or}=\frac{q_-}{q_+}=e^{2y},\qquad
 \Xi_H=\frac{8D^3\lambda_H}{g^2}.
 \label{amp-higgs-Xi}
\end{equation}
L’équation~\eqref{amp-higgs-lambda} et $D^3=e^{-6x}$ donnent
\begin{equation}
 \begin{aligned}
 \Xi_H&=e^{10y/3}=\rho_{\rm or}^{5/3},&
 \rho_{\rm or}&=\Xi_H^{3/5},\\
 q_-&=\sqrt D\,\Xi_H^{3/10},&
 q_+&=\sqrt D\,\Xi_H^{-3/10}.
 \end{aligned}
 \label{amp-higgs-canales}
\end{equation}
Les puissances sont les racines réelles positives. Le domaine exact
de la chambre orientée est
\begin{equation}
 0<D<1,\qquad 1<\Xi_H<D^{-5/3}.
 \label{amp-higgs-dominio-Xi}
\end{equation}
En effet, $0<y<x$ équivaut à $1<e^{10y/3}<e^{10x/3}$.
L’extrémité inférieure donne $q_+=q_-$ ; la supérieure donne $q_-=1$.
Pour les deux orientations ensemble, $|y|<x$ équivaut à
$D^{5/3}<\Xi_H<D^{-5/3}$.

Soit $F(q)=f(q^{30})=(1-q^{90})/(1-q^{120})$, la réponse déjà prouvée dans
\eqref{iii:eq:rchannels}--\eqref{iii:eq:monotone}.
\begin{corollary}[Confluence des lectures scalaire et constitutive]
\label{amp-higgs-impedancia-corolario}
Dans le domaine~\eqref{amp-higgs-dominio-Xi},
\begin{equation}
 \boxed{\widehat Z=
 \frac{F(\sqrt D\,\Xi_H^{3/10})}
 {F(\sqrt D\,\Xi_H^{-3/10})}.}
 \label{amp-higgs-impedancia}
\end{equation}
Les réponses complètes sont
\begin{equation}
 \begin{aligned}
 r_-&=F(\sqrt D\,\Xi_H^{3/10}),&
 r_+&=F(\sqrt D\,\Xi_H^{-3/10}),\\
 \widehat\mu&=r_-^2,& \widehat\varepsilon&=r_+^2,&
 \widehat c&=(r_-r_+)^{-1}.
 \end{aligned}
 \label{amp-higgs-vacio}
\end{equation}
De plus, les deux reconstructions satisfont
\begin{equation}
 \frac{F^{-1}(\sqrt{\widehat\mu})}
 {F^{-1}(\sqrt{\widehat\varepsilon})}
 =\rho_{\rm or}=\Xi_H^{3/5}.
 \label{amp-higgs-confluencia}
\end{equation}
\end{corollary}
\begin{proof}
Substituer les canaux de~\eqref{amp-higgs-canales} dans les quatre lectures
de~\eqref{iii:eq:four} prouve les premières identités. L’inverse
constitutif de~\eqref{iii:eq:cubic} retrouve $q_-$ et $q_+$ ; leur quotient
est $\rho_{\rm or}$. L’égalité avec $\Xi_H^{3/5}$ est déjà prouvée dans
\eqref{amp-higgs-canales}.
\end{proof}

Dans la section compatible avec $A_{\deg}=1000\alpha$, le couplage
$Q(D):=-9\log D/25$ satisfait $Q(D)=4\pi\alpha$ et
$g^2=Q(D)/(1-D)$. Par conséquent, la même coordonnée de Higgs s’écrit
\begin{equation}
 \Xi_H=\frac{8(1-D)D^3\lambda_H}{Q(D)},\qquad
 \frac1{g^2}+\frac1{g'^2}=\frac1{Q(D)}.
 \label{amp-higgs-curva}
\end{equation}
Ces expressions relient le couplage total, sa répartition rationalisée
et l’orientation scalaire dans la même section.

\subsection{Information conservée et transport d’orientation}

La paire de jauge retrouve $\alpha$ et $x$ ; le quotient
$\lambda_H/g^2$ retrouve ensuite $y$. Avec les marques $P_\pm$, on reconstruit
les opérateurs complets $T$ et $\mathcal V=r_+P_++r_-P_-$.
Les trois scalaires déterminent les valeurs propres ; les projecteurs marqués
déterminent en outre leur réalisation opératorielle. L’histoire enrichie du
chemin reste une information antérieure de résolution supérieure.

Sous $y\mapsto-y$ par rapport aux marques fixes, $D,g,g'$ restent
invariants et $\Xi_H\mapsto\Xi_H^{-1}$. Les réponses constitutives s’échangent,
$\widehat Z\mapsto\widehat Z^{-1}$ et $\widehat c$ est conservé.
Si $\lambda_H^+$ et $\lambda_H^-$ désignent les évaluations dans les deux
orientations conjuguées, on a
\begin{equation}
 \lambda_H^+\lambda_H^-
 =\frac{g^4}{64}D^{-6}.
 \label{amp-higgs-orientacion}
\end{equation}
La preuve consiste à multiplier~\eqref{amp-higgs-lambda} pour $y$ et $-y$ ;
les contributions orientées s’annulent. Cette comparaison conserve le
lecteur de signatures et change son argument angulaire. Une conjugaison simultanée
des chemins et des références doit également transporter les signatures.

Dans l’amplification $T_m=T\otimes I_{9^m}$, avec
$\tau_m=\operatorname{Tr}/(2\cdot9^m)$ et
$\mathcal R_m=\mathcal R\otimes I_{9^m}$, les coordonnées appropriées sont
\begin{equation}
 D=\exp\bigl(2\tau_m(\log T_m)\bigr),\qquad
 \rho_{\rm or}=\exp\bigl(-2\tau_m(\mathcal R_m\log T_m)\bigr).
 \label{amp-higgs-refinamiento}
\end{equation}
En effet, $\log T_m=(\log T)\otimes I_{9^m}$, de sorte que les deux traces
sont $-x$ et $-y$, respectivement. La normalisation élimine la multiplicité
de l’amplification. Les lectures de jauge, de Higgs et du vide restent ainsi
compatibles avec ces raffinements ; le déterminant ordinaire amplifié
serait $D^{9^m}$ et a une fonction distincte.

La conclusion de la composition est précise : l’invariant pair de jauge
et la coordonnée orientée scalaire reconstruisent conjointement la paire
spectrale qui gouverne le vide. Lorsque $\lambda_H$ est évalué au moyen de la
même formule, la confluence est une identité de cohérence. Une lecture
indépendante de $\lambda_H$, transportée vers des schéma et échelle identiques,
transforme cette égalité en comparaison supplémentaire. Dans les deux usages restent
explicites les signatures, les marques de feuillet et la normalisation à l’arbre.

\section{Décomposition cyclique, quart de tour et moment orienté de Catalan}
\label{iii:sec:catalan}

Le registre dodécaphasique permet de relier la réponse constitutive à
d’autres invariants du transport. Le lien s’établit sur des sous-espaces
et des applications déterminés : le défaut de rang mesure une différence
d’incidence, le déterminant conserve une réponse scalaire et un coefficient
de la résolvante retient l’orientation. La famille de moments qui culmine
en Catalan intervient à ce dernier niveau.

\subsection{Deux sous-espaces du cycle dodécaphasique}

Dans $\mathbb R^{12}$, soit $C e_j=e_{j+1\bmod12}$ et soient
\begin{equation}
 P_3^{\mathrm{cyc}}=\frac{I+C^3+C^6+C^9}{4},\qquad
 P_4^{\mathrm{cyc}}=\frac{I+C^4+C^8}{3},\qquad
 H_d=\operatorname{Ran}P_d^{\mathrm{cyc}}.
 \label{iii:eq:projectors34}
\end{equation}
Les moyennes sur les sous-groupes cycliques sont des projecteurs orthogonaux.
$H_3$ est formé des suites périodiques de période divisant $3$,
et $H_4$ de celles de période divisant $4$. Leurs dimensions sont $3$ et $4$.
La restriction de $C$ à chaque espace réalise le cycle correspondant.
L’exposant $\mathrm{cyc}$ identifie la moyenne du cycle dodécaphasique :
$P_3^{\mathrm{cyc}}$ est le projecteur sur les suites de période divisant
$3$. Son domaine et son action sont distincts de ceux du projecteur isotypique
tridimensionnel de la réalisation icosaédrique ; la coïncidence des rangs
n’identifie pas les deux opérateurs.

\begin{proposition}[Défaut de rang et réponse déterminantale]
On vérifie
\begin{equation}
 \frac{\operatorname{Tr}(P_3^{\mathrm{cyc}}-P_4^{\mathrm{cyc}})}{12}=-\frac1{12},\qquad
 \frac{\det_{H_3}(I-sC)}{\det_{H_4}(I-sC)}
 =\frac{1-s^3}{1-s^4}=f(s),\quad 0<s<1.
 \label{iii:eq:trace-det}
\end{equation}
\end{proposition}
\begin{proof}
La trace d’un projecteur orthogonal est son rang, de sorte que le numérateur
de la première expression est $3-4$. Le polynôme caractéristique d’un cycle
de longueur $d$ est $t^d-1$ ; son déterminant $\det(I-sC)$ est donc
$1-s^d$. Le quotient produit la seconde identité.
\end{proof}

La substitution $s=q^{30}$ transporte cette représentation vers le quotient
$90/120$. En $s=1$, le mode constant commun s’annule avant de prendre la
limite, qui vaut $3/4$. Cette réalisation de $-1/12$ comme défaut de trace
a un domaine fini explicite. Les réalisations régularisées et
modulaires du même scalaire conservent leurs opérateurs et domaines propres.

\subsection{Représentation fidèle du quart de tour}

Définissons le projecteur antipodal $P_{\rm ant}=(I-C^6)/2$ et
$Q=P_4^{\mathrm{cyc}}P_{\rm ant}$. Ce projecteur agit sur le cycle à douze phases et est
distinct du projecteur constitutif $P_-$ de la section~\ref{iii:sec:constitutiva}.
Les projecteurs $P_4^{\mathrm{cyc}}$ et $P_{\rm ant}$ commutent et
$Q$ est de rang $2$. Dans son image, on fixe la base orthonormée
\begin{equation}
 u=\frac1{\sqrt6}(1,0,-1,0,1,0,-1,0,1,0,-1,0)^{\mathsf T},
 \qquad v=Cu.
 \label{iii:eq:uv}
\end{equation}
On vérifie composante par composante que $Cu=v$ et $Cv=-u$. Par conséquent,
l’isométrie $W(a,b)=au+bv$ satisfait
\begin{equation}
 CW=WJ,\qquad
 J=\begin{pmatrix}0&-1\\1&0\end{pmatrix},\qquad J^2=-I.
 \label{iii:eq:quarter-intertwiner}
\end{equation}
La structure complexe du plan est réalisée par l’action effective du
cycle. Le signe dépend de l’orientation fixée : dans la même base,
$C^3W=-WJ$. Cette précision permet de composer le plan avec le lecteur orienté.

\subsection{Résolvante et profondeur illimitée}

Comme $C^2=-I$ sur $\operatorname{Ran}Q$, pour $0\leq s<1$, on a
\begin{equation}
 (I-sC)^{-1}\big|_{\operatorname{Ran}Q}=\frac{I+sC}{1+s^2},
 \qquad k_4(s):=\langle v,(I-sC)^{-1}u\rangle=\frac{s}{1+s^2}.
 \label{iii:eq:resolvent}
\end{equation}
La suite $\langle v,C^nu\rangle$ est $0,1,0,-1,\ldots$. Le lecteur
de profondeur conserve ce caractère au moyen du moment
\begin{equation}
 \mathcal C_2(s)=\sum_{k\geq0}\frac{(-1)^ks^{2k+1}}{(2k+1)^2},
 \qquad 0\leq s\leq1.
 \label{iii:eq:c2}
\end{equation}
La série converge absolument et uniformément sur cet intervalle par
comparaison avec $\sum(2k+1)^{-2}$. Pour $0<s<1$, les séries dérivées
convergent uniformément sur chaque intervalle compact intérieur. Avec
$\mathcal D=s\,d/ds$, on obtient
\begin{equation}
 \mathcal D^2\mathcal C_2(s)=k_4(s),\qquad
 \mathcal C_2(s)=\int_0^s\frac{dt}{t}\int_0^t\frac{d\xi}{\xi}\,k_4(\xi).
 \label{iii:eq:c2-integral}
\end{equation}
Les conditions $\mathcal C_2(0)=0$ et
$\lim_{s\downarrow0}\mathcal D\mathcal C_2(s)=0$ fixent les deux constantes
d’intégration. En effet, la différence de deux solutions de
$\mathcal D^2F=k_4$ a la forme $a\log s+b$ ; la seconde condition
annule $a$ et la première annule $b$.
La convergence uniforme permet d’évaluer l’extrémité $s=1$ : le résultat est
le moment de Catalan. La représentation finie du caractère conserve son
orientation ; l’indice de profondeur reste illimité.

\subsection{Composition avec la réponse constitutive}

\begin{theorem}[Relation différentielle constitutive]
Pour $0<s<1$, avec l’orientation de \eqref{iii:eq:uv},
\begin{equation}
 f(s)=\frac{1+s+s^2}{s(1+s)}\,\mathcal D^2\mathcal C_2(s).
 \label{iii:eq:catalan-vacuum}
\end{equation}
\end{theorem}
\begin{proof}
Substituer \eqref{iii:eq:resolvent} dans \eqref{iii:eq:c2-integral} et utiliser
$(1+s)(1+s^2)=1+s+s^2+s^3$ donne l’identité. Tous les dénominateurs sont
positifs sur le domaine. L’entrelaceur \eqref{iii:eq:quarter-intertwiner}
identifie le coefficient opératoriel utilisé avant son évaluation.
\end{proof}

En évaluant $s_\pm=q_\pm^{30}$, la même famille produit les deux réponses
de canal de \eqref{iii:eq:rchannels}. La limite de Catalan, le défaut de
trace et les valeurs propres constitutives maintiennent leurs fonctions différentes
au sein de cette composition. L’inversion du cycle avec les marques $u,v$
fixes change le signe de $k_4$ et conserve le déterminant $1+s^2$.
On identifie ainsi une information d’orientation que le déterminant
isolé ne distingue pas.

Le développement de Pell utilise une évaluation particulière de cette famille,
avec $\rho=2-\sqrt3$ et $\rho^{-1}=2+\sqrt3$. Sa preuve est conservée comme
unité propre d’intégration ; la définition angulaire de $q_\pm$ reste
indépendante d’une éventuelle identification de $s_\pm$ à $\rho$.

La figure~\ref{iii:fig:respuesta-momento-pell} représente la relation entre
la réponse rationnelle et le noyau orienté avant l’intégration. La section
de Pell satisfait $\rho+\rho^{-1}=4$, donc
$k_4(\rho)=(\rho+\rho^{-1})^{-1}=1/4$. Les deux valeurs constitutives
s’obtiennent, en revanche, en évaluant $f$ aux arguments $s_\pm$ déjà
construits. Le tracé inclut les prolongements continus des deux
fonctions à $[0,1]$ ; le domaine d’inversion constitutive reste
l’intervalle ouvert $(0,1)$.

\begin{figure}[htbp]
\centering
\begin{tikzpicture}[x=10cm,y=4.6cm,>=stealth]
 \pgfmathsetmacro{\pellrho}{2-sqrt(3)}
 \pgfmathsetmacro{\fpell}{(1+\pellrho+\pellrho^2)/(1+\pellrho+\pellrho^2+\pellrho^3)}
 \draw[gray!35,thin] (0,0.25)--(1,0.25);
 \draw[gray!35,thin] (0,0.5)--(1,0.5);
 \draw[gray!35,thin] (0,0.75)--(1,0.75);
 \draw[gray!35,thin] (0,1)--(1,1);
 \draw[->] (0,0)--(1.04,0) node[right] {$s$};
 \draw[->] (0,0)--(0,1.1) node[above] {valeur adimensionnelle};
 \foreach \xx/\lab in {0/0,0.5/{1/2},1/1}{
  \draw (\xx,0)--(\xx,-0.015) node[below] {$\lab$};
 }
 \foreach \yy/\lab in {0.25/{1/4},0.5/{1/2},0.75/{3/4},1/1}{
  \draw (0,\yy)--(-0.008,\yy) node[left] {$\lab$};
 }
 \draw[gray!70,densely dashed] (\pellrho,0)--(\pellrho,1.035);
 \node[below,align=center,font=\small] at (\pellrho,-0.015)
  {$\rho$\\[-1pt]$2-\sqrt3$};
 \draw[black,very thick,domain=0:1,samples=121,smooth]
  plot (\x,{(1+\x+\x^2)/(1+\x+\x^2+\x^3)});
 \draw[blue!65!black,thick,dashed,domain=0:1,samples=121,smooth]
  plot (\x,{\x/(1+\x^2)});
 \fill[black] (\pellrho,\fpell) circle[radius=1.7pt];
 \fill[blue!65!black] (\pellrho,0.25) circle[radius=1.7pt];
 \node[anchor=south west,font=\small,fill=white,inner sep=2pt] at (0.51,0.94)
  {$f(s)=\dfrac{1+s+s^2}{1+s+s^2+s^3}$};
 \node[anchor=south west,text=blue!65!black,font=\small] at (0.57,0.55)
  {$k_4(s)=\dfrac{s}{1+s^2}$};
 \node[anchor=north west,font=\small,fill=white,inner sep=2pt]
  at (0.30,0.225) {$k_4(\rho)=1/4$};
\end{tikzpicture}
\caption{Réponse constitutive et noyau du moment orienté. La courbe
en trait continu est $f$ et celle en tirets est
$k_4=\mathcal D^2\mathcal C_2$ ; toutes deux sont reliées au moyen de
\eqref{iii:eq:catalan-vacuum}. La droite verticale marque la section de Pell,
sans l’identifier aux arguments constitutifs $s_+$ ou $s_-$.}
\label{iii:fig:respuesta-momento-pell}
\end{figure}

% Delta de §10.27. Insertar al final de 05_ciclo_catalan.tex.
% La actualización afín y su balance permanecen en sus propietarios originales.
\subsection{Traces de retour et moments du quotient constitutif}
\label{amp-afin-sec}

La mise à jour du registre et la réponse du cycle partagent le
transport dodécaphasique et conservent des informations de types différents.
L’événement signé s’obtient à partir de la direction et de l’orientation de chaque
fenêtre de l’histoire APP--TRIT--TPK. Son accroissement vectoriel et le
bilan~\eqref{mem:balance} retiennent l’interaction avec le registre préalable.
La construction suivante prolonge la lecture spectrale du même cycle,
une fois ses sous-espaces d’incidence fixés. Cet ordre permet de composer
le retour, le déterminant et les moments en conservant l’histoire qui les précède.

Nous conservons l’opérateur $C$ et les projecteurs de
\eqref{iii:eq:projectors34}, et définissons
\begin{equation}
 E_{\rm cyc}=P_3^{\mathrm{cyc}}-P_4^{\mathrm{cyc}},\qquad
 b_n=\operatorname{Tr}(C^n E_{\rm cyc}),\qquad n\geq0.
 \label{amp-afin-trazas}
\end{equation}
Le symbole $b_n$ distingue cette trace spectrale de l’événement signé
$a_m$ d’une fenêtre. La trace utilise l’opérateur de transport ;
l’événement utilise en outre la trajectoire effectivement parcourue.

\begin{proposition}[Tableau des traces produit par l’incidence]
\label{amp-afin-prop-tabla}
On a
\begin{equation}
 \begin{gathered}
 b_n=3\mathbf1_{3\mid n}-4\mathbf1_{4\mid n},\\
 (b_1,\ldots,b_{12})=(0,0,3,-4,0,3,0,-4,3,0,0,-1).
 \end{gathered}
 \label{amp-afin-tabla}
\end{equation}
La période minimale est douze et la somme sur une période est nulle. Pour
$B_N=\sum_{n=1}^{N}b_n$,
\begin{equation}
 B_N=3\lfloor N/3\rfloor-4\lfloor N/4\rfloor,\qquad |B_N|\leq3.
 \label{amp-afin-parciales}
\end{equation}
\end{proposition}
\begin{proof}
Dans $H_d$, l’opérateur $C$ permute cycliquement les $d$ classes de
coordonnées. Sa puissance $n$ a $d$ points fixes si $d$ divise $n$
et aucun dans le cas contraire. Comme $P_d^{\mathrm{cyc}}$ commute avec $C$,
$\operatorname{Tr}(C^nP_d^{\mathrm{cyc}})$ est la trace de cette restriction.
Soustraire les cas $d=3,4$ prouve la formule. La période positive minimale
est douze : $b_n=-1$ se produit exactement aux multiples de douze.
Sur une période, les contributions sont $3\cdot4-4\cdot3=0$.
La somme finie compte les multiples et donne~\eqref{amp-afin-parciales}.
Ses valeurs pour $N=0,\ldots,11$ sont
$(0,0,0,3,-1,-1,2,2,-2,1,1,1)$ et se répètent lorsqu’on ajoute douze.
\end{proof}

\begin{theorem}[Reconstruction logarithmique du lecteur constitutif]
\label{amp-afin-thm-logdet}
Pour $|s|<1$, le quotient $f$ de~\eqref{iii:eq:trace-det} vérifie
\begin{align}
 -\log f(s)&=\sum_{n\geq1}\frac{b_n}{n}s^n,\label{amp-afin-logdet}\\
 -\frac{s f'(s)}{f(s)}
 &=\frac{3s^3}{1-s^3}-\frac{4s^4}{1-s^4}.
 \label{amp-afin-logder}
\end{align}
La branche holomorphe est fixée par $\log f(0)=0$.
Sur $0<s<1$, elle coïncide avec le logarithme réel.
\end{theorem}
\begin{proof}
Les polynômes $1-s^3$ et $1-s^4$ ne s’annulent pas dans le disque ouvert.
Le développement $-\log(1-w)=\sum_{k\geq1}w^k/k$, obtenu en intégrant
la série géométrique depuis zéro, converge uniformément sur les disques
compacts. Avec $w=s^3,s^4$, leur différence donne
$-\log f(s)=\sum_{k\geq1}s^{3k}/k-\sum_{k\geq1}s^{4k}/k$.
Le coefficient de $s^n$ est $b_n/n$.
La convergence locale uniforme permet de dériver et de sommer les deux séries
géométriques, ce qui prouve~\eqref{amp-afin-logder}.
La condition en zéro fixe la constante d’intégration et la branche.
\end{proof}

Ainsi, $b_0/12=-1/12$ enregistre le défaut normalisé de dimension,
tandis que toute la suite $b_n$ reconstruit une fonction.
Ce sont deux lectures de portées différentes de la même paire de sous-espaces.
Lorsqu’on substitue $s=q_\pm^{30}$, les arguments restent les sorties
de la construction angulaire préalable ; le développement ne les sélectionne pas.

\begin{theorem}[Moments de profondeur et représentation de Mellin]
\label{amp-afin-thm-momentos}
La série
\begin{equation}
 M_b(z)=\sum_{n\geq1}\frac{b_n}{n^z}
 \label{amp-afin-momento}
\end{equation}
converge localement uniformément pour $\operatorname{Re}z>0$.
Si $\sigma=\operatorname{Re}z>0$, son reste satisfait
\begin{equation}
 \left|M_b(z)-\sum_{n=1}^{N}\frac{b_n}{n^z}\right|
 \leq3\left(1+\frac{|z|}{\sigma}\right)(N+1)^{-\sigma}.
 \label{amp-afin-cota}
\end{equation}
Dans le domaine de convergence absolue $\operatorname{Re}z>1$,
\begin{equation}
 M_b(z)=(3^{1-z}-4^{1-z})\zeta(z),\qquad
 \zeta(z)=\sum_{n\geq1}n^{-z}.
 \label{amp-afin-reconocimiento}
\end{equation}
Pour $\operatorname{Re}z>0$, on a en outre
\begin{equation}
 \Gamma(z)M_b(z)=
 \int_0^\infty t^{z-1}
 \left(\frac3{e^{3t}-1}-\frac4{e^{4t}-1}\right)dt,
 \label{amp-afin-mellin}
\end{equation}
où $\Gamma(z)=\int_0^\infty u^{z-1}e^{-u}\,du$.
\end{theorem}
\begin{proof}
La sommation par parties et $|B_n|\leq3$ donnent, en passant à la limite supérieure,
\[
 \sum_{n>N}b_n n^{-z}
 =-B_N(N+1)^{-z}
   +\sum_{n>N}B_n\bigl(n^{-z}-(n+1)^{-z}\bigr).
\]
La différence de puissances est bornée en module par
$|z|\int_n^{n+1}x^{-\sigma-1}dx$. Cela prouve
\eqref{amp-afin-cota}, la convergence et son uniformité sur les compacts.
Pour $\sigma>1$, la convergence absolue permet de regrouper
les multiples de trois et de quatre séparément, ce qui donne
\eqref{amp-afin-reconocimiento}.

La série géométrique du transport, évaluée en $s=e^{-t}$, produit
\[
 \sum_{n\geq1}b_n e^{-nt}
 =\frac3{e^{3t}-1}-\frac4{e^{4t}-1}=:K_b(t).
\]
Pour $\sigma>1$, on peut intégrer terme à terme parce que
$\sum|b_n|n^{-\sigma}<\infty$. La substitution $u=nt$ donne
$\Gamma(z)n^{-z}$ et prouve~\eqref{amp-afin-mellin} dans ce demi-plan.
À l’origine, les termes $1/t$ s’annulent :
$K_b(t)=1/2+O(t)$. À l’infini, $K_b(t)=O(e^{-3t})$.
L’intégrale est donc holomorphe pour $\operatorname{Re}z>0$
par convergence locale uniforme. La série l’est aussi par la borne
précédente ; le principe d’identité prolonge l’égalité à ce domaine.
\end{proof}

En $0<\operatorname{Re}z\leq1$, l’ordre croissant de profondeur
fixe la lecture convergente de la série. La réalisation comme trace
absolue de l’opérateur diagonal correspond à $\operatorname{Re}z>1$.

\begin{corollary}[Évaluations compatibles avec l’extrémité constitutive]
\label{amp-afin-cor-extremos}
On a
\begin{equation}
 M_b(2)=\frac{\zeta(2)}{12},\qquad
 M_b(1)=\log(4/3)=-\log f(1),\qquad f(1)=\frac34.
 \label{amp-afin-extremos}
\end{equation}
\end{corollary}
\begin{proof}
La première égalité résulte de $3^{-1}-4^{-1}=1/12$.
Pour la seconde, les sommes partielles exactes sont
$H_{\lfloor N/3\rfloor}-H_{\lfloor N/4\rfloor}$, avec
$H_j=\sum_{k=1}^j1/k$ et $H_0=0$.
La comparaison intégrale de $1/x$ entre ces deux indices donne la limite
$\log(4/3)$. Annuler $1-s$ avant d’évaluer le quotient déterminantal
donne $f(1)=3/4$.
\end{proof}

Le registre affine retient les événements et leur interaction avec la mémoire
préalable ; le tableau des traces retient l’incidence des deux sous-espaces
à toutes les puissances du transport ; le moment applique une pondération
de profondeur à ce tableau déjà produit. Sa périodicité spectrale
coexiste avec une histoire d’événements qui se poursuit après le retour
de phase. Inverser $C$ conserve $b_n$, parce que $d\mid n$ est invariant
sous $n\mapsto-n$ ; le coefficient orienté de la résolvante de
\eqref{iii:eq:resolvent}, avec ses marques fixes, change de signe.
La lecture conjointe conserve une distinction que le déterminant
scalaire n’enregistre pas à lui seul. La forme quadratique et la somme des
coordonnées du registre maintiennent leurs types mathématiques ; cette
composition ne leur assigne pas d’unités d’énergie ni de charge électrique.

\subsection{Composition résiduelle et produit de moments}
\label{amp-afin-sec-productos}

La composition ultérieure des lecteurs conserve également l’opération
effectuée sur la profondeur. Dans la réalisation
$\mathcal H=\ell^2(\mathbb N_{>0})$, nous fixons
$N\delta_n=n\delta_n$. Soient $a,c$ deux tableaux résiduels périodiques
déjà produits et $Q_a\delta_n=a(n)\delta_n$,
$Q_c\delta_n=c(n)\delta_n$. Lorsqu’on les prolonge tous deux à la période commune,
les deux lecteurs agissent sur le même indice. Pour
$\operatorname{Re}z>1$, nous écrivons
$M_a(z)=\operatorname{Tr}(N^{-z}Q_a)$ et de même $M_c$.

\begin{proposition}[Composition diagonale et composition tensorielle]
\label{amp-afin-prop-productos}
On a les identités
\begin{align}
 \operatorname{Tr}(N^{-z}Q_aQ_c)
 &=\sum_{n\geq1}\frac{a(n)c(n)}{n^z},
 \label{amp-afin-diagonal}\\
 M_a(z)M_c(z)
 &=\operatorname{Tr}_{\mathcal H\otimes\mathcal H}
 \bigl((N\otimes N)^{-z}(Q_a\otimes Q_c)\bigr)\notag\\
 &=\sum_{k\geq1}\frac{(a*c)(k)}{k^z},
 \label{amp-afin-tensor}\\
 (a*c)(k)&=\sum_{d\mid k}a(d)c(k/d),
 \qquad \operatorname{Re}z>1.
 \label{amp-afin-convolucion}
\end{align}
Tous les opérateurs figurant à l’intérieur de ces traces sont à trace.
\end{proposition}
\begin{proof}
Le produit diagonal agit sur $\delta_n$ au moyen de
$a(n)c(n)$, ce qui prouve~\eqref{amp-afin-diagonal}.
Dans $\delta_m\otimes\delta_n$, l’opérateur $N\otimes N$
a pour valeur propre $mn$ et $Q_a\otimes Q_c$ agit au moyen de
$a(m)c(n)$. Si $\sigma=\operatorname{Re}z>1$, la somme des
modules de ces coefficients spectraux est bornée par
$\|a\|_\infty\|c\|_\infty(\sum n^{-\sigma})^2<\infty$.
Cette borne prouve la condition de classe de trace et permet de factoriser
la somme double ou de la regrouper par $k=mn$. Les paires de produit $k$
sont exactement $(d,k/d)$, avec $d\mid k$.
La borne analogue avec une seule somme prouve le cas diagonal.
\end{proof}

Dans les lecteurs de cette section, prendre $a(n)=b_n$ et
$c(n)=\langle v,C^nu\rangle$ donne, pour $n=3$,
$a(3)c(3)=-3$, tandis que $(a*c)(3)=3$.
Cette différence de coefficients met en évidence l’information conservée par
chaque opération. La composition diagonale maintient une profondeur commune ;
la tensorielle conserve deux profondeurs avant de les regrouper par leur produit.
La réalisation tensorielle n’attribue pas d’indépendance statistique aux
histoires HMT qui ont produit les lecteurs. Les dérivées régularisées et
les changements de normalisation d’autres moments constituent des opérations
ultérieures distinctes des deux compositions démontrées ici.

\subsection{Restitution fonctionnelle du moment orienté}
\label{iii:sec:restitucion-funcional-catalan}

La relation différentielle~\eqref{iii:eq:catalan-vacuum} admet un
inverse qui conserve la famille complète de moments. Son domaine est
la fonction constitutive produite par les secteurs du cycle, avec
l’orientation de~\eqref{iii:eq:uv}. Les incidences $90=3\cdot30$ et
$120=4\cdot30$ fixent cette fonction avant de l’évaluer dans les canaux
$s_\pm=q_\pm^{30}$. La condition à l’origine correspond à la
disparition du moment lorsque son argument contractant s’éteint.
Cette condition détermine également les constantes d’intégration.

\begin{theorem}[Restitution intégrale et unicité]
\label{iii:thm:restitucion-funcional-catalan}
Soit $f(s)=(1-s^3)/(1-s^4)$, pour $0<s<1$, la réponse
de~\eqref{iii:eq:trace-det}, et soit $\mathcal D=s\,d/ds$.
Il existe une unique fonction $F\in C^2((0,1))$ qui satisfait
\begin{equation}
 \mathcal D^2F(s)=\frac{s(1+s)}{1+s+s^2}f(s),
 \qquad \lim_{s\downarrow0}F(s)=0.
 \label{iii:eq:restitucion-funcional-problema}
\end{equation}
Elle est donnée par
\begin{align}
 F(s)&=\int_0^s\log\frac{s}{t}\,
          \frac{1+t}{1+t+t^2}f(t)\,dt
       =\int_0^s\frac{\log(s/t)}{1+t^2}\,dt
       =\mathcal C_2(s),
 \label{iii:eq:restitucion-funcional-integral}\\
 f(s)&=\frac{1+s+s^2}{s(1+s)}\mathcal D^2F(s).
 \label{iii:eq:restitucion-funcional-inversa}
\end{align}
L’extension continue à $s=1$ retrouve $F(1)=G_{\rm Cat}$.
\end{theorem}
\begin{proof}
La factorisation $1+t+t^2+t^3=(1+t)(1+t^2)$ donne
\[
 \frac{1+t}{1+t+t^2}f(t)=\frac1{1+t^2}.
\]
L’intégrande de~\eqref{iii:eq:restitucion-funcional-integral} est
intégrable en zéro, parce que
\[
 0\leq F(s)\leq\int_0^s\log(s/t)\,dt=s.
\]
Par conséquent, $F(s)\to0$. La différentiation sur l’intervalle ouvert,
d’abord sur un intervalle tronqué, puis par convergence dominée,
produit
\[
 \mathcal DF(s)=\int_0^s\frac{dt}{1+t^2},
 \qquad \mathcal D^2F(s)=\frac{s}{1+s^2}
     =\frac{s(1+s)}{1+s+s^2}f(s).
\]
En particulier, $F$ est de classe $C^2$ et satisfait l’équation.
Si $F_1,F_2$ sont deux solutions, la différence $H=F_1-F_2$
satisfait $\mathcal D^2H=0$. Au moyen de $z=\log s$, on obtient
$d^2H(e^z)/dz^2=0$, donc $H(s)=a\log s+b$.
L’existence de la limite nulle lorsque $s\downarrow0$ impose d’abord
$a=0$, puis $b=0$. La condition sur $F$ détermine donc
la solution et également la limite $\mathcal DF(s)\to0$.

Pour $s<1$, le développement géométrique de $(1+t^2)^{-1}$ et
l’intégrabilité absolue de ses termes pondérés permettent d’intégrer :
\[
 \int_0^s t^{2k}\log(s/t)\,dt
     =\frac{s^{2k+1}}{(2k+1)^2},\qquad
 F(s)=\sum_{k\geq0}\frac{(-1)^ks^{2k+1}}{(2k+1)^2}.
\]
C’est précisément la famille~\eqref{iii:eq:c2}. La convergence
absolue et uniforme de cette dernière série sur $[0,1]$ autorise le
passage à l’extrémité $s=1$. La borne
intégrable $\log(1/t)$ l’autorise également dans l’intégrale, en prolongeant l’intégrande par zéro pour
$t>s$. Enfin, isoler $f$ dans l’équation différentielle donne
\eqref{iii:eq:restitucion-funcional-inversa}, avec des dénominateurs
positifs sur $(0,1)$.
\end{proof}

La limite du moment et celle de sa dérivée seconde sont prises
dans leurs classes de convergence respectives. En particulier,
\[
 \lim_{s\uparrow1}\mathcal D^2\mathcal C_2(s)=\frac12
\]
est la limite radiale d’Abel de la série dérivée. L’évaluation
de $\mathcal C_2(1)$ utilise directement sa série absolument
convergente. Cette distinction conserve le même opérateur de profondeur
à l’intérieur et dans sa lecture de frontière.

\subsection{Compatibilité de la restitution avec les canaux et la vitesse}
\label{iii:sec:restitucion-canales-velocidad}

Le théorème restitue une fonction à partir d’une autre fonction préalablement
construite. L’inversion cubique~\eqref{iii:eq:cubic} agit, en revanche,
sur les évaluations $r_\pm$ : elle retrouve les arguments $s_\pm$
au sein de cette famille fixée. La composition des deux opérations
conserve l’orientation du quart de tour, les étiquettes $P_\pm$
et les incidences du transport. Ainsi, $G_{\rm Cat}=\mathcal C_2(1)$
est la lecture à l’extrémité d’un objet fonctionnel qui participe également
aux réponses intérieures $r_\pm=f(s_\pm)$.

Pour $\mathcal V_m=\mathcal V\otimes I_{9^m}$ et
$\tau_m=\operatorname{Tr}/(2\cdot9^m)$, la lecture de vitesse
maintient exactement la normalisation de
\eqref{iii:eq:velocidad-traza-normalizada} :
\begin{equation}
 \exp[-2\tau_m(\log\mathcal V_m)]
   =\frac1{r_+r_-}
   =(\det\mathcal V)^{-1}=\widehat c.
 \label{iii:eq:restitucion-velocidad-normalizada}
\end{equation}
En effet, chaque valeur propre $r_\pm$ se répète $9^m$ fois, de sorte
que $2\tau_m(\log\mathcal V_m)=\log r_++\log r_-$.
L’identité utilise le déterminant de la réalisation à deux feuillets
et la trace normalisée de son amplification. Les sections
dimensionnelles déjà fixées donnent alors
\[
 c_{\rm int}=c_{\rm vac}=\widehat c\,u_C,
 \qquad \ell_0=c_{\rm vac}t_0,
\]
avec le même état, la même coordinatisation et la même durée élémentaire
de~\eqref{iii:eq:longitud-elemental-constitutiva}.
La restitution fonctionnelle conserve cette section de vitesse lorsqu’on
la compose avec l’action et la longueur ; la multiplicité spectrale
est absorbée par la normalisation déclarée.

\section{Évaluation à l’échelle de Pell et fonctionnelle constitutive de Barbero--Immirzi}
\label{iii:sec:pell-barbero}

La relation entre la réponse du vide et le moment de Catalan permet de
composer deux niveaux du transport : la profondeur pondérée du caractère
de quart de tour et l’évaluation des canaux angulaires. L’échelle de
Pell distingue un argument de la famille de moments ; les canaux du
vide reçoivent, en revanche, les arguments déterminés par la paire angulaire.
La fonctionnelle de Barbero--Immirzi s’exprime sur ces derniers au moyen de
la même réponse constitutive et de son inverse. Le lien s’établit ainsi
entre des opérations définies, en conservant la position de chaque évaluation
dans la construction APP--TRIT--TPK.

\subsection{Moment orienté et évaluation dans l’unité de Pell}

Écrivons
\begin{equation}
 G_{\rm Cat}:=\mathcal C_2(1)
 =\sum_{k\geq0}\frac{(-1)^k}{(2k+1)^2},\qquad
 \lambda_{\rm P}=2+\sqrt3,\qquad
 \rho_{\rm P}=2-\sqrt3=\lambda_{\rm P}^{-1}.
 \label{iii:eq:pell-scales}
\end{equation}
La constante $G_{\rm Cat}$ est la valeur à l’extrémité du moment construit dans
\eqref{iii:eq:c2}. La paire $\lambda_{\rm P},\rho_{\rm P}$ correspond aux
orientations expansive et contractante de l’unité de Pell :
$(2+\sqrt3)(2-\sqrt3)=1$. L’identité de norme $2^2-3\cdot1^2=1$
fixe la réciprocité. La spécialisation du moment en $\rho_{\rm P}$
est postérieure à la construction de son caractère orienté.

L’équation différentielle \eqref{iii:eq:c2-integral}, avec condition nulle
à l’origine, donne
\begin{equation}
 \mathcal D\mathcal C_2(s)=\arctan s,\qquad
 \mathcal C_2(s)=\int_0^s\frac{\arctan t}{t}\,dt,
 \qquad \mathcal D=s\frac{d}{ds}.
 \label{iii:eq:pell-moment-integral}
\end{equation}
En effet, intégrer $\mathcal D^2\mathcal C_2=s/(1+s^2)$ par rapport à
$ds/s$ produit $\arctan s$ ; une seconde intégration produit l’expression
indiquée. L’intégrande se prolonge continûment en $t=0$ avec la valeur $1$.
La forme angulaire de cette intégrale permet de déterminer exactement
l’évaluation en $\rho_{\rm P}$.

\begin{lemma}[Intégrale angulaire du caractère impair]
\label{iii:lem:log-tangent}
Pour $0<\theta\leq\pi/4$,
\begin{equation}
 -\int_0^\theta\log(\tan x)\,dx
 =\sum_{\substack{n\geq1\\n\text{ impair}}}
      \frac{\sin(2n\theta)}{n^2}.
 \label{iii:eq:log-tangent}
\end{equation}
\end{lemma}
\begin{proof}
Pour $0<a<1$, le développement du logarithme à l’intérieur du disque unité implique
\begin{align*}
 L_a(x)
 &:=\log|1-ae^{2ix}|-\log|1+ae^{2ix}|\\
 &=-2\sum_{\substack{n\geq1\\n\text{ impair}}}
       \frac{a^n\cos(2nx)}{n}.
\end{align*}
La série converge uniformément sur $x$ pour chaque $a<1$ ; son intégration
terme à terme fournit
\[
 -\int_0^\theta L_a(x)\,dx
 =\sum_{\substack{n\geq1\\n\text{ impair}}}
       \frac{a^n\sin(2n\theta)}{n^2}.
\]
Lorsque $a\uparrow1$, le quotient des modules converge vers
$|1-e^{2ix}|/|1+e^{2ix}|=\tan x$ pour $0<x\leq\theta$.
La singularité du logarithme à l’origine est intégrable. Pour justifier
la limite, on peut restreindre $a\geq1/2$ et observer que
\[
 |1-ae^{2ix}|^2=(1-a)^2+4a\sin^2x
 \geq2\sin^2x,
 \qquad 1\leq|1+ae^{2ix}|\leq2
\]
sur $0\leq x\leq\pi/4$. Avec la borne supérieure
$|1-ae^{2ix}|\leq2$, ces inégalités donnent
$|L_a(x)|\leq c+|\log x|$ avec une constante indépendante de $a$.
La convergence dominée permet de passer à la limite dans l’intégrale.
Dans la série intégrée, on applique la domination par $\sum n^{-2}$.
Les deux limites prouvent \eqref{iii:eq:log-tangent}.
\end{proof}

\begin{theorem}[Évaluation Catalan--Pell]
\label{iii:thm:catalan-pell}
Le moment orienté satisfait
\begin{align}
 \mathcal C_2(\rho_{\rm P})
 &=\frac23G_{\rm Cat}-\frac{\pi}{12}\log\lambda_{\rm P},
 \label{iii:eq:catalan-pell}\\
 \mathcal D\mathcal C_2(\rho_{\rm P})&=\frac{\pi}{12},&
 \mathcal D^2\mathcal C_2(\rho_{\rm P})&=\frac14.
 \label{iii:eq:catalan-pell-derivatives}
\end{align}
\end{theorem}
\begin{proof}
La substitution $t=\tan x$ dans \eqref{iii:eq:pell-moment-integral},
suivie d’une intégration par parties, donne
\begin{equation}
 \mathcal C_2(\tan\theta)
 =\theta\log(\tan\theta)
   -\int_0^\theta\log(\tan x)\,dx.
 \label{iii:eq:angular-moment}
\end{equation}
Le terme de frontière en zéro s’annule parce que $x\log(\tan x)\to0$.
La formule de la tangente d’une différence fournit
\[
 \tan\frac{\pi}{12}
 =\tan\left(\frac\pi4-\frac\pi6\right)
 =\frac{\sqrt3-1}{\sqrt3+1}=2-\sqrt3.
\]
Il reste à évaluer la série de \eqref{iii:eq:log-tangent} en
$\theta=\pi/12$. Soit $\chi_{-4}(n)=(-1)^{(n-1)/2}$ pour $n$ impair.
Dans les six classes impaires modulo $12$, on vérifie
\begin{equation}
 2\sin\frac{n\pi}{6}
 =\chi_{-4}(n)+3\,\mathbf1_{3\mid n}\chi_{-4}(n/3).
 \label{iii:eq:character-pell}
\end{equation}
La vérification finie est
\[
\begin{array}{c|rrrrrr}
 n\bmod12&1&3&5&7&9&11\\ \hline
 2\sin(n\pi/6)&1&2&1&-1&-2&-1\\
 \chi_{-4}(n)&1&-1&1&-1&1&-1\\
 3\mathbf1_{3\mid n}\chi_{-4}(n/3)&0&3&0&0&-3&0
\end{array}
\]
et les deux membres sont périodiques modulo $12$ sur les indices impairs.
La convergence absolue permet de séparer les termes et d’introduire $n=3m$
dans le second terme :
\begin{align*}
 \sum_{\substack{n\geq1\\n\text{ impair}}}
       \frac{\sin(n\pi/6)}{n^2}
 &=\frac12\sum_{n\text{ impair}}\frac{\chi_{-4}(n)}{n^2}
  +\frac32\sum_{m\text{ impair}}\frac{\chi_{-4}(m)}{(3m)^2}\\
 &=\left(\frac12+\frac16\right)G_{\rm Cat}
 =\frac23G_{\rm Cat}.
\end{align*}
Les sommes sans borne inférieure dans cette expression parcourent les entiers
positifs. En substituant dans \eqref{iii:eq:angular-moment} et en utilisant
$\log\rho_{\rm P}=-\log\lambda_{\rm P}$, on obtient
\eqref{iii:eq:catalan-pell}. La première dérivée résulte de
$\arctan\rho_{\rm P}=\pi/12$ ; la seconde, de
\[
 \frac{\rho_{\rm P}}{1+\rho_{\rm P}^2}
 =\frac1{\rho_{\rm P}+\rho_{\rm P}^{-1}}=\frac14.
\]
\end{proof}

Le coefficient $2/3$ provient de la décomposition du caractère dans
\eqref{iii:eq:character-pell} : la restriction aux multiples de $3$
modifie le poids de profondeur par le facteur $3^{-2}$. L’évaluation
réunit donc le quart de tour et la subdivision dodécaphasique dans une
identité exacte de moments. Les coefficients sont fixés avant toute
comparaison décimale.

\subsection{Composition de phase et profondeur pondérée}

La représentation de profondeur précise ce que conserve cette évaluation.
Dans $\ell^2(\mathbb N_{>0})$, définissons
\[
 N\delta_n=n\delta_n,\qquad U_4\delta_n=i^n\delta_n,
 \qquad Q_4=\frac{U_4-U_4^*}{2i}.
\]
Pour $0<s<1$, l’opérateur $N^{-2}s^NQ_4$ est à trace, et
\begin{equation}
 \mathcal C_2(s)=\operatorname{Tr}(N^{-2}s^NQ_4).
 \label{iii:eq:pell-depth-trace}
\end{equation}
Cette égalité s’obtient en évaluant sur la base $\delta_n$ ; les
entrées diagonales sont $s^n\sin(n\pi/2)/n^2$. L’échange
$U_4\leftrightarrow U_4^*$ inverse $Q_4$ et, avec lui, le moment
orienté. L’indice $n$ et la pondération $s^n$ restent fixes.

La composition s’effectue avant de prendre la trace. Pour
$a\in\mathbb Z/4\mathbb Z$, soit $B(s,a)=s^NU_4^a$. Comme les deux
facteurs sont diagonaux dans la même base,
\begin{equation}
 B(s,a)B(t,b)=B(st,a+b),\qquad
 B(s,1)^4=s^{4N}.
 \label{iii:eq:pell-depth-composition}
\end{equation}
Le retour de phase conserve la contraction accumulée. En particulier,
$B(\rho_{\rm P},1)^4=\rho_{\rm P}^{4N}$, dont la valeur propre à la profondeur
$n$ est $\rho_{\rm P}^{4n}$. Cette réalisation donne un sens
opératoriel à l’évaluation contractante ; la somme scalaire du moment
agrège les profondeurs après avoir conservé leur indice et leur orientation.
Les arguments constitutifs $s_\pm=q_\pm^{30}$ restent déterminés par
la paire angulaire de \eqref{iii:eq:recover-angular}. Les formules précédentes
ne contiennent pas d’égalité entre ces arguments et $\rho_{\rm P}$.

\subsection{Incidence dodécaphasique et fonctionnelle de retour}

La contribution de Barbero--Immirzi utilise le transport angulaire et
l’incidence hexade--octade. Soient $\mathsf H\subset\mathsf O$ une hexade
et une octade avec la paire marquée de l’hexade conservée. Les ensembles de
drapeaux pertinents sont
\begin{equation}
 \mathcal F_6=\mathsf H\times\binom{\mathsf H}{2}
 \lhook\joinrel\longrightarrow
 \mathcal F_8=\mathsf O\times\binom{\mathsf H}{2},
 \qquad |\mathcal F_6|=90,\quad |\mathcal F_8|=120.
 \label{iii:eq:barbero-flags}
\end{equation}
Les cardinaux sont $6\binom62$ et $8\binom62$ ; la même incidence
conserve le défaut normalisé
$(90-120)/(24\binom62)=-1/12$. Ce dénombrement identifie les deux secteurs
employés par le retour, tandis que le calendrier détermine leurs
multiplicités relatives.

Pour $m\in\{90,120\}$, le retour tritique sur un canal contractant est
\begin{equation}
 S_m(q)=\sum_{j\geq0}q^{m+3mj}
       =\frac{q^m}{1-q^{3m}},\qquad 0<q<1.
 \label{iii:eq:barbero-return-series}
\end{equation}
Chaque terme correspond à la hauteur initiale $m$ et à $j$ retours de
longueur $3m$. La série géométrique conserve donc tant la hauteur
initiale que la période du retour. Le calendrier dodécaphasique
$\mathbb Z/12\mathbb Z$ apporte douze secteurs et une unique incidence
orientée de frontière, notée $\partial_{\rm ret}$. Dans le module
libre des événements, sa chaîne fondamentale est
\begin{equation}
 c_{12}^{+}
 =\sum_{j\in\mathbb Z/12\mathbb Z}[j,\mathcal F_6]
  -[\partial_{\rm ret},\mathcal F_8].
 \label{iii:eq:barbero-cycle}
\end{equation}
Le lecteur linéaire $\mathscr L_q$ assigne $S_{90}(q)$ à chaque générateur
$[j,\mathcal F_6]$ et $S_{120}(q)$ au générateur de frontière. On obtient
\begin{equation}
 \mathscr L_q(c_{12}^{+})=12S_{90}(q)-S_{120}(q).
 \label{iii:eq:barbero-weights}
\end{equation}
Les poids $12$ et $-1$ sont les multiplicités orientées de la chaîne.
L’orientation conjuguée satisfait $c_{12}^{-}=-c_{12}^{+}$, avec les
mêmes types d’incidence.

\begin{proposition}[Coefficient singulier du retour fondamental]
\label{iii:prop:barbero-pole}
Le coefficient de $(1-q)^{-1}$ de
$\mathscr L_q(c_{12}^{+})$ est $1/24$.
\end{proposition}
\begin{proof}
Pour chaque entier $m>0$, la factorisation
$1-q^{3m}=(1-q)\sum_{j=0}^{3m-1}q^j$ donne
\[
 \lim_{q\uparrow1}(1-q)S_m(q)=\frac1{3m}.
\]
Appliquer cette identité à la fonctionnelle déjà déterminée dans
\eqref{iii:eq:barbero-weights} produit
\begin{equation}
 \lim_{q\uparrow1}(1-q)\mathscr L_q(c_{12}^{+})
 =\frac{12}{270}-\frac1{360}=\frac1{24}.
 \label{iii:eq:barbero-pole}
\end{equation}
Dans la coordonnée complexe $q-1$, le résidu est $-1/24$.
\end{proof}

L’évaluation singulière vérifie une conséquence de l’incidence et du
retour tritique. L’ordre démonstratif est fixé par
\eqref{iii:eq:barbero-flags}--\eqref{iii:eq:barbero-weights} ; le
coefficient singulier se calcule après la construction de la chaîne.

\subsection{Factorisation à travers l’opérateur constitutif}

Conservons la coordinatisation angulaire relevée de la
section~\ref{iii:sec:hojas}. Pour rendre explicite la conversion angulaire,
écrivons
\begin{equation}
 A_{\deg}=\frac{180x}{\pi},\qquad
 C^*_{\deg}=\frac{180y}{\pi},\qquad
 q_\pm=e^{-x\mp y},\qquad s_\pm=q_\pm^{30}.
 \label{iii:eq:barbero-angular-lift}
\end{equation}
L’orientation considérée dans \eqref{iii:eq:rchannels} correspond à
$0<y<x$ ; l’échange de feuillets s’exprime par $y\mapsto-y$ sur le
domaine plus large $x>|y|$. Les exponentielles réelles agissent sur ces
relèvements, qui conservent les valeurs de $x\pm y$ pendant la
composition.

La fonctionnelle orientée de Barbero--Immirzi est
\begin{equation}
 \gamma=\frac{\pi A_{\deg}C^*_{\deg}}{180}
   +12\bigl[S_{90}(q_-)-S_{90}(q_+)\bigr]
   -\bigl[S_{120}(q_-)-S_{120}(q_+)\bigr].
 \label{iii:eq:barbero-functional}
\end{equation}
La première composante est l’évaluation angulaire ; la seconde est la
différence du lecteur de la chaîne fondamentale dans les canaux conjugués.
Les arguments sont reçus du transport angulaire précédemment construit.
La formule détermine un scalaire adimensionnel sur cette structure de
départ.

Soit $\psi:(3/4,1)\to(0,1)$ l’inverse de $f$ construit au moyen de
\eqref{iii:eq:cubic}. Définissons
\begin{equation}
 \mathcal H(s)=\frac{12s^3}{1-s^9}
             -\frac{s^4}{1-s^{12}}
             -\frac{(\log s)^2}{20\pi},\qquad 0<s<1.
 \label{iii:eq:barbero-h}
\end{equation}
Cette fonction est analytique réelle sur son domaine. Les exposants $3,4,9,12$
résultent de l’expression des hauteurs $90,120$ et des retours $270,360$ en
unités de $30$ ; le coefficient de la partie logarithmique est fixé lorsqu’on
transporte la contribution angulaire.

\begin{theorem}[Expression constitutive orientée de Barbero--Immirzi]
\label{iii:thm:barbero-factorization}
Dans la coordinatisation \eqref{iii:eq:barbero-angular-lift},
\begin{equation}
 \gamma=\mathcal H(s_-)-\mathcal H(s_+)
       =\mathcal H\bigl(\psi(r_-)\bigr)
        -\mathcal H\bigl(\psi(r_+)\bigr).
 \label{iii:eq:barbero-factorization}
\end{equation}
Sur la représentation bidimensionnelle des canaux,
\begin{equation}
 \gamma=-\operatorname{Tr}
    \left[\mathcal R\,\mathcal H\bigl(\psi(\mathcal V)\bigr)\right],
 \label{iii:eq:barbero-trace}
\end{equation}
où $\operatorname{Tr}$ est la trace ordinaire, avec
$\operatorname{Tr}I=2$.
\end{theorem}
\begin{proof}
La substitution $s=q^{30}$ dans \eqref{iii:eq:barbero-return-series} donne
\[
 S_{90}(q)=\frac{s^3}{1-s^9},\qquad
 S_{120}(q)=\frac{s^4}{1-s^{12}}.
\]
De plus, \eqref{iii:eq:barbero-angular-lift} implique
\[
 \log s_\pm=-30(x\pm y)
 =-\frac\pi6(A_{\deg}\pm C^*_{\deg}).
\]
Par différence de carrés,
\begin{align}
 \frac{(\log s_+)^2-(\log s_-)^2}{20\pi}
 &=\frac{900\bigl[(x+y)^2-(x-y)^2\bigr]}{20\pi}\\
 &=\frac{180xy}{\pi}
 =\frac{\pi A_{\deg}C^*_{\deg}}{180}.
 \label{iii:eq:barbero-log-term}
\end{align}
La somme de ces identités prouve la première égalité de
\eqref{iii:eq:barbero-factorization}. La seconde utilise l’inversion
$\psi(f(s_\pm))=s_\pm$ démontrée dans la
section~\ref{iii:sec:constitutiva}.

Pour l’expression opératorielle, les projecteurs de feuillet sont de rang un et
$\mathcal RP_\pm=\pm P_\pm$. Le calcul spectral fournit
\[
 \mathcal H\bigl(\psi(\mathcal V)\bigr)
 =\mathcal H(s_+)P_++\mathcal H(s_-)P_-.
\]
Multiplier par $\mathcal R$ et prendre la trace ordinaire donne
$\mathcal H(s_+)-\mathcal H(s_-)$. Son opposé est $\gamma$.
\end{proof}

La marque d’orientation $\mathcal R$ fait partie de la fonctionnelle. Si
l’on maintient cette marque fixe et que l’on échange les canaux de $\mathcal V$,
le signe de $\gamma$ change. L’évaluation du spectre sans ses étiquettes conserve
la paire non ordonnée, tandis que \eqref{iii:eq:barbero-trace} conserve la
contribution orientée. Si l’on emploie la trace normalisée
$\tau_2=\operatorname{Tr}/2$, la même formule s’écrit
$\gamma=-2\tau_2[\mathcal R\mathcal H(\psi(\mathcal V))]$.

En revanche, un changement simultané de représentation
$(\mathcal R,\mathcal V)\mapsto
(U\mathcal RU^{-1},U\mathcal VU^{-1})$ conserve la fonctionnelle :
le calcul spectral se transporte par conjugaison et la trace est
invariante. Échanger la réponse par rapport à une orientation fixe
et transporter conjointement la réponse et l’orientation sont donc
des opérations différentes.

Cette normalisation est compatible avec les amplifications de
\eqref{iii:eq:logvac}. Pour
$\mathcal V_m=\mathcal V\otimes I_{9^m}$ et
$\mathcal R_m=\mathcal R\otimes I_{9^m}$,
\begin{equation}
 \gamma=-\frac1{9^m}\operatorname{Tr}
 \left[\mathcal R_m\mathcal H\bigl(\psi(\mathcal V_m)\bigr)\right].
 \label{iii:eq:barbero-amplification}
\end{equation}
En effet, le calcul fonctionnel conserve le produit tensoriel et la trace
de $I_{9^m}$ est $9^m$. La multiplicité de la représentation se sépare
de l’évaluation orientée.

\subsection{Reconstruction à partir des coordonnées du vide}

Les identités \eqref{iii:eq:four-identities} permettent d’exprimer la même
fonctionnelle au moyen de l’impédance et de la vitesse normalisées :
\begin{align}
 \gamma={}&
 \mathcal H\!\left(\psi\!\left(
     \sqrt{\frac{\widehat Z}{\widehat c}}\right)\right)
 -\mathcal H\!\left(\psi\!\left(
     \frac1{\sqrt{\widehat Z\widehat c}}\right)\right).
 \label{iii:eq:barbero-zc}
\end{align}
Le domaine positif de cette reconstruction est décrit par
\begin{equation}
 1<\widehat Z\widehat c<\frac{16}{9},\qquad
 \frac9{16}<\frac{\widehat Z}{\widehat c}<1.
 \label{iii:eq:barbero-zc-domain}
\end{equation}
En effet, ces inégalités équivalent à
$3/4<r_\pm<1$ lorsqu’on substitue les racines positives de
\eqref{iii:eq:four-identities}. Pour l’orientation $y>0$, on conserve
en outre $r_-<r_+$, ou de manière équivalente $\widehat Z<1$. Les points du
domaine sont utilisés avec la compatibilité angulaire déjà établie par
le TPK ; l’inversion est une reconstruction ultérieure de ses canaux.
Les coordonnées sont les coordonnées internes de \eqref{iii:eq:four} : un changement
d’unités doit être défait au moyen de \eqref{iii:eq:undo-units} avant
d’appliquer $\psi$.

L’échange $y\mapsto-y$ agit comme
$(\widehat Z,\widehat c)\mapsto(\widehat Z^{-1},\widehat c)$ et
produit $\gamma\mapsto-\gamma$. Dans le cas équilibré $y=0$,
les deux canaux coïncident et \eqref{iii:eq:barbero-factorization} donne
$\gamma=0$. La vitesse normalisée conserve le produit des
canaux ; sa lecture conjointe avec l’impédance retrouve les deux
réponses et l’orientation nécessaire à l’évaluation de la fonctionnelle.
Dans ce cas équilibré, $\mathcal V=rI$ a un spectre dégénéré :
la marque $\mathcal R$ est conservée comme donnée de représentation, bien que
le scalaire $r$ ne distingue plus ses projecteurs. Pour un spectre simple,
en revanche,
$P_+=(\mathcal V-r_-I)/(r_+-r_-)$ et $P_-=I-P_+$.

Le résultat central de cette composition est la factorisation
\eqref{iii:eq:barbero-factorization}. Le moment de Catalan apporte le
coefficient orienté qui détermine $f$ ; l’inversion cubique retrouve
l’argument de chaque canal ; l’incidence dodécaphasique et la paire angulaire
déterminent $\mathcal H$. L’échelle de Pell distingue une évaluation
exacte de la famille, avec la preuve de
\eqref{iii:eq:catalan-pell}. Chaque relation conserve son argument et son
opération : la constante à l’extrémité, le moment contractant et la fonctionnelle
constitutive constituent des évaluations reliées de structures
spécifiées.

\section{Normalisation constitutive et incidences spectrales du retour}
\label{iii:sec:incidencias-espectrales}

La réponse des deux feuillets et la fonctionnelle de retour conservent deux
classes complémentaires d’information. La première comprend le produit
et le quotient des valeurs propres constitutives ; la seconde comprend les
hauteurs, les périodes et les multiplicités orientées des incidences.
Cette section réunit les opérations qui permettent de distinguer ces lectures.
Son antécédent est le transport angulaire produit par APP--TRIT--TPK,
avec ses projecteurs et sa coordinatisation relevée. Les normalisations s’appliquent
après la construction de ce transport et la conservation de son orientation.

\subsection{Séparation exacte du mode commun et du mode orienté}

Soient $r_+,r_->0$ les valeurs propres étiquetées de
\eqref{iii:eq:rchannels}. Définissons
\begin{equation}
 a=\sqrt{r_+r_-},\qquad
 \eta=\frac12\log\frac{r_+}{r_-},\qquad
 \widetilde{\mathcal V}=a^{-1}\mathcal V.
 \label{iii:eq:rec27-normalizacion}
\end{equation}
Dans la chambre $x>y>0$, on a $a\in(3/4,1)$ et $\eta>0$.
La marque $\mathcal R$ conserve $\mathcal R P_\pm=\pm P_\pm$.

\begin{proposition}[Factorisation positive avec orientation conservée]
\label{iii:prop:rec27-factorizacion}
La réponse possède la factorisation unique
\begin{equation}
 \mathcal V=a\exp(\eta\mathcal R),\qquad
 \det\widetilde{\mathcal V}=1.
 \label{iii:eq:rec27-unimodular}
\end{equation}
Ses quatre lectures constitutives satisfont
\begin{equation}
 \widehat\varepsilon=a^2e^{2\eta},\quad
 \widehat\mu=a^2e^{-2\eta},\quad
 \widehat Z=e^{-2\eta},\quad
 \widehat c=a^{-2}.
 \label{iii:eq:rec27-cuatro}
\end{equation}
Avec l’orientation fixe, l’échange de feuillets conserve $a$ et remplace
$\eta$ par $-\eta$.
\end{proposition}
\begin{proof}
Le calcul spectral donne
$\exp(\eta\mathcal R)=e^\eta P_++e^{-\eta}P_-$.
Les définitions impliquent $ae^\eta=r_+$ et $ae^{-\eta}=r_-$,
de sorte que l’on retrouve exactement $\mathcal V$.
Le produit de ses valeurs propres fixe l’unique racine positive $a$ ;
leur quotient fixe l’unique valeur réelle de $\eta$.
La substitution dans \eqref{iii:eq:four} prouve les quatre identités.
Enfin, échanger les valeurs propres inverse leur quotient et conserve
leur produit.
\end{proof}

La normalisation unimodulaire conserve toute la coordonnée orientée et
sépare la coordonnée commune. Si l’on conserve également $a$, la transformation
est réversible. Si l’on exprime uniquement $\widetilde{\mathcal V}$,
$\widehat Z$ est déterminée, tandis que $\widehat c$ nécessite le
facteur commun archivé. Cette distinction explique le rôle des feuillets
réciproques dans un lecteur de produit unitaire : un tel lecteur réalise la
composante unimodulaire. La réponse constitutive complète conserve,
en outre, son déterminant.

En particulier, les lecteurs historiques de la forme
$\mu_r^\pm=e^{\pm2u}$ et $\varepsilon_r^\pm=e^{\mp2u}$
maintiennent un produit égal à un dans chaque feuillet. Leur rapport d’impédances est
$e^{4u}$. La lecture actuelle de \eqref{iii:eq:rec27-cuatro}
maintient le produit $a^4=\widehat c^{-2}$. Le paramètre historique $u$
et la coordonnée $\eta$ ne sont identifiés que lorsqu’on déclare l’application
entre ces réalisations ; la décomposition précédente spécifie
exactement quel facteur supplémentaire permet de restituer la réponse complète.

\subsection{Fonction opératoire du moment de Catalan}

La restitution du théorème~\ref{iii:thm:restitucion-funcional-catalan}
situe la constante de Catalan à l’extrémité d’une famille fonctionnelle.
L’opération qui entre dans la réponse intérieure est
\begin{equation}
 r_\pm=
 \frac{1+s_\pm+s_\pm^2}{s_\pm(1+s_\pm)}
 \mathcal D^2\mathcal C_2(s_\pm),\qquad
 s_\pm=q_\pm^{30},\quad \mathcal D=s\frac{d}{ds}.
 \label{iii:eq:rec27-catalan-operativo}
\end{equation}
La famille et sa condition à l’origine déterminent le moment ;
l’inversion cubique détermine ses deux arguments à partir des réponses.
On conserve ainsi deux niveaux de restitution : une fonction à partir d’une fonction
et des arguments à partir d’évaluations dans une fonction déjà fixée.

La composition avec \eqref{iii:eq:rec27-normalizacion} obtient $a$ et
$\eta$ à partir de ces mêmes évaluations. La composition ultérieure avec
\eqref{iii:eq:barbero-factorization} obtient la fonctionnelle orientée
de Barbero--Immirzi. La valeur à l’extrémité $\mathcal C_2(1)$,
les réponses $\mathcal D^2\mathcal C_2(s_\pm)$ et la fonctionnelle
$\mathcal H(s_-)-\mathcal H(s_+)$ remplissent donc des fonctions
mathématiques spécifiques au sein de la même chaîne. La restitution utilise
la famille complète et les conditions d’unicité déjà démontrées.

Pour préciser quelle observation distingue les canaux, écrivons
$u_\pm=\log r_\pm$. Les variations dans une coordinatisation commune vérifient
\begin{equation}
 \begin{pmatrix}d\log\widehat c\\d\log\widehat Z\end{pmatrix}
 =\begin{pmatrix}-1&-1\\-1&1\end{pmatrix}
 \begin{pmatrix}du_+\\du_-\end{pmatrix}.
 \label{iii:eq:rec27-variaciones}
\end{equation}
Le déterminant est $-2$, de sorte que la lecture conjointe retrouve
les deux variations. Chacune séparément conserve une seule combinaison.
Cette identité exprime la fonction complémentaire de la vitesse et de
l’impédance, avec l’orientation et la coordinatisation dimensionnelle fixes.

\subsection{Trois lectures distinctes des incidences $90/120$}

Les drapeaux de \eqref{iii:eq:barbero-flags} déterminent les hauteurs
$90$ et $120$. Le défaut normalisé de leurs cardinaux est $-1/12$.
La chaîne de \eqref{iii:eq:barbero-cycle} détermine, à son tour,
les poids signés $12$ et $-1$. Chaque hauteur possède sa série de retour
$S_m(q)=q^m/(1-q^{3m})$. Conserver ces trois opérations permet de
comparer précisément les deux fonctionnelles
\begin{equation}
 L_{\rm eq}(q)=S_{90}(q)-S_{120}(q),\qquad
 L_{\rm ret}(q)=12S_{90}(q)-S_{120}(q).
 \label{iii:eq:rec27-dos-kernels}
\end{equation}
La première correspond à l’antécédent à poids égaux ; la seconde est
la lecture de la chaîne orientée actuelle.

\begin{proposition}[Normalisations du pôle de retour]
Pour $q\uparrow1$, on a
\begin{equation}
 \lim(1-q)L_{\rm eq}(q)=\frac1{1080},\qquad
 \lim(1-q)L_{\rm ret}(q)=\frac1{24},\qquad
 L_{\rm ret}-L_{\rm eq}=11S_{90}.
 \label{iii:eq:rec27-polos}
\end{equation}
\end{proposition}
\begin{proof}
La factorisation de $1-q^{3m}$ prouve
$\lim(1-q)S_m(q)=1/(3m)$.
Par linéarité, les deux coefficients sont respectivement
$1/270-1/360=1/1080$ et $12/270-1/360=1/24$.
Soustraire les deux définitions prouve la dernière égalité.
\end{proof}

Le rapport $45$ entre ces coefficients singuliers résulte des poids
spécifiés. Les fonctions conservent également leurs termes réguliers ;
le rapport des pôles n’identifie pas leurs valeurs sur tout l’intervalle.
Le défaut de cardinal, le poids de la chaîne et le coefficient singulier
sont ainsi reliés à des opérations et à des normalisations différenciées.

\clearpage
\subsection{Résolution de Fourier du motif dodécaphasique}

Sur le même cycle, le motif associé aux incidences orientées est
représenté par
\begin{equation}
 p_m=12\cos\frac{2\pi\,3m}{12}
      -\cos\frac{2\pi\,4m}{12},\qquad 0\le m<12.
 \label{iii:eq:rec27-patron}
\end{equation}
Les fréquences $3$ et $4$ réalisent les accroissements de phase $90^\circ$ et
$120^\circ$. Les arguments des fonctions trigonométriques sont
exprimés en radians. Les poids proviennent de la chaîne déjà construite ;
la représentation harmonique s’applique ensuite.

\begin{proposition}[Support et normalisation spectrale]
\label{iii:prop:rec27-fourier}
Pour la transformée non normalisée
$\widehat p_k=\sum_{m=0}^{11}p_m e^{-2\pi i km/12}$,
\begin{equation}
 \widehat p_3=\widehat p_9=72,\qquad
 \widehat p_4=\widehat p_8=-6,
 \qquad \widehat p_k=0\ \text{aux autres indices}.
 \label{iii:eq:rec27-fourier}
\end{equation}
La moyenne du motif est nulle et
$\sum_m|p_m|^2=870$.
\end{proposition}
\begin{proof}
Écrire chaque cosinus comme la demi-somme de deux exponentielles produit
les coefficients $6$ dans les caractères $3,9$ et $-1/2$ dans $4,8$.
L’orthogonalité finie
$\sum_{m=0}^{11}e^{2\pi i(j-k)m/12}=12\delta_{jk}$
donne \eqref{iii:eq:rec27-fourier}. Le coefficient d’indice zéro
s’annule. La même orthogonalité donne
$\sum|p_m|^2=(2\cdot72^2+2\cdot6^2)/12=870$.
\end{proof}

Avec la transformée divisée par $12$, les coefficients sont $6$ et
$-1/2$ ; avec la normalisation unitaire, ils sont $12\sqrt3$ et $-\sqrt3$.
Le rapport des modules $12:1$ est conservé dans les deux. La suite exacte
\begin{equation}
 2(p_0,\ldots,p_{11})
 =(22,1,-23,-2,25,1,-26,1,25,-2,-23,1)
 \label{iii:eq:rec27-patron-racional}
\end{equation}
permet de vérifier la transformée au moyen d’une arithmétique algébrique exacte.

La moyenne nulle du motif appartient à l’annulation des caractères dans le
cycle. Le défaut $-1/12$ appartient au dénombrement normalisé des incidences ;
$1/24$ appartient au pôle de la série de retour. Maintenir ces domaines
explique ce que conserve chaque lecture avant de la composer avec la paire angulaire
et avec la fonctionnelle de Barbero--Immirzi.

% BEGIN AMPLIACION_MEMORIA_COMPATIBILIDAD_20260922
\section{\HMTLang{Réciprocité différentielle et restitution des canaux constitutifs}{Réciprocité différentielle et récupération des canaux constitutifs}}
\label{delta22:iii:restitucion}
\HMTLang{Les canaux du vide, la fonctionnelle de Barbero--Immirzi et le registre dodécaphasique appartiennent à une même chaîne de constructions. Le résultat suivant réunit leurs inverses : la vitesse normalisée conserve la composante commune de la réponse, tandis que la lecture orientée de Barbero détermine sa répartition entre feuillets. La composition retrouve la paire angulaire et, sur l’image engendrée avec sa section dimensionnelle, le coefficient d’action. En conservant en outre les coordonnées régionales, elle atteint la lecture périodique de \(K\).}{Les canaux du vide, la fonctionnelle de Barbero--Immirzi et le registre dodécaphasique appartiennent à une même chaîne de constructions. Le résultat suivant réunit leurs inverses : la vitesse normalisée conserve la composante commune de la réponse, tandis que la lecture orientée de Barbero détermine sa distribution entre feuillets. La composition retrouve la paire angulaire et, sur l’image engendrée avec sa section dimensionnelle, le coefficient d’action. Conserver les coordonnées régionales donne également la lecture périodique de \(K\).}

\HMTLang{La construction précédente d’APP--TRIT--TPK et de son état enrichi demeure l’antécédent. Les paramètres employés ci-après sont les représentations angulaires déjà obtenues ; les preuves analytiques décrivent leur réalisation et les opérations de récupération. Les projecteurs étiquetés \(P_\pm\), \(\mathcal R=P_+-P_-\), l’orientation et le relèvement angulaire restent explicites.}{La construction précédente d’APP--TRIT--TPK et de son état enrichi demeure l’antécédent. Les paramètres utilisés ci-après sont les sorties angulaires déjà obtenues ; les preuves analytiques décrivent leur réalisation et leurs opérations de récupération. Les projecteurs étiquetés \(P_\pm\), \(\mathcal R=P_+-P_-\), l’orientation et le relèvement angulaire restent explicites.}

\subsection{\HMTLang{Potentiel spectral et réciprocité constitutive}{Potentiel spectral et réciprocité constitutive}}
\HMTLang{Sur le cône \(x>|y|\), soient}{Sur le cône \(x>|y|\), soient}
\[
 T=e^{-xI-y\mathcal R},\quad s_\pm=e^{-30(x\pm y)},\quad
 f(s)=\frac{1-s^3}{1-s^4},\quad
 V=(I-T^{90})(I-T^{120})^{-1}=r_+P_++r_-P_-,
\]
\[
 r_\pm=f(s_\pm),\quad
 \widehat c=(r_+r_-)^{-1},\quad \widehat Z=r_-/r_+,
 \quad c_\ell=\log\widehat c,\quad z_\ell=\log\widehat Z.
\]
\HMTLang{L’inversibilité de \(I-T^{120}\) découle de \(0<s_\pm<1\). Le quotient de polynômes prolongé est \(f(s)=(1+s+s^2)/(1+s+s^2+s^3)\) ; ses extrémités sont \(1\) et \(3/4\).}{L’inversibilité de \(I-T^{120}\) découle de \(0<s_\pm<1\). Le quotient de polynômes prolongé est \(f(s)=(1+s+s^2)/(1+s+s^2+s^3)\) ; ses valeurs aux extrémités sont \(1\) et \(3/4\).}
\begin{theorem}[\HMTLang{Potentiel de la réponse normalisée}{Potentiel de la réponse normalisée}]
\HMTLang{La série}{La série}
\begin{equation}
 \mathscr F(x,y)=\sum_{\sigma=\pm1}\sum_{n\ge1}
 \left[\frac{e^{-120n(x+\sigma y)}}{120n^2}
       -\frac{e^{-90n(x+\sigma y)}}{90n^2}\right]
 \label{delta22:iii:potencial}
\end{equation}
\HMTLang{définit un potentiel analytique réel strictement concave, avec}{définit un potentiel analytique réel strictement concave, avec}
\[
 \nabla\mathscr F=(c_\ell,z_\ell),\qquad
 \partial_y\log\widehat c=\partial_x\log\widehat Z.
\]
\HMTLang{La normalisation est \(\mathscr F\to0\) lorsque \(x-|y|\to\infty\).}{Sa normalisation est \(\mathscr F\to0\) lorsque \(x-|y|\to\infty\).}
\end{theorem}
\begin{proof}
\HMTLang{Sur chaque compact du cône, il existe \(\epsilon>0\) avec \(x\pm y\ge\epsilon\). Chaque dérivée d’ordre fixé des termes est dominée par une puissance de \(n\) multipliée par \(e^{-90n\epsilon}\) ; la somme converge uniformément. La dérivation terme à terme et \(\sum_{n\ge1}z^n/n=-\log(1-z)\) donnent les deux composantes indiquées. Pour la concavité, soient}{Sur chaque partie compacte du cône, il existe \(\epsilon>0\) avec \(x\pm y\ge\epsilon\). Toute dérivée d’ordre fixé d’un terme est bornée par une puissance de \(n\) multipliée par \(e^{-90n\epsilon}\) ; la somme converge uniformément. La dérivation terme à terme et \(\sum_{n\ge1}z^n/n=-\log(1-z)\) donnent les composantes indiquées. Pour la concavité, posons}
\[
 g(w)=\log f(e^{-w}),\quad u=30(x+y),\quad v=30(x-y),
 \quad a=g'(u),\quad b=g'(v).
\]
\[
 g'(w)=\frac3{e^{3w}-1}-\frac4{e^{4w}-1}>0,
 \qquad
 D(c_\ell,z_\ell)=-30
 \begin{pmatrix}a+b&a-b\\a-b&a+b\end{pmatrix}.
\]
\HMTLang{L’inégalité s’obtient parce que \(t/(e^{tw}-1)\) décroît strictement en \(t>0\) : sa dérivée a pour numérateur \(e^{tw}(1-tw)-1<0\). Les valeurs propres de la matrice sont \(-60a\) et \(-60b\), toutes deux négatives, et son déterminant est \(3600ab>0\). La symétrie prouve la réciprocité différentielle. La borne exponentielle donne la limite de normalisation.}{L’inégalité découle du fait que \(t/(e^{tw}-1)\) est strictement décroissante pour \(t>0\) : sa dérivée a pour numérateur \(e^{tw}(1-tw)-1<0\). Les valeurs propres de la matrice sont \(-60a\) et \(-60b\), toutes deux négatives, et son déterminant est \(3600ab>0\). La symétrie prouve la réciprocité différentielle. La borne exponentielle donne la limite de normalisation.}
\end{proof}
\HMTLang{Le même moment spectral \(\mathcal D_2(z)=\sum_{n\ge1}z^n/n^2\) permet d’écrire \eqref{delta22:iii:potencial} au moyen de ses évaluations réelles contractantes. Le caractère orienté déjà construit satisfait \(\mathcal C_2(s)=\operatorname{Im}\mathcal D_2(is)\), \(\mathcal C_2(1)=G_{\rm Cat}\). La convergence absolue atteint \(s=1\) ; pour les séries dérivées qui cessent d’y converger, on conserve la limite radiale d’Abel. Les évaluations réelles et orientées partagent l’indice de profondeur et la pondération, en conservant leurs arguments respectifs. Sous l’amplification \(T\otimes I_{9^m}\), la trace divisée par \(9^m\) conserve ces lectures, car chaque valeur propre se répète exactement \(9^m\) fois. Cette identité correspond à cette amplification déclarée.}{Le même moment spectral \(\mathcal D_2(z)=\sum_{n\ge1}z^n/n^2\) exprime \eqref{delta22:iii:potencial} au moyen d’évaluations réelles contractantes. Le caractère orienté déjà construit satisfait \(\mathcal C_2(s)=\operatorname{Im}\mathcal D_2(is)\), \(\mathcal C_2(1)=G_{\rm Cat}\). La convergence absolue s’étend jusqu’à \(s=1\) ; pour les séries dérivées qui cessent d’y converger, on conserve la limite radiale d’Abel. Les évaluations réelles et orientées partagent l’indice de profondeur et la pondération tout en conservant leurs arguments respectifs. Sous l’amplification par \(T\otimes I_{9^m}\), la trace divisée par \(9^m\) préserve ces lectures, puisque chaque valeur propre est répétée exactement \(9^m\) fois. Cette identité concerne l’amplification spécifiée.}

\begin{proposition}[\HMTLang{Sensibilité orientée et convexité}{Sensibilité orientée et convexité}]
\HMTLang{Pour \(x>0\) fixé, \(y\mapsto c_\ell(x,y)\) est paire et strictement convexe, avec un minimum unique en zéro.}{Pour \(x>0\) fixé, \(y\mapsto c_\ell(x,y)\) est paire et strictement convexe, avec son minimum unique en zéro.}
\end{proposition}
\begin{proof}
\[
 g''(w)=-\frac9{4\sinh^2(3w/2)}+\frac{16}{4\sinh^2(2w)}<0.
\]
\HMTLang{En effet, \(t/\sinh(tw/2)\) décroît : sa dérivée a le signe de \(\sinh v-v\cosh v<0\). Comme \(c_\ell=-g(u)-g(v)\), on obtient \(\partial_y^2c_\ell=-900[g''(u)+g''(v)]>0\). La symétrie sous \(y\mapsto-y\) fixe le minimum. Le développement local est}{En effet, \(t/\sinh(tw/2)\) décroît : sa dérivée a le signe de \(\sinh v-v\cosh v<0\). Puisque \(c_\ell=-g(u)-g(v)\), on obtient \(\partial_y^2c_\ell=-900[g''(u)+g''(v)]>0\). La symétrie sous \(y\mapsto-y\) fixe le minimum. Le développement local est}
\[
 c_\ell=-2g(30x)-900g''(30x)y^2+O(y^4),\qquad
 z_\ell=-60g'(30x)y+O(y^3).
\]
\end{proof}
\HMTLang{La vitesse retient la séparation des canaux au second ordre ; l’impédance conserve leur orientation au premier ordre. Évaluer le canal moyen avant le lecteur élimine précisément ces contributions.}{La vitesse retient la séparation des canaux au second ordre ; l’impédance retient leur orientation au premier ordre. Évaluer le canal moyen avant d’appliquer le lecteur élimine précisément ces contributions.}

\subsection{\HMTLang{Inversion globale au moyen de la vitesse et de Barbero--Immirzi}{Inversion globale au moyen de la vitesse et de Barbero--Immirzi}}
\HMTLang{L’inverse \(\psi=f^{-1}:(3/4,1)\to(0,1)\) est défini par l’unique racine de}{L’inverse \(\psi=f^{-1}:(3/4,1)\to(0,1)\) est défini par l’unique racine de}
\[
 rs^3+(r-1)(1+s+s^2)=0.
\]
\HMTLang{L’unicité découle de}{L’unicité découle de}
\[
 f'(s)=-\frac{s^2(s^2+2s+3)}{(1+s+s^2+s^3)^2}<0.
\]
\HMTLang{Nous conservons la fonctionnelle de retour déjà dérivée de l’incidence hexade--octade :}{Conservons la fonctionnelle de retour déjà dérivée de l’incidence hexade--octade :}
\[
 \mathcal H(s)=\frac{12s^3}{1-s^9}-\frac{s^4}{1-s^{12}}
              -\frac{(\log s)^2}{20\pi},\qquad F_{\rm B}=\mathcal H\circ\psi.
\]
\begin{theorem}[\HMTLang{Restitution constitutive orientée}{Récupération constitutive orientée}]
\HMTLang{L’application}{L’application}
\[
 (r_+,r_-)\longmapsto
 \left(\widehat c=\frac1{r_+r_-},\quad
 \gamma=F_{\rm B}(r_-)-F_{\rm B}(r_+)\right)
\]
\HMTLang{est un difféomorphisme analytique réel de \((3/4,1)^2\) sur \((1,16/9)\times\mathbb R\). Cette affirmation décrit le domaine analytique du lecteur ; l’image effectivement engendrée conserve, en outre, les compatibilités du TPK.}{est un difféomorphisme analytique réel de \((3/4,1)^2\) sur \((1,16/9)\times\mathbb R\). Cette affirmation décrit le domaine analytique du lecteur ; l’image effectivement engendrée conserve en outre les conditions de compatibilité du TPK.}
\end{theorem}
\begin{proof}
\[
 \mathcal H'(s)=\frac{36s^2(1+2s^9)}{(1-s^9)^2}
 -\frac{4s^3(1+2s^{12})}{(1-s^{12})^2}
 -\frac{\log s}{10\pi s}>0.
\]
\HMTLang{Le quotient du premier terme positif par la grandeur du second est}{Le rapport du premier terme positif à la grandeur du second est}
\[
 \frac9s\frac{1+2s^9}{1+2s^{12}}
 \left(\frac{1-s^{12}}{1-s^9}\right)^2>9,
\]
\HMTLang{et le dernier terme de la somme est également positif. Les extrémités sont \(\mathcal H(s)\sim-(\log s)^2/(20\pi)\) en zéro et \(\mathcal H(s)\sim5/[4(1-s)]\) en un. Par conséquent, \(F_{\rm B}'<0\), avec les limites \(+\infty\) en \(3/4\) et \(-\infty\) en \(1\).}{et le dernier terme de la somme est également positif. Les comportements asymptotiques aux extrémités sont \(\mathcal H(s)\sim-(\log s)^2/(20\pi)\) en zéro et \(\mathcal H(s)\sim5/[4(1-s)]\) en un. Donc \(F_{\rm B}'<0\), avec les limites \(+\infty\) en \(3/4\) et \(-\infty\) en \(1\).}

\HMTLang{Pour \(c=\widehat c\) fixé, soit \(z=\log\widehat Z\). Alors}{Pour \(c=\widehat c\) fixé, posons \(z=\log\widehat Z\). Alors}
\[
 r_+=c^{-1/2}e^{-z/2},\quad r_-=c^{-1/2}e^{z/2},\quad
 |z|<a(c)=\min\{\log c,\log(16/(9c))\}.
\]
\HMTLang{La fonction \(\mathcal G_c(z)=F_{\rm B}(r_-)-F_{\rm B}(r_+)\) satisfait}{La fonction \(\mathcal G_c(z)=F_{\rm B}(r_-)-F_{\rm B}(r_+)\) satisfait}
\[
 \mathcal G_c'(z)=\tfrac12
 [r_-F_{\rm B}'(r_-)+r_+F_{\rm B}'(r_+)]<0.
\]
\HMTLang{À l’extrémité supérieure, \(r_-\to1\), ou \(r_+\to3/4\), ou les deux, et \(\mathcal G_c\to-\infty\). À l’extrémité inférieure, elle tend vers \(+\infty\). Le théorème des valeurs intermédiaires et la monotonie donnent une unique solution pour chaque \(\gamma\). La dérivée non nulle fournit l’inverse analytique par le théorème des fonctions implicites. La description par le minimum de \(a(c)\) délimite seulement la frontière ; l’argument s’effectue sur son domaine ouvert.}{À l’extrémité supérieure, \(r_-\to1\), ou \(r_+\to3/4\), ou les deux, et \(\mathcal G_c\to-\infty\). À l’extrémité inférieure, elle tend vers \(+\infty\). Le théorème des valeurs intermédiaires et la monotonie donnent une unique solution pour chaque \(\gamma\). La dérivée non nulle donne l’inverse analytique par le théorème des fonctions implicites. Le minimum dans \(a(c)\) décrit seulement la frontière ; l’argument se déroule sur son domaine ouvert.}
\end{proof}

\subsection{\HMTLang{Récupération angulaire, action et registre périodique}{Récupération angulaire, action et registre périodique}}
\HMTLang{Une fois \(r_\pm\) obtenus, l’inversion cubique donne \(s_\pm\). Avec le relèvement et l’orientation conservés,}{Une fois \(r_\pm\) obtenus, l’inversion cubique donne \(s_\pm\). Avec le relèvement et l’orientation conservés,}
\[
 x=-\frac1{60}\log(s_+s_-),\quad y=\frac1{60}\log\frac{s_-}{s_+},
 \quad A_{\deg}=-\frac3\pi\log(s_+s_-),\quad
 C^*_{\deg}=\frac3\pi\log\frac{s_-}{s_+}.
\]
\HMTLang{Sur la section positive engendrée \(A_{\deg}=1000\alpha\), \(C^*_{\deg}=2(\eta_{\rm ret}+\alpha)\), on retrouve}{Sur la section positive engendrée \(A_{\deg}=1000\alpha\), \(C^*_{\deg}=2(\eta_{\rm ret}+\alpha)\), on retrouve}
\[
 \alpha=A_{\deg}/1000,\qquad
 \eta_{\rm ret}=C^*_{\deg}/2-\alpha,\qquad
 \hbar_{\rm ret}=s_0\eta_{\rm ret},\qquad s_0:=10^{-34}\mathcal U_S.
\]
\HMTLang{L’échelle \(s_0\) et les projecteurs sont des antécédents conservés. Dans l’orientation conjuguée, l’étiquette d’action est également transportée. La récupération est un inverse sur l’image correspondante.}{L’échelle \(s_0\) et les projecteurs sont des antécédents conservés. Sous l’orientation conjuguée, l’étiquette d’action est également transportée. La récupération est un inverse sur l’image correspondante.}

\begin{proposition}[\HMTLang{Composition avec la lecture périodique de douze blocs}{Composition avec la lecture périodique de douze blocs}]
\HMTLang{La lecture périodique prolonge les douze blocs de \(K\), tandis que les coordonnées régionales conservent leur génération coinductive complète. Avec \(\alpha_A\) comme coordonnée de la clôture arithmétique précédemment définie,}{La lecture périodique prolonge les douze blocs de \(K\), tandis que les coordonnées régionales conservent leur génération coinductive complète. Avec \(\alpha_A\) désignant la coordonnée de clôture arithmétique précédemment définie,}
\[
 \kappa_{\rm per}=\pi_{\rm HMT}+e_{\rm HMT}-\varphi_{\rm HMT}-4-\alpha_A,
 \quad B=1000,\quad N_K=(B^{12}-1)\kappa_{\rm per}.
\]
\HMTLang{Sur l’image du registre engendré, \(N_K\) est entier et ses douze blocs, avec origine fixe et zéros initiaux conservés, se retrouvent exactement par}{Sur l’image du registre engendré, \(N_K\) est un entier et ses douze blocs, avec origine fixe et zéros initiaux conservés, se retrouvent exactement par}
\[
 K_j=\left\lfloor\frac{N_K}{B^{12-j}}\right\rfloor
 -B\left\lfloor\frac{N_K}{B^{13-j}}\right\rfloor,
 \qquad 1\le j\le12.
\]
\end{proposition}
\begin{proof}
\HMTLang{La division euclidienne successive de \(N_K\) en base \(B\) détermine de manière unique \(K_j\in\{0,\ldots,B-1\}\) et produit \(N_K=\sum_{j=1}^{12}K_jB^{12-j}\). La répétition périodique de ces blocs somme la série géométrique \(N_K\sum_{m\ge1}B^{-12m}=N_K/(B^{12}-1)\). Les deux compositions sont inverses. La lecture finie \(N_K/B^{12}\) conserve une autre normalisation, distincte de la périodique. Avec des données approchées, la récupération exige un intervalle certifié pour \(N_K\) contenant un unique entier.}{La division euclidienne successive de \(N_K\) en base \(B\) détermine de manière unique \(K_j\in\{0,\ldots,B-1\}\) et donne \(N_K=\sum_{j=1}^{12}K_jB^{12-j}\). La répétition périodique de ces blocs somme la série géométrique \(N_K\sum_{m\ge1}B^{-12m}=N_K/(B^{12}-1)\). Les deux compositions sont inverses. La lecture finie \(N_K/B^{12}\) conserve une normalisation différente de celle de la lecture périodique. Avec des données approchées, la récupération nécessite un intervalle certifié pour \(N_K\) contenant un unique entier.}
\end{proof}
\HMTLang{Cette composition réunit la restitution constitutive et la récupération du registre déjà construit. La sélection exceptionnelle et temporelle détermine \(K\) au préalable ; l’inverse explique quelle information ses lecteurs permettent de restituer. La portée est le registre ordonné de douze coordonnées avec les antécédents indiqués.}{Cette composition réunit la restitution constitutive et la récupération du registre déjà construit. La sélection exceptionnelle et temporelle détermine \(K\) au préalable ; l’inverse explique quelle information ses lecteurs peuvent restituer. Sa portée est le registre ordonné de douze coordonnées avec les antécédents indiqués.}

\subsection{\HMTLang{Bornes de stabilité dans le domaine contractant}{Bornes de stabilité dans le domaine contractant}}
\HMTLang{Pour \(w>0\), écrivons \(a(w)=g'(w)\). Avec \(s=e^{-w}\),}{Pour \(w>0\), écrivons \(a(w)=g'(w)\). Avec \(s=e^{-w}\),}
\[
 a(w)=\frac{s^3(s^2+2s+3)}{(1+s+s^2)(1+s+s^2+s^3)},
 \qquad \tfrac12e^{-3w}<a(w)<3e^{-3w}.
\]
\HMTLang{Les bornes s’obtiennent en comparant les polynômes positifs du quotient ; \(g''<0\) démontre sa décroissance. Sur un domaine convexe où \(0<m_0\le30(x\pm y)\le M_0<\infty\), l’application \(\mathscr C=(c_\ell,z_\ell)\) satisfait, pour la norme euclidienne,}{Les bornes découlent de la comparaison des polynômes positifs du quotient ; \(g''<0\) prouve sa décroissance. Sur un domaine convexe où \(0<m_0\le30(x\pm y)\le M_0<\infty\), l’application \(\mathscr C=(c_\ell,z_\ell)\) satisfait, pour la norme euclidienne,}
\[
 60a(M_0)\|p-q\|\le\|\mathscr C(p)-\mathscr C(q)\|
 \le60a(m_0)\|p-q\|.
\]
\HMTLang{Pour la borne supérieure, on intègre la jacobienne sur le segment. Pour la borne inférieure, on applique le produit scalaire avec \(p-q\) à la jacobienne définie négative, puis Cauchy--Schwarz. La norme induite de la différentielle inverse est \(1/[60\min\{a(u),a(v)\}]\). En contraction profonde, elle peut croître exponentiellement, même lorsque le nombre de condition relatif reste proche de un. L’unicité exacte et la robustesse face à l’erreur sont donc spécifiées par des critères différents et compatibles.}{Pour la borne supérieure, on intègre la jacobienne le long du segment. Pour la borne inférieure, on prend le produit scalaire avec \(p-q\), on utilise le caractère défini négatif de la jacobienne, puis on applique Cauchy--Schwarz. La norme induite de la différentielle inverse est \(1/[60\min\{a(u),a(v)\}]\). En contraction profonde, elle peut croître exponentiellement même lorsque le nombre de condition relatif reste proche de un. L’unicité exacte et la robustesse face à l’erreur sont ainsi spécifiées par des critères distincts et compatibles.}

% END AMPLIACION_MEMORIA_COMPATIBILIDAD_20260922
\section{Section de charge et relations de quantification électrique}
\label{iii:sec:carga}

La réponse constitutive établit une relation entre les deux feuillets et une échelle relative de propagation. La section de charge s'obtient en composant cette réponse avec la constante de structure fine et l'action, toutes deux construites dans les chapitres précédents. L'ordre a une conséquence précise : la charge ne détermine pas rétrospectivement le registre dodécaphasique. C'est le registre, avec l'action et la réponse du vide, qui détermine la section de charge dans la droite correspondante.

\subsection{Détermination positive de la section de charge}

Considérons les sections positives $\hbar_a$, avec
$a\in\{\mathrm{pre},\mathrm{ret}\}$, et $h_a=2\pi\hbar_a$.
L’impédance $Z_{\rm vac}$ est déjà construite au moyen de
\eqref{iii:eq:dimensions}. Le quotient $h_a/Z_{\rm vac}$ appartient à
la droite quadratique de charge. Cette compatibilité dimensionnelle permet de
formuler la définition suivante sans choisir encore de coordinatisation SI.

\begin{definition}[Section positive de charge]
La section de charge correspondant à l’orientation d’action $a$ est
\begin{equation}
 e_a=\left(\frac{2\alpha h_a}{Z_{\rm vac}}\right)^{1/2}.
 \label{iii:eq:charge}
\end{equation}
La racine est prise dans la demi-droite positive de charge. Sa section conjuguée
est $-e_a$ ; le signe de cette dernière n’est pas identifié à l’indice de
retour de l’action.
\end{definition}

\begin{proposition}[Existence, unicité et compatibilité du couplage]
L’équation $Z_{\rm vac}e_a^2=2\alpha h_a$ détermine une unique section
positive. Dans la réalisation constitutive, on a aussi
\begin{equation}
 \alpha=\frac{Z_{\rm vac}e_a^2}{2h_a}
 =\frac{e_a^2}{4\pi\varepsilon_{\rm vac}\hbar_a c_{\rm vac}}.
 \label{iii:eq:coupling}
\end{equation}
\end{proposition}
\begin{proof}
La positivité de $\alpha$, $h_a$ et $Z_{\rm vac}$ garantit une unique racine positive. Les identités constitutives donnent $\varepsilon_{\rm vac}c_{\rm vac}=Z_{\rm vac}^{-1}$. Substituer cette égalité et $h_a=2\pi\hbar_a$ dans la première expression de \eqref{iii:eq:coupling} produit la seconde.
\end{proof}

La dernière expression est la forme habituelle du couplage électromagnétique
dans une coordinatisation physique. Elle apparaît ici comme identité ultérieure des sections
construites. La preuve ne recalcule pas $\alpha$ à partir d’une charge
mesurée : elle établit la cohérence algébrique du résultat précédent avec la
section qui vient d’être définie. L’identification physique exige en outre
de maintenir la coordinatisation des unités et le régime du couplage employés dans la
comparaison. Une identité entre grandeurs et une comparaison expérimentale
ont des fonctions différentes dans cette reconnaissance.

\subsection{Résistance, conductance, flux et fréquence}

Les quatre grandeurs suivantes utilisent la même section $e_a$. Leur
dépendance peut être entièrement résolue avant l’évaluation de la moindre décimale.

\begin{theorem}[Composition des quanta électriques]
Les sections
\begin{align}
 R_K&=\frac{h_a}{e_a^2}=\frac{Z_{\rm vac}}{2\alpha},
 &G_0&=\frac{2e_a^2}{h_a}=\frac{4\alpha}{Z_{\rm vac}},
 \label{iii:eq:rk-g0}\\
 \Phi_{0,a}&=\frac{h_a}{2e_a}
 =\sqrt{\frac{h_aZ_{\rm vac}}{8\alpha}},
 &K_{J,a}&=\frac{2e_a}{h_a}=\Phi_{0,a}^{-1}
 \label{iii:eq:flux-josephson}
\end{align}
satisfont exactement
\begin{equation}
 R_KG_0=2,\qquad G_0Z_{\rm vac}=4\alpha,\qquad
 K_{J,a}\Phi_{0,a}=1,\qquad
 K_{J,a}^{2}R_K=\frac4{h_a}.
 \label{iii:eq:quanta-identities}
\end{equation}
\end{theorem}
\begin{proof}
La substitution de $e_a^2=2\alpha h_a/Z_{\rm vac}$ dans les deux premières
définitions annule $h_a$ et donne \eqref{iii:eq:rk-g0}. En élevant au
carré la section positive $h_a/(2e_a)$, on obtient
$h_aZ_{\rm vac}/(8\alpha)$ ; la positivité fixe la racine dans
\eqref{iii:eq:flux-josephson}. Son inverse est $2e_a/h_a$.
Les trois premiers produits de \eqref{iii:eq:quanta-identities} se
réduisent, respectivement, à $2$, $4\alpha$ et $1$. Enfin,
$(4e_a^2/h_a^2)(h_a/e_a^2)=4/h_a$.
\end{proof}

La réalisation physique reconnaît dans ces sections la résistance de von Klitzing, le quantum de conductance, le quantum de flux et la constante de Josephson. Le facteur $2$ de la conductance appartient à la normalisation de la grandeur définie ici ; son identification à une multiplicité de canaux d'un système matériel exige de spécifier cette représentation.
Les quatre équations ne constituent pas quatre générateurs indépendants : elles partagent $\alpha$, l'action et la réponse du vide. C'est précisément cette dépendance qui produit les identités croisées et permet de vérifier une évaluation sans altérer les opérations qui l'ont engendrée.

\subsection{Retour de l’action et transport des sections}

Le quotient $R_{\rm act}=h_{\rm pre}/h_{\rm ret}$, obtenu avant le
couple angulaire, conserve la différence entre les deux sections d’action.
Il ne change pas les assignations de feuillet de la réponse constitutive. Ce sont deux
opérations distinctes : retour de l’action et échange des canaux.

\begin{proposition}[Covariance sous le retour]
Avec $\alpha$ et $Z_{\rm vac}$ communs, on a
\begin{align}
 \frac{e_{\rm pre}}{e_{\rm ret}}&=R_{\rm act}^{1/2},
 &R_K^{\rm pre}&=R_K^{\rm ret},
 &G_0^{\rm pre}&=G_0^{\rm ret},\label{iii:eq:charge-return}\\
 \frac{\Phi_{0,\rm pre}}{\Phi_{0,\rm ret}}&=R_{\rm act}^{1/2},
 &\frac{K_{J,\rm pre}}{K_{J,\rm ret}}&=R_{\rm act}^{-1/2}.
 \label{iii:eq:flux-return}
\end{align}
\end{proposition}
\begin{proof}
La première identité découle de la division des deux versions de
\eqref{iii:eq:charge} et de la prise de la racine positive. Dans
\eqref{iii:eq:rk-g0}, l’action s’est déjà annulée. Les deux autres
identités découlent des puissances $h_a^{1/2}$ et $h_a^{-1/2}$ dans
\eqref{iii:eq:flux-josephson}.
\end{proof}

La résistance et la conductance conservent le quotient constitutif, tandis que le flux et la fréquence retiennent également l'orientation de l'action. Cette séparation donne une lecture structurelle de leurs dépendances : certaines grandeurs sont des invariants du quotient des feuillets et d'autres mesurent en outre la section d'action utilisée. La comparaison métrologique doit déclarer quelle section elle évalue ; la proximité numérique des deux actions n'autorise pas à les identifier.

\section{Évaluation constitutive et réalisation dimensionnelle}
\label{iii:sec:evaluacion}

L’évaluation réunit maintenant les constructions précédentes dans leur ordre effectif.
Les développements régionaux de $\pi$, $\varphi$ et $e$ proviennent du raffinement
nonadique ; le registre $K$ conserve ses douze positions. Sa lecture périodique
détermine la section positionnelle de $\alpha$. Sur celle-ci agissent le caractère
d’action et le retour régional, puis la paire angulaire et, enfin, la
réponse constitutive. Cette évaluation est postérieure à la sélection des
règles : aucun de ses chiffres n’est utilisé pour choisir un canal.

\subsection{Évaluation de la chaîne interne}

Pour reproduire le tableau, on emploie les développements régionaux à mille
décimales engendrés par le certificat nonadique et l’entier $N_K$ de
\eqref{eq:k-numerador}. La lecture utilisée est
\[
 \alpha=\pi+e-\varphi-4-\frac{N_K}{1000^{12}-1}.
\]
La section finie de douze blocs et la section prolongée ont leurs
définitions et leur différence exacte dans le chapitre du registre. Ici, on
maintient la seconde fixe pendant toute l’évaluation. On n’échange pas
les lectures lorsque la précision augmente, et on ne remplace pas l’opérateur analytique
de la seconde voie par une coïncidence décimale avec la première.

\begin{longtable}{@{}ll@{}}
\toprule
Objeto & Évaluation décimale (24 chiffres significatifs)\\
\midrule\endhead
$\alpha$ & $0.00729735256928380099728511\ldots$ \\
$H_5$ & $1.05457181764615446962582\ldots$ \\
$\eta_{\rm ret}$ & $1.05457181691503906113369\ldots$ \\
$A_{\deg}$ & $7.29735256928380099728511\ldots$ \\
$C^*_{\deg}$ & $2.12373833896864572426195\ldots$ \\
$q_+$ & $0.848377942576404374945737\ldots$ \\
$q_-$ & $0.913660151150557285295519\ldots$ \\
$r_+$ & $0.999999628548249363178695\ldots$ \\
$r_-$ & $0.999724137189518364131517\ldots$ \\
$\widehat\varepsilon$ & $0.999999257096636702760442\ldots$ \\
$\widehat\mu$ & $0.999448350479326935089982\ldots$ \\
$\widehat Z$ & $0.999724508538937215455554\ldots$ \\
$\widehat c$ & $1.00027631048615752610639\ldots$ \\
$\gamma$ & $0.274007672694459274747255\ldots$ \\
$G_{\rm Cat}$ & $0.915965594177219015054604\ldots$ \\
$\rho$ & $0.267949192431122706472554\ldots$ \\
\bottomrule
\end{longtable}


Le calcul s’effectue avec 120 chiffres de travail ; on montre 24 chiffres
significatifs pour faciliter la comparaison entre grandeurs de différents
ordres. Les 120 chiffres sont une précision d’évaluation, non une incertitude
expérimentale ni, à eux seuls, un certificat d’erreur par intervalles.
Le paquet conserve les entrées engendrées, leurs empreintes, les opérations et
les résidus des identités vérifiées. Les preuves algébriques du
corps du texte établissent ces identités indépendamment de ce tableau.

Le sens des quatre dernières coordonnées constitutives est conservé
lors de l’évaluation : $\widehat\varepsilon$ et $\widehat\mu$ sont les réponses
quadratiques par feuillet, $\widehat Z$ est leur orientation relative et
$\widehat c$ est l’inverse de leur produit. La proximité de certains chiffres avec
l’unité exprime le petit défaut du quotient $3$--$4$ sur ces
canaux. Son origine se voit exactement dans
\[
 1-f(s)=\frac{s^3}{1+s+s^2+s^3}.
\]
La formule permet d’étudier la réponse avant de lui attribuer une grandeur
dimensionnelle : le signe du défaut, son ordre cubique et sa monotonie sont
des conséquences du même opérateur.

\subsection{Coordinatisations d’unités et covariance}

Soient $\mathcal L_S$, $\mathcal L_Q$ et $\mathcal L_C$ les droites
d’action, de charge et de vitesse. Une coordinatisation physique positive évalue les sections
$u_S$, $u_Q$ et $u_C$ en unités d’action, de charge et de vitesse. La réalisation
\eqref{iii:eq:dimensions} produit alors une impédance, une vitesse,
une perméabilité et une permittivité, avec leurs dimensions respectives.
Les coordonnées du tableau précédent restent adimensionnelles.

\begin{proposition}[Covariance de la réalisation constitutive]
Si les sections de référence changent selon
$u_S'=a u_S$, $u_Q'=b u_Q$, $u_C'=d u_C$, avec $a,b,d>0$, les
coordonnées qui représentent les mêmes grandeurs satisfont
\[
 \widehat Z'=a^{-1}b^2\widehat Z,\qquad
 \widehat c'=d^{-1}\widehat c,\qquad
 \widehat\mu'=a^{-1}db^2\widehat\mu,\qquad
 \widehat\varepsilon'=adb^{-2}\widehat\varepsilon.
\]
Les relations $\widehat\mu'\widehat\varepsilon'=(\widehat c')^{-2}$
et $\widehat\mu'/\widehat\varepsilon'=(\widehat Z')^2$ sont conservées.
\end{proposition}
\begin{proof}
On égale chaque grandeur de \eqref{iii:eq:dimensions} à son expression
dans la nouvelle base et on isole sa coordonnée. Les facteurs du produit
sont $d^2$, et ceux du quotient sont $a^{-2}b^4$, exactement ceux requis
par les coordonnées transformées de $c$ et $Z$.
\end{proof}

Cette covariance distingue deux actes : construire une section et exprimer
cette section dans des unités choisies. Le théorème spectral fixe les coordonnées
dans la normalisation interne déclarée. Une comparaison SI doit également présenter
l’évaluation des trois sections de référence. Les ajuster
au moyen des grandeurs que l’on souhaite reproduire serait un étalonnage de
la coordinatisation, non une vérification indépendante de ces grandeurs.

Pour reconstruire les canaux à partir de coordonnées exprimées dans la coordinatisation
primée, on défait ce changement avant d’utiliser la racine cubique :
\begin{equation}
 \widehat Z=ab^{-2}\widehat Z',\quad
 \widehat c=d\widehat c',\quad
 \widehat\mu=ad^{-1}b^{-2}\widehat\mu',\quad
 \widehat\varepsilon=a^{-1}d^{-1}b^2\widehat\varepsilon'.
 \label{iii:eq:undo-units}
\end{equation}
Par exemple, les entrées de l’inverse spectral sont
$r_+=(\widehat Z\widehat c)^{-1/2}$ et
$r_-=(\widehat Z/\widehat c)^{1/2}$, calculées après
\eqref{iii:eq:undo-units}. Introduire directement les coordonnées primées
dans ces expressions changerait la normalisation de l’opérateur et pourrait
les faire sortir de l’intervalle $(3/4,1)$. La covariance dimensionnelle conserve
les grandeurs physiques ; l’inversion cubique reconstruit le transport
dans la coordinatisation interne dans laquelle il a été défini.

\subsection{Contrôles d’identité et portée de la comparaison}

La comparaison reproductible s’organise par relations. En premier lieu,
$\widehat\mu\widehat\varepsilon\widehat c^2=1$ et
$\widehat\mu/(\widehat\varepsilon\widehat Z^2)=1$ vérifient la
cohérence constitutive. En deuxième lieu, l’inversion cubique retrouve
les arguments et les coordonnées angulaires de départ. Enfin,
l’évaluation de Barbero par la fonctionnelle angulaire et par le calcul fonctionnel
de la réponse coïncide par l’identité prouvée dans le chapitre précédent.
L’échange de feuillets doit inverser la contribution orientée et maintenir
le mode commun ; ce contrôle détecte une perte d’étiquetage que les
produits scalaires ne révéleraient pas.

La comparaison de la constante réduite de Planck avec les tableaux de
référence est conservée dans le développement antécédent d’action. Pour les
quatre grandeurs du vide, le présent tableau est une évaluation de leurs
coordonnées internes. Il ne se présente pas comme un tableau de coïncidences SI :
cette affirmation supplémentaire nécessite une coordinatisation numérique indépendante pour
$u_S,u_Q,u_C$. Le résultat constitutif démontré, y compris son inversion,
demeure complet sur son domaine et conserve explicitement l’interface
dimensionnelle nécessaire pour effectuer cette comparaison.

\section{Conclusions}
\label{iii:sec:conclusiones}
\begingroup
\fontsize{10.5}{12.8}\selectfont
\setlength{\parskip}{2pt plus .5pt minus .5pt}

Le résultat central est une détermination constitutive réversible dans
le domaine contractant : les deux canaux orientés d’un transport
produisent quatre coordonnées reliées et ces coordonnées permettent de
retrouver la paire angulaire. La structure fixe d’abord l’équation de
réponse ; son évaluation détermine les valeurs sans perdre l’orientation
nécessaire à la reconstruction de leurs antécédents spectraux. Cette double
direction, constructive et inverse, donne un contenu précis à la thèse
structure, équation et valeur.

\subsection*{La réponse conserve la paire angulaire}

Les incidences $90$ et $120$ déterminent une unique solution de
l’équation de complétion pour le transport fixé. Sa fonction scalaire
est une bijection strictement décroissante de $(0,1)$ sur $(3/4,1)$.
Par conséquent, l’équation cubique~\eqref{iii:eq:cubic} possède exactement
une racine admissible pour chaque réponse de canal. L’inversion n’est pas
une coïncidence numérique : elle reconstruit les deux contractions et, par
leurs logarithmes, la paire angulaire.

Les quatre coordonnées du vide satisfont
\[
\widehat\mu\widehat\varepsilon=\widehat c^{-2},\qquad
\widehat\mu/\widehat\varepsilon=\widehat Z^2.
\]
Dans les règles fixées de produit, de quotient et de positivité, ces
identités déterminent les évaluations quadratiques. Leur dépendance
commune se manifeste également sous l’échange de feuillets : la permittivité et
la perméabilité se permutent, l’impédance s’inverse et la vitesse est
conservée. Le produit symétrique et le quotient orienté transmettent ainsi
des informations complémentaires. La récupération maintient les étiquettes
de feuillet et la normalisation interne ; l’état d’origine conserve
en outre le chemin, la frontière et la mémoire.

La factorisation $\mathcal V=a\exp(\eta\mathcal R)$ identifie
ces deux informations dans l’opérateur lui-même. Le facteur positif
$a=\sqrt{\det\mathcal V}$ détermine la vitesse normalisée ;
la composante unimodulaire détermine l’impédance. Conserver les deux
permet de reconstruire la réponse. Cette séparation explique la relation
entre les feuillets réciproques de produit unitaire et les quatre lectures
du vide : la normalisation sépare le mode commun, qui demeure
nécessaire pour retrouver l’opérateur complet.

L’inverse permet de relier des réponses qui évaluent des fonctions distinctes
des mêmes canaux. Le rapport de leurs exposants reconstruit
l’anisotropie de l’ellipse d’action et l’impédance LC relative,
selon la proposition~\ref{iii:prop:forma-LC}. La section d’action
et les références de fréquence et d’impédance complètent leurs grandeurs dimensionnelles.
L’identité $c_{\rm int}=c_{\rm vac}$ établit, à son tour, que le
lecteur cinématique et le constitutif réalisent la même section de
vitesse dans une coordinatisation commune. La trace normalisée conserve cette
lecture sous amplification nonadique.

\subsection*{Composition et reconstruction entre secteurs}

La loi transportée $r\odot t=f(\psi(r)\psi(t))$ exprime la composition
du transport en coordonnées constitutives. Son élément neutre est $3/4$ dans
l’extension de frontière, et $-\log\psi(r)/30$ transforme cette composition en
somme. Les projecteurs communs préservent la réalisation opératorielle et les
étiquettes de feuillet. La vitesse d’un état unique conserve, pour sa part,
le produit ordinaire de ses deux canaux. Les deux opérations demeurent
distinguées par leurs arguments et par la preuve exacte de leur relation.

La reconstruction Higgs--impédance établit une autre connexion explicite.
Les couplages de jauge restituent la contraction commune et le
couplage électromagnétique ; l’auto-couplage de Higgs restitue la
coordonnée orientée. Dans le domaine de la réalisation à l’arbre, ces
données déterminent les deux canaux et leurs réponses du vide. Les
inégalités de domaine et les compatibilités généalogiques précisent
quels triplets appartiennent à la réalisation considérée. Conserver les
projecteurs marqués permet de restituer également l’opérateur ; l’histoire
enrichie maintient ses données de provenance supplémentaires.

\subsection*{Le cycle relie réponse, moment et orientation}

La réalisation dodécaphasique explique la provenance de la fonction
rationnelle par un quotient de déterminants en dimensions $3$ et $4$.
La différence des rangs produit $-1/12$ comme trace normalisée dans ce
domaine fini. Le plan du quart de tour conserve une autre information :
son opérateur satisfait $J^2=-I$ et son coefficient de résolvante retient
le signe d’orientation. Les deux intégrations logarithmiques de ce
coefficient, avec leurs conditions initiales, déterminent le moment
de Catalan à profondeur illimitée.

La table des traces $b_n$ reconstruit $-\log f(s)$ et produit des moments
convergents pour $\operatorname{Re}z>0$, avec une borne explicite du
reste. Sa valeur à la profondeur un est $\log(4/3)$, compatible avec
l’extrémité constitutive. La représentation de Mellin et la distinction
entre produit diagonal et tensoriel conservent les opérations qui
précèdent l’évaluation scalaire de chaque moment.

La restitution intégrale du
théorème~\ref{iii:thm:restitucion-funcional-catalan} clôt la
correspondance fonctionnelle dans les deux sens :
\[
 \mathcal C_2(s)=\int_0^s\log\frac{s}{t}\,
       \frac{1+t}{1+t+t^2}f(t)\,dt,\qquad
 f(s)=\frac{1+s+s^2}{s(1+s)}\mathcal D^2\mathcal C_2(s).
\]
La condition $\mathcal C_2(s)\to0$ à l’origine élimine toute
liberté additive et logarithmique. La valeur extrême $G_{\rm Cat}$,
le noyau orienté et la réponse constitutive se trouvent ainsi situés
dans une même collection, avec des opérations de restitution explicites.
Les arguments $s_\pm$ conservent leur détermination angulaire ; la trace
normalisée maintient la vitesse du même état en la reliant
à l’action et à la longueur. Cette correspondance spécifie le
contenu structurel qui est conservé lors du passage de la collection à ses valeurs.

L’identité Catalan--Pell précise une évaluation algébrique de cette
collection :
\[
\mathcal C_2(2-\sqrt3)
=\frac23G_{\rm Cat}-\frac{\pi}{12}\log(2+\sqrt3).
\]
Le coefficient $2/3$ provient des classes impaires du calendrier et
du poids de profondeur ; les deux premières dérivées logarithmiques dans cette section prennent les valeurs
$\pi/12$ et $1/4$. La réciprocité de l’unité de Pell relie
expansion et contraction, tandis que les deux arguments constitutifs
restent déterminés par la paire angulaire. Cette relation maintient
visibles la collection, ses lecteurs et les sections où ils sont évalués.

\Needspace{9\baselineskip}
La fonctionnelle de Barbero--Immirzi s’exprime ensuite comme différence
orientée d’une même fonction des canaux reconstruits :
\[
\gamma=\mathcal H(\psi(r_-))-\mathcal H(\psi(r_+)).
\]
L’incidence et le retour dodécaphasique fixent ses termes ;
l’inversion constitutive permet de les transporter vers la réponse du
vide. La marque d’orientation étant fixée, l’échange des canaux
change le signe de $\gamma$ ; la conjugaison conjointe conserve la
fonctionnelle. Cette distinction prouve que son contenu dépasse le spectre
sans étiquettes et se maintient sous les amplifications normalisées.

% BEGIN SINTESIS_ADITIVA_20260922
\par
La section~\ref{delta22:iii:restitucion} réunit le potentiel spectral de la réponse constitutive, la restitution globale des canaux au moyen de la vitesse normalisée et de Barbero--Immirzi, ainsi que ses bornes de stabilité. Sur l’image engendrée, avec les étiquettes et les bases conservées, la composition retrouve la paire angulaire, la section d’action et la lecture périodique du registre dodécaphasique.
% END SINTESIS_ADITIVA_20260922

La résolution harmonique du motif dodécaphasique conserve exactement
les modes $3,9,4,8$, avec les coefficients non normalisés $72,72,-6,-6$.
Le rapport des modules $12:1$ exprime les multiplicités orientées.
Le défaut d’incidence $-1/12$ et le coefficient singulier $1/24$
restent liés à leurs propres opérations. La comparaison explicite
avec le lecteur à poids égaux, dont le coefficient est $1/1080$,
précise comment la chaîne orientée intervient dans la normalisation
du retour sans remplacer les fonctionnelles par leurs pôles.

\subsection*{Du bilan arithmétique aux relations électriques}

La discrépance des vacances fournit une polarisation centrée dont la
grandeur reste inférieure à une unité et dont l’accroissement satisfait
un bilan temporel exact. Les changements entre les conventions de
plancher et de plafond sont enregistrés par un terme de frontière. Une
section de charge et une durée orientée réalisent ce bilan comme
charge centrée et courant, en conservant sa loi télescopique.

Le préfacteur électronique conserve également sa structure antécédente.
L’énumération des fibres dans $D_9^3$ fixe le mode singulier tritique
de valeur $1/2$ ; le spectre complet de son appariement prouve
\[
\|[P_{\Sigma,3},P_{\Pi,3}]\|=\sqrt3/4.
\]
Son défaut quadratique $3/4$ coïncide avec la limite du quotient des
secteurs $90/120$. Cette égalité relie des résultats d’opérateurs
spécifiés, chacun avec son espace et son action : le préfacteur est
une norme d’incompatibilité arithmétique, tandis que la réponse
intérieure dépend de ses canaux contractifs.

La section positive de charge déterminée par
$Z_{\rm vac}e_a^2=2\alpha h_a$ réunit le couplage, l’action et
l’impédance. Les relations de von Klitzing, de conductance, de flux
et de Josephson satisfont exactement
\[
R_KG_0=2,\qquad G_0Z_{\rm vac}=4\alpha,\qquad
K_{J,a}\Phi_{0,a}=1,\qquad K_{J,a}^2R_K=4/h_a.
\]
Le retour de l’action conserve résistance et conductance ; charge et
flux se transforment avec $R_{\rm act}^{1/2}$ et Josephson avec
$R_{\rm act}^{-1/2}$. Ces dépendances expliquent conjointement
quelles grandeurs retiennent l’action et lesquelles conservent uniquement leur
quotient constitutif.

La réalisation dimensionnelle préserve toutes ces identités lorsque
les bases d’action, de charge et de vitesse changent. Les valeurs
adimensionnelles obtenues appartiennent à la mise en coordonnées interne de la
construction ; leur reconnaissance métrologique exige l’évaluation
indépendante des sections dimensionnelles et du régime physique
correspondant. La portée mathématique établie est la réponse
positive et inversible, ses compositions et sa covariance dans les
domaines déclarés.

L’ensemble montre comment une même généalogie APP--TRIT--TPK conserve
les relations nécessaires pour passer de chemins et d’opérations
arithmétiques à des canaux, des moments, une réponse constitutive et des relations
électriques. Les valeurs acquièrent ainsi une provenance expliquée et
une information récupérable : elles expriment les opérations qui les
produisent et permettent de reconnaître, au sein de leur représentation, la
structure qui demeure en elles.
\par\endgroup

\appendix
\section{Dépendances, procédures et reproduction}

Les constructions utilisées sont développées dans l’article :
le registre dodécaphasique dans la section~\ref{sec:k-registro}, la clôture
positionnelle dans la section~\ref{sec:alpha}, l’action dans la
section~\ref{hmtII:sec:action} et la paire angulaire dans la
section~\ref{hmtII:sec:ii-angulos}. La réponse constitutive de la
section~\ref{iii:sec:constitutiva} compose ces sorties en conservant
les étiquettes de feuillet et d’orientation.

\subsection{Ordre de reproduction}

La reproduction distingue le producteur de coordonnées des évaluateurs
ultérieurs. Le premier énumère les états admissibles et applique les règles
d’APP, de TRIT et de TPK. Ses sorties régionales sont des entrées légitimes des
opérateurs ultérieurs du même développement. Le programme d’évaluation
du vide n’est pas présenté comme un substitut de ce producteur.

L’ordre de lecture du paquet est : définitions du noyau ; énumération
et raffinement régional ; registre signé et transformation de $K$ ;
clôture positionnelle de $\alpha$ ; section d’action et retour ; paire angulaire ;
canaux contractifs ; opérateur constitutif ; moments cycliques et fonctionnelle
de Barbero ; sections dimensionnelles et comparaison ultérieure.

Chaque étape conserve son domaine. Un fichier de chiffres régionaux porte
l’empreinte du producteur et non une attribution métrologique. Le tuple de douze
composantes de $K$ a son propriétaire dans le chapitre du registre. Les
coefficients d’action, l’incidence $90/120$ et les multiplicités
$12:-1$ se lisent dans leurs démonstrations antérieures à l’évaluation. Le
programme local reproduit ces compositions avec une précision déclarée.

\subsection{Trois vérifications de fonctions différentes}

La conservation documentaire vérifie l’identité des fichiers et la continuité
des copies. La preuve mathématique établit les identités et les limites
dans leurs domaines. L’exécution numérique vérifie une évaluation concrète,
sa stabilité et les résidus des identités. Aucune de ces trois
vérifications ne remplace les autres.

Le paquet joint comprend un manifeste des sources, les textes LaTeX,
les données régionales utilisées et les programmes de vérification. Les preuves
d’inversion, de monotonie, de moments, de transport et de covariance se trouvent
dans le corps de l’article. Les sources documentaires sont citées par
provenance ; le lecteur peut vérifier les arguments développés ici
sans utiliser un résultat numérique cible comme générateur.

La clôture positionnelle de $\alpha$ est démontrée dans la
proposition~\ref{prop:alpha-normalizacion}. Le
théorème~\ref{thm:alpha-dos-publicaciones-completas} identifie ses
préfixes à ceux de la représentation analytique corrélée du même
état, au moyen de la récurrence explicite et des cylindres compatibles.
L’inversion constitutive de
la section~\ref{iii:sec:constitutiva} est formulée dans la normalisation
interne indiquée : les changements d’unités sont annulés avant d’appliquer
cette inverse. Ces conditions délimitent les résultats qu’utilisent
les compositions ultérieures.

\section{Énumération exhaustive et détermination finie des opérateurs de transport}
\label{iii:nucleo-finito}

Cet appendice développe deux opérations finies avec des données de départ
différentes. La première énumère les conditions initiales des deux curseurs
et construit leurs émissions au moyen de la récurrence arithmétique déclarée.
La seconde détermine les opérateurs de transport à partir du registre
régional des transitions. La séparation de leurs domaines permet de préciser
ce que produit chaque démonstration : l’énumération construit le catalogue des
émissions ; la reconstruction matricielle détermine les opérateurs compatibles
avec le registre. Le théorème de reconstruction n’est pas employé comme démonstration
de la production antérieure de ce registre.

Les primitives de la première opération sont le support $I_9^2$, les quatre
orientations cardinales, le sélecteur tritique, la lecture antérieure au
déplacement, les inversions d’orientation et la règle d’émission.
Leurs définitions sont réunies ci-dessous. Celles de la seconde sont les
quatre matrices de transitions de~\eqref{iii:nf-registro}.
Leur reconstruction unique est démontrée dans le
théorème~\ref{iii:nf-unicidad} ; le calendrier qui en résulte est exprimé
dans~\eqref{iii:nf-calendario}. Ces déterminations conservent
les données de départ déclarées pour chaque opération.

\subsection{Domaine complet des conditions initiales}

Soit $I_9=\{0,\ldots,8\}$ et soit
\[
 \mathcal D=\{N,E,S,O\},\qquad
 v_N=(-1,0),\quad v_E=(0,1),\quad
 v_S=(1,0),\quad v_O=(0,-1).
\]
Une condition initiale de curseur est $s=(i,j,d)\in
\mathcal S=I_9^2\times\mathcal D$. Les deux curseurs sont indépendants
dans le domaine initial
\begin{equation}
 \Omega=\mathcal S_+\times\mathcal S_\times,
 \qquad |\mathcal S|=9^2\cdot4=324,
 \qquad |\Omega|=324^2=104\,976.
 \label{iii:nf-dominio}
\end{equation}
L’indice $i$ correspond à l’étiquette positive $i+1$ d’APP. Pour
des entiers positifs, la réduction positive est $\rho_9(n)=1+[(n-1)]_9$,
où $[a]_b$ désigne le représentant dans $\{0,\ldots,b-1\}$.
Les évaluations résiduelles des deux feuillets sont
\[
 \sigma(i,j)=\rho_9(i+j+2),\qquad
 \varpi(i,j)=\rho_9((i+1)(j+1)).
\]
Le lecteur multiplicatif fini utilise
\begin{equation}
 \begin{array}{c|rrrrrrrrr}
 a&1&2&3&4&5&6&7&8&9\\ \hline
 \ell_9(a)&0&1&0&2&5&0&4&3&0
 \end{array}.
 \label{iii:nf-log}
\end{equation}
Sur les unités, $\ell_9$ est l’exposant discret en base $2$.
Sur $3,6,9$, la valeur $0$ constitue l’extension déclarée de
l’observable d’émission ; l’appartenance à ce secteur demeure dans le
registre de trajectoire.

Pour $1\le t\le54$, on définit
\[
 r_t=1+[(t-1)]_9,
 \qquad
 \tau_t=
 \begin{cases}
 +1,&r_t\in\{1,4,7\},\\
 -1,&r_t\in\{2,5,8\},\\
 0,&r_t\in\{3,6,9\}.
 \end{cases}
\]
Chaque curseur conserve sa position $z_t$ immédiatement avant le pas
$t$ et son vecteur d’orientation $v_t$. Le curseur additif est actif
lorsque $\tau_t=+1$ ; le multiplicatif, lorsque $\tau_t=-1$. Pour chacun
d’entre eux, avec l’indicateur d’activité $a_t$, la mise à jour est
\begin{equation}
 z_{t+1}=[z_t+a_tv_t]_9,
 \qquad
 v_{t+1}=
 \begin{cases}
 -v_t,&t\in\{27,54\},\\
 v_t,&t\notin\{27,54\}.
 \end{cases}
 \label{iii:nf-actualizacion}
\end{equation}
La lecture du pas $t$ est effectuée en $z_t$, avant le déplacement.
L’événement neutre conserve les deux positions. Cette convention fixe
sans ambiguïté les extrémités des fenêtres et l’effet des
inversions cardinales.

Dans les fenêtres $W_m=\{9m-8,\ldots,9m\}$, $1\le m\le6$, soient
\begin{align}
 S_m(s_+)&=\sum_{\substack{t\in W_m\\\tau_t=+1}}
                 \sigma(z_t^+),&
 E_m(s_\times)&=\sum_{\substack{t\in W_m\\\tau_t=-1}}
                 \ell_9(\varpi(z_t^\times)),
 \label{iii:nf-firmas}\\
 s_m&=[S_m]_{10},&e_m&=[E_m]_{10}.\notag
\end{align}
Dans chaque fenêtre, il y a trois événements de chaque type. Les signatures $s$ et $e$
sont les six composantes de ces résidus, non les entiers $S_m,E_m$
antérieurs à leur réduction.

\subsection{Recensement des signatures et des émissions}

L’émission d’une paire de conditions initiales est déterminée par
\begin{equation}
 U_m=100s_m+10[s_m+e_m]_{10}+[3s_m+5e_m+1]_{10},
 \qquad 1\le m\le6.
 \label{iii:nf-emision}
\end{equation}
Le terme final $1$ est $[7\cdot3]_{10}$, correspondant aux trois
événements neutres. Les coefficients de cette règle font partie de la
récurrence déclarée. Les constantes réelles étudiées ultérieurement
n’interviennent ni dans l’énumération de son domaine ni dans son évaluation.

\begin{proposition}[Factorisation injective du catalogue fini]
\label{iii:nf-inyectividad}
L’émission distingue chaque paire de signatures $(s,e)$. Il y a $18$ signatures
additives et $26$ multiplicatives ; leurs paires produisent $468$ émissions.
Les multiplicités de ces émissions sont $144$, $432$ et $1944$,
avec $432$, $18$ et $18$ émissions respectivement.
\end{proposition}

\begin{proof}
Le chiffre des centaines de $U_m$ restitue $s_m$. Si $b_m$ est son chiffre des
dizaines, alors $e_m=[b_m-s_m]_{10}$. Le chiffre des unités vérifie
la troisième congruence de \eqref{iii:nf-emision}. Par conséquent, le
couplage est injectif.

Le recensement s’obtient en parcourant, dans l’ordre cartésien complet,
$i=0,\ldots,8$, $j=0,\ldots,8$ et $d=N,E,S,O$, et en appliquant les $54$
pas de \eqref{iii:nf-actualizacion}. Les $324$ conditions additives
se répartissent entre les $18$ signatures suivantes, chacune avec $18$
préimages. Dans les tableaux, une suite de six chiffres représente
les six composantes ordonnées d’une signature.
\begin{center}
\begin{tabular}{ccc}
$122564$&$122898$&$212645$\\
$212988$&$221456$&$221889$\\
$456221$&$465898$&$546988$\\
$564122$&$645212$&$654889$\\
$889221$&$889654$&$898122$\\
$898465$&$988212$&$988546$
\end{tabular}
\end{center}
Pour le curseur multiplicatif, les signatures et leurs multiplicités sont
\begin{center}
\begin{tabular}{rr@{\qquad}rr}
Signature&Multiplicité&Signature&Multiplicité\\ \hline
$000000$&$108$&$555555$&$24$\\
$177393$&$8$&$555717$&$8$\\
$177555$&$8$&$555771$&$8$\\
$177771$&$8$&$555933$&$8$\\
$339555$&$8$&$717555$&$8$\\
$339771$&$8$&$717717$&$8$\\
$339933$&$8$&$717933$&$8$\\
$393177$&$8$&$771177$&$8$\\
$393393$&$8$&$771339$&$8$\\
$393555$&$8$&$771555$&$8$\\
$555177$&$8$&$933339$&$8$\\
$555339$&$8$&$933555$&$8$\\
$555393$&$8$&$933717$&$8$
\end{tabular}
\end{center}
Ces tableaux résultent des sommes finies \eqref{iii:nf-firmas}, dont le
domaine, l’ordre de lecture et les inversions sont déjà spécifiés. Leurs sommes
de multiplicités sont $18\cdot18=324$ et
$24\cdot8+24+108=324$. Le domaine cartésien
\eqref{iii:nf-dominio} permet de réaliser n’importe quelle paire de signatures.
L’injectivité fournit $18\cdot26=468$ émissions et les
multiplicités
\[
 (18\cdot24,\ 18\cdot8)=(432,144),\qquad
 (18,\ 18\cdot24)=(18,432),\qquad
 (18,\ 18\cdot108)=(18,1944),
\]
où chaque paire indique le nombre d’émissions et la multiplicité de sa fibre.
Enfin,
\[
 432\cdot144+18\cdot432+18\cdot1944=104\,976.
\]
L’énumération épuise ainsi le domaine complet des conditions initiales.
\end{proof}

La réduction ternaire orientée est
\[
 \Psi(U)=([-U_1]_3,\ldots,[-U_6]_3).
\]
Appliquée aux $18\cdot26$ émissions des tableaux précédents,
elle produit le recensement des fibres suivant. Ici, $r$ compte les émissions
distinctes, non les conditions initiales.
\begin{equation}
 \begin{array}{c|rrrrrr}
 r&1&2&3&4&5&6\\ \hline
 \#\{w:\#\Psi^{-1}(w)=r\}&132&66&18&3&6&18
 \end{array}.
 \label{iii:nf-censo-ternario}
\end{equation}
La somme de la deuxième ligne est $243$, tandis que sa somme pondérée
par $r$ est $468$. On distingue ainsi l’image de $243$ mots
de l’espace ambiant $\mathbb F_3^6$, qui contient $729$ mots. Le certificat
joint conserve toutes les signatures avec leurs préimages, les $468$
émissions et leur réduction ternaire ; les formules précédentes fournissent
une procédure complète pour les reconstruire.

\subsection{Registre des transitions et reconstruction des relèvements}

Dans cette partie, on fixe le registre régional du propriétaire cité.
Soient $b_0^\chi,\ldots,b_4^\chi\in\mathbb F_3^6$, avec
$\chi\in\{\mathsf C,\mathsf P,\mathsf A\}$, ses trois suites.
Le registre organise les transitions $b_0\to b_1\to b_2$ dans
$X_0,Y_0$ et les transitions $b_2\to b_3\to b_4$ dans $X_1,Y_1$,
avec des vecteurs lignes et le même ordre régional :
\begingroup\small
\begin{equation}
\begin{aligned}
X_0&=\begin{pmatrix}
0&1&0&2&1&1\\0&1&2&2&2&2\\2&0&1&1&0&1\\
1&2&1&2&2&1\\1&2&1&2&0&0\\1&1&2&2&0&2
\end{pmatrix},&
Y_0&=\begin{pmatrix}
0&1&2&2&2&2\\0&1&0&2&1&1\\1&2&1&2&2&1\\
1&0&2&0&1&1\\1&1&2&2&0&2\\1&2&1&0&2&0
\end{pmatrix},\\[1ex]
X_1&=\begin{pmatrix}
0&1&0&2&1&1\\0&0&2&1&1&1\\1&0&2&0&1&1\\
0&1&2&2&2&2\\1&2&1&0&2&0\\0&1&0&2&1&0
\end{pmatrix},&
Y_1&=\begin{pmatrix}
0&0&2&1&1&1\\1&1&0&2&2&1\\0&1&2&2&2&2\\
1&0&2&0&1&1\\0&1&0&2&1&0\\0&1&0&2&0&0
\end{pmatrix}.
\end{aligned}
\label{iii:nf-registro}
\end{equation}
\endgroup
La production précédente s’exprime chez son propriétaire au moyen de la
composition de sélection, de transport et de mise à jour :
\[
 x_{j+1}^{\chi}
 =\mathcal U_{9j+9}\circ\cdots\circ\mathcal U_{9j+1}(x_j^{\chi}),
 \qquad b_j^\chi=\Pi_6(x_j^\chi),
 \qquad\mathcal U_t=\mathsf{Upd}_t\circ\mathsf{Tra}_t\circ\mathsf{Sel}_t.
\]
Le résultat suivant part du registre \eqref{iii:nf-registro}.
Sa démonstration permet de reconstruire les opérateurs et de vérifier la
compatibilité des transitions ; elle n’attribue pas à l’inversion matricielle
la production coordonnée des états $x_j^\chi$.

\begin{proposition}[Unicité sur le registre régional]
\label{iii:nf-unicidad}
Les systèmes $X_aL_a=Y_a$, $a=0,1$, ont une solution unique sur
$\mathbb F_3$. Les deux solutions sont
\begingroup\small
\begin{equation}
L_0=\begin{pmatrix}
2&2&2&1&2&1\\2&1&2&2&1&1\\1&0&0&1&0&2\\
0&0&1&1&1&0\\2&2&2&0&2&1\\2&1&2&1&0&0
\end{pmatrix},\qquad
L_1=\begin{pmatrix}
2&1&2&0&2&2\\1&2&1&1&2&0\\1&2&1&1&1&2\\
0&1&0&0&2&1\\2&0&2&1&0&1\\0&2&2&2&1&1
\end{pmatrix}.
\label{iii:nf-lifts}
\end{equation}
\endgroup
\end{proposition}

\begin{proof}
L’élimination dans $\mathbb F_3$ donne $\det X_0=1$ et
$\det X_1=2=-1$. Chaque application $L\mapsto X_aL$ est donc
un automorphisme de l’espace des matrices de dimension $36$ ; dans une
vectorisation par colonnes, elle est représentée par $I_6\otimes X_a$.
On obtient $L_a=X_a^{-1}Y_a$. La multiplication des matrices de
\eqref{iii:nf-lifts} par les matrices $X_a$ respectives reproduit les
deux matrices $Y_a$ de \eqref{iii:nf-registro} ; existence et unicité
sont établies sur ce registre.
\end{proof}

Les trois suites résultantes, avec une séparation explicite en blocs,
sont
\begin{align*}
 B^{\mathsf C}&=010211\mid012222\mid010211\mid002111\mid110221,\\
 B^{\mathsf P}&=201101\mid121221\mid102011\mid012222\mid102011,\\
 B^{\mathsf A}&=121200\mid112202\mid121020\mid010210\mid010200.
\end{align*}
Chaque bloc conserve ses six positions. La concaténation n’est pas
identifiée à un entier en base dix.

\subsection{Exhaustivité du calendrier et stabilité de la reconstruction}

Un calendrier binaire est $c=(c_0,c_1,c_2,c_3)\in\{0,1\}^4$.
On définit le nombre de transitions incompatibles avec le registre par
\begin{equation}
 d(c)=\sum_{\chi\in\{\mathsf C,\mathsf P,\mathsf A\}}
       \sum_{j=0}^{3}
       \mathbf1\{b_j^\chi L_{c_j}\ne b_{j+1}^\chi\}.
 \label{iii:nf-calendario}
\end{equation}
Les multiplications des suites précédentes par
\eqref{iii:nf-lifts} donnent
\[
 d(c)=3\,d_{\rm H}(c,(0,0,1,1)),
\]
où $d_{\rm H}$ est la distance de Hamming : à chacune des
quatre positions, l’opérateur alternatif ne satisfait pas les trois
transitions régionales. Le recensement complet est
\begin{center}
\begin{tabular}{cr@{\qquad}cr@{\qquad}cr@{\qquad}cr}
$c$&$d(c)$&$c$&$d(c)$&$c$&$d(c)$&$c$&$d(c)$\\ \hline
$0000$&$6$&$0100$&$9$&$1000$&$9$&$1100$&$12$\\
$0001$&$3$&$0101$&$6$&$1001$&$6$&$1101$&$9$\\
$0010$&$3$&$0110$&$6$&$1010$&$6$&$1110$&$9$\\
$0011$&$0$&$0111$&$3$&$1011$&$3$&$1111$&$6$
\end{tabular}
\end{center}
Par conséquent, $0011$ est l’unique calendrier qui reproduit les
trois suites complètes.

La sensibilité aux entrées des matrices admet également une
preuve exhaustive sans approximations. Une modification unitaire
remplace $L_a$ par $L_a+\lambda E_{ij}$, avec
$a\in\{0,1\}$, $1\le i,j\le6$ et
$\lambda\in\mathbb F_3\setminus\{0\}$. Il y a
$2\cdot6\cdot6\cdot2=144$ modifications de ce type. Comme
$X_a$ est inversible,
\[
 X_a(L_a+\lambda E_{ij})-Y_a
 =\lambda X_aE_{ij}\ne0.
\]
Toute modification détruit au moins une des transitions du
registre correspondant. Le certificat exécutable énumère en outre
les $144$ modifications et conserve, pour chacune, le nombre de
transitions non satisfaites.

La propriété démontrée est la rigidité des opérateurs par rapport au
registre régional fixé. Pour reproduire une construction antérieure
du registre lui-même, l’action coordonnée de ses opérateurs
de sélection, de transport et de mise à jour est nécessaire. Cette dépendance reste
séparée du théorème de rigidité, du catalogue fini des émissions et
des réalisations analytiques ultérieures.

\subsection{Certificat de reproduction et résidence des données}

Le programme \texttt{verificar\_nucleo\_finito.py} parcourt les $324$
conditions de chaque curseur et forme le produit des classes résultantes.
Le catalogue documentaire est chargé après cette construction pour
comparer les émissions et leurs multiplicités. Dans la seconde partie,
les quatre matrices de transitions sont extraites de la copie littérale
du propriétaire ; les solutions sont calculées par élimination exacte
modulo $3$ avant de les comparer aux matrices imprimées. Le reçu
\texttt{NUCLEO\_FINITO.json} conserve les domaines, les préimages,
les émissions, les $16$ calendriers et les $144$ modifications.

La vérification utilise des entiers et des résidus exacts. Elle ne reçoit ni valeurs
métrologiques ni développements réels des constantes. Sa portée
est celle des opérations finies et des données de départ spécifiées
dans cet appendice. Les démonstrations de compatibilité à profondeur
illimitée et les identifications physiques conservent leurs domaines
et leurs obligations d’exposition respectifs.

% Antecedentes recuperados. Fragmento autónomo en sus definiciones;
% requiere amsmath, amssymb, amsthm, booktabs y longtable del documento principal.
\section{Réalisation finie du transport avec mémoire de phase}
\label{iii:hist:antecedentes}

Les deux développements suivants réunissent des mécanismes complémentaires de la
construction : un transfert fini sur des positions APP, des modes et une phase,
et une collection de lecteurs cinématiques et constitutifs sur deux canaux
angulaires. Chaque réalisation conserve son espace d’états, ses opérations
et ses paramètres. L’opérateur fini a le domaine défini
ci-dessous ; le TPK enrichi du corps principal conserve en outre ses
propres coordonnées et prolongements. Les lecteurs réciproques sont présentés
comme des réalisations auxiliaires avec des formules et un domaine déclarés.

\subsection{Transfert fini sur les positions, les modes et la phase}
\label{iii:hist:transferencia}

Soit $\operatorname{dr}_9(n)$ le représentant dans $\{1,\ldots,9\}$ de
$n\bmod9$, avec $9$ comme représentant du résidu nul. Aux positions
$i,j\in\{0,\ldots,8\}$, on considère les deux tables
\[
 A_+(i,j)=\operatorname{dr}_9(i+j+2),\qquad
 A_\times(i,j)=\operatorname{dr}_9((i+1)(j+1)).
\]
Les entrées évaluées ici sont des entiers positifs. Ainsi, la convention
historique alternative $\operatorname{dr}_9(0)=0$ n’intervient pas dans ces tables.
L’état fini est
\[
 \Omega=(\mathbb Z/9\mathbb Z)^2\times\{+,\times\}
 \times\{N,E,S,O\}\times\mathbb Z/3\mathbb Z,
 \qquad |\Omega|=1944.
\]
Nous écrivons $x=(i,j,m,d,\tau)$. La transition conserve l’ordre suivant :
on choisit $d'$, on déplace la position, on lit APP avec le mode antérieur,
on réduit la lecture modulo $3$, on met à jour le mode et on avance la phase.
Avec
\[
 \delta_N=(-1,0),\quad\delta_E=(0,1),\quad
 \delta_S=(1,0),\quad\delta_O=(0,-1),
\]
les opérations sont
\begin{align*}
 (i',j')&=(i,j)+\delta_{d'}\pmod9,\\
 a&=A_m(i',j')\pmod3,\\
 m'&=\begin{cases}+,&a=1,\\\times,&a=2,\\m,&a=0,\end{cases}
 &\tau'&=\tau+1\pmod3.
\end{align*}
Ainsi, $F_{d'}(x)=(i',j',m',d',\tau')$. Le coût de rotation et de changement de mode est
\[
 c_\lambda(x,d')=\mathbf1_{d'\ne d}+\lambda\mathbf1_{m'\ne m},
 \qquad \beta>0,\quad\lambda\geq0,
\]
et l’opérateur agit sur la base d’états par
\begin{equation}
 T_{\beta,\lambda}e_x
 =\sum_{d'\in\{N,E,S,O\}}
 e^{-\beta c_\lambda(x,d')}e_{F_{d'}(x)}.
 \label{iii:hist:T}
\end{equation}
Les paramètres $\beta,\lambda$ sont des paramètres déclarés de ce modèle
fini. La table de référence utilise $(\beta,\lambda)=(1,1)$ ; le cas
$\lambda=0$ appartient aux tests de sensibilité.

\begin{proposition}[Équivariance de l’antécédent fini]
\label{iii:hist:equivarianza}
L’action
\[
 \Phi_{u,v,w}(i,j,m,d,\tau)
 =(i+3u,j+3v,m,d,\tau+w),\qquad
 (u,v,w)\in(\mathbb Z/3\mathbb Z)^3,
\]
commute avec $T_{\beta,\lambda}$.
\end{proposition}
\begin{proof}
On a $\operatorname{dr}_9(n)\equiv n\pmod3$. Par conséquent,
$A_+(i,j)\equiv i+j+2$ et
$A_\times(i,j)\equiv(i+1)(j+1)\pmod3$.
La translation par des multiples de trois préserve les deux lectures réduites.
Elle commute également avec le déplacement choisi et préserve le mode antérieur,
le mode postérieur et les deux directions. Elle conserve par conséquent chaque
indicateur du coût. Enfin, ajouter $w$ à la phase commute avec l’ajout d’une
unité. Chaque transition et son poids sont transportés vers la transition et le poids
correspondants ; la somme dans \eqref{iii:hist:T} conserve cette égalité.
\end{proof}

\begin{theorem}[Décomposition complète par caractères]
\label{iii:hist:descomposicion}
Pour $\omega=e^{2\pi i/3}$, l’espace $\mathbb C^\Omega$ se décompose
unitairement en vingt-sept sous-espaces de dimension $72$. Ses blocs
satisfont
\[
 T(k_a,k_b,k_t)=\omega^{k_t}T(k_a,k_b,0),
 \qquad k_a,k_b,k_t\in\{0,1,2\}.
\]
\end{theorem}
\begin{proof}
Nous écrivons $i=3a+r$, $j=3b+s$, avec $a,b,r,s\in\{0,1,2\}$.
Chaque fibre fine $u=(r,s,m,d)$ a $27$ états $(a,b,\tau)$.
Nous définissons l’isométrie $U_k:\mathbb C^{72}\to\mathbb C^\Omega$ par
\[
 U_ke_u=\frac1{\sqrt{27}}
 \sum_{a,b,\tau=0}^2
 \omega^{-(k_aa+k_bb+k_t\tau)}e_{(3a+r,3b+s,m,d,\tau)}.
\]
L’orthogonalité des caractères prouve $U_k^*U_\ell=\delta_{k\ell}I$ et
$\sum_kU_kU_k^*=I$. Si une transition franchit des frontières fines avec des accroissements
grossiers $\Delta a,\Delta b$, son coefficient dans $U_k^*TU_k$ est
\[
 e^{-\beta c_\lambda}\omega^{k_a\Delta a+k_b\Delta b+k_t}.
\]
Cela démontre $TU_k=U_kT(k)$ et la décomposition de l’opérateur. L’accroissement
temporel vaut toujours un, de sorte que le dernier facteur peut être extrait de
tout le bloc. Lorsque l’on élève au cube, ce facteur temporel disparaît, bien que
les caractères spatiaux conservent leur fonction.
\end{proof}

La décomposition précédente repose sur l’équivariance et l’orthogonalité
des caractères. La non-négativité du bloc trivial garantit en outre une
valeur propre réelle non négative égale au rayon spectral. Cette propriété
s’applique à l’espace complet, en conservant également ses états transitoires.

\subsection{Décomposition spectrale et sensibilité paramétrique}
\label{iii:hist:espectros}

Soient $\rho_1\geq\rho_2$ les deux plus grands modules, comptés avec
multiplicité, et $g=\rho_1-\rho_2$. Le tableau suivant contient les
neuf blocs spatiaux ; le théorème précédent détermine les vingt-sept
blocs temporels complets : pour chaque ligne et chaque $t=0,1,2$, la valeur propre
dominante est $\omega^t\mu_1(k_a,k_b,0)$, tandis que $\rho_1,\rho_2,g$
restent identiques. Cette forme conserve explicitement les trois phases.
Dans ce premier tableau, on fixe $\beta=\lambda=1$.

\begin{center}
\small
\begin{tabular}{ccrrrr}
\toprule
$k_a$&$k_b$&$\mu_1(k_a,k_b,0)$&$\rho_1$&$\rho_2$&$g$\\
\midrule
0&0&1.730888896&1.730888896&1.083076285&0.647812612\\
0&1&1.459962951&1.459962951&1.116843785&0.343119167\\
0&2&1.459962951&1.459962951&1.116843785&0.343119167\\
1&0&1.459962951&1.459962951&1.116843785&0.343119167\\
1&1&-1.508515719&1.508515719&1.001353790&0.507161929\\
1&2&-1.514010372&1.514010372&0.971977167&0.542033205\\
2&0&1.459962951&1.459962951&1.116843785&0.343119167\\
2&1&-1.514010372&1.514010372&0.971977167&0.542033205\\
2&2&-1.508515719&1.508515719&1.001353790&0.507161929\\
\bottomrule
\end{tabular}
\end{center}

Les signes des valeurs propres dominantes et leurs phases font partie de la
sortie : les quatre secteurs avec $k_a,k_b\ne0$ ont une valeur dominante réelle négative
pour $k_t=0$. Son module conserve le taux spectral et son signe conserve une
information supplémentaire qui intervient lors de la composition des itérations.

Les cinq applications complètes de sensibilité, avec les lignes $k_a$ et les colonnes
$k_b$, sont les suivantes ; toutes correspondent à $k_t=0$ :
\begin{align*}
\rho(1,1)&=\begin{pmatrix}
1.730888896&1.459962951&1.459962951\\
1.459962951&1.508515719&1.514010372\\
1.459962951&1.514010372&1.508515719
\end{pmatrix},\\
\rho(1,0)&=\begin{pmatrix}
2.103638324&1.685504327&1.685504327\\
1.685504327&1.853245966&1.853245966\\
1.685504327&1.853245966&1.853245966
\end{pmatrix},\\
\rho(1,2)&=\begin{pmatrix}
1.656915306&1.433869155&1.433869155\\
1.433869155&1.423885161&1.432116157\\
1.433869155&1.432116157&1.423885161
\end{pmatrix},\\
\rho(0.5,1)&=\begin{pmatrix}
2.480049445&2.124801655&2.124801655\\
2.124801655&2.281861809&2.286772914\\
2.124801655&2.286772914&2.281861809
\end{pmatrix},\\
\rho(2,1)&=\begin{pmatrix}
1.226049377&1.131390710&1.131390710\\
1.131390710&1.028041508&1.028201900\\
1.131390710&1.028201900&1.028041508
\end{pmatrix}.
\end{align*}
Ici, $\rho(\beta,\lambda)$ désigne l’application des rayons, non l’opérateur de
transfert. L’échange $i\leftrightarrow j$ accompagné de l’échange
correspondant des directions préserve APP, les transitions et le coût ;
ces applications sont donc symétriques. La variation des paramètres modifie les rayons tout en maintenant
la décomposition par caractères, qui dépend de la symétrie des
opérations. Les tableaux montrent des approximations numériques de spectres
finis. Les identités de transition, d’équivariance et de factorisation
reposent sur les preuves précédentes ; leurs conséquences numériques sont
conservées avec les données complètes des vingt-sept composantes.

\subsection{Lecteurs réciproques et tests de sensibilité}
\label{iii:hist:n6d}

Pour reproduire la réalisation angulaire antécédente, on conserve les mesures
numériques décimales en degrés
\[
 a_{\mathrm h}=7.297352569283802,\qquad
 c_{\mathrm h}=2.123738338968646.
\]
Les angles et la coordonnée adimensionnelle correspondante sont distingués par
\[
 A_{\mathrm h}=a_{\mathrm h}{}^\circ,\qquad
 C_{\mathrm h}=c_{\mathrm h}{}^\circ,\qquad
 \alpha_{\mathrm h}=a_{\mathrm h}/1000.
\]
Ces données permettent de confronter l’antécédent ; les constantes du développement
en vigueur conservent leur génération antérieure APP--TRIT--TPK et leur représentation
ultérieure. La reproduction historique ne remplace pas cette génération par les
décimales écrites ici. La conversion angulaire et les canaux utilisés sont
\[
 x_{\mathrm h}=\frac{\pi a_{\mathrm h}}{180},\qquad
 y_{\mathrm h}=\frac{\pi c_{\mathrm h}}{180},\qquad
 q_\pm=e^{-(x_{\mathrm h}\pm y_{\mathrm h})},
 \qquad0<q_+<q_-<1.
\]
Le premier lecteur exprime
\[
 \epsilon_{90}=q_-^{90}-q_+^{90},\qquad
 c_\pm=\frac{c_{\mathrm{ref}}}{1\pm\epsilon_{90}}.
\]
Ici, $c_{\mathrm{ref}}$ ne fait que normaliser le test cinématique. Pour reproduire
ses unités, la même valeur de référence historique a été utilisée,
$299792458\,\mathrm{m\,s^{-1}}$.

\begin{proposition}[Moyenne harmonique et produit de deux feuillets]
\label{iii:hist:lectores}
Si $|\epsilon_{90}|<1$, la moyenne harmonique de $c_+$ et $c_-$ est
$c_{\mathrm{ref}}$. Dans le lecteur constitutif historique
\[
 \mu_{\mathrm h}^{\pm}=e^{\pm2\alpha_{\mathrm h}},\qquad
 \varepsilon_{\mathrm h}^{\pm}=e^{\mp2\alpha_{\mathrm h}},
\]
on a $\mu_{\mathrm h}^{\pm}\varepsilon_{\mathrm h}^{\pm}=1$ et
$Z_+/Z_-=e^{4\alpha_{\mathrm h}}$.
\end{proposition}
\begin{proof}
La somme $c_+^{-1}+c_-^{-1}=2/c_{\mathrm{ref}}$ élimine le terme orienté.
Dans chaque feuillet, la somme des deux exposants constitutifs est nulle.
D’autre part, $Z_\pm=\sqrt{\mu_{\mathrm h}^{\pm}/
\varepsilon_{\mathrm h}^{\pm}}=e^{\pm2\alpha_{\mathrm h}}$.
Le quotient des impédances conserve donc l’orientation qu’élimine
le produit constitutif.
\end{proof}

Le deuxième lecteur historique provient de séries géométriques convergentes :
\begin{align*}
 S_a(q)&=\frac{q^a}{1-q^{3a}}
       =\sum_{j\geq0}q^{a(1+3j)},\\
 \delta_{a,b}(q)&=S_a(q)-S_b(q),\\
 R_{a,b}&=\frac{\delta_{a,b}(q_-)-\delta_{a,b}(q_+)}
 {q_-^a-q_+^a}.
\end{align*}
La convergence est valable pour $|q|<1$ ; le dénominateur du dernier quotient est
positif pour les canaux et les entiers positifs considérés. Les lecteurs
auxiliaires historiques sont
\[
 \zeta_{12}=1-R_{90,120}^2,\qquad
 \Delta_{4,\mathrm h}=\frac{x_{\mathrm h}-y_{\mathrm h}}{810}
 \left(1-\frac{7\pi}{12\cdot729}\right),\qquad
 \xi_{\mathrm h}=1-\frac{27}{160}\Delta_{4,\mathrm h}.
\]
L’indice $\mathrm h$ maintient explicite cette définition historique de la
torsion, en évitant de l’identifier par le symbole à un autre lecteur en vigueur.

L’évaluation à partir des chaînes décimales déclarées donne
\begin{align*}
 q_+&=0.848377942576404356\ldots,&
 q_-&=0.913660151150557274\ldots,\\
 \epsilon_{90}&=0.000295169439289827959\ldots,&
 R_{90,120}&=0.933314535237713061\ldots,\\
 \zeta_{12}&=0.128923978314011665\ldots,&
 \xi_{\mathrm h}&=0.999981235497802379\ldots.
\end{align*}
Le tableau complet des perturbations discrètes est
\begin{center}
\begin{tabular}{rrr}
\toprule
$a$&$b$&$R_{a,b}$\\\midrule
90&120&0.933314535238\\
90&119&0.927013603649\\
90&121&0.939071551912\\
89&120&0.939066204789\\
91&120&0.927019444194\\
30&40&0.569215041549\\
60&80&0.834174862952\\
\bottomrule
\end{tabular}
\end{center}
La sensibilité angulaire distingue les deux coordonnées de la paire :
\begin{center}
\begin{tabular}{lrrr}
\toprule
Variable&Variation relative&$R_{90,120}$&Variation en ppm\\\midrule
$A_{\mathrm h}$&$-0.001$&0.933059250359&$-273.525022305$\\
$A_{\mathrm h}$&$+0.001$&0.933568846564&$+272.481908792$\\
$C_{\mathrm h}$&$-0.001$&0.933388163678&$+78.889202964$\\
$C_{\mathrm h}$&$+0.001$&0.933240822366&$-78.979667757$\\
\bottomrule
\end{tabular}
\end{center}

Les évaluations ont été réalisées avec $80$ et $120$ chiffres de précision à partir
des chaînes décimales déclarées. La précision de l’évaluation conserve
la spécification finie de ces données angulaires.

Enfin, ces lecteurs maintiennent leur place dans la généalogie documentaire.
L’opérateur constitutif en vigueur
$V=(I-T^{90})(I-T^{120})^{-1}$ évalue
$(1-q_\pm^{90})/(1-q_\pm^{120})$ dans ses deux feuillets. La fonctionnelle orientée
en vigueur utilise $12S_{90}-S_{120}$. La réalisation antécédente conserve le lecteur
$S_{90}-S_{120}$ et sa normalisation par $q_-^{90}-q_+^{90}$.
La relation entre eux s’exprime en maintenant chaque opération et ses poids,
au lieu de les remplacer par un nombre unique ou un même symbole.

\begin{thebibliography}{99}
\phantomsection
\addcontentsline{toc}{section}{Références}
\bibitem{feynman1948} R. P. Feynman, ``Space-Time Approach to Non-Relativistic Quantum Mechanics'', \emph{Reviews of Modern Physics} \textbf{20}, 367--387 (1948). \href{https://doi.org/10.1103/RevModPhys.20.367}{doi:10.1103/RevModPhys.20.367}.
\bibitem{feynman1949} R. P. Feynman, ``The Theory of Positrons'', \emph{Physical Review} \textbf{76}, 749--759 (1949). \href{https://doi.org/10.1103/PhysRev.76.749}{doi:10.1103/PhysRev.76.749}.
\bibitem{feynman1965} R. P. Feynman, ``The Development of the Space-Time View of Quantum Electrodynamics'', conférence Nobel, 11 décembre 1965. \href{https://www.nobelprize.org/prizes/physics/1965/feynman/lecture/}{Texte de la conférence}.
\bibitem{mustovicary} B. Musto et J. Vicary, ``Quantum Latin Squares and Unitary Error Bases'', \emph{Quantum Information and Computation} \textbf{16}, 1318--1332 (2016). \href{https://doi.org/10.26421/QIC16.15-16-4}{doi:10.26421/QIC16.15-16-4}.
\bibitem{delascuevas2020} G. De las Cuevas, T. Drescher et T. Netzer, ``Quantum Magic Squares: Dilations and Their Limitations'', \emph{Journal of Mathematical Physics} \textbf{61}, 111704 (2020). \href{https://doi.org/10.1063/5.0022344}{doi:10.1063/5.0022344}.
\bibitem{delascuevas2023} G. De las Cuevas, T. Netzer et I. Valentiner-Branth, ``Magic Squares: Latin, Semiclassical, and Quantum'', \emph{Journal of Mathematical Physics} \textbf{64}, 022201 (2023). \href{https://doi.org/10.1063/5.0127393}{doi:10.1063/5.0127393}.
\bibitem{dirac1933} P. A. M. Dirac, ``The Lagrangian in Quantum Mechanics'', \emph{Physikalische Zeitschrift der Sowjetunion} \textbf{3}, 64--72 (1933). \href{https://www.informationphilosopher.com/solutions/scientists/dirac/Lagrangian_1933.pdf}{Reproduction de l'article original}.
\bibitem{pauli1946} W. Pauli, ``Exclusion Principle and Quantum Mechanics'', conférence Nobel, 13 décembre 1946, dans \emph{Nobel Lectures, Physics 1942--1962}, Elsevier, 1964, pp. 27--43, notamment p. 42. \href{https://www.nobelprize.org/uploads/2018/06/pauli-lecture.pdf}{Texte de la conférence}.
\bibitem{codata2018} E. Tiesinga, P. J. Mohr, D. B. Newell et B. N. Taylor, ``CODATA Recommended Values of the Fundamental Physical Constants: 2018'', \emph{Reviews of Modern Physics} \textbf{93}, 025010 (2021). \href{https://doi.org/10.1103/RevModPhys.93.025010}{doi:10.1103/RevModPhys.93.025010}. \href{https://physics.nist.gov/cuu/Constants/ArchiveASCII/allascii_2018.txt}{Tableaux archivés de 2018}.
\bibitem{codata2022} P. J. Mohr, D. B. Newell, B. N. Taylor et E. Tiesinga, ``CODATA Recommended Values of the Fundamental Physical Constants: 2022'', \emph{Reviews of Modern Physics} \textbf{97}, 025002 (2025). \href{https://doi.org/10.1103/RevModPhys.97.025002}{doi:10.1103/RevModPhys.97.025002}. \href{https://physics.nist.gov/cuu/Constants/Table/allascii.txt}{Tableau des constantes}.
\bibitem{flm} I. B. Frenkel, J. Lepowsky et A. Meurman, \emph{Vertex Operator Algebras and the Monster}, Pure and Applied Mathematics, vol. 134, Academic Press, 1988.
\bibitem{borcherds} R. E. Borcherds, ``Monstrous Moonshine and Monstrous Lie Superalgebras'', \emph{Inventiones Mathematicae} \textbf{109}, 405--444 (1992). \href{https://doi.org/10.1007/BF01232032}{doi:10.1007/BF01232032}.
\bibitem{conwaysloane} J. H. Conway et N. J. A. Sloane, \emph{Sphere Packings, Lattices and Groups}, 3\textsuperscript{e} éd., Grundlehren der mathematischen Wissenschaften 290, Springer, 1999. \href{https://doi.org/10.1007/978-1-4757-6568-7}{doi:10.1007/978-1-4757-6568-7}.
\bibitem{bipm2018} Conférence générale des poids et mesures, \emph{Resolution 1 of the 26th CGPM: On the revision of the International System of Units}, 2018. \href{https://www.bipm.org/en/committees/cg/cgpm/26-2018/resolution-1}{Résolution officielle}. \href{https://doi.org/10.59161/CGPM2018RES1E}{doi:10.59161/CGPM2018RES1E}.
\bibitem{dirac1931} P. A. M. Dirac, ``Quantised singularities in the electromagnetic field'', \emph{Proceedings of the Royal Society of London. Series A} \textbf{133}, 60--72 (1931). \href{https://doi.org/10.1098/rspa.1931.0130}{doi:10.1098/rspa.1931.0130}.
\bibitem{lep2006} ALEPH, DELPHI, L3, OPAL et SLD Collaborations; LEP Electroweak Working Group; SLD Electroweak Group et SLD Heavy Flavour Group, ``Precision electroweak measurements on the Z resonance'', \emph{Physics Reports} \textbf{427}, 257--454 (2006). \href{https://doi.org/10.1016/j.physrep.2005.12.006}{doi:10.1016/j.physrep.2005.12.006}.
\end{thebibliography}

\end{document}
